% !TeX program = xelatex
% !TeX encoding = UTF-8
\documentclass[UTF8,10pt,a4paper,fontset=fandol]{ctexart}
\usepackage{amsmath,amssymb,graphicx,tabularx,array,multirow,enumitem,titlesec,xcolor,ragged2e,etoolbox}
\usepackage[a4paper,top=1.6cm,bottom=1.6cm,left=1.6cm,right=1.6cm]{geometry}
\setlength{\parindent}{2em}
\setlength{\parskip}{0.15em}
\setlength{\tabcolsep}{4pt}
\renewcommand{\arraystretch}{1.15}
\setlist[enumerate]{leftmargin=2.6em,itemsep=0.1em,topsep=0.2em}
\titleformat{\section}{\large\bfseries}{\thesection}{0.6em}{}
\titleformat{\subsection}{\normalsize\bfseries}{\thesubsection}{0.6em}{}
\titleformat{\subsubsection}{\normalsize\bfseries}{\thesubsubsection}{0.6em}{}
\titlespacing*{\section}{0pt}{0.7em}{0.25em}
\titlespacing*{\subsection}{0pt}{0.55em}{0.2em}
\titlespacing*{\subsubsection}{0pt}{0.4em}{0.15em}
\raggedbottom
\BeforeBeginEnvironment{tabularx}{\par\noindent}
\newcommand{\fieldvalue}[1]{#1}
\newcommand{\handwritten}[1]{#1}
\newcommand{\checkboxfield}[1]{\ifstrequal{#1}{checked}{☑}{\ifstrequal{#1}{unclear}{?}{☐}}}
\begin{document}

\begin{center}
{\Large\bfseries 永水性滤器使用记录}
\end{center}

\noindent
\begin{tabular}{@{}p{0.23\linewidth}p{0.73\linewidth}@{}}
文件编号： & REC-YS-ZOP-PD-02-012-a \\
\end{tabular}

\small
\noindent
\begin{tabular}{|p{0.18\linewidth}|p{0.25\linewidth}|p{0.18\linewidth}|p{0.29\linewidth}|}
\hline
过滤器进厂批号 & & 过滤器序列号 &
% #VALUE_ID: LEX-P0001-V0001
% #FIELD_VALUE: 过滤器序列号
% #HANDWRITTEN: A1-2\quad A4-8\quad A1-6 / A4-? \quad A8-? \quad A2-7 \quad A3-2
\fieldvalue{\handwritten{A1-2\quad A4-8\quad A1-6 / A4-?\quad A8-?\quad A2-7\quad A3-2}} \\
\hline
生产厂家 &
% #VALUE_ID: LEX-P0001-V0002
% #FIELD_VALUE: 生产厂家--单胞生物
\fieldvalue{\checkboxfield{unchecked}} 单胞生物\quad
% #VALUE_ID: LEX-P0001-V0003
% #FIELD_VALUE: 生产厂家--Cobetter
\fieldvalue{\checkboxfield{unchecked}} Cobetter\quad
% #VALUE_ID: LEX-P0001-V0004
% #FIELD_VALUE: 生产厂家--其他
\fieldvalue{\checkboxfield{checked}} 其他 &
% #VALUE_ID: LEX-P0001-V0005
% #FIELD_VALUE: 其他生产厂家
% #HANDWRITTEN: 瑞普生物
\fieldvalue{\handwritten{瑞普生物}} & \\
\hline
是否需要组装 &
% #VALUE_ID: LEX-P0001-V0006
% #FIELD_VALUE: 是否需要组装--是
\fieldvalue{\checkboxfield{unchecked}} 是\quad
% #VALUE_ID: LEX-P0001-V0007
% #FIELD_VALUE: 是否需要组装--否
\fieldvalue{\checkboxfield{unchecked}} 否
& 滤前完整性测试结果 & 泄点值：\par
% #VALUE_ID: LEX-P0001-V0008
% #FIELD_VALUE: 泄点值
% #HANDWRITTEN: 0.5
\fieldvalue{\handwritten{0.5}} mbar \\
\hline
使用前灭菌批号 & & 灭菌前完整性 & \\
\hline
操作人/日期 & & 复核人/日期 & \\
\hline
测试设备 &
% #VALUE_ID: LEX-P0001-V0009
% #FIELD_VALUE: 测试设备--完整性测试仪
\fieldvalue{\checkboxfield{checked}} 完整性测试仪\quad
% #VALUE_ID: LEX-P0001-V0010
% #FIELD_VALUE: 测试设备--其他
\fieldvalue{\checkboxfield{unchecked}} 其他
& 设备编号 &
% #VALUE_ID: LEX-P0001-V0011
% #FIELD_VALUE: 设备编号
% #HANDWRITTEN: 112230
\fieldvalue{\handwritten{112230}} \\
\hline
\multirow{5}{*}{过滤产品名称} &
% #VALUE_ID: LEX-P0001-V0012
% #FIELD_VALUE: 过滤产品名称--1\%人自袋冲液
\fieldvalue{\checkboxfield{unchecked}}1\%人自袋冲液 &
% #VALUE_ID: LEX-P0001-V0013
% #FIELD_VALUE: 过滤产品名称--2\%人自袋冲液
\fieldvalue{\checkboxfield{unchecked}}2\%人自袋冲液 &
% #VALUE_ID: LEX-P0001-V0014
% #FIELD_VALUE: 过滤产品名称--干细胞冻存液
\fieldvalue{\checkboxfield{unchecked}}干细胞冻存液 \\
\cline{2-4}
&
% #VALUE_ID: LEX-P0001-V0015
% #FIELD_VALUE: 过滤产品名称--裂解液
\fieldvalue{\checkboxfield{unchecked}}裂解液 &
% #VALUE_ID: LEX-P0001-V0016
% #FIELD_VALUE: 过滤产品名称--完全培养基
\fieldvalue{\checkboxfield{unchecked}}完全培养基 &
% #VALUE_ID: LEX-P0001-V0017
% #FIELD_VALUE: 过滤产品名称--成品冻存液
\fieldvalue{\checkboxfield{unchecked}}成品冻存液 \\
\cline{2-4}
&
% #VALUE_ID: LEX-P0001-V0018
% #FIELD_VALUE: 过滤产品名称--10\% DMSO
\fieldvalue{\checkboxfield{unchecked}}10\% DMSO &
% #VALUE_ID: LEX-P0001-V0019
% #FIELD_VALUE: 过滤产品名称--0.9\%氯化钠溶液
\fieldvalue{\checkboxfield{unchecked}}0.9\%氯化钠溶液 &
% #VALUE_ID: LEX-P0001-V0020
% #FIELD_VALUE: 过滤产品名称--其他
\fieldvalue{\checkboxfield{checked}}其他 \\
\cline{2-4}
&
% #VALUE_ID: LEX-P0001-V0021
% #FIELD_VALUE: 其他过滤产品名称
% #HANDWRITTEN: 2\%血浆
\fieldvalue{\handwritten{2\%血浆}} & & \\
\hline
过滤产品批号 &
\multicolumn{3}{l|}{
% #VALUE_ID: LEX-P0001-V0022
% #FIELD_VALUE: 过滤产品批号
% #HANDWRITTEN: H? A?3222025/12/?
\fieldvalue{\handwritten{H? A?3222025/12/?}}} \\
\hline
\multirow{4}{*}{测试程序} &
% #VALUE_ID: LEX-P0001-V0023
% #FIELD_VALUE: 测试程序--C02BBSLENAI1P
\fieldvalue{\checkboxfield{unchecked}}C02BBSLENAI1P &
% #VALUE_ID: LEX-P0001-V0024
% #FIELD_VALUE: 测试程序--L05SSLENAI1P
\fieldvalue{\checkboxfield{unchecked}}L05SSLENAI1P &
% #VALUE_ID: LEX-P0001-V0025
% #FIELD_VALUE: 测试程序--C02BBSHMN1P
\fieldvalue{\checkboxfield{unchecked}}C02BBSHMN1P \\
\cline{2-4}
&
% #VALUE_ID: LEX-P0001-V0026
% #FIELD_VALUE: 测试程序--L10SSLENAI1P
\fieldvalue{\checkboxfield{unchecked}}L10SSLENAI1P &
% #VALUE_ID: LEX-P0001-V0027
% #FIELD_VALUE: 测试程序--L20SSLENA1P
\fieldvalue{\checkboxfield{unchecked}}L20SSLENA1P &
% #VALUE_ID: LEX-P0001-V0028
% #FIELD_VALUE: 测试程序--L05BBSHMNA1P
\fieldvalue{\checkboxfield{unchecked}}L05BBSHMNA1P \\
\cline{2-4}
&
% #VALUE_ID: LEX-P0001-V0029
% #FIELD_VALUE: 测试程序--其他
\fieldvalue{\checkboxfield{checked}}其他：
% #VALUE_ID: LEX-P0001-V0030
% #FIELD_VALUE: 其他测试程序
% #HANDWRITTEN: SFMPEF-2224
\fieldvalue{\handwritten{SFMPEF-2224}} & & \\
\hline
滤器状态 &
\multicolumn{3}{l|}{
% #VALUE_ID: LEX-P0001-V0031
% #FIELD_VALUE: 滤器状态--第一组--滤前
\fieldvalue{\checkboxfield{unchecked}}滤前\quad
% #VALUE_ID: LEX-P0001-V0032
% #FIELD_VALUE: 滤器状态--第一组--滤后
\fieldvalue{\checkboxfield{checked}}滤后\qquad
% #VALUE_ID: LEX-P0001-V0033
% #FIELD_VALUE: 滤器状态--第二组--滤前
\fieldvalue{\checkboxfield{unchecked}}滤前\quad
% #VALUE_ID: LEX-P0001-V0034
% #FIELD_VALUE: 滤器状态--第二组--滤后
\fieldvalue{\checkboxfield{checked}}滤后\qquad
% #VALUE_ID: LEX-P0001-V0035
% #FIELD_VALUE: 滤器状态--第三组--滤前
\fieldvalue{\checkboxfield{unchecked}}滤前\quad
% #VALUE_ID: LEX-P0001-V0036
% #FIELD_VALUE: 滤器状态--第三组--滤后
\fieldvalue{\checkboxfield{checked}}滤后} \\
\hline
清洗时长 &
% #VALUE_ID: LEX-P0001-V0037
% #FIELD_VALUE: 清洗时长--第一组
% #HANDWRITTEN: 3
\fieldvalue{\handwritten{3}} min &
% #VALUE_ID: LEX-P0001-V0038
% #FIELD_VALUE: 清洗时长--第二组
% #HANDWRITTEN: 3
\fieldvalue{\handwritten{3}} min &
% #VALUE_ID: LEX-P0001-V0039
% #FIELD_VALUE: 清洗时长--第三组
% #HANDWRITTEN: 3
\fieldvalue{\handwritten{3}} min \\
\hline
测试日期 &
% #VALUE_ID: LEX-P0001-V0040
% #FIELD_VALUE: 测试日期--第一组
% #HANDWRITTEN: 2025.12.12
\fieldvalue{\handwritten{2025.12.12}} &
% #VALUE_ID: LEX-P0001-V0041
% #FIELD_VALUE: 测试日期--第二组
% #HANDWRITTEN: 2025.12.12
\fieldvalue{\handwritten{2025.12.12}} &
% #VALUE_ID: LEX-P0001-V0042
% #FIELD_VALUE: 测试日期--第三组
% #HANDWRITTEN: 2025.12.12
\fieldvalue{\handwritten{2025.12.12}} \\
\hline
\multirow{3}{*}{测试结果} &
% #VALUE_ID: LEX-P0001-V0043
% #FIELD_VALUE: 测试结果--第一组--通过
\fieldvalue{\checkboxfield{checked}}通过，实际测点值：
% #VALUE_ID: LEX-P0001-V0046
% #FIELD_VALUE: 第一组实际测点值
% #HANDWRITTEN: 450mbar
\fieldvalue{\handwritten{450mbar}} &
% #VALUE_ID: LEX-P0001-V0044
% #FIELD_VALUE: 测试结果--第二组--通过
\fieldvalue{\checkboxfield{checked}}通过，实际测点值：
% #VALUE_ID: LEX-P0001-V0047
% #FIELD_VALUE: 第二组实际测点值
% #HANDWRITTEN: 400mbar
\fieldvalue{\handwritten{400mbar}} &
% #VALUE_ID: LEX-P0001-V0045
% #FIELD_VALUE: 测试结果--第三组--通过
\fieldvalue{\checkboxfield{checked}}通过，实际测点值：
% #VALUE_ID: LEX-P0001-V0048
% #FIELD_VALUE: 第三组实际测点值
% #HANDWRITTEN: 450mbar
\fieldvalue{\handwritten{450mbar}} \\
\cline{2-4}
&
% #VALUE_ID: LEX-P0001-V0049
% #FIELD_VALUE: 测试结果--第一组--失败
\fieldvalue{\checkboxfield{unchecked}}失败，失败原因：
% #VALUE_ID: LEX-P0001-V0052
% #FIELD_VALUE: 第一组失败原因
% #HANDWRITTEN: /
\fieldvalue{\handwritten{/}} &
% #VALUE_ID: LEX-P0001-V0050
% #FIELD_VALUE: 测试结果--第二组--失败
\fieldvalue{\checkboxfield{unchecked}}失败，失败原因：
% #VALUE_ID: LEX-P0001-V0053
% #FIELD_VALUE: 第二组失败原因
% #HANDWRITTEN: /
\fieldvalue{\handwritten{/}} &
% #VALUE_ID: LEX-P0001-V0051
% #FIELD_VALUE: 测试结果--第三组--失败
\fieldvalue{\checkboxfield{unchecked}}失败，失败原因：
% #VALUE_ID: LEX-P0001-V0054
% #FIELD_VALUE: 第三组失败原因
% #HANDWRITTEN: /
\fieldvalue{\handwritten{/}} \\
\hline
测试人/日期 &
% #VALUE_ID: LEX-P0001-V0055
% #FIELD_VALUE: 第一组测试人
% #HANDWRITTEN: 李霞
\fieldvalue{\handwritten{李霞}}\quad
% #VALUE_ID: LEX-P0001-V0056
% #FIELD_VALUE: 第一组测试日期
% #HANDWRITTEN: 2025.12.12
\fieldvalue{\handwritten{2025.12.12}} &
% #VALUE_ID: LEX-P0001-V0057
% #FIELD_VALUE: 第二组测试人
% #HANDWRITTEN: 李霞
\fieldvalue{\handwritten{李霞}}\quad
% #VALUE_ID: LEX-P0001-V0058
% #FIELD_VALUE: 第二组测试日期
% #HANDWRITTEN: 2025.12.12
\fieldvalue{\handwritten{2025.12.12}} &
% #VALUE_ID: LEX-P0001-V0059
% #FIELD_VALUE: 第三组测试人
% #HANDWRITTEN: 李霞
\fieldvalue{\handwritten{李霞}}\quad
% #VALUE_ID: LEX-P0001-V0060
% #FIELD_VALUE: 第三组测试日期
% #HANDWRITTEN: 2025.12.12
\fieldvalue{\handwritten{2025.12.12}} \\
\hline
复核人/日期 &
% #VALUE_ID: LEX-P0001-V0061
% #FIELD_VALUE: 第一组复核人
% #HANDWRITTEN: 白俊玲
\fieldvalue{\handwritten{白俊玲}}\quad
% #VALUE_ID: LEX-P0001-V0062
% #FIELD_VALUE: 第一组复核日期
% #HANDWRITTEN: 2025.12.12
\fieldvalue{\handwritten{2025.12.12}} &
% #VALUE_ID: LEX-P0001-V0063
% #FIELD_VALUE: 第二组复核人
% #HANDWRITTEN: 白俊玲
\fieldvalue{\handwritten{白俊玲}}\quad
% #VALUE_ID: LEX-P0001-V0064
% #FIELD_VALUE: 第二组复核日期
% #HANDWRITTEN: 2025.12.12
\fieldvalue{\handwritten{2025.12.12}} &
% #VALUE_ID: LEX-P0001-V0065
% #FIELD_VALUE: 第三组复核人
% #HANDWRITTEN: 白俊玲
\fieldvalue{\handwritten{白俊玲}}\quad
% #VALUE_ID: LEX-P0001-V0066
% #FIELD_VALUE: 第三组复核日期
% #HANDWRITTEN: 2025.12.12
\fieldvalue{\handwritten{2025.12.12}} \\
\hline
备注： &
\multicolumn{3}{l|}{
% #VALUE_ID: LEX-P0001-V0067
% #FIELD_VALUE: 备注
% #HANDWRITTEN: /
\fieldvalue{\handwritten{/}}} \\
\hline
\end{tabular}

\normalsize
\begin{center}
第 1 页 共 1 页
\end{center}

% TODO: Several handwritten names, serial-number characters, and product-batch characters are partly illegible in the source image and should be visually reviewed.
% LEXOID_PAGE_COMPLETED: 1/70

\newpage

\begin{center}
\textbf{永水性滤器使用记录}
\end{center}

\noindent
\begin{tabular}{|p{0.18\linewidth}|p{0.30\linewidth}|p{0.20\linewidth}|p{0.25\linewidth}|}
\hline
过滤器进厂批号 & & 过滤器序列号 &
% #VALUE_ID: LEX-P0002-V0001
% #FIELD_VALUE: 过滤器序列号
% #TODO #HANDWRITTEN: B1-4; handwritten serial number
\fieldvalue{\handwritten{B1-4}}\newline
% #VALUE_ID: LEX-P0002-V0002
% #FIELD_VALUE: 过滤器序列号
% #TODO #HANDWRITTEN: B3-2; handwritten serial number
\fieldvalue{\handwritten{B3-2}}\newline
% #VALUE_ID: LEX-P0002-V0003
% #FIELD_VALUE: 过滤器序列号
% #TODO #HANDWRITTEN: A?-2; final character unclear
\fieldvalue{\handwritten{A?-2}}\\
\cline{4-4}
&&&
% #VALUE_ID: LEX-P0002-V0004
% #FIELD_VALUE: 过滤器序列号
% #TODO #HANDWRITTEN: B1-4; handwritten serial number
\fieldvalue{\handwritten{B1-4}}\newline
% #VALUE_ID: LEX-P0002-V0005
% #FIELD_VALUE: 过滤器序列号
% #TODO #HANDWRITTEN: B3-1; handwritten serial number
\fieldvalue{\handwritten{B3-1}}\newline
% #VALUE_ID: LEX-P0002-V0006
% #FIELD_VALUE: 过滤器序列号
% #HANDWRITTEN: A4-7
\fieldvalue{\handwritten{A4-7}}\\
\hline
生产厂家 &
$\square$ 华生物\quad $\square$ Cobetter\quad
% #VALUE_ID: LEX-P0002-V0007
% #FIELD_VALUE: 生产厂家--其他
% #HANDWRITTEN: check mark
\fieldvalue{\handwritten{\checkboxfield{checked}}}\ 其他 &
& 
% #VALUE_ID: LEX-P0002-V0008
% #FIELD_VALUE: 生产厂家--其他
% #TODO #HANDWRITTEN: 赛特; handwriting partially unclear
\fieldvalue{\handwritten{赛特}}\\
\hline
是否需要组装 &
% #VALUE_ID: LEX-P0002-V0009
% #FIELD_VALUE: 是否需要组装--是
\fieldvalue{\checkboxfield{unchecked}}\ 是\quad
% #VALUE_ID: LEX-P0002-V0010
% #FIELD_VALUE: 是否需要组装--否
\fieldvalue{\checkboxfield{unchecked}}\ 否
& 滤前完整性测试结果 & 泡点：\underline{\hspace{1.5cm}} mbar\\
\hline
使用前灭菌批号 & & 灭菌有效期至 &
% #VALUE_ID: LEX-P0002-V0011
% #FIELD_VALUE: 灭菌有效期至
% #HANDWRITTEN: 2025-12-12
\fieldvalue{\handwritten{2025-12-12}}\\
\hline
操作人/日期 & & 复核人/日期 &\\
\hline
测试设备 &
% #VALUE_ID: LEX-P0002-V0012
% #FIELD_VALUE: 测试设备--完整性测试仪
% #HANDWRITTEN: check mark
\fieldvalue{\handwritten{\checkboxfield{checked}}} 完整性测试仪\quad
% #VALUE_ID: LEX-P0002-V0013
% #FIELD_VALUE: 测试设备--其他
\fieldvalue{\checkboxfield{unchecked}} 其他
& 设备编号 &
% #VALUE_ID: LEX-P0002-V0014
% #FIELD_VALUE: 设备编号
% #HANDWRITTEN: 1122307
\fieldvalue{\handwritten{1122307}}\\
\hline
过滤产品名称 &
% #VALUE_ID: LEX-P0002-V0015
% #FIELD_VALUE: 过滤产品名称--1\%人白缓冲液
\fieldvalue{\checkboxfield{unchecked}}1\%人白缓冲液\quad
% #VALUE_ID: LEX-P0002-V0016
% #FIELD_VALUE: 过滤产品名称--2\%人白缓冲液
\fieldvalue{\checkboxfield{unchecked}}2\%人白缓冲液
&
% #VALUE_ID: LEX-P0002-V0017
% #FIELD_VALUE: 过滤产品名称--干细胞冻存液
\fieldvalue{\checkboxfield{unchecked}}干细胞冻存液
&\\
&
% #VALUE_ID: LEX-P0002-V0018
% #FIELD_VALUE: 过滤产品名称--裂解液
\fieldvalue{\checkboxfield{unchecked}}裂解液\quad
% #VALUE_ID: LEX-P0002-V0019
% #FIELD_VALUE: 过滤产品名称--完全培养基
\fieldvalue{\checkboxfield{unchecked}}完全培养基\quad
% #VALUE_ID: LEX-P0002-V0020
% #FIELD_VALUE: 过滤产品名称--成品冻存液
\fieldvalue{\checkboxfield{unchecked}}成品冻存液
&&\\
&
% #VALUE_ID: LEX-P0002-V0021
% #FIELD_VALUE: 过滤产品名称--10\%DMSO
\fieldvalue{\checkboxfield{unchecked}}10\%DMSO\quad
% #VALUE_ID: LEX-P0002-V0022
% #FIELD_VALUE: 过滤产品名称--0.9\%氯化钠溶液
\fieldvalue{\checkboxfield{unchecked}}0.9\%氯化钠溶液
&
% #VALUE_ID: LEX-P0002-V0023
% #FIELD_VALUE: 过滤产品名称--其他
% #HANDWRITTEN: check mark
\fieldvalue{\handwritten{\checkboxfield{checked}}}其他\quad
% #VALUE_ID: LEX-P0002-V0024
% #FIELD_VALUE: 过滤产品名称--其他
% #TODO #HANDWRITTEN: 2?use?e; handwriting unclear
\fieldvalue{\handwritten{2?use?e}} &\\
\hline
过滤产品批号 &
% #VALUE_ID: LEX-P0002-V0025
% #FIELD_VALUE: 过滤产品批号
% #TODO #HANDWRITTEN: PF-H?332?0025-1209; handwriting partially unclear
\fieldvalue{\handwritten{PF-H?332?0025-1209}} &&\\
\hline
测试程序 &
% #VALUE_ID: LEX-P0002-V0026
% #FIELD_VALUE: 测试程序--C02BBSLENAP1P
\fieldvalue{\checkboxfield{unchecked}}C02BBSLENAP1P\quad
% #VALUE_ID: LEX-P0002-V0027
% #FIELD_VALUE: 测试程序--C05SLENAP1P
\fieldvalue{\checkboxfield{unchecked}}C05SLENAP1P\quad
% #VALUE_ID: LEX-P0002-V0028
% #FIELD_VALUE: 测试程序--C02BBSHMN1P
\fieldvalue{\checkboxfield{unchecked}}C02BBSHMN1P
&&\\
&
% #VALUE_ID: LEX-P0002-V0029
% #FIELD_VALUE: 测试程序--L10SSLENAP1P
\fieldvalue{\checkboxfield{unchecked}}L10SSLENAP1P\quad
% #VALUE_ID: LEX-P0002-V0030
% #FIELD_VALUE: 测试程序--L20SSSLENAP1P
\fieldvalue{\checkboxfield{unchecked}}L20SSSLENAP1P\quad
% #VALUE_ID: LEX-P0002-V0031
% #FIELD_VALUE: 测试程序--L05BSHMNA1P
\fieldvalue{\checkboxfield{unchecked}}L05BSHMNA1P
&&\\
&
% #VALUE_ID: LEX-P0002-V0032
% #FIELD_VALUE: 测试程序--其他
% #HANDWRITTEN: check mark
\fieldvalue{\handwritten{\checkboxfield{checked}}}其他\quad
% #VALUE_ID: LEX-P0002-V0033
% #FIELD_VALUE: 测试程序--其他
% #HANDWRITTEN: SEMP?ES-222E
\fieldvalue{\handwritten{SEMP?ES-222E}} &&\\
\hline
过滤器状态 &
% #VALUE_ID: LEX-P0002-V0034
% #FIELD_VALUE: 第一组过滤器状态--滤前
\fieldvalue{\checkboxfield{unchecked}}滤前\quad
% #VALUE_ID: LEX-P0002-V0035
% #FIELD_VALUE: 第一组过滤器状态--滤后
% #HANDWRITTEN: check mark
\fieldvalue{\handwritten{\checkboxfield{checked}}}滤后
&
% #VALUE_ID: LEX-P0002-V0036
% #FIELD_VALUE: 第二组过滤器状态--滤前
\fieldvalue{\checkboxfield{unchecked}}滤前\quad
% #VALUE_ID: LEX-P0002-V0037
% #FIELD_VALUE: 第二组过滤器状态--滤后
% #HANDWRITTEN: check mark
\fieldvalue{\handwritten{\checkboxfield{checked}}}滤后
&
% #VALUE_ID: LEX-P0002-V0038
% #FIELD_VALUE: 第三组过滤器状态--滤前
\fieldvalue{\checkboxfield{unchecked}}滤前\quad
% #VALUE_ID: LEX-P0002-V0039
% #FIELD_VALUE: 第三组过滤器状态--滤后
% #HANDWRITTEN: check mark
\fieldvalue{\handwritten{\checkboxfield{checked}}}滤后\\
\hline
消洗时长 &
% #VALUE_ID: LEX-P0002-V0040
% #FIELD_VALUE: 第一组消洗时长
% #HANDWRITTEN: 3
\fieldvalue{\handwritten{3}}\ min
&
% #VALUE_ID: LEX-P0002-V0041
% #FIELD_VALUE: 第二组消洗时长
% #HANDWRITTEN: 3
\fieldvalue{\handwritten{3}}\ min
&
% #VALUE_ID: LEX-P0002-V0042
% #FIELD_VALUE: 第三组消洗时长
% #HANDWRITTEN: 3
\fieldvalue{\handwritten{3}}\ min\\
\hline
测试日期 &
% #VALUE_ID: LEX-P0002-V0043
% #FIELD_VALUE: 第一组测试日期
% #HANDWRITTEN: 2025.12.12
\fieldvalue{\handwritten{2025.12.12}}
&
% #VALUE_ID: LEX-P0002-V0044
% #FIELD_VALUE: 第二组测试日期
% #HANDWRITTEN: 2025.12.12
\fieldvalue{\handwritten{2025.12.12}}
&
% #VALUE_ID: LEX-P0002-V0045
% #FIELD_VALUE: 第三组测试日期
% #HANDWRITTEN: 2025.12.12
\fieldvalue{\handwritten{2025.12.12}}\\
\hline
测试结果 &
% #VALUE_ID: LEX-P0002-V0046
% #FIELD_VALUE: 第一组测试结果--通过
% #HANDWRITTEN: check mark
\fieldvalue{\handwritten{\checkboxfield{checked}}}通过，实际泡点值：
% #VALUE_ID: LEX-P0002-V0047
% #FIELD_VALUE: 第一组实际泡点值
% #HANDWRITTEN: 6000 mbar
\fieldvalue{\handwritten{6000\ mbar}}&&\\
&
% #VALUE_ID: LEX-P0002-V0048
% #FIELD_VALUE: 第一组测试结果--失败
\fieldvalue{\checkboxfield{unchecked}}失败，失败原因：\underline{\hspace{1.5cm}}
&
% #VALUE_ID: LEX-P0002-V0049
% #FIELD_VALUE: 第二组测试结果--通过
% #HANDWRITTEN: check mark
\fieldvalue{\handwritten{\checkboxfield{checked}}}通过，实际泡点值：
% #VALUE_ID: LEX-P0002-V0050
% #FIELD_VALUE: 第二组实际泡点值
% #HANDWRITTEN: 5000 mbar
\fieldvalue{\handwritten{5000\ mbar}}&\\
&
% #VALUE_ID: LEX-P0002-V0051
% #FIELD_VALUE: 第二组测试结果--失败
\fieldvalue{\checkboxfield{unchecked}}失败，失败原因：\underline{\hspace{1.5cm}}
&
% #VALUE_ID: LEX-P0002-V0052
% #FIELD_VALUE: 第三组测试结果--通过
% #HANDWRITTEN: check mark
\fieldvalue{\handwritten{\checkboxfield{checked}}}通过，实际泡点值：
% #VALUE_ID: LEX-P0002-V0053
% #FIELD_VALUE: 第三组实际泡点值
% #HANDWRITTEN: 5000 mbar
\fieldvalue{\handwritten{5000\ mbar}}&\\
&
% #VALUE_ID: LEX-P0002-V0054
% #FIELD_VALUE: 第三组测试结果--失败
\fieldvalue{\checkboxfield{unchecked}}失败，失败原因：\underline{\hspace{1.5cm}} &&\\
\hline
测试人/日期 &
% #VALUE_ID: LEX-P0002-V0055
% #FIELD_VALUE: 第一组测试人/日期
% #TODO #HANDWRITTEN: 李艳2025.12.12; name partially unclear
\fieldvalue{\handwritten{李艳\ 2025.12.12}}
&
% #VALUE_ID: LEX-P0002-V0056
% #FIELD_VALUE: 第二组测试人/日期
% #TODO #HANDWRITTEN: 李艳2025.12.12; name partially unclear
\fieldvalue{\handwritten{李艳\ 2025.12.12}}
&
% #VALUE_ID: LEX-P0002-V0057
% #FIELD_VALUE: 第三组测试人/日期
% #TODO #HANDWRITTEN: 李艳2025.12.12; name partially unclear
\fieldvalue{\handwritten{李艳\ 2025.12.12}}\\
\hline
复核人/日期 &
% #VALUE_ID: LEX-P0002-V0058
% #FIELD_VALUE: 第一组复核人/日期
% #TODO #HANDWRITTEN: handwriting illegible, 2025.12.12
\fieldvalue{\handwritten{[姓名不清]\ 2025.12.12}}
&
% #VALUE_ID: LEX-P0002-V0059
% #FIELD_VALUE: 第二组复核人/日期
% #TODO #HANDWRITTEN: handwriting illegible, 2025.12.12
\fieldvalue{\handwritten{[姓名不清]\ 2025.12.12}}
&
% #VALUE_ID: LEX-P0002-V0060
% #FIELD_VALUE: 第三组复核人/日期
% #TODO #HANDWRITTEN: handwriting illegible, 2025.12.12
\fieldvalue{\handwritten{[姓名不清]\ 2025.12.12}}\\
\hline
备注：&
% #VALUE_ID: LEX-P0002-V0061
% #FIELD_VALUE: 备注
% #HANDWRITTEN: /
\fieldvalue{\handwritten{/}} &&\\
\hline
\end{tabular}

\noindent
% TODO: The form header contains printed company/product identifiers and a red stamp-like box; no editable value was clearly distinguishable beyond the handwritten entries annotated above.
% LEXOID_PAGE_COMPLETED: 2/70

\newpage

\begin{center}
\textbf{永永性滤器使用记录}
\end{center}

\noindent
\begin{tabular}{|p{0.16\linewidth}|p{0.25\linewidth}|p{0.16\linewidth}|p{0.31\linewidth}|}
\hline
过滤器进厂批号 & & 过滤器序列号 &
% #VALUE_ID: LEX-P0003-V0001
% #FIELD_VALUE: 过滤器序列号
\fieldvalue{\handwritten{A3-5\quad A2-6\quad B2-4\\A4-4\quad A4-6\quad B?\\A?\quad H1-7\quad AH-5}}\\
\hline
生产厂家 &
\begin{tabular}{@{}l@{}}
$\square$ 华生物\\
$\square$ Cobetter\\
其他\quad
% #VALUE_ID: LEX-P0003-V0002
% #FIELD_VALUE: 生产厂家--其他
\fieldvalue{\checkboxfield{checked}}
\end{tabular}
&
\multicolumn{2}{l|}{}\\
\hline
是否需要组装 &
\begin{tabular}{@{}l@{}}
是\quad
% #VALUE_ID: LEX-P0003-V0003
% #FIELD_VALUE: 是否需要组装--是
\fieldvalue{\checkboxfield{unchecked}}
\qquad 否\quad
% #VALUE_ID: LEX-P0003-V0004
% #FIELD_VALUE: 是否需要组装--否
\fieldvalue{\checkboxfield{unchecked}}
\end{tabular}
& 滤前完整性测试结果 &
泄点值：
% #VALUE_ID: LEX-P0003-V0005
% #FIELD_VALUE: 泄点值
% #TODO #HANDWRITTEN: 12-12; handwriting is partially unclear
\fieldvalue{\handwritten{12-12}}\ mbar\\
\hline
使用前灭菌批号 & & 灭菌有效期至 &
% #VALUE_ID: LEX-P0003-V0006
% #FIELD_VALUE: 灭菌有效期至
% #TODO #HANDWRITTEN: 2025.12.25; handwriting is partially unclear
\fieldvalue{\handwritten{2025.12.25}}\\
\hline
操作人/日期 & & 复核人/日期 & \\
\hline
测试设备 &
\begin{tabular}{@{}l@{}}
$\square$ 完整性测试仪\quad
% #VALUE_ID: LEX-P0003-V0007
% #FIELD_VALUE: 测试设备--完整性测试仪
\fieldvalue{\checkboxfield{checked}}\\
$\square$ 其他\quad
% #VALUE_ID: LEX-P0003-V0008
% #FIELD_VALUE: 测试设备--其他
\fieldvalue{\checkboxfield{unchecked}}
\end{tabular}
& 设备编号 &
% #VALUE_ID: LEX-P0003-V0009
% #FIELD_VALUE: 设备编号
% #TODO #HANDWRITTEN: 1122307; final digit is unclear
\fieldvalue{\handwritten{1122307}}\\
\hline
过滤产品名称 &
\begin{tabular}{@{}l@{}}
$\square$1\%人白蛋白稀释液\quad
% #VALUE_ID: LEX-P0003-V0010
% #FIELD_VALUE: 过滤产品名称--1\%人白蛋白稀释液
\fieldvalue{\checkboxfield{unchecked}}\\
$\square$裂解液\quad
% #VALUE_ID: LEX-P0003-V0011
% #FIELD_VALUE: 过滤产品名称--裂解液
\fieldvalue{\checkboxfield{unchecked}}\\
$\square$1\%DMSO\quad
% #VALUE_ID: LEX-P0003-V0012
% #FIELD_VALUE: 过滤产品名称--1\%DMSO
\fieldvalue{\checkboxfield{unchecked}}\\
$\square$其他\quad
% #VALUE_ID: LEX-P0003-V0013
% #FIELD_VALUE: 过滤产品名称--其他
\fieldvalue{\checkboxfield{checked}}\\
% #VALUE_ID: LEX-P0003-V0014
% #FIELD_VALUE: 过滤产品名称--其他填写内容
% #TODO #HANDWRITTEN: 2?sh?eye; handwriting is illegible
\fieldvalue{\handwritten{2?sh?eye}}
\end{tabular}
&
\begin{tabular}{@{}l@{}}
$\square$2\%人白蛋白稀释液\quad
% #VALUE_ID: LEX-P0003-V0015
% #FIELD_VALUE: 过滤产品名称--2\%人白蛋白稀释液
\fieldvalue{\checkboxfield{unchecked}}\\
$\square$完全培养基\quad
% #VALUE_ID: LEX-P0003-V0016
% #FIELD_VALUE: 过滤产品名称--完全培养基
\fieldvalue{\checkboxfield{unchecked}}\\
$\square$0.9\%氯化钠溶液\quad
% #VALUE_ID: LEX-P0003-V0017
% #FIELD_VALUE: 过滤产品名称--0.9\%氯化钠溶液
\fieldvalue{\checkboxfield{unchecked}}
\end{tabular}&\\
\hline
过滤产品批号 &
\multicolumn{3}{l|}{
% #VALUE_ID: LEX-P0003-V0018
% #FIELD_VALUE: 过滤产品批号
% #HANDWRITTEN: A4S3222025/12/19
\fieldvalue{\handwritten{A4S3222025/12/19}}}\\
\hline
测试程序 &
\multicolumn{3}{l|}{
\begin{tabular}{@{}lll@{}}
$\square$C02BBSLENA1P\quad
% #VALUE_ID: LEX-P0003-V0019
% #FIELD_VALUE: 测试程序--C02BBSLENA1P
\fieldvalue{\checkboxfield{unchecked}}
&
$\square$L05SSLENA1P\quad
% #VALUE_ID: LEX-P0003-V0020
% #FIELD_VALUE: 测试程序--L05SSLENA1P
\fieldvalue{\checkboxfield{unchecked}}
&
$\square$C02BBSHMN1P\quad
% #VALUE_ID: LEX-P0003-V0021
% #FIELD_VALUE: 测试程序--C02BBSHMN1P
\fieldvalue{\checkboxfield{unchecked}}
\\
$\square$L01SSLENA1P\quad
% #VALUE_ID: LEX-P0003-V0022
% #FIELD_VALUE: 测试程序--L01SSLENA1P
\fieldvalue{\checkboxfield{unchecked}}
&
$\square$L20SSLENA1P\quad
% #VALUE_ID: LEX-P0003-V0023
% #FIELD_VALUE: 测试程序--L20SSLENA1P
\fieldvalue{\checkboxfield{unchecked}}
&
$\square$L05BBSHMNA1P\quad
% #VALUE_ID: LEX-P0003-V0024
% #FIELD_VALUE: 测试程序--L05BBSHMNA1P
\fieldvalue{\checkboxfield{unchecked}}
\\
$\square$其他\quad
% #VALUE_ID: LEX-P0003-V0025
% #FIELD_VALUE: 测试程序--其他
\fieldvalue{\checkboxfield{checked}}
&
% #VALUE_ID: LEX-P0003-V0026
% #FIELD_VALUE: 测试程序--其他填写内容
% #HANDWRITTEN: SFM?ES-2225
\fieldvalue{\handwritten{SFM?ES-2225}}
&
\end{tabular}}\\
\hline
滤器状态 &
\begin{tabular}{@{}l@{}}
$\square$滤前\quad
% #VALUE_ID: LEX-P0003-V0027
% #FIELD_VALUE: 滤器状态--第一列滤前
\fieldvalue{\checkboxfield{unchecked}}
\quad 滤后\quad
% #VALUE_ID: LEX-P0003-V0028
% #FIELD_VALUE: 滤器状态--第一列滤后
\fieldvalue{\checkboxfield{checked}}
\end{tabular}
&
\begin{tabular}{@{}l@{}}
$\square$滤前\quad
% #VALUE_ID: LEX-P0003-V0029
% #FIELD_VALUE: 滤器状态--第二列滤前
\fieldvalue{\checkboxfield{unchecked}}
\quad 滤后\quad
% #VALUE_ID: LEX-P0003-V0030
% #FIELD_VALUE: 滤器状态--第二列滤后
\fieldvalue{\checkboxfield{checked}}
\end{tabular}
&
\begin{tabular}{@{}l@{}}
$\square$滤前\quad
% #VALUE_ID: LEX-P0003-V0031
% #FIELD_VALUE: 滤器状态--第三列滤前
\fieldvalue{\checkboxfield{unchecked}}
\quad 滤后\quad
% #VALUE_ID: LEX-P0003-V0032
% #FIELD_VALUE: 滤器状态--第三列滤后
\fieldvalue{\checkboxfield{checked}}
\end{tabular}\\
\hline
消滤时长 &
% #VALUE_ID: LEX-P0003-V0033
% #FIELD_VALUE: 消滤时长--第一列
% #HANDWRITTEN: 3
\fieldvalue{\handwritten{3}}\ min &
% #VALUE_ID: LEX-P0003-V0034
% #FIELD_VALUE: 消滤时长--第二列
% #HANDWRITTEN: 3
\fieldvalue{\handwritten{3}}\ min &
% #VALUE_ID: LEX-P0003-V0035
% #FIELD_VALUE: 消滤时长--第三列
% #HANDWRITTEN: 3
\fieldvalue{\handwritten{3}}\ min\\
\hline
测试日期 &
% #VALUE_ID: LEX-P0003-V0036
% #FIELD_VALUE: 测试日期--第一列
% #HANDWRITTEN: 2025.12.12
\fieldvalue{\handwritten{2025.12.12}} &
% #VALUE_ID: LEX-P0003-V0037
% #FIELD_VALUE: 测试日期--第二列
% #HANDWRITTEN: 2025.12.12
\fieldvalue{\handwritten{2025.12.12}} &
% #VALUE_ID: LEX-P0003-V0038
% #FIELD_VALUE: 测试日期--第三列
% #HANDWRITTEN: 2025.12.12
\fieldvalue{\handwritten{2025.12.12}}\\
\hline
测试结果 &
\begin{tabular}{@{}l@{}}
$\square$通过，实际测点值：
% #VALUE_ID: LEX-P0003-V0039
% #FIELD_VALUE: 测试结果--第一列通过
\fieldvalue{\checkboxfield{checked}}\\
% #VALUE_ID: LEX-P0003-V0040
% #FIELD_VALUE: 测试结果--第一列实际测点值
% #HANDWRITTEN: 6000mbar
\fieldvalue{\handwritten{6000mbar}}\\
$\square$失败，失败原因：
% #VALUE_ID: LEX-P0003-V0041
% #FIELD_VALUE: 测试结果--第一列失败
\fieldvalue{\checkboxfield{unchecked}}
\end{tabular}
&
\begin{tabular}{@{}l@{}}
$\square$通过，实际测点值：
% #VALUE_ID: LEX-P0003-V0042
% #FIELD_VALUE: 测试结果--第二列通过
\fieldvalue{\checkboxfield{checked}}\\
% #VALUE_ID: LEX-P0003-V0043
% #FIELD_VALUE: 测试结果--第二列实际测点值
% #HANDWRITTEN: 5000mbar
\fieldvalue{\handwritten{5000mbar}}\\
$\square$失败，失败原因：
% #VALUE_ID: LEX-P0003-V0044
% #FIELD_VALUE: 测试结果--第二列失败
\fieldvalue{\checkboxfield{unchecked}}
\end{tabular}
&
\begin{tabular}{@{}l@{}}
$\square$通过，实际测点值：
% #VALUE_ID: LEX-P0003-V0045
% #FIELD_VALUE: 测试结果--第三列通过
\fieldvalue{\checkboxfield{checked}}\\
% #VALUE_ID: LEX-P0003-V0046
% #FIELD_VALUE: 测试结果--第三列实际测点值
% #HANDWRITTEN: 4050mbar
\fieldvalue{\handwritten{4050mbar}}\\
$\square$失败，失败原因：
% #VALUE_ID: LEX-P0003-V0047
% #FIELD_VALUE: 测试结果--第三列失败
\fieldvalue{\checkboxfield{unchecked}}
\end{tabular}\\
\hline
测试人/日期 &
% #VALUE_ID: LEX-P0003-V0048
% #FIELD_VALUE: 测试人/日期--第一列
% #HANDWRITTEN: 李超 2025.12.12
\fieldvalue{\handwritten{李超\ 2025.12.12}} &
% #VALUE_ID: LEX-P0003-V0049
% #FIELD_VALUE: 测试人/日期--第二列
% #HANDWRITTEN: 李超 2025.12.12
\fieldvalue{\handwritten{李超\ 2025.12.12}} &
% #VALUE_ID: LEX-P0003-V0050
% #FIELD_VALUE: 测试人/日期--第三列
% #HANDWRITTEN: 李超 2025.12.12
\fieldvalue{\handwritten{李超\ 2025.12.12}}\\
\hline
复核人/日期 &
% #VALUE_ID: LEX-P0003-V0051
% #FIELD_VALUE: 复核人/日期--第一列
% #TODO #HANDWRITTEN: 赵?强 2025.12.12; signature is unclear
\fieldvalue{\handwritten{赵?强\ 2025.12.12}} &
% #VALUE_ID: LEX-P0003-V0052
% #FIELD_VALUE: 复核人/日期--第二列
% #TODO #HANDWRITTEN: 赵?强 2025.12.12; signature is unclear
\fieldvalue{\handwritten{赵?强\ 2025.12.12}} &
% #VALUE_ID: LEX-P0003-V0053
% #FIELD_VALUE: 复核人/日期--第三列
% #TODO #HANDWRITTEN: 赵?强 2025.12.12; signature is unclear
\fieldvalue{\handwritten{赵?强\ 2025.12.12}}\\
\hline
备注 &
\multicolumn{3}{l|}{
% #VALUE_ID: LEX-P0003-V0054
% #FIELD_VALUE: 备注
% #TODO #HANDWRITTEN: /; handwritten slash-like mark
\fieldvalue{\handwritten{/}}}\\
\hline
\end{tabular}

\begin{center}
第1页 共1页
\end{center}
% LEXOID_PAGE_COMPLETED: 3/70

\newpage

\begin{center}
\textbf{永永性滤器使用记录}
\end{center}

\noindent
文件编码：REC-YZ-SOP-PD-02-012-2\hfill 版本编号：1.5

\small
\begin{tabular}{|p{0.17\linewidth}|p{0.26\linewidth}|p{0.17\linewidth}|p{0.30\linewidth}|}
\hline
过滤器进厂批号 & & 过滤器序列号 &
% #VALUE_ID: LEX-P0004-V0001
% #TODO #HANDWRITTEN: A3-2; handwriting in serial-number area
\handwritten{A3-2}\quad
% #VALUE_ID: LEX-P0004-V0002
% #TODO #HANDWRITTEN: B3-2; handwriting in serial-number area
\handwritten{B3-2}\quad
% #VALUE_ID: LEX-P0004-V0003
% #TODO #HANDWRITTEN: A4-6; handwriting in serial-number area
\handwritten{A4-6}\quad
% #VALUE_ID: LEX-P0004-V0004
% #TODO #HANDWRITTEN: B4-2; handwriting in serial-number area
\handwritten{B4-2}\quad
% #VALUE_ID: LEX-P0004-V0005
% #TODO #HANDWRITTEN: A2-?; final character unclear
\handwritten{A2-?}\quad
% #VALUE_ID: LEX-P0004-V0006
% #HANDWRITTEN: B2-1
\handwritten{B2-1}\quad
% #VALUE_ID: LEX-P0004-V0007
% #HANDWRITTEN: A2-4
\handwritten{A2-4}\quad
% #VALUE_ID: LEX-P0004-V0008
% #HANDWRITTEN: B2-3
\handwritten{B2-3}\quad
% #VALUE_ID: LEX-P0004-V0009
% #HANDWRITTEN: A3-4
\handwritten{A3-4}
\\
\hline
生产厂家 &
\begin{tabular}{l}
$\square$ 牵牛生物\\
$\square$ Cobbeter\\
% #VALUE_ID: LEX-P0004-V0010
% #FIELD_VALUE: 生产厂家—其他
% #HANDWRITTEN: ✓
\fieldvalue{\handwritten{\checkboxfield{checked}}} 其他
\end{tabular}
&
是否需要组装 &
\begin{tabular}{l}
% #VALUE_ID: LEX-P0004-V0011
% #FIELD_VALUE: 是否需要组装—是
\fieldvalue{\checkboxfield{unchecked}} 是
\quad
% #VALUE_ID: LEX-P0004-V0012
% #FIELD_VALUE: 是否需要组装—否
\fieldvalue{\checkboxfield{unchecked}} 否
\end{tabular}
\\
\hline
使用前灭菌批号 & & 灭菌有效期 &
% #VALUE_ID: LEX-P0004-V0013
% #FIELD_VALUE: 灭菌有效期
% #HANDWRITTEN: 2025.12.12
\fieldvalue{\handwritten{2025.12.12}}
\\
\hline
操作人/日期 & & 复核人/日期 & \\
\hline
测试设备 &
\begin{tabular}{l}
% #VALUE_ID: LEX-P0004-V0014
% #FIELD_VALUE: 测试设备—完整性测试仪
% #HANDWRITTEN: ✓
\fieldvalue{\handwritten{\checkboxfield{checked}}} 完整性测试仪\\
% #VALUE_ID: LEX-P0004-V0015
% #FIELD_VALUE: 测试设备—其他
\fieldvalue{\checkboxfield{unchecked}} 其他
\end{tabular}
&
设备编号 &
% #VALUE_ID: LEX-P0004-V0016
% #FIELD_VALUE: 设备编号
% #HANDWRITTEN: 1122307
\fieldvalue{\handwritten{1122307}}
\\
\hline
过滤产品名称 &
\begin{tabular}{l}
% #VALUE_ID: LEX-P0004-V0017
% #FIELD_VALUE: 过滤产品名称—1\%人白蛋白液
\fieldvalue{\checkboxfield{unchecked}}1\%人白蛋白液\\
% #VALUE_ID: LEX-P0004-V0018
% #FIELD_VALUE: 过滤产品名称—裂解液
\fieldvalue{\checkboxfield{unchecked}}裂解液\\
% #VALUE_ID: LEX-P0004-V0019
% #FIELD_VALUE: 过滤产品名称—1\% DMSO
\fieldvalue{\checkboxfield{unchecked}}1\%DMSO\\
% #VALUE_ID: LEX-P0004-V0020
% #FIELD_VALUE: 过滤产品名称—其他
% #HANDWRITTEN: ✓
\fieldvalue{\handwritten{\checkboxfield{checked}}}其他
\end{tabular}
&
&
\begin{tabular}{l}
% #VALUE_ID: LEX-P0004-V0021
% #FIELD_VALUE: 过滤产品名称—2\%人白蛋白液
\fieldvalue{\checkboxfield{unchecked}}2\%人白蛋白液\\
% #VALUE_ID: LEX-P0004-V0022
% #FIELD_VALUE: 过滤产品名称—完全培养基
\fieldvalue{\checkboxfield{unchecked}}完全培养基\\
% #VALUE_ID: LEX-P0004-V0023
% #FIELD_VALUE: 过滤产品名称—0.9\%氯化钠溶液
\fieldvalue{\checkboxfield{unchecked}}0.9\%氯化钠溶液
\end{tabular}
\\
\hline
& \multicolumn{3}{l|}{
% #VALUE_ID: LEX-P0004-V0024
% #FIELD_VALUE: 过滤产品名称—其他具体名称
% #HANDWRITTEN: 2 Whe?e/e
\fieldvalue{\handwritten{2 Whe?e/e}}
}\\
\hline
过滤产品批号 &
\multicolumn{3}{l|}{
% #VALUE_ID: LEX-P0004-V0025
% #FIELD_VALUE: 过滤产品批号
% #HANDWRITTEN: A432220251204
\fieldvalue{\handwritten{A432220251204}}
}\\
\hline
测试程序 &
\begin{tabular}{l}
% #VALUE_ID: LEX-P0004-V0026
% #FIELD_VALUE: 测试程序—C02BSSLENA1P
\fieldvalue{\checkboxfield{unchecked}}C02BSSLENA1P\\
% #VALUE_ID: LEX-P0004-V0027
% #FIELD_VALUE: 测试程序—L10SSSLENA1P
\fieldvalue{\checkboxfield{unchecked}}L10SSSLENA1P\\
% #VALUE_ID: LEX-P0004-V0028
% #FIELD_VALUE: 测试程序—其他
% #HANDWRITTEN: ✓
\fieldvalue{\handwritten{\checkboxfield{checked}}}其他：
% #VALUE_ID: LEX-P0004-V0029
% #FIELD_VALUE: 测试程序—其他程序名称
% #HANDWRITTEN: SMPES-2025
\fieldvalue{\handwritten{SMPES-2025}}
\end{tabular}
&
&
\begin{tabular}{l}
% #VALUE_ID: LEX-P0004-V0030
% #FIELD_VALUE: 测试程序—I05SSSLENA1P
\fieldvalue{\checkboxfield{unchecked}}I05SSSLENA1P\\
% #VALUE_ID: LEX-P0004-V0031
% #FIELD_VALUE: 测试程序—L20SSSLENA1P
\fieldvalue{\checkboxfield{unchecked}}L20SSSLENA1P\\
% #VALUE_ID: LEX-P0004-V0032
% #FIELD_VALUE: 测试程序—C02BBSHMN1P
\fieldvalue{\checkboxfield{unchecked}}C02BBSHMN1P\\
% #VALUE_ID: LEX-P0004-V0033
% #FIELD_VALUE: 测试程序—L05BBSHMNA1P
\fieldvalue{\checkboxfield{unchecked}}L05BBSHMNA1P
\end{tabular}
\\
\hline
仪器状态 &
\begin{tabular}{c}
$\square$滤前\quad
% #VALUE_ID: LEX-P0004-V0034
% #FIELD_VALUE: 仪器状态第一组—滤后
% #HANDWRITTEN: ✓
\fieldvalue{\handwritten{\checkboxfield{checked}}}滤后
\end{tabular}
&
&
\begin{tabular}{c}
$\square$滤前\quad
% #VALUE_ID: LEX-P0004-V0035
% #FIELD_VALUE: 仪器状态第二组—滤后
% #HANDWRITTEN: ✓
\fieldvalue{\handwritten{\checkboxfield{checked}}}滤后
\end{tabular}
\\
\hline
& \multicolumn{3}{l|}{
% #VALUE_ID: LEX-P0004-V0036
% #FIELD_VALUE: 仪器状态第三组—滤前
\fieldvalue{\checkboxfield{unchecked}}滤前\quad
% #VALUE_ID: LEX-P0004-V0037
% #FIELD_VALUE: 仪器状态第三组—滤后
% #HANDWRITTEN: ✓
\fieldvalue{\handwritten{\checkboxfield{checked}}}滤后
}\\
\hline
消洗时长 &
% #VALUE_ID: LEX-P0004-V0038
% #FIELD_VALUE: 第一组消洗时长
% #HANDWRITTEN: 2
\fieldvalue{\handwritten{2}} min
&
% #VALUE_ID: LEX-P0004-V0039
% #FIELD_VALUE: 第二组消洗时长
% #HANDWRITTEN: 3
\fieldvalue{\handwritten{3}} min
&
% #VALUE_ID: LEX-P0004-V0040
% #FIELD_VALUE: 第三组消洗时长
% #HANDWRITTEN: 3
\fieldvalue{\handwritten{3}} min
\\
\hline
测试日期 &
% #VALUE_ID: LEX-P0004-V0041
% #FIELD_VALUE: 第一组测试日期
% #HANDWRITTEN: 2025.12.12
\fieldvalue{\handwritten{2025.12.12}}
&
% #VALUE_ID: LEX-P0004-V0042
% #FIELD_VALUE: 第二组测试日期
% #HANDWRITTEN: 2025.12.12
\fieldvalue{\handwritten{2025.12.12}}
&
% #VALUE_ID: LEX-P0004-V0043
% #FIELD_VALUE: 第三组测试日期
% #HANDWRITTEN: 2025.12.12
\fieldvalue{\handwritten{2025.12.12}}
\\
\hline
测试结果 &
\begin{tabular}{l}
% #VALUE_ID: LEX-P0004-V0044
% #FIELD_VALUE: 第一组测试结果—通过
% #HANDWRITTEN: ✓
\fieldvalue{\handwritten{\checkboxfield{checked}}}通过，实测滤点值：\\
% #VALUE_ID: LEX-P0004-V0045
% #FIELD_VALUE: 第一组实测滤点值
% #HANDWRITTEN: 450 mbar
\fieldvalue{\handwritten{450\,mbar}}\\
% #VALUE_ID: LEX-P0004-V0046
% #FIELD_VALUE: 第一组测试结果—失败
\fieldvalue{\checkboxfield{unchecked}}失败，失败原因：
\end{tabular}
&
\begin{tabular}{l}
% #VALUE_ID: LEX-P0004-V0047
% #FIELD_VALUE: 第二组测试结果—通过
% #HANDWRITTEN: ✓
\fieldvalue{\handwritten{\checkboxfield{checked}}}通过，实测滤点值：\\
% #VALUE_ID: LEX-P0004-V0048
% #FIELD_VALUE: 第二组实测滤点值
% #HANDWRITTEN: 450 mbar
\fieldvalue{\handwritten{450\,mbar}}\\
% #VALUE_ID: LEX-P0004-V0049
% #FIELD_VALUE: 第二组测试结果—失败
\fieldvalue{\checkboxfield{unchecked}}失败，失败原因：
\end{tabular}
&
\begin{tabular}{l}
% #VALUE_ID: LEX-P0004-V0050
% #FIELD_VALUE: 第三组测试结果—通过
% #HANDWRITTEN: ✓
\fieldvalue{\handwritten{\checkboxfield{checked}}}通过，实测滤点值：\\
% #VALUE_ID: LEX-P0004-V0051
% #FIELD_VALUE: 第三组实测滤点值
% #HANDWRITTEN: 5050 mbar
\fieldvalue{\handwritten{5050\,mbar}}\\
% #VALUE_ID: LEX-P0004-V0052
% #FIELD_VALUE: 第三组测试结果—失败
\fieldvalue{\checkboxfield{unchecked}}失败，失败原因：
\end{tabular}
\\
\hline
测试人/日期 &
% #VALUE_ID: LEX-P0004-V0053
% #FIELD_VALUE: 第一组测试人/日期
% #HANDWRITTEN: 陈强 2025.12.12
\fieldvalue{\handwritten{陈强\ 2025.12.12}}
&
% #VALUE_ID: LEX-P0004-V0054
% #FIELD_VALUE: 第二组测试人/日期
% #HANDWRITTEN: 陈强 2025.12.12
\fieldvalue{\handwritten{陈强\ 2025.12.12}}
&
% #VALUE_ID: LEX-P0004-V0055
% #FIELD_VALUE: 第三组测试人/日期
% #HANDWRITTEN: 陈强 2025.12.12
\fieldvalue{\handwritten{陈强\ 2025.12.12}}
\\
\hline
复核人/日期 &
% #VALUE_ID: LEX-P0004-V0056
% #FIELD_VALUE: 第一组复核人/日期
% #HANDWRITTEN: 刘艳玲 2025.12.12
\fieldvalue{\handwritten{刘艳玲\ 2025.12.12}}
&
% #VALUE_ID: LEX-P0004-V0057
% #FIELD_VALUE: 第二组复核人/日期
% #HANDWRITTEN: 刘艳玲 2025.12.12
\fieldvalue{\handwritten{刘艳玲\ 2025.12.12}}
&
% #VALUE_ID: LEX-P0004-V0058
% #FIELD_VALUE: 第三组复核人/日期
% #HANDWRITTEN: 刘艳玲 2025.12.12
\fieldvalue{\handwritten{刘艳玲\ 2025.12.12}}
\\
\hline
备注 &
\multicolumn{3}{l|}{
% #VALUE_ID: LEX-P0004-V0059
% #HANDWRITTEN: ✓
\handwritten{✓}
}\\
\hline
\end{tabular}

\normalsize
\begin{center}
第 1 页 共 1 页
\end{center}
% LEXOID_PAGE_COMPLETED: 4/70

\newpage

\section*{Palltronic Flowstar IV 测试报告}

\begin{tabular}{@{}p{0.28\linewidth}p{0.62\linewidth}@{}}
日期 &
% #VALUE_ID: LEX-P0005-V0001
% #FIELD_VALUE: 日期
\fieldvalue{12/20/2025 10:06:34}
\\
页码 &
% #VALUE_ID: LEX-P0005-V0002
% #FIELD_VALUE: 页码
\fieldvalue{1 of 2}
\\
Serial/issue number &
% #VALUE_ID: LEX-P0005-V0003
% #FIELD_VALUE: Serial/issue number
\fieldvalue{238807424?41}
\\
测试人员 &
% #VALUE_ID: LEX-P0005-V0004
% #FIELD_VALUE: 测试人员
\fieldvalue{Air}
\\
测试程序 &
% #VALUE_ID: LEX-P0005-V0005
% #FIELD_VALUE: 测试程序
\fieldvalue{S?F?E?S-2225}
\\
测试结果 &
% #VALUE_ID: LEX-P0005-V0006
% #FIELD_VALUE: 测试结果
\fieldvalue{Q2}
\\
装置 &
% #VALUE_ID: LEX-P0005-V0007
% #FIELD_VALUE: 装置
\fieldvalue{Shiyingzhou}
\\
设备编号 &
% #VALUE_ID: LEX-P0005-V0008
% #FIELD_VALUE: 设备编号
\fieldvalue{ZN?5?}
\\
程序编号 &
% #VALUE_ID: LEX-P0005-V0009
% #FIELD_VALUE: 程序编号
\fieldvalue{A?322?202?19}
\\
软件版本 &
% #VALUE_ID: LEX-P0005-V0010
% #FIELD_VALUE: 软件版本
\fieldvalue{S?F?E?S-2225}
\\
测试气体 &
% #VALUE_ID: LEX-P0005-V0011
% #FIELD_VALUE: 测试气体
\fieldvalue{Air}
\\
测试压力 &
% #VALUE_ID: LEX-P0005-V0012
% #FIELD_VALUE: 测试压力
\fieldvalue{4500 mbar}
\\
最大压力 &
% #VALUE_ID: LEX-P0005-V0013
% #FIELD_VALUE: 最大压力
\fieldvalue{7000 mbar}
\end{tabular}

\vspace{1em}

测试结果：\underline{气密性测试结果：4500 mbar}

\vspace{1em}

% #VALUE_ID: LEX-P0005-V0014
% #HANDWRITTEN: 审核签字 2026/12/2
\handwritten{审核签字 2026/12/2}

\vspace{2em}

% #VALUE_ID: LEX-P0005-V0015
% #HANDWRITTEN: [签名]
\handwritten{[签名]}

\section*{Palltronic Flowstar IV 测试报告}

\begin{tabular}{@{}p{0.28\linewidth}p{0.62\linewidth}@{}}
日期 &
% #VALUE_ID: LEX-P0005-V0016
% #FIELD_VALUE: 日期
\fieldvalue{12/20/2025 10:06:34}
\\
页码 &
% #VALUE_ID: LEX-P0005-V0017
% #FIELD_VALUE: 页码
\fieldvalue{2 of 2}
\end{tabular}

\vspace{1em}

\begin{center}
\begin{tabular}{@{}c@{}}
\textit{大气压差 (mbar/min)}\\[-0.3em]
\begin{tabular}{c|c|c|c|c}
1800 & 1350 & 900 & 450 & 0
\end{tabular}\\[0.5em]
\fbox{\parbox[c][0.30\textheight][c]{0.72\linewidth}{%
\centering
% TODO: Reproduce the visible pressure/time graph; the plotted curve descends sharply near the right side.
}}\\[-0.2em]
\begin{tabular}{c|c|c|c|c|c}
0 & 1400 & 2700 & 4200 & 5600 & 7000
\end{tabular}\\
时间（秒）
\end{tabular}
\end{center}

\vspace{1em}

% #VALUE_ID: LEX-P0005-V0018
% #HANDWRITTEN: [签名] 2026/12/2
\handwritten{[签名] 2026/12/2}

\vspace{1em}

% #VALUE_ID: LEX-P0005-V0019
% #HANDWRITTEN: 6/3
\handwritten{6/3}
% LEXOID_PAGE_COMPLETED: 5/70

\newpage

\section{Palltronic Flowstar IV 测试结果}

\begin{tabular}{@{}p{0.28\linewidth}p{0.62\linewidth}@{}}
日期 &
% #VALUE_ID: LEX-P0006-V0001
% #FIELD_VALUE: 日期
\fieldvalue{12/Dec/2025 10:21:19}
\\
页码 &
% #VALUE_ID: LEX-P0006-V0002
% #FIELD_VALUE: 页码
\fieldvalue{1 of 2}
\\
测试号 &
% #VALUE_ID: LEX-P0006-V0003
% #FIELD_VALUE: 测试号
\fieldvalue{258804241512}
\\
测试程序 &
% #VALUE_ID: LEX-P0006-V0004
% #FIELD_VALUE: 测试程序
\fieldvalue{SMP?E5225}
\\
测试结果 &
% #VALUE_ID: LEX-P0006-V0005
% #FIELD_VALUE: 测试结果
\fieldvalue{合格}
\\
测试介质 &
% #VALUE_ID: LEX-P0006-V0006
% #FIELD_VALUE: 测试介质
\fieldvalue{空气}
\\
测试气体 &
% #VALUE_ID: LEX-P0006-V0007
% #FIELD_VALUE: 测试气体
\fieldvalue{Air}
\\
测试压力 &
% #VALUE_ID: LEX-P0006-V0008
% #FIELD_VALUE: 测试压力
\fieldvalue{1000 mbar}
\\
最大压力 &
% #VALUE_ID: LEX-P0006-V0009
% #FIELD_VALUE: 最大压力
\fieldvalue{4500 mbar}
\\
测试大压力 &
% #VALUE_ID: LEX-P0006-V0010
% #FIELD_VALUE: 测试大压力
\fieldvalue{7000 mbar}
\end{tabular}

\subsection{结果}

\begin{tabular}{@{}p{0.28\linewidth}p{0.62\linewidth}@{}}
结果 &
% #VALUE_ID: LEX-P0006-V0011
% #FIELD_VALUE: 结果
\fieldvalue{通过}
\\
高压范围 &
% #VALUE_ID: LEX-P0006-V0012
% #FIELD_VALUE: 高压范围
\fieldvalue{5000 mbar}
\\
测试点 &
% #VALUE_ID: LEX-P0006-V0013
% #FIELD_VALUE: 测试点
\fieldvalue{5000 mbar}
\\
测试完成时间 &
% #VALUE_ID: LEX-P0006-V0014
% #FIELD_VALUE: 测试完成时间
\fieldvalue{12/Dec/2025 10:31:51}
\end{tabular}

\subsection{注释}

\begin{enumerate}
\item 整体测试在高压范围内。
\end{enumerate}

\subsection{签名}

% #VALUE_ID: LEX-P0006-V0015
% #FIELD_VALUE: 签名
% #HANDWRITTEN: [签名]
\fieldvalue{\handwritten{[签名]}}

\bigskip

\section{Palltronic Flowstar IV 测试结果}

\begin{tabular}{@{}p{0.28\linewidth}p{0.62\linewidth}@{}}
日期 &
% #VALUE_ID: LEX-P0006-V0016
% #FIELD_VALUE: 日期
\fieldvalue{12/Dec/2025 10:21:19}
\\
页码 &
% #VALUE_ID: LEX-P0006-V0017
% #FIELD_VALUE: 页码
\fieldvalue{2 of 2}
\end{tabular}

\subsection{测试曲线}

% TODO: Reproduce the visible pressure-flow graph, including the plotted curve and axis labels.
\begin{center}
\fbox{\parbox[c][0.30\textheight][c]{0.82\linewidth}{\centering
\textit{测试曲线图（原图位置）}
}}
\end{center}

\subsection{签名}

% #VALUE_ID: LEX-P0006-V0018
% #FIELD_VALUE: 签名
% #HANDWRITTEN: [签名]
\fieldvalue{\handwritten{[签名]}}

% #VALUE_ID: LEX-P0006-V0019
% #HANDWRITTEN: 2025/12/12
\handwritten{2025/12/12}

% TODO #HANDWRITTEN: 7/3?; partially legible handwritten notation at the right edge.
% #VALUE_ID: LEX-P0006-V0020
\handwritten{7/3?}
% LEXOID_PAGE_COMPLETED: 6/70

\newpage

\section{Palltronic Flowstar IV 测试结果}

\begin{tabular}{@{}p{0.28\linewidth}p{0.62\linewidth}@{}}
日期 &
% #VALUE_ID: LEX-P0007-V0001
% #FIELD_VALUE: 日期
\fieldvalue{12/Dec/2025 10:35:49}
\\[2pt]
时间 &
% #VALUE_ID: LEX-P0007-V0002
% #FIELD_VALUE: 时间
\fieldvalue{10:42}
\\[2pt]
Serial/Result number &
% #TODO: Serial/Result number is partially unclear in the scan.
% #VALUE_ID: LEX-P0007-V0003
% #FIELD_VALUE: Serial/Result number
\fieldvalue{23880744-41616}
\\[2pt]
测试程序 &
% #VALUE_ID: LEX-P0007-V0004
% #FIELD_VALUE: 测试程序
\fieldvalue{SFM?ES-2255}
\\[2pt]
测试类型 &
% #VALUE_ID: LEX-P0007-V0005
% #FIELD_VALUE: 测试类型
\fieldvalue{QZ}
\\[2pt]
生产区域 &
% #VALUE_ID: LEX-P0007-V0006
% #FIELD_VALUE: 生产区域
\fieldvalue{C}
\\[2pt]
过滤器规格 &
% #VALUE_ID: LEX-P0007-V0007
% #FIELD_VALUE: 过滤器规格
\fieldvalue{2\textsuperscript{''}}
\\[2pt]
过滤器编号 &
% #TODO: Filter number is partially unclear in the scan.
% #VALUE_ID: LEX-P0007-V0008
% #FIELD_VALUE: 过滤器编号
\fieldvalue{A4222025-21019}
\\[2pt]
产品编号 &
% #VALUE_ID: LEX-P0007-V0009
% #FIELD_VALUE: 产品编号
\fieldvalue{SFM?ES-2255}
\\[2pt]
产品名称 &
% #TODO: Product name is unclear in the scan.
% #VALUE_ID: LEX-P0007-V0010
% #FIELD_VALUE: 产品名称
\fieldvalue{Shu?}
\\[2pt]
过滤器材质 &
% #VALUE_ID: LEX-P0007-V0011
% #FIELD_VALUE: 过滤器材质
\fieldvalue{Air}
\\[2pt]
最大测试压力 &
% #VALUE_ID: LEX-P0007-V0012
% #FIELD_VALUE: 最大测试压力
\fieldvalue{4000\ \mathrm{mbar}}
\\[2pt]
最小测试压力 &
% #VALUE_ID: LEX-P0007-V0013
% #FIELD_VALUE: 最小测试压力
\fieldvalue{1000\ \mathrm{mbar}}
\\[2pt]
最小测试点 &
% #VALUE_ID: LEX-P0007-V0014
% #FIELD_VALUE: 最小测试点
\fieldvalue{4500\ \mathrm{mbar}}
\\[2pt]
结果 &
% #VALUE_ID: LEX-P0007-V0015
% #FIELD_VALUE: 结果
\fieldvalue{测试合格，4950\ \mathrm{mbar}}
\\[2pt]
测试完成时间 &
% #VALUE_ID: LEX-P0007-V0016
% #FIELD_VALUE: 测试完成时间
\fieldvalue{12/Dec/2025 10:46:14}
\\[2pt]
签名 &
% #VALUE_ID: LEX-P0007-V0017
% #FIELD_VALUE: 签名
% #TODO #HANDWRITTEN: illegible signature; blue handwritten characters are unclear.
\fieldvalue{\handwritten{A?B?-?4-?30192}}
\end{tabular}

\vspace{1em}

% #VALUE_ID: LEX-P0007-V0018
% #HANDWRITTEN: 2026/12/2
\handwritten{2026/12/2}

% #VALUE_ID: LEX-P0007-V0019
% #HANDWRITTEN: 签名（字迹不清）
\handwritten{签名}

\section{测试曲线}

\begin{center}
\fbox{\parbox{0.86\linewidth}{%
\centering
% TODO: Figure/graph visible on this page has no available filename.
测试曲线图\\[4pt]
纵轴：大气压差（mbar/min）\\
横轴：时间（分钟）\\[4pt]
曲线刻度可见：0、450、900、1350、1800；\\
横轴刻度可见：0、1400、2800、4200、5600、7000。
}}
\end{center}

\begin{tabular}{@{}p{0.28\linewidth}p{0.62\linewidth}@{}}
日期 &
% #VALUE_ID: LEX-P0007-V0020
% #FIELD_VALUE: 日期
\fieldvalue{12/Dec/2025 10:35:49}
\\[2pt]
时间 &
% #VALUE_ID: LEX-P0007-V0021
% #FIELD_VALUE: 时间
\fieldvalue{10:42}
\\[2pt]
签名 &
% #VALUE_ID: LEX-P0007-V0022
% #FIELD_VALUE: 签名
% #TODO #HANDWRITTEN: illegible signature; blue handwritten characters are unclear.
\fieldvalue{\handwritten{签名}}
\end{tabular}

% #VALUE_ID: LEX-P0007-V0023
% #HANDWRITTEN: 8/35
\handwritten{8/35}

% #VALUE_ID: LEX-P0007-V0024
% #HANDWRITTEN: 2026/12/2
\handwritten{2026/12/2}
% LEXOID_PAGE_COMPLETED: 7/70

\newpage

\section{Palltronic Flowstar IV 测漏合格结果}

\begin{tabular}{@{}p{0.34\linewidth}p{0.56\linewidth}@{}}
日期 / Date &
% #VALUE_ID: LEX-P0008-V0001
% #FIELD_VALUE: 日期 / Date
\fieldvalue{12/Dec/2025 10:50:36}
\\[2pt]
页码 / Page &
% #VALUE_ID: LEX-P0008-V0002
% #FIELD_VALUE: 页码 / Page
\fieldvalue{1 of 2}
\\[2pt]
Serial/instrument number &
% #VALUE_ID: LEX-P0008-V0003
% #FIELD_VALUE: Serial/instrument number
\fieldvalue{238804746?67}
\\[2pt]
测试气体 &
% #VALUE_ID: LEX-P0008-V0004
% #FIELD_VALUE: 测试气体
\fieldvalue{Air}
\\[2pt]
测试温度 &
% #VALUE_ID: LEX-P0008-V0005
% #FIELD_VALUE: 测试温度
\fieldvalue{21.?\,^\circ\mathrm{C}}
\\[2pt]
测试压力 &
% #VALUE_ID: LEX-P0008-V0006
% #FIELD_VALUE: 测试压力
\fieldvalue{4500\,mbar}
\\[2pt]
测试气体 &
% #VALUE_ID: LEX-P0008-V0007
% #FIELD_VALUE: 测试气体
\fieldvalue{Air}
\\[2pt]
过滤器型号 &
% #VALUE_ID: LEX-P0008-V0008
% #FIELD_VALUE: 过滤器型号
\fieldvalue{Sh?ngy?lou}
\\[2pt]
过滤器序列号 &
% #VALUE_ID: LEX-P0008-V0009
% #FIELD_VALUE: 过滤器序列号
\fieldvalue{A423?025012019}
\\[2pt]
程序编号 &
% #VALUE_ID: LEX-P0008-V0010
% #FIELD_VALUE: 程序编号
\fieldvalue{SFMP?S-225}
\\[2pt]
产品名称 &
% #VALUE_ID: LEX-P0008-V0011
% #FIELD_VALUE: 产品名称
\fieldvalue{---}
\\[2pt]
产品编号 &
% #VALUE_ID: LEX-P0008-V0012
% #FIELD_VALUE: 产品编号
\fieldvalue{---}
\\[2pt]
测试次数 &
% #VALUE_ID: LEX-P0008-V0013
% #FIELD_VALUE: 测试次数
\fieldvalue{1.000}
\\[2pt]
最大允许流量 &
% #VALUE_ID: LEX-P0008-V0014
% #FIELD_VALUE: 最大允许流量
\fieldvalue{4500\,mbar}
\\[2pt]
最小允许流量 &
% #VALUE_ID: LEX-P0008-V0015
% #FIELD_VALUE: 最小允许流量
\fieldvalue{7200\,mbar}
\\[2pt]
流量稳定时间 &
% #VALUE_ID: LEX-P0008-V0016
% #FIELD_VALUE: 流量稳定时间
\fieldvalue{1:00}
\\[2pt]
测试完成时间 &
% #VALUE_ID: LEX-P0008-V0017
% #FIELD_VALUE: 测试完成时间
\fieldvalue{19/Dec/2025 11:01:09}
\end{tabular}

\vspace{1em}

\noindent\textbf{注释 / Notes:}

\vspace{1.5em}

\noindent\textbf{签名 / Signature:}

% #VALUE_ID: LEX-P0008-V0018
% #FIELD_VALUE: 签名 / Signature
% #TODO #HANDWRITTEN: [签名不清]; 蓝色手写签名难以辨认
\fieldvalue{\handwritten{[签名不清]}} \hfill
% #VALUE_ID: LEX-P0008-V0019
% #FIELD_VALUE: 日期
% #HANDWRITTEN: 2025/12/12
\fieldvalue{\handwritten{2025/12/12}}

\vfill

\section{Palltronic Flowstar IV 测漏合格结果}

\begin{tabular}{@{}p{0.34\linewidth}p{0.56\linewidth}@{}}
日期 / Date &
% #VALUE_ID: LEX-P0008-V0020
% #FIELD_VALUE: 日期 / Date
\fieldvalue{12/Dec/2025 10:50:36}
\\[2pt]
页码 / Page &
% #VALUE_ID: LEX-P0008-V0021
% #FIELD_VALUE: 页码 / Page
\fieldvalue{2 of 2}
\end{tabular}

\vspace{1em}

% TODO: The following graph is visible on this page; reproduce the plotted curve if vector extraction is available.
\begin{center}
\fbox{\parbox[c][0.34\textheight][c]{0.86\linewidth}{\centering
\textit{图表：流量变化（mbar/min）}\\[0.5em]
纵轴刻度：0、450、900、1350、1800\\
横轴刻度：0、1400、2800、4200、5600、7000\\[0.5em]
\textit{页面中显示一条在约 2800--4200 区间下降的曲线}
}}
\end{center}

\noindent\textbf{签名 / Signature:}

% #VALUE_ID: LEX-P0008-V0022
% #FIELD_VALUE: 签名 / Signature
% #TODO #HANDWRITTEN: [签名不清]; 蓝色手写签名难以辨认
\fieldvalue{\handwritten{[签名不清]}} \hfill
% #VALUE_ID: LEX-P0008-V0023
% #FIELD_VALUE: 日期
% #HANDWRITTEN: 2025/12/12
\fieldvalue{\handwritten{2025/12/12}}

% #VALUE_ID: LEX-P0008-V0024
% #HANDWRITTEN: 9/35
\handwritten{9/35}
% LEXOID_PAGE_COMPLETED: 8/70

\newpage

\section{Palltronic Flowstar IV 流量结果}

\begin{tabular}{@{}p{0.32\linewidth}p{0.58\linewidth}@{}}
日期 &
% #VALUE_ID: LEX-P0009-V0001
% #FIELD_VALUE: 日期
\fieldvalue{12/Dec/2025 11:04:18}
\\[2pt]
仪号 &
% #VALUE_ID: LEX-P0009-V0002
% #FIELD_VALUE: 仪号
\fieldvalue{FES084R2-21?1}
\\[2pt]
Serial/Result number &
% #VALUE_ID: LEX-P0009-V0003
% #FIELD_VALUE: Serial/Result number
\fieldvalue{28807442618}
\\[2pt]
测量序号 &
% #VALUE_ID: LEX-P0009-V0004
% #FIELD_VALUE: 测量序号
\fieldvalue{SFMP?S2225}
\\[2pt]
滤器型号 &
% #VALUE_ID: LEX-P0009-V0005
% #FIELD_VALUE: 滤器型号
\fieldvalue{C}
\\[2pt]
滤器序列号 &
% #VALUE_ID: LEX-P0009-V0006
% #FIELD_VALUE: 滤器序列号
\fieldvalue{Q?}
\\[2pt]
测量人员 &
% #VALUE_ID: LEX-P0009-V0007
% #FIELD_VALUE: 测量人员
\fieldvalue{Shiyinglou}
\\[2pt]
校验编号 &
% #VALUE_ID: LEX-P0009-V0008
% #FIELD_VALUE: 校验编号
\fieldvalue{A342?2025-2019}
\\[2pt]
产品型号 &
% #VALUE_ID: LEX-P0009-V0009
% #FIELD_VALUE: 产品型号
\fieldvalue{A?}
\\[2pt]
产品编号 &
% #VALUE_ID: LEX-P0009-V0010
% #FIELD_VALUE: 产品编号
\fieldvalue{SFMP?S-2225}
\\[2pt]
测试模式 &
% #VALUE_ID: LEX-P0009-V0011
% #FIELD_VALUE: 测试模式
\fieldvalue{C?}
\\[2pt]
测量气体 &
% #VALUE_ID: LEX-P0009-V0012
% #FIELD_VALUE: 测量气体
\fieldvalue{Air}
\\[2pt]
流量单位 &
% #VALUE_ID: LEX-P0009-V0013
% #FIELD_VALUE: 流量单位
\fieldvalue{Sm\textsuperscript{3}/h}
\\[2pt]
最大流量 &
% #VALUE_ID: LEX-P0009-V0014
% #FIELD_VALUE: 最大流量
\fieldvalue{1.000}
\\[2pt]
最大压力 &
% #VALUE_ID: LEX-P0009-V0015
% #FIELD_VALUE: 最大压力
\fieldvalue{4500 mbar}
\\[2pt]
最小压力 &
% #VALUE_ID: LEX-P0009-V0016
% #FIELD_VALUE: 最小压力
\fieldvalue{2000 mbar}
\\[2pt]
结果 &
% #VALUE_ID: LEX-P0009-V0017
% #FIELD_VALUE: 结果
\fieldvalue{流量范围内：5000 mbar}
\\[2pt]
校验时间 &
% #VALUE_ID: LEX-P0009-V0018
% #FIELD_VALUE: 校验时间
\fieldvalue{12/Dec/2025 11:44:49}
\end{tabular}

\vspace{1em}

签名：%
% #VALUE_ID: LEX-P0009-V0019
% #FIELD_VALUE: 签名
% #TODO #HANDWRITTEN: 手写中文签名；字迹难以辨认
\fieldvalue{\handwritten{手写中文签名}}

日期：%
% #VALUE_ID: LEX-P0009-V0020
% #FIELD_VALUE: 日期
% #HANDWRITTEN: 2025/12/12
\fieldvalue{\handwritten{2025/12/12}}

\section{Palltronic Flowstar IV 流量结果}

\begin{tabular}{@{}p{0.32\linewidth}p{0.58\linewidth}@{}}
日期 &
% #VALUE_ID: LEX-P0009-V0021
% #FIELD_VALUE: 日期
\fieldvalue{12/Dec/2025 11:04:18}
\\[2pt]
测量点 &
% #VALUE_ID: LEX-P0009-V0022
% #FIELD_VALUE: 测量点
\fieldvalue{2.02}
\end{tabular}

\vspace{1em}

\begin{center}
\setlength{\fboxsep}{8pt}
\fbox{\begin{minipage}{0.78\linewidth}
\centering
% TODO: Reproduce the visible pressure--flow graph, including the axes,
% tick labels, and plotted curve.
\textit{压力--流量曲线图（原图可见）}

\vspace{0.5em}
纵轴：大气压力（mbar）\qquad 横轴：流量（mbar）
\end{minipage}}
\end{center}

\vspace{1em}

签名：%
% #VALUE_ID: LEX-P0009-V0023
% #FIELD_VALUE: 签名
% #TODO #HANDWRITTEN: 手写中文签名；字迹难以辨认
\fieldvalue{\handwritten{手写中文签名}}

日期：%
% #VALUE_ID: LEX-P0009-V0024
% #FIELD_VALUE: 日期
% #HANDWRITTEN: 2025/12/12
\fieldvalue{\handwritten{2025/12/12}}

\vspace{0.5em}

% #VALUE_ID: LEX-P0009-V0025
% #FIELD_VALUE: 右下角手写标记
% #TODO #HANDWRITTEN: 10/35；字符可能为手写页码或标记
\fieldvalue{\handwritten{10/35}}
% LEXOID_PAGE_COMPLETED: 9/70

\newpage

\section{Palltronic Flowstar IV 测试结果}

\begin{tabular}{@{}p{0.24\linewidth}p{0.68\linewidth}@{}}
日期 &
% #VALUE_ID: LEX-P0010-V0001
% #FIELD_VALUE: 日期
\fieldvalue{12/Dec/2025 11:19:02}
\\
页码 &
% #VALUE_ID: LEX-P0010-V0002
% #FIELD_VALUE: 页码
\fieldvalue{1 of 1}
\\
Serial number &
% #VALUE_ID: LEX-P0010-V0003
% #FIELD_VALUE: Serial number
\fieldvalue{\textit{[部分字符不清]}}
\\
功能 &
% #VALUE_ID: LEX-P0010-V0004
% #FIELD_VALUE: 功能
\fieldvalue{Bubble Point}
\\
测试程序 &
% #VALUE_ID: LEX-P0010-V0005
% #FIELD_VALUE: 测试程序
\fieldvalue{SMP?FS-2225}
\\
滤器编号 &
% #VALUE_ID: LEX-P0010-V0006
% #FIELD_VALUE: 滤器编号
\fieldvalue{Q}
\\
滤器类型 &
% #VALUE_ID: LEX-P0010-V0007
% #FIELD_VALUE: 滤器类型
\fieldvalue{\textit{[字符不清]}}
\\
产品名称 &
% #VALUE_ID: LEX-P0010-V0008
% #FIELD_VALUE: 产品名称
\fieldvalue{\textit{[字符不清]}}
\\
批号 &
% #VALUE_ID: LEX-P0010-V0009
% #FIELD_VALUE: 批号
\fieldvalue{A222?}
\\
测试压力 &
% #VALUE_ID: LEX-P0010-V0010
% #FIELD_VALUE: 测试压力
\fieldvalue{4 bar}
\\
测试介质 &
% #VALUE_ID: LEX-P0010-V0011
% #FIELD_VALUE: 测试介质
\fieldvalue{Air}
\\
测试气体 &
% #VALUE_ID: LEX-P0010-V0012
% #FIELD_VALUE: 测试气体
\fieldvalue{Air}
\\
测试流量 &
% #VALUE_ID: LEX-P0010-V0013
% #FIELD_VALUE: 测试流量
\fieldvalue{5000}
\\
最小压力 &
% #VALUE_ID: LEX-P0010-V0014
% #FIELD_VALUE: 最小压力
\fieldvalue{1500 mbar}
\\
最大压力 &
% #VALUE_ID: LEX-P0010-V0015
% #FIELD_VALUE: 最大压力
\fieldvalue{7000 mbar}
\\
测试点 &
% #VALUE_ID: LEX-P0010-V0016
% #FIELD_VALUE: 测试点
\fieldvalue{4950 mbar}
\\
测试完成时间 &
% #VALUE_ID: LEX-P0010-V0017
% #FIELD_VALUE: 测试完成时间
\fieldvalue{12/Dec/2025 11:29:24}
\\
\end{tabular}

\vspace{1em}

\noindent 注释

\begin{enumerate}
\item 签名：

% #VALUE_ID: LEX-P0010-V0018
% #FIELD_VALUE: 签名
% #HANDWRITTEN: A2-?A?A2-?A?（难以辨认）
\fieldvalue{\handwritten{A2-?A?A2-?A?}}
\end{enumerate}

% #VALUE_ID: LEX-P0010-V0019
% #HANDWRITTEN: [中文签名]（难以辨认）
\handwritten{[中文签名]}

% #VALUE_ID: LEX-P0010-V0020
% #HANDWRITTEN: 2025.12.12
\handwritten{2025.12.12}

\section{Palltronic Flowstar IV 测试结果}

\begin{tabular}{@{}p{0.24\linewidth}p{0.68\linewidth}@{}}
日期 &
% #VALUE_ID: LEX-P0010-V0021
% #FIELD_VALUE: 日期
\fieldvalue{12/Dec/2025 11:19:02}
\\
页码 &
% #VALUE_ID: LEX-P0010-V0022
% #FIELD_VALUE: 页码
\fieldvalue{2 of 2}
\\
\end{tabular}

\vspace{1em}

\begin{figure}[ht]
\centering
\fbox{\parbox[c][0.34\textheight][c]{0.82\linewidth}{\centering
% TODO: 此处为测试压力/时间曲线图；原页面显示纵轴刻度约为
% 0、450、900、1350、1800 mbar，横轴刻度约为
% 0、1400、2800、4200、5600、7000 ms。
测试曲线图
}}
\caption{测试压力曲线}
\end{figure}

\noindent
纵轴：气体流量（mbar/min）；横轴：时间（msec）。

\vspace{1em}

% #VALUE_ID: LEX-P0010-V0023
% #HANDWRITTEN: [中文签名]（难以辨认）
\handwritten{[中文签名]}

% #VALUE_ID: LEX-P0010-V0024
% #HANDWRITTEN: 2025.12.12
\handwritten{2025.12.12}

% #VALUE_ID: LEX-P0010-V0025
% #HANDWRITTEN: 11/35
\handwritten{11/35}
% LEXOID_PAGE_COMPLETED: 10/70

\newpage

\section*{Palltronic Flowstar IV 测试报告}

\begin{tabular}{@{}p{0.30\linewidth}p{0.60\linewidth}@{}}
日期 &
% #VALUE_ID: LEX-P0011-V0001
% #FIELD_VALUE: 日期
\fieldvalue{12/Dec/2025 11:35:0} \\

页码 &
% #VALUE_ID: LEX-P0011-V0002
% #FIELD_VALUE: 页码
\fieldvalue{1 of 2} \\

Serial/Result number &
% #VALUE_ID: LEX-P0011-V0003
% #FIELD_VALUE: Serial/Result number
% #TODO: Transcription uncertain in the scan.
\fieldvalue{A2880/424?260 / 12?T} \\

测试名称 &
% #VALUE_ID: LEX-P0011-V0004
% #FIELD_VALUE: 测试名称
% #TODO: Transcription uncertain in the scan.
\fieldvalue{S?M?FS-2225} \\

测试 &
% #VALUE_ID: LEX-P0011-V0005
% #FIELD_VALUE: 测试
% #TODO: Transcription uncertain in the scan.
\fieldvalue{C?} \\

操作员 &
% #VALUE_ID: LEX-P0011-V0006
% #FIELD_VALUE: 操作员
\fieldvalue{Shiyonghou} \\

过滤器编号 &
% #VALUE_ID: LEX-P0011-V0007
% #FIELD_VALUE: 过滤器编号
% #TODO: Value is faint/unclear in the scan.
\fieldvalue{?} \\

产品名称 &
% #VALUE_ID: LEX-P0011-V0008
% #FIELD_VALUE: 产品名称
% #TODO: Transcription uncertain in the scan.
\fieldvalue{Zh?sh?hou} \\

产品编号 &
% #VALUE_ID: LEX-P0011-V0009
% #FIELD_VALUE: 产品编号
% #TODO: Value is faint/unclear in the scan.
\fieldvalue{?} \\

压力 &
% #VALUE_ID: LEX-P0011-V0010
% #FIELD_VALUE: 压力
\fieldvalue{5000 mbar} \\

测试介质 &
% #VALUE_ID: LEX-P0011-V0011
% #FIELD_VALUE: 测试介质
\fieldvalue{Air} \\

过滤器型号 &
% #VALUE_ID: LEX-P0011-V0012
% #FIELD_VALUE: 过滤器型号
% #TODO: Transcription uncertain in the scan.
\fieldvalue{S?M?FS-2225} \\

过滤器序列号 &
% #VALUE_ID: LEX-P0011-V0013
% #FIELD_VALUE: 过滤器序列号
% #TODO: Transcription uncertain in the scan.
\fieldvalue{S?M?FS-2225} \\

最大压力 &
% #VALUE_ID: LEX-P0011-V0014
% #FIELD_VALUE: 最大压力
\fieldvalue{7000 mbar} \\

最小压力 &
% #VALUE_ID: LEX-P0011-V0015
% #FIELD_VALUE: 最小压力
\fieldvalue{4500 mbar} \\

测试时间 &
% #VALUE_ID: LEX-P0011-V0016
% #FIELD_VALUE: 测试时间
\fieldvalue{12/Dec/2025 11:43:13} \\

结果 &
% #VALUE_ID: LEX-P0011-V0017
% #FIELD_VALUE: 结果
\fieldvalue{Passed} \\

测试完成时间 &
% #VALUE_ID: LEX-P0011-V0018
% #FIELD_VALUE: 测试完成时间
\fieldvalue{12/Dec/2025 11:43:13} \\
\end{tabular}

\vspace{1em}

签名：%
% #VALUE_ID: LEX-P0011-V0019
% #FIELD_VALUE: 签名
% #TODO #HANDWRITTEN: illegible signature; handwritten Chinese characters.
\fieldvalue{\handwritten{签名}}

\hspace{2em}日期：%
% #VALUE_ID: LEX-P0011-V0020
% #FIELD_VALUE: 日期
% #HANDWRITTEN: 2025/12/12
\fieldvalue{\handwritten{2025/12/12}}

\vspace{2em}

\section*{Palltronic Flowstar IV 测试报告}

\begin{tabular}{@{}p{0.30\linewidth}p{0.60\linewidth}@{}}
日期 &
% #VALUE_ID: LEX-P0011-V0021
% #FIELD_VALUE: 日期
\fieldvalue{12/Dec/2025 11:35:50} \\

页码 &
% #VALUE_ID: LEX-P0011-V0022
% #FIELD_VALUE: 页码
\fieldvalue{2 of 2} \\
\end{tabular}

\vspace{1em}

\begin{center}
\begin{tikzpicture}
  \draw (0,0) rectangle (0.78\linewidth,4.2cm);
  \draw[->] (0,0) -- (0.78\linewidth,0)
    node[below right]{时间 (test time)};
  \draw[->] (0,0) -- (0,4.2cm)
    node[above left]{气体流量 (mbar/min)};
  \draw[thick] (0.02\linewidth,0.35cm)
    -- (0.38\linewidth,0.35cm)
    .. controls (0.47\linewidth,0.35cm) and (0.48\linewidth,0.85cm)
    .. (0.49\linewidth,1.45cm)
    -- (0.49\linewidth,3.95cm);
  \node[below] at (0.10\linewidth,0){1400};
  \node[below] at (0.28\linewidth,0){2000};
  \node[below] at (0.42\linewidth,0){2400};
  \node[below] at (0.60\linewidth,0){5600};
  \node[below] at (0.74\linewidth,0){7000};
  \node[left] at (0,0.4cm){0};
  \node[left] at (0,1.4cm){450};
  \node[left] at (0,2.4cm){900};
  \node[left] at (0,3.4cm){1350};
  \node[left] at (0,4.0cm){1800};
\end{tikzpicture}
\end{center}

签名：%
% #VALUE_ID: LEX-P0011-V0023
% #FIELD_VALUE: 签名
% #TODO #HANDWRITTEN: illegible signature; handwritten Chinese characters.
\fieldvalue{\handwritten{签名}}

\hspace{2em}日期：%
% #VALUE_ID: LEX-P0011-V0024
% #FIELD_VALUE: 日期
% #HANDWRITTEN: 2025/12/12
\fieldvalue{\handwritten{2025/12/12}}

% #VALUE_ID: LEX-P0011-V0025
% #HANDWRITTEN: 12/13
\handwritten{12/13}
% LEXOID_PAGE_COMPLETED: 11/70

\newpage

\section*{Palltronic Flowstar IV 测漏结果}

\begin{tabular}{@{}p{0.30\linewidth}p{0.60\linewidth}@{}}
日期 &
% #VALUE_ID: LEX-P0012-V0001
% #FIELD_VALUE: 日期
\fieldvalue{12/Dec/2025 13:32:24}
\\
页码 &
% #VALUE_ID: LEX-P0012-V0002
% #FIELD_VALUE: 页码
\fieldvalue{1 of 2}
\\
Serial/Result number &
% #VALUE_ID: LEX-P0012-V0003
% #FIELD_VALUE: Serial/Result number
\fieldvalue{288074462621}
\\
测试序号 &
% #VALUE_ID: LEX-P0012-V0004
% #FIELD_VALUE: 测试序号
\fieldvalue{SFMPes-2225}
\\
测试程序 &
% #VALUE_ID: LEX-P0012-V0005
% #FIELD_VALUE: 测试程序
\fieldvalue{C2}
\\
测试模式 &
% #VALUE_ID: LEX-P0012-V0006
% #FIELD_VALUE: 测试模式
\fieldvalue{C2}
\\
生产区域 &
% #VALUE_ID: LEX-P0012-V0007
% #FIELD_VALUE: 生产区域
\fieldvalue{Shiyonghou}
\\
过滤器类型 &
% #VALUE_ID: LEX-P0012-V0008
% #FIELD_VALUE: 过滤器类型
\fieldvalue{2}
\\
产品名称 &
% #VALUE_ID: LEX-P0012-V0009
% #FIELD_VALUE: 产品名称
\fieldvalue{AA3220251029}
\\
过滤器编号 &
% #VALUE_ID: LEX-P0012-V0010
% #FIELD_VALUE: 过滤器编号
\fieldvalue{SFMPes-2225}
\\
测试气体 &
% #VALUE_ID: LEX-P0012-V0011
% #FIELD_VALUE: 测试气体
\fieldvalue{Air}
\\
测试压力 &
% #VALUE_ID: LEX-P0012-V0012
% #FIELD_VALUE: 测试压力
\fieldvalue{1000 mbar}
\\
最小测试压力 &
% #VALUE_ID: LEX-P0012-V0013
% #FIELD_VALUE: 最小测试压力
\fieldvalue{4500 mbar}
\\
最小泡点 &
% #VALUE_ID: LEX-P0012-V0014
% #FIELD_VALUE: 最小泡点
\fieldvalue{2000 mbar}
\\
结果 &
% #VALUE_ID: LEX-P0012-V0015
% #FIELD_VALUE: 结果
\fieldvalue{通过}
\\
泡点结果 &
% #VALUE_ID: LEX-P0012-V0016
% #FIELD_VALUE: 泡点结果
\fieldvalue{泡点结果区间：5000 mbar}
\\
测试总时间 &
% #VALUE_ID: LEX-P0012-V0017
% #FIELD_VALUE: 测试总时间
\fieldvalue{1:20/Dec/2025 15:27:06}
\end{tabular}

\vspace{1em}

审核签名：%
% #VALUE_ID: LEX-P0012-V0018
% #FIELD_VALUE: 审核签名
% #HANDWRITTEN: [签名，难以辨认]
\fieldvalue{\handwritten{[签名，难以辨认]}}
\qquad
日期：%
% #VALUE_ID: LEX-P0012-V0019
% #FIELD_VALUE: 日期
% #HANDWRITTEN: 2025.12.12
\fieldvalue{\handwritten{2025.12.12}}

\vspace{2em}

\section*{Palltronic Flowstar IV 测漏结果}

\begin{tabular}{@{}p{0.30\linewidth}p{0.60\linewidth}@{}}
日期 &
% #VALUE_ID: LEX-P0012-V0020
% #FIELD_VALUE: 日期
\fieldvalue{12/Dec/2025 13:32:24}
\\
页码 &
% #VALUE_ID: LEX-P0012-V0021
% #FIELD_VALUE: 页码
\fieldvalue{2 of 2}
\end{tabular}

\vspace{1em}

\noindent
% TODO: The visible graph is reproduced schematically; exact plotted curve and axis typography may require the source figure.
\begin{center}
\begin{tabular}{@{}c@{}}
\framebox[0.72\linewidth][c]{%
\begin{minipage}[c][0.28\textheight][c]{0.68\linewidth}
\centering
大气流量（mbar/min）\\[0.5em]
\makebox[\linewidth]{\rule{0pt}{0.18\textheight}\hfill\raisebox{0.02\textheight}{\rule{0.55\linewidth}{0.4pt}}\hfill}\\[-0.8em]
\hfill 流量曲线图 \hfill
\end{minipage}}
\end{tabular}
\end{center}

\vspace{1em}

审核签名：%
% #VALUE_ID: LEX-P0012-V0022
% #FIELD_VALUE: 审核签名
% #HANDWRITTEN: [签名，难以辨认]
\fieldvalue{\handwritten{[签名，难以辨认]}}
\qquad
日期：%
% #VALUE_ID: LEX-P0012-V0023
% #FIELD_VALUE: 日期
% #HANDWRITTEN: 2025.12.12
\fieldvalue{\handwritten{2025.12.12}}

\vspace{1em}

% #VALUE_ID: LEX-P0012-V0024
% #HANDWRITTEN: 3/35
\handwritten{3/35}
% LEXOID_PAGE_COMPLETED: 12/70

\newpage

\begin{center}
\textbf{PALL\quad Paltronic FlowStart IV 流量站结果}
\end{center}

\begin{tabular}{@{}p{0.28\linewidth}p{0.62\linewidth}@{}}
日期 &
% #VALUE_ID: LEX-P0013-V0001
% #FIELD_VALUE: 日期
\fieldvalue{12/Dec/2025 13:27:38}
\\[2pt]
产品 &
% #VALUE_ID: LEX-P0013-V0002
% #FIELD_VALUE: 产品
\fieldvalue{FFS04LGR/2/24/1T}
\\[2pt]
Serial number &
% #VALUE_ID: LEX-P0013-V0003
% #FIELD_VALUE: Serial number
\fieldvalue{230287044262}
\\[2pt]
功能 &
% #VALUE_ID: LEX-P0013-V0004
% #FIELD_VALUE: 功能
\fieldvalue{FlowStart}
\\[2pt]
测试模式 &
% #VALUE_ID: LEX-P0013-V0005
% #FIELD_VALUE: 测试模式
\fieldvalue{Q2}
\\[2pt]
测试程序 &
% #VALUE_ID: LEX-P0013-V0006
% #FIELD_VALUE: 测试程序
\fieldvalue{SFMPE5-22SS}
\\[2pt]
气体 &
% #VALUE_ID: LEX-P0013-V0007
% #FIELD_VALUE: 气体
\fieldvalue{C}
\\[2pt]
黏度 &
% #VALUE_ID: LEX-P0013-V0008
% #FIELD_VALUE: 黏度
\fieldvalue{Viscosity}
\\[2pt]
流量范围 &
% #VALUE_ID: LEX-P0013-V0009
% #FIELD_VALUE: 流量范围
\fieldvalue{Zishingle}
\\[2pt]
测试流量 &
% #VALUE_ID: LEX-P0013-V0010
% #FIELD_VALUE: 测试流量
\fieldvalue{3422}
\\[2pt]
测试压力 &
% #VALUE_ID: LEX-P0013-V0011
% #FIELD_VALUE: 测试压力
\fieldvalue{Air}
\\[2pt]
最小压力 &
% #VALUE_ID: LEX-P0013-V0012
% #FIELD_VALUE: 最小压力
\fieldvalue{1000 mbar}
\\[2pt]
最大压力 &
% #VALUE_ID: LEX-P0013-V0013
% #FIELD_VALUE: 最大压力
\fieldvalue{7000 mbar}
\end{tabular}

\vspace{0.5em}

\begin{tabular}{@{}p{0.38\linewidth}p{0.52\linewidth}@{}}
测试压力范围 &
% #VALUE_ID: LEX-P0013-V0014
% #FIELD_VALUE: 测试压力范围
\fieldvalue{5000 mbar}
\\[2pt]
实际压力 &
% #VALUE_ID: LEX-P0013-V0015
% #FIELD_VALUE: 实际压力
\fieldvalue{5000 mbar}
\\[2pt]
测试完成时间 &
% #VALUE_ID: LEX-P0013-V0016
% #FIELD_VALUE: 测试完成时间
\fieldvalue{12/Dec/2025 13:38:20}
\end{tabular}

\vspace{1em}

\noindent
1.\ 签名\hfill
% #VALUE_ID: LEX-P0013-V0017
% #FIELD_VALUE: 签名
% #HANDWRITTEN: [签名，难以辨认]
\fieldvalue{\handwritten{[签名，难以辨认]}}
\\[1em]
\hrule

\vspace{1em}

\begin{center}
\textbf{PALL\quad Paltronic FlowStart IV 流量站结果}
\end{center}

\begin{tabular}{@{}p{0.25\linewidth}p{0.65\linewidth}@{}}
日期 &
% #VALUE_ID: LEX-P0013-V0018
% #FIELD_VALUE: 日期
\fieldvalue{12/Dec/2025 13:27:38}
\\[2pt]
页码 &
% #VALUE_ID: LEX-P0013-V0019
% #FIELD_VALUE: 页码
\fieldvalue{2 of 4}
\end{tabular}

\vspace{1em}

\begin{center}
\setlength{\fboxsep}{4pt}
\fbox{%
\begin{minipage}{0.78\linewidth}
\centering
\textbf{大气流量（mbar/min）}

\vspace{0.5em}
\begin{tabular}{@{}c@{\hspace{1em}}c@{}}
1800 & \rule{0pt}{1.2em}\\
1350 & \\
900 & \\
450 & \\
0 & \rule{0pt}{1.2em}
\end{tabular}

\vspace{-1em}
\noindent\rule{\linewidth}{0.4pt}

\begin{tabular}{@{}p{0.16\linewidth}p{0.16\linewidth}p{0.16\linewidth}p{0.16\linewidth}p{0.16\linewidth}p{0.16\linewidth}@{}}
0 & 1300 & 2700 & 4000 & 5600 & 7000
\end{tabular}

\centering
时间（min）
\end{minipage}%
}
\end{center}

% TODO: The graph trace is visible but cannot be represented accurately without the source figure.
% TODO: Handwritten annotations appear over the report and graph; portions are illegible.
% #VALUE_ID: LEX-P0013-V0020
% #HANDWRITTEN: [蓝色手写签名及日期，难以辨认]
\handwritten{[蓝色手写签名及日期，难以辨认]}
% LEXOID_PAGE_COMPLETED: 13/70

\newpage

\section{Panasonic Flowstart IV 流量报告}

\begin{tabular}{@{}p{0.28\linewidth}p{0.62\linewidth}@{}}
日期 &
% #VALUE_ID: LEX-P0014-V0001
% #FIELD_VALUE: 日期
\fieldvalue{12/20/2025 13:42:19} \\

页码 &
% #VALUE_ID: LEX-P0014-V0002
% #FIELD_VALUE: 页码
\fieldvalue{1 of 2} \\

产品型号 &
% #TODO: 产品型号字符部分难以辨认
% #VALUE_ID: LEX-P0014-V0003
% #FIELD_VALUE: 产品型号
\fieldvalue{FSF40?G?R-22/41} \\

Serial/Pass number &
% #TODO: Serial/Pass number 字符部分难以辨认
% #VALUE_ID: LEX-P0014-V0004
% #FIELD_VALUE: Serial/Pass number
\fieldvalue{288?014?423?} \\

实验员 &
% #VALUE_ID: LEX-P0014-V0005
% #FIELD_VALUE: 实验员
\fieldvalue{王?} \\

测试介质 &
% #VALUE_ID: LEX-P0014-V0006
% #FIELD_VALUE: 测试介质
\fieldvalue{S?ME?F?S-2225} \\

设备 &
% #VALUE_ID: LEX-P0014-V0007
% #FIELD_VALUE: 设备
\fieldvalue{QZ} \\

环境温度 &
% #VALUE_ID: LEX-P0014-V0008
% #FIELD_VALUE: 环境温度
\fieldvalue{Shenyang} \\

测试设备 &
% #VALUE_ID: LEX-P0014-V0009
% #FIELD_VALUE: 测试设备
\fieldvalue{Shenyang} \\

设备编号 &
% #VALUE_ID: LEX-P0014-V0010
% #FIELD_VALUE: 设备编号
\fieldvalue{FS?A?202?5019} \\

压力 &
% #VALUE_ID: LEX-P0014-V0011
% #FIELD_VALUE: 压力
\fieldvalue{5MPa FS-2225} \\

流量 &
% #VALUE_ID: LEX-P0014-V0012
% #FIELD_VALUE: 流量
\fieldvalue{Air} \\

测试压力 &
% #VALUE_ID: LEX-P0014-V0013
% #FIELD_VALUE: 测试压力
\fieldvalue{1500 mbar} \\

最大允许流量 &
% #VALUE_ID: LEX-P0014-V0014
% #FIELD_VALUE: 最大允许流量
\fieldvalue{4500 mbar} \\

最大压力 &
% #VALUE_ID: LEX-P0014-V0015
% #FIELD_VALUE: 最大压力
\fieldvalue{7200 mbar} \\

结论 &
% #VALUE_ID: LEX-P0014-V0016
% #FIELD_VALUE: 结论
\fieldvalue{通过，合格} \\

安全阀设定 &
% #VALUE_ID: LEX-P0014-V0017
% #FIELD_VALUE: 安全阀设定
\fieldvalue{5000 mbar} \\

测试完成时间 &
% #VALUE_ID: LEX-P0014-V0018
% #FIELD_VALUE: 测试完成时间
\fieldvalue{12/20/2025 13:25:22} \\

签名编号 &
% #TODO: 签名编号字符部分难以辨认
% #VALUE_ID: LEX-P0014-V0019
% #FIELD_VALUE: 签名编号
\fieldvalue{A3-5A-62-4?01?912} \\
\end{tabular}

% #VALUE_ID: LEX-P0014-V0020
% #FIELD_VALUE: 签名
% #HANDWRITTEN: 手写签名；字迹难以辨认
\fieldvalue{\handwritten{手写签名}}

% #VALUE_ID: LEX-P0014-V0021
% #FIELD_VALUE: 日期
% #HANDWRITTEN: 2025/12/12
\fieldvalue{\handwritten{2025/12/12}}

\section{Panasonic Flowstart IV 流量报告}

\begin{tabular}{@{}p{0.28\linewidth}p{0.62\linewidth}@{}}
日期 &
% #VALUE_ID: LEX-P0014-V0022
% #FIELD_VALUE: 日期
\fieldvalue{12/20/2025 13:42:19} \\

页码 &
% #VALUE_ID: LEX-P0014-V0023
% #FIELD_VALUE: 页码
\fieldvalue{2 of 2} \\
\end{tabular}

% TODO: 页面下半部分显示一幅流量测试曲线图；未提供可用图像文件，需后续补入图形。
\begin{figure}[ht]
\centering
\fbox{\parbox{0.78\linewidth}{\centering
流量测试曲线图（页面可见部分）\\[2em]
\textit{TODO：根据原图重绘曲线及坐标轴}
}}
\caption{流量测试曲线}
\end{figure}

图中可见坐标标注：

\begin{itemize}
    \item 纵轴：大气流量（mbar/min）；刻度约为 0、450、900、1350、1800。
    \item 横轴：压差（mbar）；刻度约为 0、1400、2800、4200、5600、7000。
\end{itemize}

% #VALUE_ID: LEX-P0014-V0024
% #FIELD_VALUE: 签名
% #HANDWRITTEN: 手写签名；字迹难以辨认
\fieldvalue{\handwritten{手写签名}}

% #VALUE_ID: LEX-P0014-V0025
% #FIELD_VALUE: 日期
% #HANDWRITTEN: 2025/12/12
\fieldvalue{\handwritten{2025/12/12}}

% #VALUE_ID: LEX-P0014-V0026
% #FIELD_VALUE: 时间
% #HANDWRITTEN: 13:35
\fieldvalue{\handwritten{13:35}}
% LEXOID_PAGE_COMPLETED: 14/70

\newpage

\section*{PALL Patronic Flowstar IV 测试结果}

\begin{tabular}{@{}p{0.28\linewidth}p{0.62\linewidth}@{}}
日期 &
% #VALUE_ID: LEX-P0015-V0001
% #FIELD_VALUE: 日期
\fieldvalue{12/Dec/2025 13:56:44}
\\
项目 &
% #VALUE_ID: LEX-P0015-V0002
% #FIELD_VALUE: 项目
\fieldvalue{FFS04CA, 21211}
\\
Serial/Result number &
% #VALUE_ID: LEX-P0015-V0003
% #FIELD_VALUE: Serial/Result number
\fieldvalue{288?04?412?84}
\\
测试序号 &
% #VALUE_ID: LEX-P0015-V0004
% #FIELD_VALUE: 测试序号
\fieldvalue{2}
\\
测试名称 &
% #VALUE_ID: LEX-P0015-V0005
% #FIELD_VALUE: 测试名称
\fieldvalue{SFMP?S-225}
\\
介质 &
% #VALUE_ID: LEX-P0015-V0006
% #FIELD_VALUE: 介质
\fieldvalue{O2}
\\
测试设备 &
% #VALUE_ID: LEX-P0015-V0007
% #FIELD_VALUE: 测试设备
\fieldvalue{Shy?olnou}
\\
测试气体 &
% #VALUE_ID: LEX-P0015-V0008
% #FIELD_VALUE: 测试气体
\fieldvalue{SF6}
\\
产品名称 &
% #VALUE_ID: LEX-P0015-V0009
% #FIELD_VALUE: 产品名称
\fieldvalue{A?4?S?2025?019}
\\
产品型号 &
% #VALUE_ID: LEX-P0015-V0010
% #FIELD_VALUE: 产品型号
\fieldvalue{434?}
\\
测试流量 &
% #VALUE_ID: LEX-P0015-V0011
% #FIELD_VALUE: 测试流量
\fieldvalue{Air}
\\
测试温度 &
% #VALUE_ID: LEX-P0015-V0012
% #FIELD_VALUE: 测试温度
\fieldvalue{20?}
\\
最大压力 &
% #VALUE_ID: LEX-P0015-V0013
% #FIELD_VALUE: 最大压力
\fieldvalue{7000 mbar}
\\
最小压力 &
% #VALUE_ID: LEX-P0015-V0014
% #FIELD_VALUE: 最小压力
\fieldvalue{4500 mbar}
\\
压力变化 &
% #VALUE_ID: LEX-P0015-V0015
% #FIELD_VALUE: 压力变化
\fieldvalue{1000}
\end{tabular}

\vspace{1em}

\textbf{结果}

\begin{tabular}{@{}p{0.42\linewidth}p{0.48\linewidth}@{}}
流量在目标压力下 &
% #VALUE_ID: LEX-P0015-V0016
% #FIELD_VALUE: 流量在目标压力下
\fieldvalue{5000 mbar}
\\
测试完成时间 &
% #VALUE_ID: LEX-P0015-V0017
% #FIELD_VALUE: 测试完成时间
\fieldvalue{12/Dec/2025 14:07:12}
\\
1. 签名 &
% #VALUE_ID: LEX-P0015-V0018
% #FIELD_VALUE: 1. 签名
% #HANDWRITTEN: 不合格?（字迹不清）
\fieldvalue{\handwritten{不合格?}}
\end{tabular}

\vspace{1em}

% #VALUE_ID: LEX-P0015-V0019
% #HANDWRITTEN: 袁瑾
\handwritten{袁瑾}
\qquad
% #VALUE_ID: LEX-P0015-V0020
% #HANDWRITTEN: 2025/12/2
\handwritten{2025/12/2}

\section*{PALL Patronic Flowstar IV 测试结果}

\begin{tabular}{@{}p{0.28\linewidth}p{0.62\linewidth}@{}}
日期 &
% #VALUE_ID: LEX-P0015-V0021
% #FIELD_VALUE: 日期
\fieldvalue{12/Dec/2025 13:56:44}
\\
页码 &
% #VALUE_ID: LEX-P0015-V0022
% #FIELD_VALUE: 页码
\fieldvalue{2 of 2}
\end{tabular}

\vspace{1em}

\begin{center}
\begin{tikzpicture}
\draw (0,0) rectangle (0.82\linewidth,5cm);
\draw[->] (0,0) -- (0,5.4cm)
  node[above] {大气流量 (mbar/min)};
\draw[->] (0,0) -- (0.86\linewidth,0)
  node[right] {压力 (mbar)};

\node[left] at (0,0) {0};
\node[left] at (0,1.35cm) {450};
\node[left] at (0,2.7cm) {900};
\node[left] at (0,4.05cm) {1350};
\node[left] at (0,5cm) {1800};

\node[below] at (0,0) {0};
\node[below] at (0.16\linewidth,0) {1400};
\node[below] at (0.33\linewidth,0) {2800};
\node[below] at (0.50\linewidth,0) {4200};
\node[below] at (0.66\linewidth,0) {5600};
\node[below] at (0.82\linewidth,0) {7000};

\draw[gray] (0.57\linewidth,5cm) -- (0.57\linewidth,0.8cm);
\draw[gray] plot[smooth] coordinates {
  (0.57\linewidth,0.8cm)
  (0.585\linewidth,0.48cm)
  (0.61\linewidth,0.35cm)
  (0.64\linewidth,0.27cm)
  (0.67\linewidth,0.20cm)
};
\end{tikzpicture}
\end{center}

% #VALUE_ID: LEX-P0015-V0023
% #HANDWRITTEN: 袁瑾
\handwritten{袁瑾}
\qquad
% #VALUE_ID: LEX-P0015-V0024
% #HANDWRITTEN: 2025/12/2
\handwritten{2025/12/2}
\qquad
% #VALUE_ID: LEX-P0015-V0025
% #HANDWRITTEN: 16/35
\handwritten{16/35}

% TODO: Several printed field values in the rotated source image are partially illegible; transcriptions marked with question marks require review.
% LEXOID_PAGE_COMPLETED: 15/70

\newpage

\section{Palltronic Flowstar IV 测试结果}

\begin{tabular}{@{}p{0.28\linewidth}p{0.62\linewidth}@{}}
日期 &
% #VALUE_ID: LEX-P0016-V0001
% #FIELD_VALUE: 日期
\fieldvalue{12/20/2025 13:40:47}
\\
页码 &
% #VALUE_ID: LEX-P0016-V0002
% #FIELD_VALUE: 页码
\fieldvalue{1 of 2}
\\
产品 &
% #VALUE_ID: LEX-P0016-V0003
% #FIELD_VALUE: 产品
\fieldvalue{FFS40G2R-2171}
\\
Serial/part number &
% #VALUE_ID: LEX-P0016-V0004
% #FIELD_VALUE: Serial/part number
\fieldvalue{238804742825}
\\
测试程序 &
% #VALUE_ID: LEX-P0016-V0005
% #FIELD_VALUE: 测试程序
\fieldvalue{SMPA-ES225}
\\
测试介质 &
% #VALUE_ID: LEX-P0016-V0006
% #FIELD_VALUE: 测试介质
\fieldvalue{N2}
\\
测试气体 &
% #VALUE_ID: LEX-P0016-V0007
% #FIELD_VALUE: 测试气体
\fieldvalue{C}
\\
测试模块 &
% #VALUE_ID: LEX-P0016-V0008
% #FIELD_VALUE: 测试模块
\fieldvalue{SMV?}
\\
测试标准 &
% #VALUE_ID: LEX-P0016-V0009
% #FIELD_VALUE: 测试标准
\fieldvalue{N2?}
\\
程序名称 &
% #VALUE_ID: LEX-P0016-V0010
% #FIELD_VALUE: 程序名称
\fieldvalue{Z?}
\\
程序编号 &
% #VALUE_ID: LEX-P0016-V0011
% #FIELD_VALUE: 程序编号
\fieldvalue{434?}
\\
设备编号 &
% #VALUE_ID: LEX-P0016-V0012
% #FIELD_VALUE: 设备编号
\fieldvalue{SMPA-ES225}
\\
测试压力 &
% #VALUE_ID: LEX-P0016-V0013
% #FIELD_VALUE: 测试压力
\fieldvalue{4950\ mbar}
\\
测试开始压力 &
% #VALUE_ID: LEX-P0016-V0014
% #FIELD_VALUE: 测试开始压力
\fieldvalue{-2000\ mbar}
\\
测试结束压力 &
% #VALUE_ID: LEX-P0016-V0015
% #FIELD_VALUE: 测试结束压力
\fieldvalue{-4500\ mbar}
\\
\end{tabular}

\vspace{1em}

\begin{tabular}{@{}p{0.28\linewidth}p{0.62\linewidth}@{}}
结果 &
% #VALUE_ID: LEX-P0016-V0016
% #FIELD_VALUE: 结果
\fieldvalue{通过，在标准范围内：4950\ mbar}
\\
测试完成时间 &
% #VALUE_ID: LEX-P0016-V0017
% #FIELD_VALUE: 测试完成时间
\fieldvalue{12/20/2025 14:21:08}
\\
\end{tabular}

\vspace{1em}

\textbf{注释}

\begin{enumerate}
\item
% #VALUE_ID: LEX-P0016-V0018
% #FIELD_VALUE: 签名
% #TODO #HANDWRITTEN: A17-25-54-5?90?; handwriting is partially illegible
\fieldvalue{\handwritten{A17-25-54-5?90?}}
\end{enumerate}

% #VALUE_ID: LEX-P0016-V0019
% #FIELD_VALUE: 签名及日期
% #TODO #HANDWRITTEN: [签名] 2025.12.12; Chinese signature is illegible
\fieldvalue{\handwritten{[签名]\quad 2025.12.12}}

\section{测试曲线}

\begin{tabular}{@{}p{0.28\linewidth}p{0.62\linewidth}@{}}
日期 &
% #VALUE_ID: LEX-P0016-V0020
% #FIELD_VALUE: 日期
\fieldvalue{12/20/2025 13:40:47}
\\
页码 &
% #VALUE_ID: LEX-P0016-V0021
% #FIELD_VALUE: 页码
\fieldvalue{2 of 2}
\\
\end{tabular}

\vspace{1em}

% TODO: Reproduce the visible test curve. The curve is plotted against time
% (approximately 0--7200 s) and flow rate (0--1800 mbar/min).
\begin{figure}[ht]
\centering
\fbox{\parbox{0.82\linewidth}{\centering
Test curve plot\\[1ex]
纵轴：升压速率（mbar/min），刻度 0、450、900、1350、1800\\
横轴：时间（s），刻度约为 0、1000、2000、4200、5600、7200
}}
\end{figure}

% #VALUE_ID: LEX-P0016-V0022
% #FIELD_VALUE: 签名及日期
% #TODO #HANDWRITTEN: [签名] 2025/12/12; Chinese signature is illegible
\fieldvalue{\handwritten{[签名]\quad 2025/12/12}}

% #VALUE_ID: LEX-P0016-V0023
% #FIELD_VALUE: 页码标记
% #HANDWRITTEN: 17/35
\fieldvalue{\handwritten{17/35}}
% LEXOID_PAGE_COMPLETED: 16/70

\newpage

\section{Palltronic Flowstar IV 测试站结果}

\begin{tabular}{@{}p{0.24\linewidth}p{0.68\linewidth}@{}}
日期 &
% #VALUE_ID: LEX-P0017-V0001
% #FIELD_VALUE: 日期
\fieldvalue{12/Dec/2025 15:24:33}
\\
页码 &
% #VALUE_ID: LEX-P0017-V0002
% #FIELD_VALUE: 页码
\fieldvalue{1 of 2}
\\
设备号 &
% #VALUE_ID: LEX-P0017-V0003
% #FIELD_VALUE: 设备号
\fieldvalue{FTS04G/R-2/2/1}
\\
序列号 &
% #VALUE_ID: LEX-P0017-V0004
% #FIELD_VALUE: 序列号
\fieldvalue{230807442626}
\\
功能 &
% #VALUE_ID: LEX-P0017-V0005
% #FIELD_VALUE: 功能
\fieldvalue{流量}
\\
型号 &
% #VALUE_ID: LEX-P0017-V0006
% #FIELD_VALUE: 型号
\fieldvalue{SFM?PS-22ZS}
\\
压力单位 &
% #VALUE_ID: LEX-P0017-V0007
% #FIELD_VALUE: 压力单位
\fieldvalue{Q2}
\\
生产区段 &
% #VALUE_ID: LEX-P0017-V0008
% #FIELD_VALUE: 生产区段
\fieldvalue{C}
\\
生产商 &
% #VALUE_ID: LEX-P0017-V0009
% #FIELD_VALUE: 生产商
\fieldvalue{Shinyoung}
\\
产品名称 &
% #VALUE_ID: LEX-P0017-V0010
% #FIELD_VALUE: 产品名称
\fieldvalue{Zhu?}
\\
产品编号 &
% #VALUE_ID: LEX-P0017-V0011
% #FIELD_VALUE: 产品编号
\fieldvalue{A?322?0252019}
\\
过滤器编号 &
% #VALUE_ID: LEX-P0017-V0012
% #FIELD_VALUE: 过滤器编号
\fieldvalue{SFM?PS-22ZS}
\\
测试气体 &
% #VALUE_ID: LEX-P0017-V0013
% #FIELD_VALUE: 测试气体
\fieldvalue{Air}
\\
测试类型 &
% #VALUE_ID: LEX-P0017-V0014
% #FIELD_VALUE: 测试类型
\fieldvalue{泄漏测试}
\\
最低测试压力 &
% #VALUE_ID: LEX-P0017-V0015
% #FIELD_VALUE: 最低测试压力
\fieldvalue{1000}
\\
最大测试压力 &
% #VALUE_ID: LEX-P0017-V0016
% #FIELD_VALUE: 最大测试压力
\fieldvalue{4500 mbar}
\end{tabular}

\subsection{结果}

\begin{tabular}{@{}p{0.32\linewidth}p{0.60\linewidth}@{}}
结果 &
% #VALUE_ID: LEX-P0017-V0017
% #FIELD_VALUE: 结果
\fieldvalue{通过}
\\
最低测试压力 &
% #VALUE_ID: LEX-P0017-V0018
% #FIELD_VALUE: 最低测试压力
\fieldvalue{4950 mbar}
\\
测试日期 &
% #VALUE_ID: LEX-P0017-V0019
% #FIELD_VALUE: 测试日期
\fieldvalue{12/Dec/2025 15:35:15}
\\
签字 &
% #VALUE_ID: LEX-P0017-V0020
% #FIELD_VALUE: 签字
% #TODO #HANDWRITTEN: illegible signature; blue handwritten signature
\fieldvalue{\handwritten{[签名]}}
\end{tabular}

% #VALUE_ID: LEX-P0017-V0021
% #HANDWRITTEN: 康校 2026/1/12
\handwritten{康校 2026/1/12}

\section{Palltronic Flowstar IV 测试站结果}

\begin{tabular}{@{}p{0.24\linewidth}p{0.68\linewidth}@{}}
日期 &
% #VALUE_ID: LEX-P0017-V0022
% #FIELD_VALUE: 日期
\fieldvalue{12/Dec/2025 15:24:33}
\\
页码 &
% #VALUE_ID: LEX-P0017-V0023
% #FIELD_VALUE: 页码
\fieldvalue{1 of 2}
\end{tabular}

\subsection{流量曲线}

\begin{center}
\begin{tabular}{@{}c@{}}
\begin{picture}(300,180)
\put(25,20){\line(1,0){245}}
\put(25,20){\line(0,1){135}}
\put(25,155){\line(1,0){245}}
\put(270,20){\line(0,1){135}}
\put(25,20){\line(1,1){0}}
\put(25,20){\makebox(0,0)[r]{0}}
\put(65,20){\makebox(0,-12)[c]{1400}}
\put(115,20){\makebox(0,-12)[c]{2800}}
\put(165,20){\makebox(0,-12)[c]{4200}}
\put(215,20){\makebox(0,-12)[c]{5600}}
\put(270,20){\makebox(0,-12)[c]{7000}}
\put(25,155){\makebox(-8,0)[r]{1600}}
\put(25,110){\makebox(-8,0)[r]{1350}}
\put(25,65){\makebox(-8,0)[r]{900}}
\put(25,20){\makebox(-8,0)[r]{450}}
\put(130,0){\makebox(0,0)[c]{大气流量 (mbar/min)}}
\put(0,88){\rotatebox{90}{\makebox(0,0)[c]{压力 (mbar)}}}
\qbezier(25,20)(40,20)(52,22)
\qbezier(52,22)(60,25)(65,42)
\qbezier(65,42)(70,70)(82,85)
\qbezier(82,85)(90,95)(105,96)
\qbezier(105,96)(150,96)(270,96)
\end{picture}
\end{tabular}
\end{center}

% #VALUE_ID: LEX-P0017-V0024
% #HANDWRITTEN: 康校 2026/1/12
\handwritten{康校 2026/1/12}

% #VALUE_ID: LEX-P0017-V0025
% #HANDWRITTEN: 18/35
\handwritten{18/35}
% LEXOID_PAGE_COMPLETED: 17/70

\newpage

\rotatebox{90}{%
\begin{minipage}{0.92\textheight}
\small

\begin{center}
\textbf{Palltronic Flowstar IV 测试结果}
\end{center}

\begin{tabular}{@{}p{0.25\linewidth}p{0.68\linewidth}@{}}
日期 &
% #VALUE_ID: LEX-P0018-V0001
% #FIELD_VALUE: 日期
\fieldvalue{12/Dec/2025 13:49:01}
\\
设备 &
% #VALUE_ID: LEX-P0018-V0002
% #FIELD_VALUE: 设备
\fieldvalue{FFS04LC (2 24?)}
\\
Serial/Serial number &
% #VALUE_ID: LEX-P0018-V0003
% #FIELD_VALUE: Serial/Serial number
\fieldvalue{238804244?27}
\\
测试类型 &
% #VALUE_ID: LEX-P0018-V0004
% #FIELD_VALUE: 测试类型
\fieldvalue{SFMPES-2225}
\\
测试程序 &
% #VALUE_ID: LEX-P0018-V0005
% #FIELD_VALUE: 测试程序
\fieldvalue{SFMPES-2225}
\\
测试气体 &
% #VALUE_ID: LEX-P0018-V0006
% #FIELD_VALUE: 测试气体
\fieldvalue{O2}
\\
测试压力单位 &
% #VALUE_ID: LEX-P0018-V0007
% #FIELD_VALUE: 测试压力单位
\fieldvalue{C}
\\
测试方法 &
% #VALUE_ID: LEX-P0018-V0008
% #FIELD_VALUE: 测试方法
\fieldvalue{Single}
\\
测试持续时间 &
% #VALUE_ID: LEX-P0018-V0009
% #FIELD_VALUE: 测试持续时间
\fieldvalue{2?}
\\
泄漏流量 &
% #VALUE_ID: LEX-P0018-V0010
% #FIELD_VALUE: 泄漏流量
\fieldvalue{1.000}
\\
最小压力 &
% #VALUE_ID: LEX-P0018-V0011
% #FIELD_VALUE: 最小压力
\fieldvalue{4500 mbar}
\\
最大压力 &
% #VALUE_ID: LEX-P0018-V0012
% #FIELD_VALUE: 最大压力
\fieldvalue{7000 mbar}
\end{tabular}

\vspace{0.5em}
\hrule
\vspace{0.5em}

\begin{tabular}{@{}p{0.25\linewidth}p{0.68\linewidth}@{}}
结果 &
% #VALUE_ID: LEX-P0018-V0013
% #FIELD_VALUE: 结果
\fieldvalue{测试在规定范围内：5050 mbar}
\\
测试完成时间 &
% #VALUE_ID: LEX-P0018-V0014
% #FIELD_VALUE: 测试完成时间
\fieldvalue{12/Dec/2025 14:49:54}
\end{tabular}

\vspace{1em}

% #VALUE_ID: LEX-P0018-V0015
% #HANDWRITTEN: 林? 2025/12/12
\handwritten{林? 2025/12/12}

\vspace{1em}
\hrule
\vspace{0.5em}

\begin{center}
\textbf{Palltronic Flowstar IV 测试结果}
\end{center}

\begin{tabular}{@{}p{0.25\linewidth}p{0.68\linewidth}@{}}
日期 &
% #VALUE_ID: LEX-P0018-V0016
% #FIELD_VALUE: 日期
\fieldvalue{12/Dec/2025 13:49:01}
\\
时间 &
% #VALUE_ID: LEX-P0018-V0017
% #FIELD_VALUE: 时间
\fieldvalue{2:012}
\end{tabular}

\vspace{1em}

\begin{center}
\begin{picture}(300,150)
\put(0,0){\framebox(300,150){}}
\put(8,138){\makebox(0,0)[l]{大气压差 (mbar/min)}}
\put(8,8){\makebox(0,0)[l]{时间 (msec)}}
\put(18,130){\makebox(0,0)[r]{0}}
\put(54,130){\makebox(0,0)[r]{450}}
\put(105,130){\makebox(0,0)[r]{900}}
\put(155,130){\makebox(0,0)[r]{1350}}
\put(205,130){\makebox(0,0)[r]{1800}}
\put(45,-3){\makebox(0,0)[c]{1400}}
\put(105,-3){\makebox(0,0)[c]{2800}}
\put(165,-3){\makebox(0,0)[c]{4200}}
\put(225,-3){\makebox(0,0)[c]{5600}}
\put(285,-3){\makebox(0,0)[r]{7000}}
\put(120,82){\line(1,0){0}}
\put(0,35){\line(1,0){300}}
\end{picture}
\end{center}

% TODO: The plotted curve is visible on the page; reproduce graph geometry if required.

% #VALUE_ID: LEX-P0018-V0018
% #HANDWRITTEN: 林? 2025/12/12
\handwritten{林? 2025/12/12}

% #VALUE_ID: LEX-P0018-V0019
% #HANDWRITTEN: 19/3?
\handwritten{19/3?}

\end{minipage}%
}
% LEXOID_PAGE_COMPLETED: 18/70

\newpage

\section*{Palltronic Flowstar IV 测试结果}

\begin{tabular}{@{}p{0.28\linewidth}p{0.62\linewidth}@{}}
日期 &
% #VALUE_ID: LEX-P0019-V0001
% #FIELD_VALUE: 日期
\fieldvalue{12/Dec/2025 15:40:41}
\\
项目 &
% #VALUE_ID: LEX-P0019-V0002
% #FIELD_VALUE: 项目
\fieldvalue{1.0}
\\
页码 &
% #VALUE_ID: LEX-P0019-V0003
% #FIELD_VALUE: 页码
\fieldvalue{1 of 2}
\\
设备 &
% #VALUE_ID: LEX-P0019-V0004
% #FIELD_VALUE: 设备
\fieldvalue{Palltronic Flowstar IV}
\\
序列号 &
% #VALUE_ID: LEX-P0019-V0005
% #FIELD_VALUE: 序列号
\fieldvalue{2288742628}
\\
功能 &
% #VALUE_ID: LEX-P0019-V0006
% #FIELD_VALUE: 功能
\fieldvalue{泡点}
\\
测试程序 &
% #VALUE_ID: LEX-P0019-V0007
% #FIELD_VALUE: 测试程序
\fieldvalue{SM-PFS-2225}
\\
测试气体 &
% #VALUE_ID: LEX-P0019-V0008
% #FIELD_VALUE: 测试气体
\fieldvalue{Air}
\\
过滤器型号 &
% #VALUE_ID: LEX-P0019-V0009
% #FIELD_VALUE: 过滤器型号
\fieldvalue{QZ}
\\
最大压力 &
% #VALUE_ID: LEX-P0019-V0010
% #FIELD_VALUE: 最大压力
\fieldvalue{4500 mbar}
\\
最小压力 &
% #VALUE_ID: LEX-P0019-V0011
% #FIELD_VALUE: 最小压力
\fieldvalue{7000 mbar}
\\
结果 &
% #VALUE_ID: LEX-P0019-V0012
% #FIELD_VALUE: 结果
\fieldvalue{通过}
\\
测试滤点 &
% #VALUE_ID: LEX-P0019-V0013
% #FIELD_VALUE: 测试滤点
\fieldvalue{5050 mbar}
\\
测试完成时间 &
% #VALUE_ID: LEX-P0019-V0014
% #FIELD_VALUE: 测试完成时间
\fieldvalue{12/Dec/2025 15:40:44}
\\
签名 &
% #VALUE_ID: LEX-P0019-V0015
% #FIELD_VALUE: 签名
% #TODO #HANDWRITTEN: [签名]; handwriting is unclear
\fieldvalue{\handwritten{[签名]}}
\\
日期（签名旁） &
% #VALUE_ID: LEX-P0019-V0016
% #FIELD_VALUE: 日期（签名旁）
% #HANDWRITTEN: 2025.12.12
\fieldvalue{\handwritten{2025.12.12}}
\end{tabular}

\vspace{1em}

% #TODO: Several printed field labels and values in the source image are partially obscured by the page rotation and low resolution.

\section*{Palltronic Flowstar IV 测试结果}

\begin{tabular}{@{}p{0.28\linewidth}p{0.62\linewidth}@{}}
日期 &
% #VALUE_ID: LEX-P0019-V0017
% #FIELD_VALUE: 日期
\fieldvalue{12/Dec/2025 15:45:40}
\\
页码 &
% #VALUE_ID: LEX-P0019-V0018
% #FIELD_VALUE: 页码
\fieldvalue{2 of 2}
\end{tabular}

\vspace{1em}

\begin{center}
\begin{tabular}{@{}c@{}}
\text{气体流量 (mbar/min)}\\[-0.25em]
\begin{picture}(280,150)
\put(25,10){\line(0,1){125}}
\put(25,10){\line(1,0){255}}
\put(25,135){\line(1,0){0}}
\put(25,135){\makebox(0,0)[r]{1800}}
\put(25,104){\makebox(0,0)[r]{1350}}
\put(25,73){\makebox(0,0)[r]{900}}
\put(25,42){\makebox(0,0)[r]{450}}
\put(25,10){\makebox(0,0)[r]{0}}
\put(25,10){\line(1,0){255}}
\put(25,10){\makebox(0,-14)[c]{1400}}
\put(88,10){\makebox(0,-14)[c]{2800}}
\put(152,10){\makebox(0,-14)[c]{4200}}
\put(216,10){\makebox(0,-14)[c]{5600}}
\put(280,10){\makebox(0,-14)[c]{7000}}
\put(152,10){\line(0,1){72}}
\put(152,82){\line(1,0){15}}
\put(167,82){\line(0,-1){13}}
\put(167,69){\line(1,0){8}}
\put(175,69){\line(0,-1){6}}
\put(183,63){\line(1,0){8}}
\put(191,63){\line(0,-1){4}}
\end{picture}\\[-0.25em]
\text{时间 (s)}
\end{tabular}
\end{center}

% #VALUE_ID: LEX-P0019-V0019
% #FIELD_VALUE: 图表签名
% #HANDWRITTEN: [签名]
\noindent 图表签名：\fieldvalue{\handwritten{[签名]}}

\hfill
% #VALUE_ID: LEX-P0019-V0020
% #FIELD_VALUE: 图表日期
% #HANDWRITTEN: 2025.12.12
\fieldvalue{\handwritten{2025.12.12}}

\hfill
% #VALUE_ID: LEX-P0019-V0021
% #FIELD_VALUE: 右侧日期
% #TODO #HANDWRITTEN: 2025; final digit is unclear
\fieldvalue{\handwritten{2025}}
% LEXOID_PAGE_COMPLETED: 19/70

\newpage

\section{Panasonic Flowstar 测试报告}

\begin{tabular}{@{}p{0.28\linewidth}p{0.62\linewidth}@{}}
日期 &
% #VALUE_ID: LEX-P0020-V0001
% #FIELD_VALUE: 日期
\fieldvalue{2020/02/05 10:06:34}
\\
页码 &
% #VALUE_ID: LEX-P0020-V0002
% #FIELD_VALUE: 页码
\fieldvalue{1 of 2}
\\
项目名称 &
% #VALUE_ID: LEX-P0020-V0003
% #FIELD_VALUE: 项目名称
% TODO: The project name is partially unclear in the scan.
\fieldvalue{SFM?C?}
\\
Serial/Result number &
% #VALUE_ID: LEX-P0020-V0004
% #FIELD_VALUE: Serial/Result number
% TODO: The serial/result number is partially illegible.
\fieldvalue{288?0?4-4261}
\\
客户 &
% #VALUE_ID: LEX-P0020-V0005
% #FIELD_VALUE: 客户
% TODO: Customer name is unclear in the scan.
\fieldvalue{?}
\\
测试程序 &
% #VALUE_ID: LEX-P0020-V0006
% #FIELD_VALUE: 测试程序
\fieldvalue{SFM?ES-2225}
\\
样品名 &
% #VALUE_ID: LEX-P0020-V0007
% #FIELD_VALUE: 样品名
\fieldvalue{CZ}
\\
测试人员 &
% #VALUE_ID: LEX-P0020-V0008
% #FIELD_VALUE: 测试人员
\fieldvalue{Siyongzhou}
\\
测试气体 &
% #VALUE_ID: LEX-P0020-V0009
% #FIELD_VALUE: 测试气体
% TODO: The value is unclear in the scan.
\fieldvalue{?}
\\
测试温度 &
% #VALUE_ID: LEX-P0020-V0010
% #FIELD_VALUE: 测试温度
% TODO: The value is unclear in the scan.
\fieldvalue{?}
\\
测试压力 &
% #VALUE_ID: LEX-P0020-V0011
% #FIELD_VALUE: 测试压力
\fieldvalue{4000}
\\
最大测试压力 &
% #VALUE_ID: LEX-P0020-V0012
% #FIELD_VALUE: 最大测试压力
\fieldvalue{4500 mbar}
\\
最大压力 &
% #VALUE_ID: LEX-P0020-V0013
% #FIELD_VALUE: 最大压力
\fieldvalue{7000 mbar}
\end{tabular}

\vspace{1em}

% #VALUE_ID: LEX-P0020-V0014
% #HANDWRITTEN: 李?霞
\handwritten{李?霞}
\qquad
% #VALUE_ID: LEX-P0020-V0015
% #HANDWRITTEN: 2025/12/2
\handwritten{2025/12/2}

\section{测试结果}

\begin{tabular}{@{}p{0.28\linewidth}p{0.62\linewidth}@{}}
日期 &
% #VALUE_ID: LEX-P0020-V0016
% #FIELD_VALUE: 日期
\fieldvalue{2020/02/05 10:06:34}
\\
页码 &
% #VALUE_ID: LEX-P0020-V0017
% #FIELD_VALUE: 页码
\fieldvalue{2 of 2}
\end{tabular}

\vspace{1em}

\begin{figure}[h]
\centering
\setlength{\fboxsep}{8pt}
\fbox{\begin{minipage}{0.82\linewidth}
\centering
\textbf{大气体积 (mbar/min)}

\vspace{0.5em}
\begin{tabular}{c|ccccc}
1800 & & & & & \\
1350 & & & & & \\
900  & & & & & \\
450  & & & & & \\
0    &\multicolumn{5}{c}{\rule{0pt}{1.2em}}\\[-0.3em]
\hline
& 1400 & 2800 & 4200 & 5600 & 7000
\end{tabular}

\vspace{0.3em}
\textit{大气压 (mbar)}
\end{minipage}}
\caption{测试曲线}
\end{figure}

% TODO: The plotted curve is visible on the page but is not reproduced precisely here.

% #VALUE_ID: LEX-P0020-V0018
% #HANDWRITTEN: ????????? 2025/12/2
\handwritten{????????? 2025/12/2}

% #VALUE_ID: LEX-P0020-V0019
% #FIELD_VALUE: 姓名
% #HANDWRITTEN: 祝?文
\fieldvalue{\handwritten{祝?文}}

% #VALUE_ID: LEX-P0020-V0020
% #HANDWRITTEN: 2/3
\handwritten{2/3}
% LEXOID_PAGE_COMPLETED: 20/70

\newpage

\section*{Palltronic Flowstar IV 测试结果}

\begin{tabular}{@{}p{0.24\linewidth}p{0.68\linewidth}@{}}
日期 &
% #VALUE_ID: LEX-P0021-V0001
% #FIELD_VALUE: 日期
\fieldvalue{12/Dec/2025 10:21:19}
\\
页码 &
% #VALUE_ID: LEX-P0021-V0002
% #FIELD_VALUE: 页码
\fieldvalue{1 of 2}
\\
测试号 &
% #VALUE_ID: LEX-P0021-V0003
% #FIELD_VALUE: 测试号
\fieldvalue{2588042515}
\\
功能 &
% #VALUE_ID: LEX-P0021-V0004
% #FIELD_VALUE: 功能
\fieldvalue{Air}
\\
测试结果 &
% #VALUE_ID: LEX-P0021-V0005
% #FIELD_VALUE: 测试结果
\fieldvalue{Passed}
\\
测试程序 &
% #VALUE_ID: LEX-P0021-V0006
% #FIELD_VALUE: 测试程序
\fieldvalue{FMSPES-225}
\\
测试区域 &
% #VALUE_ID: LEX-P0021-V0007
% #FIELD_VALUE: 测试区域
\fieldvalue{C2}
\\
生产线名称 &
% #VALUE_ID: LEX-P0021-V0008
% #FIELD_VALUE: 生产线名称
\fieldvalue{Shijingshan}
\\
产品名称 &
% #VALUE_ID: LEX-P0021-V0009
% #FIELD_VALUE: 产品名称
\fieldvalue{Zhu...}
\\
产品编号 &
% #VALUE_ID: LEX-P0021-V0010
% #FIELD_VALUE: 产品编号
\fieldvalue{Z...}
\\
程序名称 &
% #VALUE_ID: LEX-P0021-V0011
% #FIELD_VALUE: 程序名称
\fieldvalue{SMF...}
\\
程序版本号 &
% #VALUE_ID: LEX-P0021-V0012
% #FIELD_VALUE: 程序版本号
\fieldvalue{SMF...}
\\
测试气体 &
% #VALUE_ID: LEX-P0021-V0013
% #FIELD_VALUE: 测试气体
\fieldvalue{Air}
\\
测试压力 &
% #VALUE_ID: LEX-P0021-V0014
% #FIELD_VALUE: 测试压力
\fieldvalue{1000 mbar}
\\
最大压力 &
% #VALUE_ID: LEX-P0021-V0015
% #FIELD_VALUE: 最大压力
\fieldvalue{4500 mbar}
\\
最小压力 &
% #VALUE_ID: LEX-P0021-V0016
% #FIELD_VALUE: 最小压力
\fieldvalue{7000 mbar}
\end{tabular}

\bigskip

\noindent\textbf{结果}

\begin{tabular}{@{}p{0.28\linewidth}p{0.64\linewidth}@{}}
泄漏测试结果 &
% #VALUE_ID: LEX-P0021-V0017
% #FIELD_VALUE: 泄漏测试结果
\fieldvalue{通过}
\\
测试压力 &
% #VALUE_ID: LEX-P0021-V0018
% #FIELD_VALUE: 测试压力
\fieldvalue{5000 mbar}
\\
测试开始时间 &
% #VALUE_ID: LEX-P0021-V0019
% #FIELD_VALUE: 测试开始时间
\fieldvalue{12/Dec/2025 10:35:31}
\\
\end{tabular}

\bigskip

\noindent\textbf{注释}

\begin{tabular}{@{}p{0.28\linewidth}p{0.64\linewidth}@{}}
注释内容 &
% #VALUE_ID: LEX-P0021-V0020
% #FIELD_VALUE: 注释内容
\fieldvalue{B1-B1-B4-...}
\end{tabular}

\bigskip

% #VALUE_ID: LEX-P0021-V0021
% #FIELD_VALUE: 签字
% #HANDWRITTEN: 2025/12/12
\noindent 签字：\fieldvalue{\handwritten{2025/12/12}}

\bigskip

\section*{Palltronic Flowstar IV 测试结果}

\begin{tabular}{@{}p{0.24\linewidth}p{0.68\linewidth}@{}}
日期 &
% #VALUE_ID: LEX-P0021-V0022
% #FIELD_VALUE: 日期
\fieldvalue{12/Dec/2025 10:21:19}
\\
页码 &
% #VALUE_ID: LEX-P0021-V0023
% #FIELD_VALUE: 页码
\fieldvalue{2 of 2}
\end{tabular}

\bigskip

\noindent
\begin{picture}(300,170)
\put(35,15){\line(1,0){245}}
\put(35,15){\line(0,1){140}}
\put(35,155){\line(1,0){245}}
\put(280,15){\line(0,1){140}}
\put(32,12){\makebox(0,0)[r]{0}}
\put(32,85){\makebox(0,0)[r]{450}}
\put(32,120){\makebox(0,0)[r]{900}}
\put(32,140){\makebox(0,0)[r]{1350}}
\put(32,155){\makebox(0,0)[r]{1800}}
\put(70,5){\makebox(0,0){1400}}
\put(115,5){\makebox(0,0){2800}}
\put(165,5){\makebox(0,0){4200}}
\put(225,5){\makebox(0,0){5600}}
\put(270,5){\makebox(0,0){7000}}
\put(145,0){\makebox(0,0){时间 (min)}}
\put(5,105){\makebox(0,0)[l]{大气压强 (mbar/min)}}
\put(205,82){\line(0,-1){58}}
\put(205,24){\line(1,0){12}}
\end{picture}

\bigskip

% #VALUE_ID: LEX-P0021-V0024
% #FIELD_VALUE: 签字
% #HANDWRITTEN: 2025/12/12
\noindent 签字：\fieldvalue{\handwritten{2025/12/12}}

\bigskip

% #VALUE_ID: LEX-P0021-V0025
% #HANDWRITTEN: 2435
\noindent\textcolor{blue}{\handwritten{2435}}
% LEXOID_PAGE_COMPLETED: 21/70

\newpage

\section{Palltronic Flowstar IV 测流结果}

\begin{tabular}{@{}p{0.27\linewidth}p{0.65\linewidth}@{}}
日期 &
% #VALUE_ID: LEX-P0022-V0001
% #FIELD_VALUE: 日期
\fieldvalue{12/Dec/2025 10:35:49}
\\
序号 &
% #VALUE_ID: LEX-P0022-V0002
% #FIELD_VALUE: 序号
% TODO: 部分字符不清晰。
\fieldvalue{FFSW?4? (2.21)}
\\
Serial/results number &
% #VALUE_ID: LEX-P0022-V0003
% #FIELD_VALUE: Serial/results number
% TODO: 字符辨识不完全确定。
\fieldvalue{288?424?216}
\\
测试程序 &
% #VALUE_ID: LEX-P0022-V0004
% #FIELD_VALUE: 测试程序
% TODO: 字符辨识不完全确定。
\fieldvalue{FFSW?S-225}
\\
测试气体 &
% #VALUE_ID: LEX-P0022-V0005
% #FIELD_VALUE: 测试气体
\fieldvalue{N2}
\\
气体温度 &
% #VALUE_ID: LEX-P0022-V0006
% #FIELD_VALUE: 气体温度
\fieldvalue{20\,${}^{\circ}$C}
\\
气体压力 &
% #VALUE_ID: LEX-P0022-V0007
% #FIELD_VALUE: 气体压力
\fieldvalue{2}
\\
测试设备 &
% #VALUE_ID: LEX-P0022-V0008
% #FIELD_VALUE: 测试设备
\fieldvalue{Symphogon}
\\
产品编号 &
% #VALUE_ID: LEX-P0022-V0009
% #FIELD_VALUE: 产品编号
% TODO: 字符辨识不完全确定。
\fieldvalue{A242?025?2010}
\\
产品名称 &
% #VALUE_ID: LEX-P0022-V0010
% #FIELD_VALUE: 产品名称
% TODO: 字符辨识不完全确定。
\fieldvalue{?}
\\
测试模式 &
% #VALUE_ID: LEX-P0022-V0011
% #FIELD_VALUE: 测试模式
\fieldvalue{Simul}
\\
测试比例 &
% #VALUE_ID: LEX-P0022-V0012
% #FIELD_VALUE: 测试比例
\fieldvalue{1:1000}
\\
测试压力 &
% #VALUE_ID: LEX-P0022-V0013
% #FIELD_VALUE: 测试压力
\fieldvalue{4500\,mbar}
\\
最大压力 &
% #VALUE_ID: LEX-P0022-V0014
% #FIELD_VALUE: 最大压力
\fieldvalue{1000\,mbar}
\end{tabular}

\subsection{结果}

\begin{tabular}{@{}p{0.27\linewidth}p{0.65\linewidth}@{}}
测试结果 &
% #VALUE_ID: LEX-P0022-V0015
% #FIELD_VALUE: 测试结果
\fieldvalue{通过；测试压降：4950\,mbar}
\\
测试时间 &
% #VALUE_ID: LEX-P0022-V0016
% #FIELD_VALUE: 测试时间
\fieldvalue{12/Dec/2025 10:54:14}
\end{tabular}

\subsection{签名}

% #VALUE_ID: LEX-P0022-V0017
% #FIELD_VALUE: 签名
% TODO #HANDWRITTEN: （签名）2025.12.12；字迹部分难以辨识。
\fieldvalue{\handwritten{（签名）2025.12.12}}

\subsection{流量曲线}

\begin{center}
\textbf{Palltronic Flowstar IV 测流结果}

\begin{tabular}{@{}p{0.25\linewidth}p{0.65\linewidth}@{}}
日期 &
% #VALUE_ID: LEX-P0022-V0018
% #FIELD_VALUE: 日期
\fieldvalue{12/Dec/2025 10:35:49}
\\
页码 &
% #VALUE_ID: LEX-P0022-V0019
% #FIELD_VALUE: 页码
\fieldvalue{2 of 2}
\end{tabular}

\vspace{1em}

\begin{picture}(0,0)
\put(0,0){\framebox(250,150){}}
\put(5,145){\makebox(0,0)[l]{1800}}
\put(5,120){\makebox(0,0)[l]{1350}}
\put(5,95){\makebox(0,0)[l]{900}}
\put(5,70){\makebox(0,0)[l]{450}}
\put(5,45){\makebox(0,0)[l]{0}}
\put(20,155){\makebox(0,0)[l]{灭菌空气 (mbar/min)}}
\put(35,-8){\makebox(0,0)[l]{0}}
\put(75,-8){\makebox(0,0)[l]{1000}}
\put(115,-8){\makebox(0,0)[l]{2800}}
\put(155,-8){\makebox(0,0)[l]{4200}}
\put(195,-8){\makebox(0,0)[l]{5600}}
\put(235,-8){\makebox(0,0)[l]{7000}}
\put(120,-20){\makebox(0,0)[c]{压差 (mbar)}}
\end{picture}

\vspace{8em}

% #VALUE_ID: LEX-P0022-V0020
% #HANDWRITTEN: （签名）2025/12/12
\handwritten{（签名）2025/12/12}

\end{center}

% #VALUE_ID: LEX-P0022-V0021
% #HANDWRITTEN: 23/35
\handwritten{23/35}
% LEXOID_PAGE_COMPLETED: 22/70

\newpage

\begin{center}
\rotatebox{90}{%
\begin{minipage}{0.96\textheight}
\section*{PALL Palltronic Flowstar IV 测试结果}

\begin{tabular}{@{}p{0.30\linewidth}p{0.58\linewidth}@{}}
日期 &
% #VALUE_ID: LEX-P0023-V0001
% #FIELD_VALUE: 日期
\fieldvalue{12/Dec/2025 10:50:36}
\\[2pt]
页号 &
% #VALUE_ID: LEX-P0023-V0002
% #FIELD_VALUE: 页号
\fieldvalue{1 of 2}
\\[2pt]
测试/Result number &
% #VALUE_ID: LEX-P0023-V0003
% #FIELD_VALUE: 测试/Result number
\fieldvalue{23880442671}
\\[2pt]
操作员/Operator &
% #VALUE_ID: LEX-P0023-V0004
% #FIELD_VALUE: 操作员/Operator
\fieldvalue{SF}
\\[2pt]
测试程序/Test program &
% #VALUE_ID: LEX-P0023-V0005
% #FIELD_VALUE: 测试程序/Test program
\fieldvalue{SAMPES-2525}
\\[2pt]
生产区域 &
% #VALUE_ID: LEX-P0023-V0006
% #FIELD_VALUE: 生产区域
\fieldvalue{C}
\\[2pt]
过滤器规格 &
% #VALUE_ID: LEX-P0023-V0007
% #FIELD_VALUE: 过滤器规格
\fieldvalue{Shenglong}
\\[2pt]
过滤器编号 &
% #VALUE_ID: LEX-P0023-V0008
% #FIELD_VALUE: 过滤器编号
\fieldvalue{2}
\\[2pt]
产品编号 &
% #VALUE_ID: LEX-P0023-V0009
% #FIELD_VALUE: 产品编号
\fieldvalue{4A?202?}
\\[2pt]
产品名称 &
% #VALUE_ID: LEX-P0023-V0010
% #FIELD_VALUE: 产品名称
\fieldvalue{SAMPES-2525}
\\[2pt]
测试类型 &
% #VALUE_ID: LEX-P0023-V0011
% #FIELD_VALUE: 测试类型
\fieldvalue{---}
\\[2pt]
测试压力 &
% #VALUE_ID: LEX-P0023-V0012
% #FIELD_VALUE: 测试压力
\fieldvalue{1000}
\\[2pt]
最小测试压力 &
% #VALUE_ID: LEX-P0023-V0013
% #FIELD_VALUE: 最小测试压力
\fieldvalue{4500\,mbar}
\\[2pt]
最大测试压力 &
% #VALUE_ID: LEX-P0023-V0014
% #FIELD_VALUE: 最大测试压力
\fieldvalue{7000\,mbar}
\\[2pt]
测试结束时间 &
% #VALUE_ID: LEX-P0023-V0015
% #FIELD_VALUE: 测试结束时间
\fieldvalue{12/Dec/2025 10:49:30}
\\[2pt]
测试开始时间 &
% #VALUE_ID: LEX-P0023-V0016
% #FIELD_VALUE: 测试开始时间
\fieldvalue{12/Dec/2025 10:49:30}
\\
\end{tabular}

\vspace{0.5em}
\noindent\textbf{备注：}

\vspace{1em}
\noindent\textbf{签名：}

% #VALUE_ID: LEX-P0023-V0017
% #FIELD_VALUE: 签名
% #HANDWRITTEN: 2025.12.12
\fieldvalue{\handwritten{2025.12.12}}

\vspace{1em}
% #VALUE_ID: LEX-P0023-V0018
% #HANDWRITTEN: 签名及日期；字迹部分难以辨认
\handwritten{签名，2025.12.12}

\bigskip
\hrule
\bigskip

\section*{PALL Palltronic Flowstar IV 测试结果}

\begin{tabular}{@{}p{0.30\linewidth}p{0.58\linewidth}@{}}
日期 &
% #VALUE_ID: LEX-P0023-V0019
% #FIELD_VALUE: 日期
\fieldvalue{12/Dec/2025 10:50:36}
\\[2pt]
页号 &
% #VALUE_ID: LEX-P0023-V0020
% #FIELD_VALUE: 页号
\fieldvalue{2 of 2}
\\
\end{tabular}

\vspace{1em}

\begin{center}
\begin{tikzpicture}[x=0.0012\linewidth,y=0.0012\linewidth]
  \draw (0,0) rectangle (650,300);
  \draw[->] (0,0) -- (0,325);
  \draw[->] (0,0) -- (680,0);
  \node[rotate=90] at (-42,150) {大气流量 (mbar/min)};
  \node at (330,-35) {时间 (msec)};
  \node[above] at (0,300) {450};
  \node[above] at (215,300) {900};
  \node[above] at (430,300) {1350};
  \node[above] at (640,300) {1800};
  \node[below] at (0,-2) {1400};
  \node[below] at (215,-2) {2000};
  \node[below] at (430,-2) {4200};
  \node[below] at (540,-2) {5600};
  \node[below] at (640,-2) {7000};
  \draw[thick] (0,3) -- (430,3)
    .. controls (455,3) and (470,12) .. (478,75)
    .. controls (485,160) and (490,250) .. (492,295);
\end{tikzpicture}
\end{center}

% #VALUE_ID: LEX-P0023-V0021
% #HANDWRITTEN: 2025.12.12
\handwritten{2025.12.12}

% #VALUE_ID: LEX-P0023-V0022
% #HANDWRITTEN: 24/3?; blue handwriting is partly unclear
\handwritten{24/3?}

\end{minipage}%
}
\end{center}
% LEXOID_PAGE_COMPLETED: 23/70

\newpage

\section*{Palltronic FlowStar IV 流量结果}

\begin{tabular}{@{}p{0.30\linewidth}p{0.60\linewidth}@{}}
日期 &
% #VALUE_ID: LEX-P0024-V0001
% #FIELD_VALUE: 日期
\fieldvalue{12/Dec/2025 11:04:18}
\\
页码 &
% #VALUE_ID: LEX-P0024-V0002
% #FIELD_VALUE: 页码
\fieldvalue{2 of 2}
\\
编号 &
% #VALUE_ID: LEX-P0024-V0003
% #FIELD_VALUE: 编号
\fieldvalue{F5?G? / 12?1}
\\
Serial/Result number &
% #VALUE_ID: LEX-P0024-V0004
% #FIELD_VALUE: Serial/Result number
\fieldvalue{28807442618}
\\
测试结果 &
% #VALUE_ID: LEX-P0024-V0005
% #FIELD_VALUE: 测试结果
\fieldvalue{合格}
\\
测试型号 &
% #VALUE_ID: LEX-P0024-V0006
% #FIELD_VALUE: 测试型号
\fieldvalue{SFM?S-2235}
\\
测试模式 &
% #VALUE_ID: LEX-P0024-V0007
% #FIELD_VALUE: 测试模式
\fieldvalue{02}
\\
过滤器型号 &
% #VALUE_ID: LEX-P0024-V0008
% #FIELD_VALUE: 过滤器型号
\fieldvalue{2N?}
\\
过滤器序列号 &
% #VALUE_ID: LEX-P0024-V0009
% #FIELD_VALUE: 过滤器序列号
\fieldvalue{A?}
\\
产品名称 &
% #VALUE_ID: LEX-P0024-V0010
% #FIELD_VALUE: 产品名称
\fieldvalue{A?}
\\
过滤器类型 &
% #VALUE_ID: LEX-P0024-V0011
% #FIELD_VALUE: 过滤器类型
\fieldvalue{SFM?S-2235}
\\
测试介质 &
% #VALUE_ID: LEX-P0024-V0012
% #FIELD_VALUE: 测试介质
\fieldvalue{Simu}
\\
测试流量 &
% #VALUE_ID: LEX-P0024-V0013
% #FIELD_VALUE: 测试流量
\fieldvalue{1000}
\\
测试压力 &
% #VALUE_ID: LEX-P0024-V0014
% #FIELD_VALUE: 测试压力
\fieldvalue{4500\ \mathrm{mbar}}
\\
最大允许压差 &
% #VALUE_ID: LEX-P0024-V0015
% #FIELD_VALUE: 最大允许压差
\fieldvalue{5000\ \mathrm{mbar}}
\\
扩散流量 &
% #VALUE_ID: LEX-P0024-V0016
% #FIELD_VALUE: 扩散流量
\fieldvalue{120\ \mathrm{mbar}}
\end{tabular}

\vspace{1em}

% #VALUE_ID: LEX-P0024-V0017
% #FIELD_VALUE: 签名及日期
% #HANDWRITTEN: 李健 2025.12.12
\fieldvalue{\handwritten{李健 2025.12.12}}

\vfill

\section*{Palltronic FlowStar IV 流量结果}

\begin{tabular}{@{}p{0.30\linewidth}p{0.60\linewidth}@{}}
日期 &
% #VALUE_ID: LEX-P0024-V0018
% #FIELD_VALUE: 日期
\fieldvalue{12/Dec/2025 11:04:18}
\\
页码 &
% #VALUE_ID: LEX-P0024-V0019
% #FIELD_VALUE: 页码
\fieldvalue{2 of 2}
\end{tabular}

\vspace{0.5em}

\begin{center}
\begin{tabular}{|p{0.12\linewidth}|p{0.66\linewidth}|}
\hline
\multicolumn{2}{|c|}{大气流量（mbar/min）}\\
\hline
0 & 450 \hfill 900 \hfill 1350 \hfill 1800\\
\hline
\multicolumn{2}{c}{}
\\[-1.5ex]
\multicolumn{2}{c}{\rule{0pt}{7em}
\hfill
\begin{picture}(0,0)
% TODO: Reproduce the plotted curve visible in the graph.
\end{picture}}
\\
\hline
\multicolumn{2}{c}{}
\\[-1.5ex]
\multicolumn{2}{c}{0 \hfill 1000 \hfill 2000 \hfill 4000 \hfill 5000 \hfill 7000}\\
\multicolumn{2}{c}{测试时间（mbar）}\\
\end{tabular}
\end{center}

% #VALUE_ID: LEX-P0024-V0020
% #FIELD_VALUE: 签名及日期
% #HANDWRITTEN: 李健 2025.12.12
\fieldvalue{\handwritten{李健 2025.12.12}}

% #VALUE_ID: LEX-P0024-V0021
% #HANDWRITTEN: 25/35
\handwritten{25/35}
% LEXOID_PAGE_COMPLETED: 24/70

\newpage

\section*{Palltronic Flowstar IV 测试结果}

\begin{tabular}{@{}p{0.25\linewidth}p{0.65\linewidth}@{}}
日期 &
% #VALUE_ID: LEX-P0025-V0001
% #FIELD_VALUE: 日期
\fieldvalue{12/Dec/2025 11:19:02}
\\[2pt]
项目 &
% #VALUE_ID: LEX-P0025-V0002
% #FIELD_VALUE: 项目
\fieldvalue{1.02}
\\[2pt]
设备号 &
% #VALUE_ID: LEX-P0025-V0003
% #FIELD_VALUE: 设备号
% TODO: 设备号部分字符辨认不清。
\fieldvalue{FFSQ?LCR/22?1T}
\\[2pt]
测试结果 &
% #VALUE_ID: LEX-P0025-V0004
% #FIELD_VALUE: 测试结果
% TODO: 测试结果数值部分辨认不清。
\fieldvalue{28.???}
\\[2pt]
功能 &
% #VALUE_ID: LEX-P0025-V0005
% #FIELD_VALUE: 功能
\fieldvalue{测漏}
\\[2pt]
序列号 &
% #VALUE_ID: LEX-P0025-V0006
% #FIELD_VALUE: 序列号
\fieldvalue{SFMPES2-225}
\\[2pt]
测试类型 &
% #VALUE_ID: LEX-P0025-V0007
% #FIELD_VALUE: 测试类型
\fieldvalue{Q2}
\\[2pt]
生产线 &
% #VALUE_ID: LEX-P0025-V0008
% #FIELD_VALUE: 生产线
\fieldvalue{C}
\\[2pt]
产品名称 &
% #VALUE_ID: LEX-P0025-V0009
% #FIELD_VALUE: 产品名称
\fieldvalue{Syringe/10}
\\[2pt]
产品编号 &
% #VALUE_ID: LEX-P0025-V0010
% #FIELD_VALUE: 产品编号
\fieldvalue{2N/sleeve}
\\[2pt]
过滤器型号 &
% #VALUE_ID: LEX-P0025-V0011
% #FIELD_VALUE: 过滤器型号
% TODO: 过滤器型号部分字符辨认不清。
\fieldvalue{A?220?202?2019}
\\[2pt]
过滤器序列号 &
% #VALUE_ID: LEX-P0025-V0012
% #FIELD_VALUE: 过滤器序列号
% TODO: 过滤器序列号部分字符辨认不清。
\fieldvalue{2N/sleeve}
\\[2pt]
测试压力 &
% #VALUE_ID: LEX-P0025-V0013
% #FIELD_VALUE: 测试压力
\fieldvalue{450\ mbar}
\\[2pt]
测试气体 &
% #VALUE_ID: LEX-P0025-V0014
% #FIELD_VALUE: 测试气体
\fieldvalue{Air}
\\[2pt]
测试时长 &
% #VALUE_ID: LEX-P0025-V0015
% #FIELD_VALUE: 测试时长
\fieldvalue{1:00}
\\[2pt]
最大压力 &
% #VALUE_ID: LEX-P0025-V0016
% #FIELD_VALUE: 最大压力
\fieldvalue{4500\ mbar}
\\[2pt]
最小压力 &
% #VALUE_ID: LEX-P0025-V0017
% #FIELD_VALUE: 最小压力
\fieldvalue{<2000\ mbar}
\\[2pt]
结果 &
% #VALUE_ID: LEX-P0025-V0018
% #FIELD_VALUE: 结果
\fieldvalue{通过}
\\[2pt]
测量终点 &
% #VALUE_ID: LEX-P0025-V0019
% #FIELD_VALUE: 测量终点
\fieldvalue{测试终点：4950\ mbar}
\\[2pt]
测试开始时间 &
% #VALUE_ID: LEX-P0025-V0020
% #FIELD_VALUE: 测试开始时间
\fieldvalue{12/Dec/2025 11:29:24}
\\[2pt]
注释 &
% #VALUE_ID: LEX-P0025-V0021
% #FIELD_VALUE: 注释
% TODO: 注释中的字母和数字辨认不清。
\fieldvalue{A2A?3?2?109?12?9?}
\\[2pt]
签名 &
% #VALUE_ID: LEX-P0025-V0022
% #FIELD_VALUE: 签名
% #TODO #HANDWRITTEN: 蒋华 2025.12.12; 手写签名及日期辨认不完全清楚。
\fieldvalue{\handwritten{蒋华 2025.12.12}}
\end{tabular}

\vspace{1em}
\noindent
\textbf{Palltronic Flowstar IV 测试结果}

\begin{tabular}{@{}p{0.25\linewidth}p{0.65\linewidth}@{}}
日期 &
% #VALUE_ID: LEX-P0025-V0023
% #FIELD_VALUE: 日期
\fieldvalue{12/Dec/2025 11:19:02}
\\[2pt]
页码 &
% #VALUE_ID: LEX-P0025-V0024
% #FIELD_VALUE: 页码
\fieldvalue{2 of 2}
\end{tabular}

\vspace{1em}
\noindent
% TODO: 此处为测试压力--时间曲线图；原图显示纵轴“大气绝对压力 (mbar/min)”刻度约为
% 0、450、900、1350、1800，横轴“时间 (min:sec)”刻度约为
% 0、1400、2800、4200、5600、7000。曲线在中后段出现明显下降后保持近似水平。
\begin{center}
\setlength{\unitlength}{1mm}
\begin{picture}(125,65)
\put(5,5){\framebox(110,52){}}
\put(5,5){\line(1,0){110}}
\put(5,5){\line(0,1){52}}
\put(5,57){\makebox(30,4)[l]{\scriptsize 大气绝对压力 (mbar/min)}}
\put(45,1){\makebox(45,4)[c]{\scriptsize 时间 (min:sec)}}
\put(6,51){\makebox(12,4)[r]{\scriptsize 1800}}
\put(6,38){\makebox(12,4)[r]{\scriptsize 1350}}
\put(6,25){\makebox(12,4)[r]{\scriptsize 900}}
\put(6,12){\makebox(12,4)[r]{\scriptsize 450}}
\put(7,2){\makebox(8,4)[c]{\scriptsize 0}}
\put(27,2){\makebox(14,4)[c]{\scriptsize 1400}}
\put(48,2){\makebox(14,4)[c]{\scriptsize 2800}}
\put(69,2){\makebox(14,4)[c]{\scriptsize 4200}}
\put(90,2){\makebox(14,4)[c]{\scriptsize 5600}}
\put(108,2){\makebox(14,4)[c]{\scriptsize 7000}}
\put(8,51){\line(1,-1){48}}
\put(56,3){\line(0,1){30}}
\put(56,33){\line(1,-1){8}}
\put(64,25){\line(1,0){42}}
\end{picture}
\end{center}

\noindent
签名：%
% #VALUE_ID: LEX-P0025-V0025
% #FIELD_VALUE: 签名
% #TODO #HANDWRITTEN: 蒋华 2025.12.12; 手写签名及日期辨认不完全清楚。
\fieldvalue{\handwritten{蒋华 2025.12.12}}

\vspace{1em}
\noindent
% #VALUE_ID: LEX-P0025-V0026
% #HANDWRITTEN: 20/33
\handwritten{20/33}
% LEXOID_PAGE_COMPLETED: 25/70

\newpage

\section{Palltronic Flowstar 测试报告}

\begin{tabular}{@{}p{0.28\linewidth}p{0.62\linewidth}@{}}
日期 &
% #VALUE_ID: LEX-P0026-V0001
% #FIELD_VALUE: 日期
\fieldvalue{12/20/2025 13:15:50}
\\
序号 &
% #VALUE_ID: LEX-P0026-V0002
% #FIELD_VALUE: 序号
\fieldvalue{1 of 2}
\\
产品型号 &
% #VALUE_ID: LEX-P0026-V0003
% #FIELD_VALUE: 产品型号
\fieldvalue{\texttt{S-4042}}
\\
Seal test number &
% #VALUE_ID: LEX-P0026-V0004
% #FIELD_VALUE: Seal test number
\fieldvalue{\texttt{238804-42620}}
\\
结果 &
% #VALUE_ID: LEX-P0026-V0005
% #FIELD_VALUE: 结果
\fieldvalue{Pass}
\\
测试程序 &
% #VALUE_ID: LEX-P0026-V0006
% #FIELD_VALUE: 测试程序
\fieldvalue{\texttt{SFM PES-225}}
\\
测试气体 &
% #VALUE_ID: LEX-P0026-V0007
% #FIELD_VALUE: 测试气体
\fieldvalue{\texttt{O2}}
\\
测试操作者 &
% #VALUE_ID: LEX-P0026-V0008
% #FIELD_VALUE: 测试操作者
\fieldvalue{\texttt{Shiyonglou}}
\\
产品名称 &
% #VALUE_ID: LEX-P0026-V0009
% #FIELD_VALUE: 产品名称
\fieldvalue{\texttt{A2?}}
\\
产品编号 &
% #VALUE_ID: LEX-P0026-V0010
% #FIELD_VALUE: 产品编号
\fieldvalue{\texttt{A?}}
\\
测试模式 &
% #VALUE_ID: LEX-P0026-V0011
% #FIELD_VALUE: 测试模式
\fieldvalue{\texttt{A?}}
\\
泄漏流量单位 &
% #VALUE_ID: LEX-P0026-V0012
% #FIELD_VALUE: 泄漏流量单位
\fieldvalue{\texttt{sccm}}
\\
测试压力 &
% #VALUE_ID: LEX-P0026-V0013
% #FIELD_VALUE: 测试压力
\fieldvalue{1000}
\\
最大流量 &
% #VALUE_ID: LEX-P0026-V0014
% #FIELD_VALUE: 最大流量
\fieldvalue{4500\ \text{mbar}}
\\
最大压力 &
% #VALUE_ID: LEX-P0026-V0015
% #FIELD_VALUE: 最大压力
\fieldvalue{7000\ \text{mbar}}
\end{tabular}

\vspace{1em}
\hrule
\vspace{0.5em}

\begin{tabular}{@{}p{0.28\linewidth}p{0.62\linewidth}@{}}
测试结果 &
% #VALUE_ID: LEX-P0026-V0016
% #FIELD_VALUE: 测试结果
\fieldvalue{通过，5000\ \text{mbar}}
\\
泄漏测试时间 &
% #VALUE_ID: LEX-P0026-V0017
% #FIELD_VALUE: 泄漏测试时间
\fieldvalue{5000\ \text{msec}}
\\
测试完成时间 &
% #VALUE_ID: LEX-P0026-V0018
% #FIELD_VALUE: 测试完成时间
\fieldvalue{12/20/2025 13:41:13}
\end{tabular}

\vspace{1em}

% #VALUE_ID: LEX-P0026-V0019
% #FIELD_VALUE: 签名
% #HANDWRITTEN: [签名]
\fieldvalue{\handwritten{[签名]}}

% #VALUE_ID: LEX-P0026-V0020
% #FIELD_VALUE: 日期
% #HANDWRITTEN: 2025/12/12
\fieldvalue{\handwritten{2025/12/12}}

\section{Palltronic Flowstar 测试报告}

\begin{tabular}{@{}p{0.28\linewidth}p{0.62\linewidth}@{}}
日期 &
% #VALUE_ID: LEX-P0021-V0021
% #FIELD_VALUE: 日期
\fieldvalue{12/20/2025 13:15:50}
\\
序号 &
% #VALUE_ID: LEX-P0022-V0022
% #FIELD_VALUE: 序号
\fieldvalue{2 of 2}
\end{tabular}

\vspace{1em}

\begin{center}
\setlength{\fboxsep}{8pt}
\fbox{\parbox{0.82\linewidth}{%
\centering
% TODO: The page contains a plotted test graph; reproduce the plotted curve if graphical reconstruction is required.

\textbf{大气压差 (mbar/min)}

\vspace{0.5em}
\begin{tabular}{c|ccccc}
1800 & & & & & \\
1350 & & & & & \\
900  & & & & & \\
450  & & & & & \\
0    & & & & & \\
\hline
     & 0 & 1400 & 2800 & 4200 & 5600 \quad 7000
\end{tabular}

\vspace{0.5em}
测试时间
}}
\end{center}

% #VALUE_ID: LEX-P0026-V0023
% #FIELD_VALUE: 签名
% #HANDWRITTEN: [签名]
\fieldvalue{\handwritten{[签名]}}

% #VALUE_ID: LEX-P0026-V0024
% #FIELD_VALUE: 日期
% #HANDWRITTEN: 2025/12/12
\fieldvalue{\handwritten{2025/12/12}}

% #VALUE_ID: LEX-P0026-V0025
% #HANDWRITTEN: 21/35
\handwritten{21/35}
% LEXOID_PAGE_COMPLETED: 26/70

\newpage

\section*{Palltronic Flowstar IV 测试报告}

\begin{tabular}{@{}p{0.25\linewidth}p{0.65\linewidth}@{}}
日期 &
% #VALUE_ID: LEX-P0027-V0001
% #FIELD_VALUE: 日期
\fieldvalue{12/Dec/2025 13:42:24}
\\
页号 & 1 of 2 \\
设备 &
% #VALUE_ID: LEX-P0027-V0002
% #FIELD_VALUE: 设备
\fieldvalue{FMS04L CR? /2?1}
\\
Serial/Asset number &
% #VALUE_ID: LEX-P0027-V0003
% #FIELD_VALUE: Serial/Asset number
\fieldvalue{3288024/42621?}
\\
测试模式 &
% #VALUE_ID: LEX-P0027-V0004
% #FIELD_VALUE: 测试模式
\fieldvalue{Test}
\\
测试程序 &
% #VALUE_ID: LEX-P0027-V0005
% #FIELD_VALUE: 测试程序
\fieldvalue{SPMES-2225}
\\
生产区域 &
% #VALUE_ID: LEX-P0027-V0006
% #FIELD_VALUE: 生产区域
\fieldvalue{QZ}
\\
生产设备 &
% #VALUE_ID: LEX-P0027-V0007
% #FIELD_VALUE: 生产设备
\fieldvalue{Sheng?jing?hua}
\\
过滤器型号 &
% #VALUE_ID: LEX-P0027-V0008
% #FIELD_VALUE: 过滤器型号
\fieldvalue{2?}
\\
产品编号 &
% #VALUE_ID: LEX-P0027-V0009
% #FIELD_VALUE: 产品编号
\fieldvalue{A4?220?25109}
\\
产品名称 &
% #VALUE_ID: LEX-P0027-V0010
% #FIELD_VALUE: 产品名称
\fieldvalue{Filter?}
\\
过滤器批号 &
% #VALUE_ID: LEX-P0027-V0011
% #FIELD_VALUE: 过滤器批号
\fieldvalue{SPMES-2225}
\\
过滤器名称 &
% #VALUE_ID: LEX-P0027-V0012
% #FIELD_VALUE: 过滤器名称
\fieldvalue{S?}
\\
测试介质 &
% #VALUE_ID: LEX-P0027-V0013
% #FIELD_VALUE: 测试介质
\fieldvalue{Air}
\\
测试压力 &
% #VALUE_ID: LEX-P0027-V0014
% #FIELD_VALUE: 测试压力
\fieldvalue{1.000}
\\
最小测试流量 &
% #VALUE_ID: LEX-P0027-V0015
% #FIELD_VALUE: 最小测试流量
\fieldvalue{4500 mbar}
\end{tabular}

\vspace{1em}

\begin{tabular}{@{}p{0.25\linewidth}p{0.65\linewidth}@{}}
结果 &
% #VALUE_ID: LEX-P0027-V0016
% #FIELD_VALUE: 结果
\fieldvalue{通过}
\\
测试点 &
% #VALUE_ID: LEX-P0027-V0017
% #FIELD_VALUE: 测试点
\fieldvalue{5000 mbar}
\\
测试结束时间 &
% #VALUE_ID: LEX-P0027-V0018
% #FIELD_VALUE: 测试结束时间
\fieldvalue{10/Dec/2025 13:42:06}
\\
签名 &
% #VALUE_ID: LEX-P0027-V0019
% #FIELD_VALUE: 签名
% TODO #HANDWRITTEN: [签名不清]; handwritten signature is illegible
\fieldvalue{\handwritten{[签名不清]}}
\\
日期 &
% #VALUE_ID: LEX-P0027-V0020
% #FIELD_VALUE: 日期
% #HANDWRITTEN: 2025.12.12
\fieldvalue{\handwritten{2025.12.12}}
\end{tabular}

\vspace{2em}

\section*{Palltronic Flowstar IV 测试报告}

\begin{tabular}{@{}p{0.25\linewidth}p{0.65\linewidth}@{}}
日期 &
% #VALUE_ID: LEX-P0021-V0021
% #FIELD_VALUE: 日期
\fieldvalue{12/Dec/2025 13:42:24}
\\
页号 & 2 of 2
\end{tabular}

\vspace{1em}

\begin{center}
\begin{tabular}{c}
\text{大气压强 (mbar/min)}\\[-0.25em]
\begin{picture}(260,150)
\put(25,10){\vector(1,0){220}}
\put(25,10){\vector(0,1){125}}
\put(15,5){\makebox(0,0)[r]{0}}
\put(25,0){\makebox(0,0)[t]{1400}}
\put(90,0){\makebox(0,0)[t]{2000}}
\put(125,0){\makebox(0,0)[t]{2400}}
\put(175,0){\makebox(0,0)[t]{5000}}
\put(245,0){\makebox(0,0)[t]{7000}}
\put(25,48){\makebox(0,0)[r]{450}}
\put(25,82){\makebox(0,0)[r]{900}}
\put(25,108){\makebox(0,0)[r]{1350}}
\put(25,135){\makebox(0,0)[r]{1800}}
\put(135,-10){\makebox(0,0)[c]{时间 (mbar)}}
\put(25,10){\line(1,0){150}}
\put(175,10){\line(0,1){115}}
\put(175,125){\line(1,0){70}}
\put(175,10){\line(1,0){70}}
\put(175,10){\curve(0,0)(4,0)(8,2)(12,10)(15,35)(17,70)(19,95)(22,108)}
\end{picture}
\end{tabular}
\end{center}

\begin{tabular}{@{}p{0.25\linewidth}p{0.65\linewidth}@{}}
签名 &
% #VALUE_ID: LEX-P0027-V0022
% #FIELD_VALUE: 签名
% TODO #HANDWRITTEN: [签名不清]; handwritten signature is illegible
\fieldvalue{\handwritten{[签名不清]}}
\\
日期 &
% #VALUE_ID: LEX-P0027-V0023
% #FIELD_VALUE: 日期
% #HANDWRITTEN: 2025.12.12
\fieldvalue{\handwritten{2025.12.12}}
\end{tabular}

% TODO: The page contains two rotated report panels; graph curve and several metadata values are partially obscured or difficult to read in the source image.
% LEXOID_PAGE_COMPLETED: 27/70

\newpage

\section*{Paltronic FlowStart IV 流量记录}

\begin{tabular}{@{}p{0.28\linewidth}p{0.65\linewidth}@{}}
日期 &
% #VALUE_ID: LEX-P0028-V0001
% #FIELD_VALUE: 日期
\fieldvalue{12/Dec/2025 13:27:38}
\\
页码 &
% #VALUE_ID: LEX-P0028-V0002
% #FIELD_VALUE: 页码
\fieldvalue{1 of 2}
\\
序列号 &
% #VALUE_ID: LEX-P0028-V0003
% #FIELD_VALUE: 序列号
% TODO: 序列号部分字符难以辨认。
\fieldvalue{FFS04LGR/242?7T}
\\
设备 Resource number &
% #VALUE_ID: LEX-P0028-V0004
% #FIELD_VALUE: 设备 Resource number
% TODO: 设备编号部分字符难以辨认。
\fieldvalue{988?047?422?2}
\\
功能编号 &
% #VALUE_ID: LEX-P0028-V0005
% #FIELD_VALUE: 功能编号
% TODO: 功能编号部分字符难以辨认。
\fieldvalue{2A?}
\\
灭菌批号 &
% #VALUE_ID: LEX-P0028-V0006
% #FIELD_VALUE: 灭菌批号
% TODO: 灭菌批号部分字符难以辨认。
\fieldvalue{S?M?S-2225}
\\
操作员 &
% #VALUE_ID: LEX-P0028-V0007
% #FIELD_VALUE: 操作员
\fieldvalue{QZ}
\\
检验品 &
% #VALUE_ID: LEX-P0028-V0008
% #FIELD_VALUE: 检验品
\fieldvalue{SFMEPS-2225}
\\
生产名称 &
% #VALUE_ID: LEX-P0028-V0009
% #FIELD_VALUE: 生产名称
% TODO: 生产名称字符难以辨认。
\fieldvalue{2A?}
\\
生产批号 &
% #VALUE_ID: LEX-P0028-V0010
% #FIELD_VALUE: 生产批号
% TODO: 生产批号字符难以辨认。
\fieldvalue{2?N?}
\\
过滤器型号 &
% #VALUE_ID: LEX-P0028-V0011
% #FIELD_VALUE: 过滤器型号
\fieldvalue{S?ME?S-2225}
\\
过滤器序列号 &
% #VALUE_ID: LEX-P0028-V0012
% #FIELD_VALUE: 过滤器序列号
\fieldvalue{-}
\\
过滤器尺寸 &
% #VALUE_ID: LEX-P0028-V0013
% #FIELD_VALUE: 过滤器尺寸
\fieldvalue{5\,mm}
\\
测量单位 &
% #VALUE_ID: LEX-P0028-V0014
% #FIELD_VALUE: 测量单位
\fieldvalue{mbar}
\\
最大压力 &
% #VALUE_ID: LEX-P0028-V0015
% #FIELD_VALUE: 最大压力
\fieldvalue{4500\,mbar}
\\
结果 &
% #VALUE_ID: LEX-P0028-V0016
% #FIELD_VALUE: 结果
\fieldvalue{通过}
\\
实施地点 &
% #VALUE_ID: LEX-P0028-V0017
% #FIELD_VALUE: 实施地点
\fieldvalue{洁净室}
\\
实施压力 &
% #VALUE_ID: LEX-P0028-V0018
% #FIELD_VALUE: 实施压力
\fieldvalue{5000\,mbar}
\\
测量数据时间 &
% #VALUE_ID: LEX-P0028-V0019
% #FIELD_VALUE: 测量数据时间
\fieldvalue{12/Dec/2025 13:38:20}
\\
\end{tabular}

\vspace{1em}

% #VALUE_ID: LEX-P0028-V0020
% #HANDWRITTEN: 康静
\handwritten{康静}
\qquad
% #VALUE_ID: LEX-P0028-V0021
% #HANDWRITTEN: 2025.12.12
\handwritten{2025.12.12}

\section*{Paltronic FlowStart IV 流量记录}

\begin{tabular}{@{}p{0.28\linewidth}p{0.65\linewidth}@{}}
日期 &
% #VALUE_ID: LEX-P0028-V0022
% #FIELD_VALUE: 日期
\fieldvalue{12/Dec/2025 13:27:38}
\\
页码 &
% #VALUE_ID: LEX-P0028-V0023
% #FIELD_VALUE: 页码
\fieldvalue{2 of 2}
\\
\end{tabular}

\vspace{1em}

\begin{center}
\begin{tikzpicture}
  \draw (0,0) rectangle (0.82\linewidth,4.0cm);
  \draw[->] (0,0) -- (0.82\linewidth,0)
    node[below, midway] {时间（mbar）};
  \draw[->] (0,0) -- (0,4.0cm)
    node[above, midway, rotate=90] {大气体流量（mbar/min）};

  \foreach \x/\lab in {0/0,0.16\linewidth/1400,0.32\linewidth/2800,
                        0.48\linewidth/4200,0.64\linewidth/5600,
                        0.80\linewidth/7000}
    \draw (\x,0.08cm) -- (\x,-0.08cm)
      node[below] {\scriptsize \lab};

  \foreach \y/\lab in {0/0,1.0cm/450,2.0cm/900,3.0cm/1350,4.0cm/1800}
    \draw (0.08cm,\y) -- (-0.08cm,\y)
      node[left] {\scriptsize \lab};

  \draw[thin] (0.02\linewidth,3.92cm)
    -- (0.48\linewidth,3.55cm)
    -- (0.57\linewidth,3.35cm)
    -- (0.61\linewidth,2.65cm)
    -- (0.64\linewidth,0.30cm)
    -- (0.70\linewidth,0.10cm)
    -- (0.80\linewidth,0.08cm);
\end{tikzpicture}
\end{center}

% TODO: The graph is reproduced schematically from the visible page; exact plotted coordinates are approximate.

% #VALUE_ID: LEX-P0028-V0024
% #HANDWRITTEN: 康静
\handwritten{康静}
\qquad
% #VALUE_ID: LEX-P0028-V0025
% #HANDWRITTEN: 2025.12.12
\handwritten{2025.12.12}

\vfill

% #VALUE_ID: LEX-P0028-V0026
% #HANDWRITTEN: 29/35
\handwritten{29/35}
% LEXOID_PAGE_COMPLETED: 28/70

\newpage

\section{Palltronic Flowstar IV 测试结果}

\begin{tabular}{@{}p{0.30\linewidth}p{0.62\linewidth}@{}}
日期 &
% #VALUE_ID: LEX-P0029-V0001
% #FIELD_VALUE: 日期
\fieldvalue{12/Dec/2025 13:42:19} \\[2pt]
页码 &
% #VALUE_ID: LEX-P0029-V0002
% #FIELD_VALUE: 页码
\fieldvalue{1 of 2} \\[2pt]
Serial/Result number &
% #VALUE_ID: LEX-P0029-V0003
% #FIELD_VALUE: Serial/Result number
\fieldvalue{238804?4223} \\[2pt]
测试单元 &
% #VALUE_ID: LEX-P0029-V0004
% #FIELD_VALUE: 测试单元
\fieldvalue{Air?} \\[2pt]
测试程序 &
% #VALUE_ID: LEX-P0029-V0005
% #FIELD_VALUE: 测试程序
\fieldvalue{SFMPES-2225} \\[2pt]
测试结果 &
% #VALUE_ID: LEX-P0029-V0006
% #FIELD_VALUE: 测试结果
\fieldvalue{C: 02} \\[2pt]
过滤器 &
% #VALUE_ID: LEX-P0029-V0007
% #FIELD_VALUE: 过滤器
\fieldvalue{Shiyonglou} \\[2pt]
产品编号 &
% #VALUE_ID: LEX-P0029-V0008
% #FIELD_VALUE: 产品编号
\fieldvalue{A?222?} \\[2pt]
产品名称 &
% #VALUE_ID: LEX-P0029-V0009
% #FIELD_VALUE: 产品名称
\fieldvalue{A?222?} \\[2pt]
过滤器型号 &
% #VALUE_ID: LEX-P0029-V0010
% #FIELD_VALUE: 过滤器型号
\fieldvalue{SFMPES-2225} \\[2pt]
测试气体 &
% #VALUE_ID: LEX-P0029-V0011
% #FIELD_VALUE: 测试气体
\fieldvalue{Air} \\[2pt]
测试压力 &
% #VALUE_ID: LEX-P0029-V0012
% #FIELD_VALUE: 测试压力
\fieldvalue{1.000} \\[2pt]
最大压力 &
% #VALUE_ID: LEX-P0029-V0013
% #FIELD_VALUE: 最大压力
\fieldvalue{4500 mbar} \\[2pt]
最小压力 &
% #VALUE_ID: LEX-P0029-V0014
% #FIELD_VALUE: 最小压力
\fieldvalue{7000 mbar}
\end{tabular}

\medskip
\noindent\rule{\linewidth}{0.4pt}

\begin{tabular}{@{}p{0.30\linewidth}p{0.62\linewidth}@{}}
结果 &
% #VALUE_ID: LEX-P0029-V0015
% #FIELD_VALUE: 结果
\fieldvalue{通过} \\[2pt]
最大压力 &
% #VALUE_ID: LEX-P0029-V0016
% #FIELD_VALUE: 最大压力
\fieldvalue{4500 mbar} \\[2pt]
最小压力 &
% #VALUE_ID: LEX-P0029-V0017
% #FIELD_VALUE: 最小压力
\fieldvalue{7000 mbar} \\[2pt]
测试结束时间 &
% #VALUE_ID: LEX-P0029-V0018
% #FIELD_VALUE: 测试结束时间
\fieldvalue{12/Dec/2025 13:52:52}
\end{tabular}

% TODO: The handwritten signature/name in the center of the report is partly illegible.
% #VALUE_ID: LEX-P0029-V0019
% #HANDWRITTEN: [illegible signature/name]
\handwritten{[illegible signature/name]}

% #VALUE_ID: LEX-P0029-V0020
% #HANDWRITTEN: 2025/12/12
\handwritten{2025/12/12}

\section{Palltronic Flowstar IV 测试结果}

\begin{tabular}{@{}p{0.30\linewidth}p{0.62\linewidth}@{}}
日期 &
% #VALUE_ID: LEX-P0029-V0021
% #FIELD_VALUE: 日期
\fieldvalue{12/Dec/2025 13:42:19} \\[2pt]
页码 &
% #VALUE_ID: LEX-P0029-V0022
% #FIELD_VALUE: 页码
\fieldvalue{2 of 2}
\end{tabular}

\medskip

\begin{center}
\begin{tikzpicture}
  \draw (0,0) rectangle (0.78\linewidth,4.5cm);
  % TODO: Reproduce the plotted curve and axis labels from the visible graph.
  \draw[->] (0.35cm,0.35cm) -- (0.35cm,4.15cm);
  \draw[->] (0.35cm,0.35cm) -- (0.73\linewidth,0.35cm);
  \node[rotate=90] at (0.08cm,2.25cm) {\scriptsize 大气补偿量 (mbar/min)};
  \node at (0.40\linewidth,0.05cm) {\scriptsize 压力 (mbar)};
\end{tikzpicture}
\end{center}

% TODO: The dark handwritten signature/name near the graph is partly illegible.
% #VALUE_ID: LEX-P0029-V0023
% #HANDWRITTEN: [illegible signature/name]
\handwritten{[illegible signature/name]}

% #VALUE_ID: LEX-P0029-V0024
% #HANDWRITTEN: 2025/12/12
\handwritten{2025/12/12}

% TODO: The blue handwritten entry at the lower right is partly illegible.
% #VALUE_ID: LEX-P0029-V0025
% #HANDWRITTEN: 30/?
\handwritten{30/?}

% TODO: The second blue handwritten entry at the lower right is partly illegible.
% #VALUE_ID: LEX-P0029-V0026
% #HANDWRITTEN: 45/?
\handwritten{45/?}

% TODO: The final blue handwritten Chinese note is illegible.
% #VALUE_ID: LEX-P0029-V0027
% #HANDWRITTEN: [illegible Chinese note]
\handwritten{[illegible Chinese note]}
% LEXOID_PAGE_COMPLETED: 29/70

\newpage

\section{Pailtronic Flowstar IV 测试结果}

\begin{tabular}{@{}p{0.28\linewidth}p{0.62\linewidth}@{}}
日期 &
% #VALUE_ID: LEX-P0030-V0001
% #FIELD_VALUE: 日期
\fieldvalue{12/Dec/2025 13:56:44}
\\
型号 &
% #VALUE_ID: LEX-P0030-V0002
% #FIELD_VALUE: 型号
\fieldvalue{FS06L4R2-2127}
\\
序列号（Serial/result number） &
% #VALUE_ID: LEX-P0030-V0003
% #FIELD_VALUE: 序列号（Serial/result number）
\fieldvalue{238804742464}
\\
测试编号 &
% #VALUE_ID: LEX-P0030-V0004
% #FIELD_VALUE: 测试编号
\fieldvalue{SFM?ES-5225}
\\
测试介质 &
% #VALUE_ID: LEX-P0030-V0005
% #FIELD_VALUE: 测试介质
\fieldvalue{空气}
\\
生产厂家 &
% #VALUE_ID: LEX-P0030-V0006
% #FIELD_VALUE: 生产厂家
\fieldvalue{Shiny?}
\\
生产日期 &
% #VALUE_ID: LEX-P0030-V0007
% #FIELD_VALUE: 生产日期
\fieldvalue{[不清晰]}
\\
工作压力 &
% #VALUE_ID: LEX-P0030-V0008
% #FIELD_VALUE: 工作压力
\fieldvalue{[不清晰]}
\\
工作温度 &
% #VALUE_ID: LEX-P0030-V0009
% #FIELD_VALUE: 工作温度
\fieldvalue{[不清晰]}
\\
阀体材料 &
% #VALUE_ID: LEX-P0030-V0010
% #FIELD_VALUE: 阀体材料
\fieldvalue{[不清晰]}
\\
测试台编号 &
% #VALUE_ID: LEX-P0030-V0011
% #FIELD_VALUE: 测试台编号
\fieldvalue{SFM?ES-5225}
\\
测试压力 &
% #VALUE_ID: LEX-P0030-V0012
% #FIELD_VALUE: 测试压力
\fieldvalue{[不清晰]}
\\
测试介质温度 &
% #VALUE_ID: LEX-P0030-V0013
% #FIELD_VALUE: 测试介质温度
\fieldvalue{[不清晰]}
\end{tabular}

\vspace{1em}

\begin{tabular}{@{}p{0.28\linewidth}p{0.62\linewidth}@{}}
流量测试范围 &
% #VALUE_ID: LEX-P0030-V0014
% #FIELD_VALUE: 流量测试范围
\fieldvalue{5000 mbar}
\\
测试时间 &
% #VALUE_ID: LEX-P0030-V0015
% #FIELD_VALUE: 测试时间
\fieldvalue{12/Dec/2025 14:02:15}
\\
\end{tabular}

\vspace{1em}

签名：

% #VALUE_ID: LEX-P0030-V0016
% #FIELD_VALUE: 签名
% #HANDWRITTEN: 康健
\fieldvalue{\handwritten{康健}}

日期：

% #VALUE_ID: LEX-P0030-V0017
% #FIELD_VALUE: 日期
% #HANDWRITTEN: 2025/12/12
\fieldvalue{\handwritten{2025/12/12}}

% TODO: The report contains several additional small printed field labels and values that are illegible in the supplied page image.

\section{Pailtronic Flowstar IV 测试结果}

\begin{tabular}{@{}p{0.28\linewidth}p{0.62\linewidth}@{}}
日期 &
% #VALUE_ID: LEX-P0030-V0018
% #FIELD_VALUE: 日期
\fieldvalue{12/Dec/2025 13:56:44}
\\
时间 &
% #VALUE_ID: LEX-P0030-V0019
% #FIELD_VALUE: 时间
\fieldvalue{02}
\\
\end{tabular}

\vspace{1em}

% TODO: Figure visible on this page: a flow/pressure chart with vertical scale from 0 to 1800 mbar/min and horizontal scale from 0 to 7000. No source filename is available.
\begin{center}
\fbox{\parbox{0.9\linewidth}{\centering
TODO: Reproduce the visible Flowstar IV test graph here.\\
The graph shows a mostly flat trace followed by a sharp change near the right side.
}}
\end{center}

\vspace{1em}

签名：

% #VALUE_ID: LEX-P0030-V0020
% #FIELD_VALUE: 签名
% #HANDWRITTEN: 康健
\fieldvalue{\handwritten{康健}}

日期：

% #VALUE_ID: LEX-P0030-V0021
% #FIELD_VALUE: 日期
% #HANDWRITTEN: 2025/12/12
\fieldvalue{\handwritten{2025/12/12}}

% #VALUE_ID: LEX-P0030-V0022
% #HANDWRITTEN: 31/35
\handwritten{31/35}
% LEXOID_PAGE_COMPLETED: 30/70

\newpage

\section*{Palltronic Flowstar IV 测试结果}

\begin{tabular}{@{}p{0.24\linewidth}p{0.68\linewidth}@{}}
日期 &
% #VALUE_ID: LEX-P0031-V0001
% #FIELD_VALUE: 日期
\fieldvalue{12/2/2025 14:10:47}
\\
页码 &
% #VALUE_ID: LEX-P0031-V0002
% #FIELD_VALUE: 页码
\fieldvalue{1 of 2}
\\
产品 &
% #VALUE_ID: LEX-P0031-V0003
% #FIELD_VALUE: 产品
% TODO: Product identifier is partially unclear in the scan.
\fieldvalue{FF?CG2}
\\
序列号 &
% #VALUE_ID: LEX-P0031-V0004
% #FIELD_VALUE: 序列号
% TODO: Serial number is partially illegible.
\fieldvalue{23880?A?425?T}
\\
设备 &
% #VALUE_ID: LEX-P0031-V0005
% #FIELD_VALUE: 设备
% TODO: Device value is unclear.
\fieldvalue{??}
\\
测试编号 &
% #VALUE_ID: LEX-P0031-V0006
% #FIELD_VALUE: 测试编号
\fieldvalue{S?F?E?S-2325}
\\
滤芯型号 &
% #VALUE_ID: LEX-P0031-V0007
% #FIELD_VALUE: 滤芯型号
\fieldvalue{Q2}
\\
测试方法 &
% #VALUE_ID: LEX-P0031-V0008
% #FIELD_VALUE: 测试方法
\fieldvalue{Sil?inghou}
\\
开始压力 &
% #VALUE_ID: LEX-P0031-V0009
% #FIELD_VALUE: 开始压力
\fieldvalue{4500 mbar}
\\
终止压力 &
% #VALUE_ID: LEX-P0031-V0010
% #FIELD_VALUE: 终止压力
\fieldvalue{1000 mbar}
\\
测试气体 &
% #VALUE_ID: LEX-P0031-V0011
% #FIELD_VALUE: 测试气体
\fieldvalue{空气}
\\
测试压力 &
% #VALUE_ID: LEX-P0031-V0012
% #FIELD_VALUE: 测试压力
\fieldvalue{7000 mbar}
\end{tabular}

\vspace{1em}
\noindent\textbf{结果}

\begin{tabular}{@{}p{0.30\linewidth}p{0.62\linewidth}@{}}
测试压力范围 &
% #VALUE_ID: LEX-P0031-V0013
% #FIELD_VALUE: 测试压力范围
\fieldvalue{4950 mbar}
\\
测试气体 &
% #VALUE_ID: LEX-P0031-V0014
% #FIELD_VALUE: 测试气体
\fieldvalue{空气}
\\
测试持续时间 &
% #VALUE_ID: LEX-P0031-V0015
% #FIELD_VALUE: 测试持续时间
\fieldvalue{12/2/2025 14:21:08}
\\
签名 &
% #VALUE_ID: LEX-P0031-V0016
% #FIELD_VALUE: 签名
% TODO #HANDWRITTEN: illegible signature; handwritten signature is not clearly readable.
\fieldvalue{\handwritten{[签名]}}
\end{tabular}

\vspace{0.5em}
% #VALUE_ID: LEX-P0031-V0017
% #FIELD_VALUE: 签名日期
% #HANDWRITTEN: 2025/12/12
\fieldvalue{\handwritten{2025/12/12}}

\vspace{2em}
\section*{Palltronic Flowstar IV 测试结果}

\begin{tabular}{@{}p{0.22\linewidth}p{0.70\linewidth}@{}}
日期 &
% #VALUE_ID: LEX-P0031-V0018
% #FIELD_VALUE: 日期
\fieldvalue{12/2/2025 14:10:47}
\\
页码 &
% #VALUE_ID: LEX-P0031-V0019
% #FIELD_VALUE: 页码
\fieldvalue{2 of 2}
\end{tabular}

\vspace{1em}
\noindent
\begin{center}
\setlength{\fboxsep}{8pt}
\fbox{%
\begin{minipage}{0.78\linewidth}
\centering
% TODO: Reproduce the visible flow graph; no source image filename is available.
\begin{picture}(260,150)
\put(35,20){\line(1,0){200}}
\put(35,20){\line(0,1){110}}
\put(35,20){\makebox(0,0)[r]{0}}
\put(35,75){\makebox(0,0)[r]{1400}}
\put(35,130){\makebox(0,0)[r]{2800}}
\put(35,20){\makebox(0,-15)[t]{0}}
\put(85,20){\makebox(0,-15)[t]{1400}}
\put(135,20){\makebox(0,-15)[t]{2800}}
\put(185,20){\makebox(0,-15)[t]{4200}}
\put(235,20){\makebox(0,-15)[t]{5600}}
\put(130,0){\makebox(0,0)[t]{压降（mbar）}}
\put(0,85){\rotatebox{90}{\makebox(0,0)[c]{大气流量（mbar/min）}}}
\put(35,20){\line(1,0){190}}
\put(35,20){\line(1,1){90}}
\end{picture}

\vspace{0.5em}
大气流量（mbar/min）：4.50,\ 900,\ 1350,\ 1800
\end{minipage}%
}
\end{center}

\vspace{1em}
\noindent
% #VALUE_ID: LEX-P0031-V0020
% #FIELD_VALUE: 签名
% TODO #HANDWRITTEN: illegible signature; handwritten signature is not clearly readable.
\fieldvalue{\handwritten{[签名]}}
\hspace{1em}
% #VALUE_ID: LEX-P0031-V0021
% #FIELD_VALUE: 签名日期
% #HANDWRITTEN: 2025/12/12
\fieldvalue{\handwritten{2025/12/12}}

\vfill
% #VALUE_ID: LEX-P0031-V0022
% #HANDWRITTEN: 37/35
\handwritten{37/35}
% LEXOID_PAGE_COMPLETED: 31/70

\newpage

\section{Palltronic Flowstar IV 测试结论}

\begin{tabular}{@{}p{0.30\linewidth}p{0.62\linewidth}@{}}
日期 &
% #VALUE_ID: LEX-P0032-V0001
% #FIELD_VALUE: 日期
\fieldvalue{12/Dec/2025 14:24:33}
\\
页码 &
% #VALUE_ID: LEX-P0032-V0002
% #FIELD_VALUE: 页码
\fieldvalue{1 of 2}
\\
测试结果编号 &
% #VALUE_ID: LEX-P0032-V0003
% #FIELD_VALUE: 测试结果编号
\fieldvalue{FPS041GR/2426 22/17}
\\
测试结果 &
% #VALUE_ID: LEX-P0032-V0004
% #FIELD_VALUE: 测试结果
\fieldvalue{Pass}
\\
测试编号 &
% #VALUE_ID: LEX-P0032-V0005
% #FIELD_VALUE: 测试编号
\fieldvalue{S-MFPS-2225}
\\
测试仪器 &
% #VALUE_ID: LEX-P0032-V0006
% #FIELD_VALUE: 测试仪器
\fieldvalue{SQ}
\\
生产区域 &
% #VALUE_ID: LEX-P0032-V0007
% #FIELD_VALUE: 生产区域
\fieldvalue{Shiyinglou}
\\
生产批次 &
% #VALUE_ID: LEX-P0032-V0008
% #FIELD_VALUE: 生产批次
\fieldvalue{2}
\\
产品名称 &
% #VALUE_ID: LEX-P0032-V0009
% #FIELD_VALUE: 产品名称
\fieldvalue{A?422020512019}
\\
过滤器型号 &
% #VALUE_ID: LEX-P0032-V0010
% #FIELD_VALUE: 过滤器型号
\fieldvalue{A?422020512019}
\\
过滤器序列号 &
% #VALUE_ID: LEX-P0032-V0011
% #FIELD_VALUE: 过滤器序列号
\fieldvalue{?}
\\
过滤器尺寸 &
% #VALUE_ID: LEX-P0032-V0012
% #FIELD_VALUE: 过滤器尺寸
\fieldvalue{S?}
\\
测试流量 &
% #VALUE_ID: LEX-P0032-V0013
% #FIELD_VALUE: 测试流量
\fieldvalue{1000}
\\
测试压力 &
% #VALUE_ID: LEX-P0032-V0014
% #FIELD_VALUE: 测试压力
\fieldvalue{4500 mbar}
\\
最大压差 &
% #VALUE_ID: LEX-P0032-V0015
% #FIELD_VALUE: 最大压差
\fieldvalue{7000 mbar}
\end{tabular}

\bigskip
\hrule
\bigskip

\begin{tabular}{@{}p{0.30\linewidth}p{0.62\linewidth}@{}}
结果 &
% #VALUE_ID: LEX-P0032-V0016
% #FIELD_VALUE: 结果
\fieldvalue{通过}
\\
测试压力 &
% #VALUE_ID: LEX-P0032-V0017
% #FIELD_VALUE: 测试压力
\fieldvalue{4950 mbar}
\\
测试持续时间 &
% #VALUE_ID: LEX-P0032-V0018
% #FIELD_VALUE: 测试持续时间
\fieldvalue{12/Dec/2025 14:35:15}
\end{tabular}

\subsection{注释}

\begin{enumerate}
\item 签名：

% #VALUE_ID: LEX-P0032-V0019
% #FIELD_VALUE: 签名标识
\fieldvalue{A3 34A-42-190g129e}

% #VALUE_ID: LEX-P0032-V0020
% #HANDWRITTEN: 陈保
\handwritten{陈保}

% #VALUE_ID: LEX-P0032-V0021
% #HANDWRITTEN: 2025/1/12
\handwritten{2025/1/12}
\end{enumerate}

\section{Palltronic Flowstar IV 测试曲线}

\begin{tabular}{@{}p{0.30\linewidth}p{0.62\linewidth}@{}}
日期 &
% #VALUE_ID: LEX-P0032-V0022
% #FIELD_VALUE: 日期
\fieldvalue{12/Dec/2025 14:24:33}
\\
页码 &
% #VALUE_ID: LEX-P0032-V0023
% #FIELD_VALUE: 页码
\fieldvalue{2 of 2}
\end{tabular}

\bigskip

% TODO: Recreate the visible test graph with axes and curve; the source image contains no separate figure filename.
\begin{center}
\fbox{\parbox[c][0.28\textheight][c]{0.88\linewidth}{\centering
测试曲线图\\[1ex]
横轴：压力（mbar）\\
纵轴：大气流量（mbar/min）\\[1ex]
（图形及曲线见原页面）
}}
\end{center}

% #VALUE_ID: LEX-P0032-V0024
% #HANDWRITTEN: 陈保
\handwritten{陈保}

% #VALUE_ID: LEX-P0032-V0025
% #HANDWRITTEN: 2025/1/12
\handwritten{2025/1/12}

% #VALUE_ID: LEX-P0032-V0026
% #HANDWRITTEN: 33/35
\handwritten{33/35}
% LEXOID_PAGE_COMPLETED: 32/70

\newpage

\section{Platform FlowStar IV 流量结果}

\textbf{PALL Platform FlowStar IV 流量结果}

\begin{description}
  \item[日期：]
  % #VALUE_ID: LEX-P0033-V0001
  % #FIELD_VALUE: 日期
  \fieldvalue{12/Dec/2025 14:39:01}

  \item[项目号：]
  % #VALUE_ID: LEX-P0033-V0002
  % #FIELD_VALUE: 项目号
  % TODO: The final characters of the project number are unclear in the scan.
  \fieldvalue{FF5016/GF-2/241}

  \item[产品号：]
  % #VALUE_ID: LEX-P0033-V0003
  % #FIELD_VALUE: 产品号
  \fieldvalue{238807442627}

  \item[批号：]
  % #VALUE_ID: LEX-P0033-V0004
  % #FIELD_VALUE: 批号
  % TODO: The batch-number value is unclear in the scan.
  \fieldvalue{[不清晰]}

  \item[流量测试：]
  % #VALUE_ID: LEX-P0033-V0005
  % #FIELD_VALUE: 流量测试
  \fieldvalue{SMFES-2225}

  \item[压力：]
  % #VALUE_ID: LEX-P0033-V0006
  % #FIELD_VALUE: 压力
  \fieldvalue{0.2}

  \item[介质区域：]
  % #VALUE_ID: LEX-P0033-V0007
  % #FIELD_VALUE: 介质区域
  \fieldvalue{C}

  \item[介质：]
  % #VALUE_ID: LEX-P0033-V0008
  % #FIELD_VALUE: 介质
  % TODO: The value is partially illegible.
  \fieldvalue{Shinyou}

  \item[测试项目：]
  % #VALUE_ID: LEX-P0033-V0009
  % #FIELD_VALUE: 测试项目
  % TODO: The field label and value are partially unclear.
  \fieldvalue{7?}

  \item[设备编号：]
  % #VALUE_ID: LEX-P0033-V0010
  % #FIELD_VALUE: 设备编号
  % TODO: The equipment-number value is partially illegible.
  \fieldvalue{A?322?022?019}

  \item[过压测试：]
  % #VALUE_ID: LEX-P0033-V0011
  % #FIELD_VALUE: 过压测试
  \fieldvalue{SMFES-2225}

  \item[结果：]
  % #VALUE_ID: LEX-P0033-V0012
  % #FIELD_VALUE: 结果
  \fieldvalue{通过}

  \item[流量：]
  % #VALUE_ID: LEX-P0033-V0013
  % #FIELD_VALUE: 流量
  \fieldvalue{1000}

  \item[测试压力：]
  % #VALUE_ID: LEX-P0033-V0014
  % #FIELD_VALUE: 测试压力
  \fieldvalue{450~mbar}

  \item[最大压力：]
  % #VALUE_ID: LEX-P0033-V0015
  % #FIELD_VALUE: 最大压力
  \fieldvalue{1000~mbar}
\end{description}

% #VALUE_ID: LEX-P0033-V0016
% #HANDWRITTEN: 承梅 2025.12.12
\fieldvalue{\handwritten{承梅 2025.12.12}}

\section{签字记录}

\begin{description}
  \item[签名：]
  % #VALUE_ID: LEX-P0033-V0017
  % #FIELD_VALUE: 签名
  % #HANDWRITTEN: [签名]
  \fieldvalue{\handwritten{[签名]}}

  \item[签名日期：]
  % #VALUE_ID: LEX-P0033-V0018
  % #FIELD_VALUE: 签名日期
  \fieldvalue{29/Dec/2025 14:49:51}
\end{description}

\section{流量曲线}

\textbf{PALL Platform FlowStar IV 流量结果}

\begin{description}
  \item[日期：]
  % #VALUE_ID: LEX-P0033-V0019
  % #FIELD_VALUE: 日期
  \fieldvalue{12/Dec/2025 14:39:01}

  \item[项目编号：]
  % #VALUE_ID: LEX-P0033-V0020
  % #FIELD_VALUE: 项目编号
  % TODO: The project-number value is unclear in the scan.
  \fieldvalue{2 02?}
\end{description}

\begin{center}
\setlength{\fboxsep}{6pt}
\begin{minipage}{0.78\linewidth}
\centering
\textbf{大气体流量（mbar/min）}

\vspace{0.3em}
\begin{tabular}{c|c}
\hline
\multicolumn{2}{c}{流量曲线} \\
\hline
\multicolumn{2}{c}{纵轴：0,\ 450,\ 900,\ 1350,\ 1800} \\
\hline
\multicolumn{2}{c}{横轴：0,\ 1000,\ 2000,\ 3000,\ 4000,\ 5000,\ 6000,\ 7000} \\
\hline
\end{tabular}

\vspace{0.5em}
% TODO: Reconstruct the plotted curve from the visible graph.
\end{minipage}
\end{center}

% #VALUE_ID: LEX-P0033-V0021
% #HANDWRITTEN: 承梅 2025.12.12
\fieldvalue{\handwritten{承梅 2025.12.12}}
% LEXOID_PAGE_COMPLETED: 33/70

\newpage

\section*{Palltronic Flowstar IV 测试结果}

\begin{tabular}{@{}p{0.28\linewidth}p{0.62\linewidth}@{}}
日期 &
% #VALUE_ID: LEX-P0034-V0001
% #FIELD_VALUE: 日期
\fieldvalue{12/20/2025 13:45:01}
\\[2pt]
页码 &
% #VALUE_ID: LEX-P0034-V0002
% #FIELD_VALUE: 页码
\fieldvalue{1 of 2}
\\[2pt]
Serial/Result number &
% #VALUE_ID: LEX-P0034-V0003
% #FIELD_VALUE: Serial/Result number
% #TODO: Serial/result number is partially illegible in the scan.
\fieldvalue{2B?804?4268}
\\[2pt]
测试产品 &
% #VALUE_ID: LEX-P0034-V0004
% #FIELD_VALUE: 测试产品
\fieldvalue{Air}
\\[2pt]
过滤器型号 &
% #VALUE_ID: LEX-P0034-V0005
% #FIELD_VALUE: 过滤器型号
\fieldvalue{SFPES-225}
\\[2pt]
过滤器类型 &
% #VALUE_ID: LEX-P0034-V0006
% #FIELD_VALUE: 过滤器类型
% #TODO: Filter type is partially illegible.
\fieldvalue{Zun?e}
\\[2pt]
产品编号 &
% #VALUE_ID: LEX-P0034-V0007
% #FIELD_VALUE: 产品编号
% #TODO: Product number is partially illegible.
\fieldvalue{A3?}
\\[2pt]
测试气体 &
% #VALUE_ID: LEX-P0034-V0008
% #FIELD_VALUE: 测试气体
\fieldvalue{Air}
\\[2pt]
测试流量 &
% #VALUE_ID: LEX-P0034-V0009
% #FIELD_VALUE: 测试流量
\fieldvalue{1000}
\\[2pt]
测试压力 &
% #VALUE_ID: LEX-P0034-V0010
% #FIELD_VALUE: 测试压力
\fieldvalue{5050 mbar}
\\[2pt]
润湿液 &
% #VALUE_ID: LEX-P0034-V0011
% #FIELD_VALUE: 润湿液
% #TODO: Wetting-liquid value is partially illegible.
\fieldvalue{SF?}
\\[2pt]
最大压力 &
% #VALUE_ID: LEX-P0034-V0012
% #FIELD_VALUE: 最大压力
\fieldvalue{7000 mbar}
\\[2pt]
测试结束时间 &
% #VALUE_ID: LEX-P0034-V0013
% #FIELD_VALUE: 测试结束时间
\fieldvalue{12/20/2025 15:04:44}
\end{tabular}

\vspace{1em}

\noindent 签名：%
% #VALUE_ID: LEX-P0034-V0014
% #FIELD_VALUE: 签名
% #HANDWRITTEN: 手写签名
\fieldvalue{\handwritten{手写签名}}

\noindent 日期：%
% #VALUE_ID: LEX-P0034-V0015
% #FIELD_VALUE: 日期
% #HANDWRITTEN: 2025/12/12
\fieldvalue{\handwritten{2025/12/12}}

\vspace{1em}

\noindent
% #TODO: The lower portion of the first report contains a handwritten entry beside the signature/date area; the exact printed field label is unclear.
% #VALUE_ID: LEX-P0034-V0016
% #FIELD_VALUE: 未辨识的手写字段
% #TODO #HANDWRITTEN: - A2?B?3?; handwriting is unclear.
\fieldvalue{\handwritten{- A2?B?3?}}

\noindent
% #VALUE_ID: LEX-P0034-V0017
% #FIELD_VALUE: 日期
% #HANDWRITTEN: 2025/12/12
\fieldvalue{\handwritten{2025/12/12}}

\bigskip
\hrule
\bigskip

\section*{Palltronic Flowstar IV 测试结果}

\begin{tabular}{@{}p{0.28\linewidth}p{0.62\linewidth}@{}}
日期 &
% #VALUE_ID: LEX-P0034-V0018
% #FIELD_VALUE: 日期
\fieldvalue{12/20/2025 14:54:01}
\\[2pt]
页码 &
% #VALUE_ID: LEX-P0034-V0019
% #FIELD_VALUE: 页码
\fieldvalue{2 of 2}
\end{tabular}

\vspace{1em}

\begin{center}
\setlength{\fboxsep}{6pt}
\fbox{%
\begin{minipage}{0.82\linewidth}
\centering
\textbf{大气流量 (mbar/min)}

\vspace{0.5em}
\begin{picture}(260,125)
\put(28,12){\line(1,0){210}}
\put(28,12){\line(0,1){95}}
\put(28,107){\line(1,0){210}}
\put(238,12){\line(0,1){95}}
\put(22,7){\makebox(0,0)[r]{0}}
\put(72,7){\makebox(0,0)[c]{1400}}
\put(120,7){\makebox(0,0)[c]{2800}}
\put(168,7){\makebox(0,0)[c]{4200}}
\put(216,7){\makebox(0,0)[c]{5600}}
\put(238,7){\makebox(0,0)[c]{7000}}
\put(18,107){\makebox(0,0)[r]{1800}}
\put(18,83){\makebox(0,0)[r]{1350}}
\put(18,59){\makebox(0,0)[r]{900}}
\put(18,35){\makebox(0,0)[r]{450}}
\put(130,-3){\makebox(0,0)[c]{流量 (mbar/min)}}
\put(4,60){\rotatebox{90}{\makebox(0,0)[c]{大气流量}}}
% TODO: The plotted curve is faint; only its approximate descending shape is represented.
\qbezier(28,105)(100,105)(150,103)
\qbezier(150,103)(180,101)(185,82)
\qbezier(185,82)(188,68)(194,62)
\end{picture}
\end{minipage}}
\end{center}

\vspace{1em}

\noindent 签名：%
% #VALUE_ID: LEX-P0034-V0020
% #FIELD_VALUE: 签名
% #HANDWRITTEN: 手写签名
\fieldvalue{\handwritten{手写签名}}

\noindent 日期：%
% #VALUE_ID: LEX-P0034-V0021
% #FIELD_VALUE: 日期
% #HANDWRITTEN: 2025/12/12
\fieldvalue{\handwritten{2025/12/12}}

\hfill
% #VALUE_ID: LEX-P0034-V0022
% #FIELD_VALUE: 手写页码标记
% #HANDWRITTEN: 35/35
\fieldvalue{\handwritten{35/35}}
% LEXOID_PAGE_COMPLETED: 34/70

\newpage

\begin{center}
{\Large 艾米迈托赛注射液（Z2）无菌工艺模拟试验\par}
\vspace{0.4em}
{\Large 生产记录\par}
\end{center}

\vspace{1.5em}

\noindent
\begin{tabular}{@{}l l@{}}
工序名称：&
% #VALUE_ID: LEX-P0035-V0001
% #FIELD_VALUE: 工序名称
\fieldvalue{细胞分裂}
\\[1.2em]
规格：&
% #VALUE_ID: LEX-P0035-V0002
% #FIELD_VALUE: 规格
\fieldvalue{\textbf{12.5$\pm$0.5ml}}
\\[1.2em]
生产批号：&
% #VALUE_ID: LEX-P0035-V0003
% #FIELD_VALUE: 生产批号
% #HANDWRITTEN: APS2025D016 (Z2)
\fieldvalue{\handwritten{APS2025D016 (Z2)}}
\\[1.2em]
批量：&
% #VALUE_ID: LEX-P0035-V0004
% #FIELD_VALUE: 批量
% #HANDWRITTEN: $>$240
\fieldvalue{\handwritten{$>$240}}\ \ 瓶
\\[1.2em]
生产日期：&
% #VALUE_ID: LEX-P0035-V0005
% #FIELD_VALUE: 生产日期（起始日期）
% #HANDWRITTEN: 2025年12月13日
\fieldvalue{\handwritten{2025年12月13日}}
\quad 至 \quad
% #VALUE_ID: LEX-P0035-V0006
% #FIELD_VALUE: 生产日期（结束日期）
% #HANDWRITTEN: 2027年12月13日
\fieldvalue{\handwritten{2027年12月13日}}
\end{tabular}

\vfill

\begin{center}
第 1 页 共 23 页
\end{center}
% LEXOID_PAGE_COMPLETED: 35/70

\newpage

\noindent
\begin{tabular}{|p{0.29\linewidth}|p{0.20\linewidth}|p{0.25\linewidth}|p{0.12\linewidth}|}
\hline
品名 & 规格 & 生产批号 & 批量 \\
\hline
% #VALUE_ID: LEX-P0036-V0001
% #FIELD_VALUE: 品名
\fieldvalue{艾米迈托霉钠注射液(2乙)无菌工艺模拟试验}
&
% #VALUE_ID: LEX-P0036-V0002
% #FIELD_VALUE: 规格
\fieldvalue{12.5$\pm$0.5ml}
&
% #VALUE_ID: LEX-P0036-V0003
% #FIELD_VALUE: 生产批号
% #TODO #HANDWRITTEN: APS250501(062); handwriting is partially illegible
\fieldvalue{\handwritten{APS250501(062)}}
&
% #VALUE_ID: LEX-P0036-V0004
% #FIELD_VALUE: 批量
% #HANDWRITTEN: 2740
\fieldvalue{\handwritten{2740}} 袋
\\
\hline
\end{tabular}

\section{生产工序五：细胞分装}

\subsection{1.\quad 生产前确认}

\noindent
\begin{tabular}{|p{0.60\linewidth}|p{0.30\linewidth}|}
\hline
项目 & 操作记录 \\
\hline
操作人员按《人员进出洁净生产区管理规程》要求进入洁净生产车间； &
操作员名称及编码：

% #VALUE_ID: LEX-P0036-V0005
% #FIELD_VALUE: 操作员名称及编码（第一人）
% #TODO #HANDWRITTEN: 张桂芳 5226; handwriting is partially illegible
\fieldvalue{\handwritten{张桂芳 5226}}

% #VALUE_ID: LEX-P0036-V0006
% #FIELD_VALUE: 操作员名称及编码（第二人）
% #TODO #HANDWRITTEN: 张桂娟 5230; handwriting is partially illegible
\fieldvalue{\handwritten{张桂娟 5230}}
\\
\hline
确认生产区操作间在清场有效期内（C级日清洁效期3天）； &
是否在清洁有效期内：

% #VALUE_ID: LEX-P0036-V0007
% #FIELD_VALUE: 清洁有效期内--是
% #HANDWRITTEN: 2025.12.15
\fieldvalue{\handwritten{2025.12.15}}

% #VALUE_ID: LEX-P0036-V0008
% #FIELD_VALUE: 清洁有效期内--是
\fieldvalue{\checkboxfield{checked}} 是

% #VALUE_ID: LEX-P0036-V0009
% #FIELD_VALUE: 清洁有效期内--否
\fieldvalue{\checkboxfield{unchecked}} 否
\\
\hline
确认生产区内无上批产品遗留物、文件及与本产品生产无关物料； &
是否符合要求：

% #VALUE_ID: LEX-P0036-V0010
% #FIELD_VALUE: 无遗留物--是
\fieldvalue{\checkboxfield{checked}} 是

% #VALUE_ID: LEX-P0036-V0011
% #FIELD_VALUE: 无遗留物--否
\fieldvalue{\checkboxfield{unchecked}} 否
\\
\hline
检查发放的批记录及其它相关配套文件均与生产指令相对应，且均受控； &
是否符合要求：

% #VALUE_ID: LEX-P0036-V0012
% #FIELD_VALUE: 文件受控--是
\fieldvalue{\checkboxfield{checked}} 是

% #VALUE_ID: LEX-P0036-V0013
% #FIELD_VALUE: 文件受控--否
\fieldvalue{\checkboxfield{unchecked}} 否
\\
\hline
确认仪器完好，且在验证有效期内，计量器具有“合格证”且在有效期内，详见《主要设备/仪表确认表》； &
是否符合要求：

% #VALUE_ID: LEX-P0036-V0014
% #FIELD_VALUE: 仪器设备符合要求--是
\fieldvalue{\checkboxfield{checked}} 是

% #VALUE_ID: LEX-P0036-V0015
% #FIELD_VALUE: 仪器设备符合要求--否
\fieldvalue{\checkboxfield{unchecked}} 否
\\
\hline
将操作间状态标识更换为“生产加工状态卡”，写明品名、规格、批号等，并确认与本批生产指令一致； &
是否符合要求：

% #VALUE_ID: LEX-P0036-V0016
% #FIELD_VALUE: 状态标识符合要求--是
\fieldvalue{\checkboxfield{checked}} 是

% #VALUE_ID: LEX-P0036-V0017
% #FIELD_VALUE: 状态标识符合要求--否
\fieldvalue{\checkboxfield{unchecked}} 否
\\
\hline
核对本次生产用物料耗材、溶液和试剂的品名、规格、批号、数量，检查物料外包装的完整性以及是否在有效期/复验期内，并记录相关信息，详见《物料确认表》和《溶液/试剂确认表》； &
是否符合要求：

% #VALUE_ID: LEX-P0036-V0018
% #FIELD_VALUE: 物料耗材符合要求--是
\fieldvalue{\checkboxfield{checked}} 是

% #VALUE_ID: LEX-P0036-V0019
% #FIELD_VALUE: 物料耗材符合要求--否
\fieldvalue{\checkboxfield{unchecked}} 否
\\
\hline
备注：

% #VALUE_ID: LEX-P0036-V0020
% #FIELD_VALUE: 备注
% #HANDWRITTEN: /
\fieldvalue{\handwritten{/}}
&
\\
\hline
\end{tabular}

\vspace{0.5em}

\noindent
确认人/日期：

% #VALUE_ID: LEX-P0036-V0021
% #FIELD_VALUE: 确认人
% #TODO #HANDWRITTEN: 赵俊峰; handwriting is partially illegible
\fieldvalue{\handwritten{赵俊峰}}

% #VALUE_ID: LEX-P0036-V0022
% #FIELD_VALUE: 确认日期
% #HANDWRITTEN: 2025.12.13
\fieldvalue{\handwritten{2025.12.13}}
\hfill
复核人/日期：

% #VALUE_ID: LEX-P0036-V0023
% #FIELD_VALUE: 复核人
% #TODO #HANDWRITTEN: 赵影; handwriting is partially illegible
\fieldvalue{\handwritten{赵影}}

% #VALUE_ID: LEX-P0036-V0024
% #FIELD_VALUE: 复核日期
% #HANDWRITTEN: 2025.12.15
\fieldvalue{\handwritten{2025.12.15}}
% LEXOID_PAGE_COMPLETED: 36/70

\newpage

\begin{center}
\begin{tabular}{|p{0.34\linewidth}|p{0.20\linewidth}|p{0.20\linewidth}|p{0.12\linewidth}|}
\hline
品名 & 规格 & 生产批号 & 批量 \\
\hline
% #VALUE_ID: LEX-P0037-V0001
% #FIELD_VALUE: 品名
\fieldvalue{艾米迈克霉素注射液(2Z)无菌工艺模拟试验}
&
% #VALUE_ID: LEX-P0037-V0002
% #FIELD_VALUE: 规格
\fieldvalue{12.5$\pm$0.5 ml}
&
% #VALUE_ID: LEX-P0037-V0003
% #FIELD_VALUE: 生产批号
% #TODO #HANDWRITTEN: AP?25?01b(??)24?0; handwriting unclear
\fieldvalue{\handwritten{AP?25?01b(??)24?0}}
&
% #VALUE_ID: LEX-P0037-V0004
% #FIELD_VALUE: 批量
\fieldvalue{袋}
\\
\hline
\end{tabular}
\end{center}

\subsection*{1.1\quad 主要设备/仪器确认}

\begin{center}
\small
\begin{tabular}{|c|p{0.18\linewidth}|p{0.17\linewidth}|p{0.16\linewidth}|p{0.18\linewidth}|p{0.13\linewidth}|}
\hline
序号 & 设备名称 & 型号 & 设备编号 & 检查项 & 确认结果 \\
\hline
1 & 电容性检测仪 & CB-AirDetecter &
% #VALUE_ID: LEX-P0037-V0005
% #FIELD_VALUE: 设备编号
% #HANDWRITTEN: 112363?
\fieldvalue{\handwritten{112363?}} &
&
% #VALUE_ID: LEX-P0037-V0006
% #FIELD_VALUE: 确认结果--是
\fieldvalue{\checkboxfield{checked}}是\quad
% #VALUE_ID: LEX-P0037-V0007
% #FIELD_VALUE: 确认结果--否
\fieldvalue{\checkboxfield{unchecked}}否
\\
\hline
2 & 电容性检测仪 & CB-AirDetecter &
% #VALUE_ID: LEX-P0037-V0008
% #FIELD_VALUE: 设备编号
% #HANDWRITTEN: 1123624
\fieldvalue{\handwritten{1123624}} &
&
% #VALUE_ID: LEX-P0037-V0009
% #FIELD_VALUE: 确认结果--是
\fieldvalue{\checkboxfield{checked}}是\quad
% #VALUE_ID: LEX-P0037-V0010
% #FIELD_VALUE: 确认结果--否
\fieldvalue{\checkboxfield{unchecked}}否
\\
\hline
3 & 高精度快速分装工作站（简称 Pack Pro） & Gentle P-Pac Pro &
% #VALUE_ID: LEX-P0037-V0011
% #FIELD_VALUE: 设备编号
% #HANDWRITTEN: 1123715
\fieldvalue{\handwritten{1123715}} &
&
% #VALUE_ID: LEX-P0037-V0012
% #FIELD_VALUE: 确认结果--是
\fieldvalue{\checkboxfield{checked}}是\quad
% #VALUE_ID: LEX-P0037-V0013
% #FIELD_VALUE: 确认结果--否
\fieldvalue{\checkboxfield{unchecked}}否
\\
\hline
4 & 接管机 & TSCD II &
% #VALUE_ID: LEX-P0037-V0014
% #FIELD_VALUE: 设备编号
% #HANDWRITTEN: 1123704
\fieldvalue{\handwritten{1123704}} &
&
% #VALUE_ID: LEX-P0037-V0015
% #FIELD_VALUE: 确认结果--是
\fieldvalue{\checkboxfield{checked}}是\quad
% #VALUE_ID: LEX-P0037-V0016
% #FIELD_VALUE: 确认结果--否
\fieldvalue{\checkboxfield{unchecked}}否
\\
\hline
5 & 接管机 & TSCD II &
% #VALUE_ID: LEX-P0037-V0017
% #FIELD_VALUE: 设备编号
% #HANDWRITTEN: 112?306
\fieldvalue{\handwritten{112?306}} &
&
% #VALUE_ID: LEX-P0037-V0018
% #FIELD_VALUE: 确认结果--是
\fieldvalue{\checkboxfield{checked}}是\quad
% #VALUE_ID: LEX-P0037-V0019
% #FIELD_VALUE: 确认结果--否
\fieldvalue{\checkboxfield{unchecked}}否
\\
\hline
6 & 接管机 & TSCD II &
% #VALUE_ID: LEX-P0037-V0020
% #FIELD_VALUE: 设备编号
% #HANDWRITTEN: 1123627
\fieldvalue{\handwritten{1123627}} &
&
% #VALUE_ID: LEX-P0037-V0021
% #FIELD_VALUE: 确认结果--是
\fieldvalue{\checkboxfield{checked}}是\quad
% #VALUE_ID: LEX-P0037-V0022
% #FIELD_VALUE: 确认结果--否
\fieldvalue{\checkboxfield{unchecked}}否
\\
\hline
7 & 封管机 & Biosealer CR 6AA &
% #VALUE_ID: LEX-P0037-V0023
% #FIELD_VALUE: 设备编号
% #HANDWRITTEN: 1120006
\fieldvalue{\handwritten{1120006}} &
\multirow{6}{*}{\raggedright
1.设备应有“完好”标识。\\
2.设备在验证有效期内。\\
3.计量器具有“合格证”且在有效期内。}
&
% #VALUE_ID: LEX-P0037-V0024
% #FIELD_VALUE: 确认结果--是
\fieldvalue{\checkboxfield{checked}}是\quad
% #VALUE_ID: LEX-P0037-V0025
% #FIELD_VALUE: 确认结果--否
\fieldvalue{\checkboxfield{unchecked}}否
\\
\hline
8 & 封管机 & Biosealer CR 6AA &
% #VALUE_ID: LEX-P0037-V0026
% #FIELD_VALUE: 设备编号
% #HANDWRITTEN: 112?009
\fieldvalue{\handwritten{112?009}} &
&
% #VALUE_ID: LEX-P0037-V0027
% #FIELD_VALUE: 确认结果--是
\fieldvalue{\checkboxfield{checked}}是\quad
% #VALUE_ID: LEX-P0037-V0028
% #FIELD_VALUE: 确认结果--否
\fieldvalue{\checkboxfield{unchecked}}否
\\
\hline
9 & 封管机 & Biosealer CR 6AA &
% #VALUE_ID: LEX-P0037-V0029
% #FIELD_VALUE: 设备编号
% #HANDWRITTEN: 112?002
\fieldvalue{\handwritten{112?002}} &
&
% #VALUE_ID: LEX-P0037-V0030
% #FIELD_VALUE: 确认结果--是
\fieldvalue{\checkboxfield{checked}}是\quad
% #VALUE_ID: LEX-P0037-V0031
% #FIELD_VALUE: 确认结果--否
\fieldvalue{\checkboxfield{unchecked}}否
\\
\hline
10 & 封管机 & Biosealer CR 6AA &
% #VALUE_ID: LEX-P0037-V0032
% #FIELD_VALUE: 设备编号
% #HANDWRITTEN: 112000?
\fieldvalue{\handwritten{112000?}} &
&
% #VALUE_ID: LEX-P0037-V0033
% #FIELD_VALUE: 确认结果--是
\fieldvalue{\checkboxfield{checked}}是\quad
% #VALUE_ID: LEX-P0037-V0034
% #FIELD_VALUE: 确认结果--否
\fieldvalue{\checkboxfield{unchecked}}否
\\
\hline
11 & 封管机 & Biosealer CR 6AA &
% #VALUE_ID: LEX-P0037-V0035
% #FIELD_VALUE: 设备编号
% #HANDWRITTEN: 112?07
\fieldvalue{\handwritten{112?07}} &
&
% #VALUE_ID: LEX-P0037-V0036
% #FIELD_VALUE: 确认结果--是
\fieldvalue{\checkboxfield{checked}}是\quad
% #VALUE_ID: LEX-P0037-V0037
% #FIELD_VALUE: 确认结果--否
\fieldvalue{\checkboxfield{unchecked}}否
\\
\hline
12 & 封管机 & Biosealer CR 6AA &
% #VALUE_ID: LEX-P0037-V0038
% #FIELD_VALUE: 设备编号
% #HANDWRITTEN: 112?0?
\fieldvalue{\handwritten{112?0?}} &
&
% #VALUE_ID: LEX-P0037-V0039
% #FIELD_VALUE: 确认结果--是
\fieldvalue{\checkboxfield{checked}}是\quad
% #VALUE_ID: LEX-P0037-V0040
% #FIELD_VALUE: 确认结果--否
\fieldvalue{\checkboxfield{unchecked}}否
\\
\hline
13 & 自动配液仪 & Gentle Mixer &
% #VALUE_ID: LEX-P0037-V0041
% #FIELD_VALUE: 设备编号
% #HANDWRITTEN: 123714
\fieldvalue{\handwritten{123714}} &
&
% #VALUE_ID: LEX-P0037-V0042
% #FIELD_VALUE: 确认结果--是
\fieldvalue{\checkboxfield{checked}}是\quad
% #VALUE_ID: LEX-P0037-V0043
% #FIELD_VALUE: 确认结果--否
\fieldvalue{\checkboxfield{unchecked}}否
\\
\hline
14 & 自动配液仪 & Gentle Mixer & & &
% #VALUE_ID: LEX-P0037-V0044
% #FIELD_VALUE: 确认结果--是
\fieldvalue{\checkboxfield{checked}}是\quad
% #VALUE_ID: LEX-P0037-V0045
% #FIELD_VALUE: 确认结果--否
\fieldvalue{\checkboxfield{unchecked}}否
\\
\hline
15 & & & & &
% #VALUE_ID: LEX-P0037-V0046
% #FIELD_VALUE: 确认结果--是
\fieldvalue{\checkboxfield{unchecked}}是\quad
% #VALUE_ID: LEX-P0037-V0047
% #FIELD_VALUE: 确认结果--否
\fieldvalue{\checkboxfield{unchecked}}否
\\
\hline
16 & & & & &
% #VALUE_ID: LEX-P0037-V0048
% #FIELD_VALUE: 确认结果--是
\fieldvalue{\checkboxfield{unchecked}}是\quad
% #VALUE_ID: LEX-P0037-V0049
% #FIELD_VALUE: 确认结果--否
\fieldvalue{\checkboxfield{unchecked}}否
\\
\hline
17 & & & & &
% #VALUE_ID: LEX-P0037-V0050
% #FIELD_VALUE: 确认结果--是
\fieldvalue{\checkboxfield{unchecked}}是\quad
% #VALUE_ID: LEX-P0037-V0051
% #FIELD_VALUE: 确认结果--否
\fieldvalue{\checkboxfield{unchecked}}否
\\
\hline
18 & & & & &
% #VALUE_ID: LEX-P0037-V0052
% #FIELD_VALUE: 确认结果--是
\fieldvalue{\checkboxfield{unchecked}}是\quad
% #VALUE_ID: LEX-P0037-V0053
% #FIELD_VALUE: 确认结果--否
\fieldvalue{\checkboxfield{unchecked}}否
\\
\hline
\end{tabular}
\end{center}

备注：\hrulefill

% #VALUE_ID: LEX-P0037-V0054
% #FIELD_VALUE: 确认人
% #TODO #HANDWRITTEN: 赵?玲; handwriting unclear
\text{确认人：}\fieldvalue{\handwritten{赵?玲}}
\qquad
% #VALUE_ID: LEX-P0037-V0055
% #FIELD_VALUE: 确认日期
% #HANDWRITTEN: 2025.12.13
\fieldvalue{\handwritten{2025.12.13}}

% #VALUE_ID: LEX-P0037-V0056
% #FIELD_VALUE: 复核人
% #TODO #HANDWRITTEN: 孙?昊; handwriting unclear
\text{复核人：}\fieldvalue{\handwritten{孙?昊}}
\qquad
% #VALUE_ID: LEX-P0037-V0057
% #FIELD_VALUE: 复核日期
% #HANDWRITTEN: 2025.12.13
\fieldvalue{\handwritten{2025.12.13}}

\subsection*{1.2\quad 物料确认表}

\begin{center}
\small
\begin{tabular}{|c|p{0.25\linewidth}|p{0.13\linewidth}|p{0.12\linewidth}|p{0.12\linewidth}|c|p{0.16\linewidth}|}
\hline
序号 & 名称 & 品牌 & 规格 & 批号 & 数量 & 是否符合要求 \\
\hline
1 & 细胞冻存袋 & Origen & 4 个/包 & & &
% #VALUE_ID: LEX-P0037-V0058
% #FIELD_VALUE: 是否符合要求--是
\fieldvalue{\checkboxfield{checked}}是\quad
% #VALUE_ID: LEX-P0037-V0059
% #FIELD_VALUE: 是否符合要求--否
\fieldvalue{\checkboxfield{unchecked}}否
\\
\hline
2 & 焊接片 & 泰尔康 & 70 个/包 & & &
% #VALUE_ID: LEX-P0037-V0060
% #FIELD_VALUE: 是否符合要求--是
\fieldvalue{\checkboxfield{unchecked}}是\quad
% #VALUE_ID: LEX-P0037-V0061
% #FIELD_VALUE: 是否符合要求--否
\fieldvalue{\checkboxfield{unchecked}}否
\\
\hline
3 & 产品标签 & & & & &
% #VALUE_ID: LEX-P0037-V0062
% #FIELD_VALUE: 是否符合要求--是
\fieldvalue{\checkboxfield{unchecked}}是\quad
% #VALUE_ID: LEX-P0037-V0063
% #FIELD_VALUE: 是否符合要求--否
\fieldvalue{\checkboxfield{unchecked}}否
\\
\hline
4 & 一次性分装套件（PacPro-20） & 泰析 &
% #VALUE_ID: LEX-P0037-V0064
% #FIELD_VALUE: 规格
% #HANDWRITTEN: 10/箱
\fieldvalue{\handwritten{10/箱}} &
% #TODO #HANDWRITTEN: PacPro-20; blue handwritten correction/notation across the row
\handwritten{PacPro-20} &
% #VALUE_ID: LEX-P0037-V0065
% #FIELD_VALUE: 数量
% #HANDWRITTEN: 24
\fieldvalue{\handwritten{24}} &
% #VALUE_ID: LEX-P0037-V0066
% #FIELD_VALUE: 是否符合要求--是
\fieldvalue{\checkboxfield{checked}}是\quad
% #VALUE_ID: LEX-P0037-V0067
% #FIELD_VALUE: 是否符合要求--否
\fieldvalue{\checkboxfield{unchecked}}否
\\
\hline
5 & 一次性分装套件（PacPro-50） & 泰析 & & & &
% #VALUE_ID: LEX-P0037-V0068
% #FIELD_VALUE: 是否符合要求--是
\fieldvalue{\checkboxfield{unchecked}}是\quad
% #VALUE_ID: LEX-P0037-V0069
% #FIELD_VALUE: 是否符合要求--否
\fieldvalue{\checkboxfield{unchecked}}否
\\
\hline
6 & 一次性润高套件（10L） & 泰析 & 1 个/盒 &
% #VALUE_ID: LEX-P0037-V0070
% #FIELD_VALUE: 批号
% #HANDWRITTEN: H1080/20250709
\fieldvalue{\handwritten{H1080/20250709}} &
% #VALUE_ID: LEX-P0037-V0071
% #FIELD_VALUE: 数量
% #HANDWRITTEN: 1
\fieldvalue{\handwritten{1}} &
% #VALUE_ID: LEX-P0037-V0072
% #FIELD_VALUE: 是否符合要求--是
\fieldvalue{\checkboxfield{checked}}是\quad
% #VALUE_ID: LEX-P0037-V0073
% #FIELD_VALUE: 是否符合要求--否
\fieldvalue{\checkboxfield{unchecked}}否
\\
\hline
\end{tabular}
\end{center}

\begin{center}
第 3 页 共 23 页
\end{center}
% LEXOID_PAGE_COMPLETED: 37/70

\newpage

\begin{center}
\begin{tabular}{|p{0.18\linewidth}|p{0.20\linewidth}|p{0.25\linewidth}|p{0.17\linewidth}|}
\hline
品名 & 规格 & 生产批号 & 批量\\
\hline
艾迈托托套针注液(Z2)无菌工艺模拟试验 & 12.5$\pm$0.5\,ml &
% #VALUE_ID: LEX-P0038-V0001
% #FIELD_VALUE: 生产批号
% #TODO #HANDWRITTEN: A?2025 201?62?; handwriting partially illegible
\fieldvalue{\handwritten{A?2025 201?62?}} &
% #VALUE_ID: LEX-P0038-V0002
% #FIELD_VALUE: 批量
% #TODO #HANDWRITTEN: 2?40; handwriting partially illegible
\fieldvalue{\handwritten{2?40}} 袋\\
\hline
\end{tabular}
\end{center}

\vspace{0.5em}

\begin{center}
\small
\begin{tabular}{|p{0.06\linewidth}|p{0.30\linewidth}|p{0.06\linewidth}|p{0.14\linewidth}|p{0.12\linewidth}|p{0.11\linewidth}|p{0.14\linewidth}|}
\hline
序号 & 品名 & 规格 & 包装规格 & 生产批号 & 批量 & 判定\\
\hline
7 & 一次性分装套件 (PacPro-50) & / & 1Pcs/包 & & &
% #VALUE_ID: LEX-P0038-V0003
% #FIELD_VALUE: 是
\fieldvalue{\checkboxfield{unchecked}} 是\quad
% #VALUE_ID: LEX-P0038-V0004
% #FIELD_VALUE: 否
\fieldvalue{\checkboxfield{checked}} 否\\
\hline
8 & 一次性分装套件 (PacPro-50) & / & 1Pcs/包 & & &
% #VALUE_ID: LEX-P0038-V0005
% #FIELD_VALUE: 是
\fieldvalue{\checkboxfield{unchecked}} 是\quad
% #VALUE_ID: LEX-P0038-V0006
% #FIELD_VALUE: 否
\fieldvalue{\checkboxfield{unchecked}} 否\\
\hline
9 & 一次性分装套件 (PacPro-50) & / & 1Pcs/包 & & &
% #VALUE_ID: LEX-P0038-V0007
% #FIELD_VALUE: 是
\fieldvalue{\checkboxfield{unchecked}} 是\quad
% #VALUE_ID: LEX-P0038-V0008
% #FIELD_VALUE: 否
\fieldvalue{\checkboxfield{unchecked}} 否\\
\hline
10 & 一次性分装套件 (PacPro-50) & / & 1Pcs/包 & & &
% #VALUE_ID: LEX-P0038-V0009
% #FIELD_VALUE: 是
\fieldvalue{\checkboxfield{unchecked}} 是\quad
% #VALUE_ID: LEX-P0038-V0010
% #FIELD_VALUE: 否
\fieldvalue{\checkboxfield{unchecked}} 否\\
\hline
11 & 一次性分装套件 (PacPro-20) & / & 1Pcs/包 & & &
% #VALUE_ID: LEX-P0038-V0011
% #FIELD_VALUE: 是
\fieldvalue{\checkboxfield{unchecked}} 是\quad
% #VALUE_ID: LEX-P0038-V0012
% #FIELD_VALUE: 否
\fieldvalue{\checkboxfield{unchecked}} 否\\
\hline
12 & 一次性分装套件 (PacPro-20) & / & 1Pcs/包 &
% #VALUE_ID: LEX-P0038-V0013
% #FIELD_VALUE: 生产批号/备注
% #TODO #HANDWRITTEN: A?0?12; handwriting partially illegible
\fieldvalue{\handwritten{A?0?12}} & &
% #VALUE_ID: LEX-P0038-V0014
% #FIELD_VALUE: 是
\fieldvalue{\checkboxfield{unchecked}} 是\quad
% #VALUE_ID: LEX-P0038-V0015
% #FIELD_VALUE: 否
\fieldvalue{\checkboxfield{unchecked}} 否\\
\hline
13 & 一次性分装套件 (PacPro-20) & / & 1Pcs/包 &
% #VALUE_ID: LEX-P0038-V0016
% #FIELD_VALUE: 生产批号/备注
% #TODO #HANDWRITTEN: 数字; handwritten text overlaps the printed row
\fieldvalue{\handwritten{数字}} & &
% #VALUE_ID: LEX-P0038-V0017
% #FIELD_VALUE: 是
\fieldvalue{\checkboxfield{unchecked}} 是\quad
% #VALUE_ID: LEX-P0038-V0018
% #FIELD_VALUE: 否
\fieldvalue{\checkboxfield{unchecked}} 否\\
\hline
14 & 一次性分装套件 (PacPro-20) & / & 1Pcs/包 &
% #VALUE_ID: LEX-P0038-V0019
% #FIELD_VALUE: 生产批号/备注
% #TODO #HANDWRITTEN: 手写字样不清; handwriting overlaps the row
\fieldvalue{\handwritten{手写字样不清}} & &
% #VALUE_ID: LEX-P0038-V0020
% #FIELD_VALUE: 是
\fieldvalue{\checkboxfield{unchecked}} 是\quad
% #VALUE_ID: LEX-P0038-V0021
% #FIELD_VALUE: 否
\fieldvalue{\checkboxfield{unchecked}} 否\\
\hline
15 & 一次性分装套件 (PacPro-20) & / & 1Pcs/包 & & &
% #VALUE_ID: LEX-P0038-V0022
% #FIELD_VALUE: 是
\fieldvalue{\checkboxfield{unchecked}} 是\quad
% #VALUE_ID: LEX-P0038-V0023
% #FIELD_VALUE: 否
\fieldvalue{\checkboxfield{unchecked}} 否\\
\hline
16 & 一次性分装套件 (PacPro-20) & / & 1Pcs/包 & & &
% #VALUE_ID: LEX-P0038-V0024
% #FIELD_VALUE: 是
\fieldvalue{\checkboxfield{unchecked}} 是\quad
% #VALUE_ID: LEX-P0038-V0025
% #FIELD_VALUE: 否
\fieldvalue{\checkboxfield{unchecked}} 否\\
\hline
17 & 一次性分装套件 (PacPro-20) & / & 1Pcs/包 & & &
% #VALUE_ID: LEX-P0038-V0026
% #FIELD_VALUE: 是
\fieldvalue{\checkboxfield{unchecked}} 是\quad
% #VALUE_ID: LEX-P0038-V0027
% #FIELD_VALUE: 否
\fieldvalue{\checkboxfield{unchecked}} 否\\
\hline
18 & 一次性分装套件 (PacPro-20) & / & 1Pcs/包 & & &
% #VALUE_ID: LEX-P0038-V0028
% #FIELD_VALUE: 是
\fieldvalue{\checkboxfield{unchecked}} 是\quad
% #VALUE_ID: LEX-P0038-V0029
% #FIELD_VALUE: 否
\fieldvalue{\checkboxfield{unchecked}} 否\\
\hline
19 & 一次性分装套件 (PacPro-20) & / & 1Pcs/包 & & &
% #VALUE_ID: LEX-P0038-V0030
% #FIELD_VALUE: 是
\fieldvalue{\checkboxfield{unchecked}} 是\quad
% #VALUE_ID: LEX-P0038-V0031
% #FIELD_VALUE: 否
\fieldvalue{\checkboxfield{unchecked}} 否\\
\hline
20 &
% #VALUE_ID: LEX-P0038-V0032
% #FIELD_VALUE: 品名/备注
% #TODO #HANDWRITTEN: 一次性分装套件……; handwriting partially illegible
\fieldvalue{\handwritten{一次性分装套件……}} &
% #VALUE_ID: LEX-P0038-V0033
% #FIELD_VALUE: 规格/备注
% #HANDWRITTEN: 数字
\fieldvalue{\handwritten{数字}} &
% #VALUE_ID: LEX-P0038-V0034
% #FIELD_VALUE: 包装规格/备注
% #HANDWRITTEN: 1/2
\fieldvalue{\handwritten{1/2}} &
% #VALUE_ID: LEX-P0038-V0035
% #FIELD_VALUE: 生产批号/备注
% #TODO #HANDWRITTEN: H?86?260?301?; handwriting partially illegible
\fieldvalue{\handwritten{H?86?260?301?}} &
% #VALUE_ID: LEX-P0036
% #FIELD_VALUE: 批量/备注
% #HANDWRITTEN: 2
\fieldvalue{\handwritten{2}} &
% #VALUE_ID: LEX-P0038-V0037
% #FIELD_VALUE: 是
\fieldvalue{\checkboxfield{checked}} 是\quad
% #VALUE_ID: LEX-P0038-V0038
% #FIELD_VALUE: 否
\fieldvalue{\checkboxfield{unchecked}} 否\\
\hline
\end{tabular}
\end{center}

备注：
% #VALUE_ID: LEX-P0038-V0039
% #FIELD_VALUE: 备注
% #TODO #HANDWRITTEN: 短横线; handwritten mark in the remarks field
\fieldvalue{\handwritten{短横线}}

\vspace{0.5em}

确认人/日期：
% #VALUE_ID: LEX-P0038-V0040
% #FIELD_VALUE: 确认人
% #HANDWRITTEN: 刘晓华
\fieldvalue{\handwritten{刘晓华}}
\qquad
% #VALUE_ID: LEX-P0038-V0041
% #FIELD_VALUE: 确认日期
% #HANDWRITTEN: 2025.12.13
\fieldvalue{\handwritten{2025.12.13}}

\qquad
复核人/日期：
% #VALUE_ID: LEX-P0038-V0042
% #FIELD_VALUE: 复核人
% #HANDWRITTEN: 赵多
\fieldvalue{\handwritten{赵多}}
\qquad
% #VALUE_ID: LEX-P0038-V0043
% #FIELD_VALUE: 复核日期
% #HANDWRITTEN: 2025.12.13
\fieldvalue{\handwritten{2025.12.13}}

\subsection{2.\ 一次性分装套件与细胞冻存袋焊接}

\begin{center}
\small
\begin{tabular}{|p{0.46\linewidth}|p{0.44\linewidth}|}
\hline
项目 & 操作记录\\
\hline
1）使用气密性检测仪对一次性分装套件进行气密性检测，检测合格后方可使用。 &
一次性分装套件焊接数量：
% #VALUE_ID: LEX-P0038-V0044
% #FIELD_VALUE: 一次性分装套件焊接数量
% #HANDWRITTEN: 24
\fieldvalue{\handwritten{24}} 个/套，共24套\\
\hline
2）使用管路机将细胞冻存袋与气密性检测通过的一次性分装套件焊接。 &
是否符合要求：
% #VALUE_ID: LEX-P0038-V0045
% #FIELD_VALUE: 是
\fieldvalue{\checkboxfield{checked}} 是\quad
% #VALUE_ID: LEX-P0038-V0046
% #FIELD_VALUE: 否
\fieldvalue{\checkboxfield{unchecked}} 否\\
\hline
3）焊接后的一次性分装套件各接口检查是否无漏点，并填写焊接批号、规格、日期等信息；焊接批号填写规则：A年（4位）+月（2位）+日（2位）+日流水号（2位），如：A2024062501……A2024062508； &
\\
\hline
\multicolumn{2}{|c|}{一次性分装套件焊接记录}\\
\hline
序号 & 焊接批号 \quad 一次性分装套件批号 \quad 细胞冻存袋批号 \quad 细胞冻存袋批号 \quad 贴签数量（个）\\
\hline
\end{tabular}
\end{center}

% TODO: The welding-record table continues on the following page; only its visible header is included here.

\begin{center}
第 4 页 共 23 页
\end{center}
% LEXOID_PAGE_COMPLETED: 38/70

\newpage

\begin{center}
\small
记录编号：YZ-VP-PD-APS-012(2502)-07 \quad 版本号：1.0

\vspace{2pt}
\begin{tabular}{|p{0.28\linewidth}|p{0.20\linewidth}|p{0.20\linewidth}|p{0.14\linewidth}|}
\hline
\centering 品名 & \centering 规格 & \centering 生产批号 & \centering 批量\\
\hline
艾米沃托莨注射液（Z2）无菌工艺模拟试验
&
\centering 12.5$\pm$0.5 ml
&
% #VALUE_ID: LEX-P0039-V0001
% #FIELD_VALUE: 生产批号
% #TODO #HANDWRITTEN: APS?520?8(622); 字迹不清
\fieldvalue{\handwritten{APS?520?8(622)}}
&
% #VALUE_ID: LEX-P0039-V0002
% #FIELD_VALUE: 批量
\fieldvalue{7240}\ 袋
\\
\hline
\end{tabular}

\vspace{10pt}

\begin{tabular}{|p{0.07\linewidth}|p{0.18\linewidth}|p{0.25\linewidth}|p{0.22\linewidth}|p{0.14\linewidth}|}
\hline
& & \multicolumn{2}{c|}{数量（个）} &\\
\hline
1. & & & &\\
\hline
2. & & & &\\
\hline
3. & & & &\\
\hline
4. & & & &\\
\hline
5. & & & &\\
\hline
6. & &
\multicolumn{2}{c|}{
% #VALUE_ID: LEX-P0039-V0003
% #HANDWRITTEN: 无菌模拟 2025/12
% #TODO #HANDWRITTEN: 无菌模拟 2025/12; 斜向手写字迹部分难以辨认
\handwritten{无菌模拟 2025/12}
} &\\
\hline
7. & & & &\\
\hline
8. & & & &\\
\hline
9. & & & &\\
\hline
10. & & & &\\
\hline
\end{tabular}

\vspace{2pt}

\begin{tabular}{|p{0.60\linewidth}|p{0.40\linewidth}|}
\hline
4）使用气密性检测仪对焊接后的 一次性分装袋进行气密性检测，检测合格后进行贴签：
&
一次性分装袋存检验合格数量：
\\[-2pt]
&
% #VALUE_ID: LEX-P0039-V0004
% #FIELD_VALUE: 一次性分装袋存检验合格数量（个/套）
% #HANDWRITTEN: 10
\fieldvalue{\handwritten{10}}\ 个/套，共
% #VALUE_ID: LEX-P0039-V0005
% #FIELD_VALUE: 一次性分装袋存检验合格数量（套）
% #HANDWRITTEN: 24
\fieldvalue{\handwritten{24}}\ 套
\\[-2pt]
&
% #VALUE_ID: LEX-P0039-V0006
% #FIELD_VALUE: 一次性分装袋存检验合格数量（个/套）
% #HANDWRITTEN: 1
\fieldvalue{\handwritten{1}}\ 个/套，共
% #VALUE_ID: LEX-P0039-V0007
% #FIELD_VALUE: 一次性分装袋存检验合格数量（套）
% #HANDWRITTEN: 2
\fieldvalue{\handwritten{2}}\ 套
\\
\hline
5）撕掉和贴标签签字的一次性分装袋存放于分装间。
&
是否符合要求：
是
% #VALUE_ID: LEX-P0039-V0008
% #FIELD_VALUE: 是否符合要求：是
% #HANDWRITTEN: checked
\fieldvalue{\handwritten{\checkboxfield{checked}}}
\quad 否
% #VALUE_ID: LEX-P0039-V0009
% #FIELD_VALUE: 是否符合要求：否
\fieldvalue{\checkboxfield{unchecked}}
\\
\hline
备注： &\\[12pt]
\hline
\end{tabular}

\vspace{5pt}

操作人/日期：
% #VALUE_ID: LEX-P0039-V0010
% #FIELD_VALUE: 操作人
% #TODO #HANDWRITTEN: 刘?斌; 姓名字迹不清
% #HANDWRITTEN: 刘?斌
\fieldvalue{\handwritten{刘?斌}}
\quad
% #VALUE_ID: LEX-P0039-V0011
% #FIELD_VALUE: 操作日期
% #HANDWRITTEN: 2025.12.13
\fieldvalue{\handwritten{2025.12.13}}
\qquad
复核人/日期：
% #VALUE_ID: LEX-P0039-V0012
% #FIELD_VALUE: 复核人
% #TODO #HANDWRITTEN: 韩?; 姓名字迹不清
% #HANDWRITTEN: 韩?
\fieldvalue{\handwritten{韩?}}
\quad
% #VALUE_ID: LEX-P0039-V0013
% #FIELD_VALUE: 复核日期
% #HANDWRITTEN: 2025.12.13
\fieldvalue{\handwritten{2025.12.13}}

\end{center}

\vfill

\begin{center}
第 5 页 共 23 页
\end{center}
% LEXOID_PAGE_COMPLETED: 39/70

\newpage

\begin{center}
\small
\begin{tabular}{|p{0.34\linewidth}|p{0.20\linewidth}|p{0.19\linewidth}|p{0.13\linewidth}|}
\hline
品名 & 规格 & 生产批号 & 批量 \\ \hline
% #VALUE_ID: LEX-P0040-V0001
% #FIELD_VALUE: 品名
\fieldvalue{艾尔近托克塞注射液(ZZ)无菌工艺模拟试验}
&
% #VALUE_ID: LEX-P0040-V0002
% #FIELD_VALUE: 规格
\fieldvalue{12.5$\pm$0.5 ml}
&
% #VALUE_ID: LEX-P0040-V0003
% #FIELD_VALUE: 生产批号
% #HANDWRITTEN: APS 2025/201（?）/12
\fieldvalue{\handwritten{APS 2025/201（?）/12}}
&
% #VALUE_ID: LEX-P0040-V0004
% #FIELD_VALUE: 批量
% #HANDWRITTEN: 240
\fieldvalue{\handwritten{240}} 瓶
\\ \hline
\end{tabular}
\end{center}

\section*{3.\quad 耗材安装}

\subsection*{3.1\quad 耗材安装1}

\begin{center}
\small
\begin{tabular}{|p{0.67\linewidth}|p{0.29\linewidth}|}
\hline
\centering 项目 & \centering 操作记录 \\ \hline
1）设备开机：Pack Pro 开机，登录操作账户，选择需要的工艺；
&
设备编号：
% #VALUE_ID: LEX-P0040-V0005
% #FIELD_VALUE: 设备编号
% #HANDWRITTEN: 1123/15
\fieldvalue{\handwritten{1123/15}}

工艺名称：
% #VALUE_ID: LEX-P0040-V0006
% #FIELD_VALUE: 工艺名称
% #HANDWRITTEN: P-PacPro
\fieldvalue{\handwritten{P-PacPro}}
\\ \hline
\begin{minipage}[t]{\linewidth}
2）耗材安装

① 耗材安装前检查：确认当前设备上没有任何耗材。确认无误后，点击“确认”进入下一步。

② 耗材安装复位：根据连接的分装模块，选择对应的耗材安装序号，点击耗材安装，使分装模块复位至安装状态。

③ 耗材安装：套入卡板并安装1--13号病毒侧饱管路和注射泵，对准卡位放入四个角的卡槽，并旋转金属压座固定在耗材卡槽中。

④ 安装传感器管路：安装传感器位置的管路，并将管路塞入管槽中。

⑤ 连接排气阀：将分装模块耗材管路末端的排气阀与设备自带的管路连接。
\end{minipage}
&
是否符合要求：
% #VALUE_ID: LEX-P0040-V0007
% #FIELD_VALUE: 是
\fieldvalue{\checkboxfield{checked}} 是
\quad
% #VALUE_ID: LEX-P0040-V0008
% #FIELD_VALUE: 否
\fieldvalue{\checkboxfield{unchecked}} 否

模块1耗材连接号：
% #VALUE_ID: LEX-P0040-V0009
% #FIELD_VALUE: 模块1耗材连接号
% #TODO #HANDWRITTEN: illegible; handwritten serial number is unclear
\fieldvalue{\handwritten{\text{illegible}}}

模块2耗材连接号：
% #VALUE_ID: LEX-P0040-V0010
% #FIELD_VALUE: 模块2耗材连接号
% #TODO #HANDWRITTEN: illegible; handwritten serial number is unclear
\fieldvalue{\handwritten{\text{illegible}}}

模块3耗材连接号：
% #VALUE_ID: LEX-P0040-V0011
% #FIELD_VALUE: 模块3耗材连接号
% #TODO #HANDWRITTEN: illegible; handwritten serial number is unclear
\fieldvalue{\handwritten{\text{illegible}}}

模块4耗材连接号：
% #VALUE_ID: LEX-P0040-V0012
% #FIELD_VALUE: 模块4耗材连接号
% #TODO #HANDWRITTEN: illegible; handwritten serial number is unclear
\fieldvalue{\handwritten{\text{illegible}}}
\\ \hline
备注：
% #VALUE_ID: LEX-P0040-V0013
% #FIELD_VALUE: 备注
% #HANDWRITTEN: /
\fieldvalue{\handwritten{/}}
&
\\ \hline
\end{tabular}
\end{center}

操作人/日期：
% #VALUE_ID: LEX-P0040-V0014
% #FIELD_VALUE: 操作人
% #HANDWRITTEN: 赵海娜
\fieldvalue{\handwritten{赵海娜}}
\quad
% #VALUE_ID: LEX-P0040-V0015
% #FIELD_VALUE: 操作日期
% #HANDWRITTEN: 2025.12.13
\fieldvalue{\handwritten{2025.12.13}}
\hfill
复核人/日期：
% #VALUE_ID: LEX-P0040-V0016
% #FIELD_VALUE: 复核人
% #HANDWRITTEN: 张宁
\fieldvalue{\handwritten{张宁}}
\quad
% #VALUE_ID: LEX-P0040-V0017
% #FIELD_VALUE: 复核日期
% #HANDWRITTEN: 2025.12.13
\fieldvalue{\handwritten{2025.12.13}}

\subsection*{3.2\quad 耗材安装2}

\begin{center}
\small
\begin{tabular}{|p{0.67\linewidth}|p{0.29\linewidth}|}
\hline
\centering 项目 & \centering 操作记录 \\ \hline
1）设备开机：Pack Pro 开机，登录操作账户，选择需要的工艺；
&
设备编号：
% #VALUE_ID: LEX-P0040-V0018
% #FIELD_VALUE: 设备编号
% #HANDWRITTEN: 1137/15
\fieldvalue{\handwritten{1137/15}}

工艺名称：
% #VALUE_ID: LEX-P0040-V0019
% #FIELD_VALUE: 工艺名称
% #HANDWRITTEN: P-PacPro
\fieldvalue{\handwritten{P-PacPro}}
\\ \hline
\begin{minipage}[t]{\linewidth}
2）耗材安装

① 耗材安装复位：根据连接的分装模块，选择对应的耗材安装序号，点击耗材安装，使分装模块复位至安装状态。

② 耗材安装：套入卡板并安装1--13号病毒侧饱管路和注射泵，对准卡位放入四个角的卡槽，并旋转金属压座固定在耗材卡槽中。

③ 安装气泡传感器管路：安装气泡传感器位置的管路，并将管路塞入管槽。

④ 连接排气阀：将分装模块耗材管路末端的排气阀与设备自带的管路连接。
\end{minipage}
&
是否符合要求：
% #VALUE_ID: LEX-P0040-V0020
% #FIELD_VALUE: 是
\fieldvalue{\checkboxfield{checked}} 是
\quad
% #VALUE_ID: LEX-P0040-V0021
% #FIELD_VALUE: 否
\fieldvalue{\checkboxfield{unchecked}} 否

模块1耗材连接号：
% #VALUE_ID: LEX-P0040-V0022
% #FIELD_VALUE: 模块1耗材连接号
% #TODO #HANDWRITTEN: illegible; handwritten serial number is unclear
\fieldvalue{\handwritten{\text{illegible}}}

模块2耗材连接号：
% #VALUE_ID: LEX-P0040-V0023
% #FIELD_VALUE: 模块2耗材连接号
% #TODO #HANDWRITTEN: illegible; handwritten serial number is unclear
\fieldvalue{\handwritten{\text{illegible}}}

模块3耗材连接号：
% #VALUE_ID: LEX-P0040-V0024
% #FIELD_VALUE: 模块3耗材连接号
% #TODO #HANDWRITTEN: illegible; handwritten serial number is unclear
\fieldvalue{\handwritten{\text{illegible}}}

模块4耗材连接号：
% #VALUE_ID: LEX-P0040-V0025
% #FIELD_VALUE: 模块4耗材连接号
% #TODO #HANDWRITTEN: illegible; handwritten serial number is unclear
\fieldvalue{\handwritten{\text{illegible}}}
\\ \hline
备注： & \\ \hline
\end{tabular}
\end{center}

操作人/日期：
% #VALUE_ID: LEX-P0040-V0026
% #FIELD_VALUE: 操作人
% #HANDWRITTEN: 赵海娜
\fieldvalue{\handwritten{赵海娜}}
\quad
% #VALUE_ID: LEX-P0040-V0027
% #FIELD_VALUE: 操作日期
% #HANDWRITTEN: 2025.12.13
\fieldvalue{\handwritten{2025.12.13}}
\hfill
复核人/日期：
% #VALUE_ID: LEX-P0040-V0028
% #FIELD_VALUE: 复核人
% #HANDWRITTEN: 张宁
\fieldvalue{\handwritten{张宁}}
\quad
% #VALUE_ID: LEX-P0040-V0029
% #FIELD_VALUE: 复核日期
% #HANDWRITTEN: 2025.12.13
\fieldvalue{\handwritten{2025.12.13}}

\begin{center}
第 6 页 共 23 页
\end{center}
% LEXOID_PAGE_COMPLETED: 40/70

\newpage

\noindent
\begin{tabular}{|p{0.22\linewidth}|p{0.20\linewidth}|p{0.25\linewidth}|p{0.16\linewidth}|}
\hline
品名 &
% #VALUE_ID: LEX-P0041-V0001
% #FIELD_VALUE: 品名
\fieldvalue{艾尔迈克莘注射液(22)无菌工艺模拟试验} &
规格 &
% #VALUE_ID: LEX-P0041-V0002
% #FIELD_VALUE: 规格
\fieldvalue{12.5$\pm$0.5\,ml}
\\
\hline
生产批号 &
% #VALUE_ID: LEX-P0041-V0003
% #FIELD_VALUE: 生产批号
% #TODO #HANDWRITTEN: APS20210106(??); 部分字符难以辨认
\fieldvalue{\handwritten{APS20210106(??)}} &
批量 &
% #VALUE_ID: LEX-P0041-V0004
% #FIELD_VALUE: 批量
% #HANDWRITTEN: \(\geq\)240
\fieldvalue{\handwritten{\(\geq\)240}} 袋
\\
\hline
\end{tabular}

\vspace{0.5em}

\section*{3.3.\quad 细胞分装液制备}

\noindent
\begin{tabular}{|p{0.56\linewidth}|p{0.36\linewidth}|}
\hline
\centering 项目 & \centering 操作 \tabularnewline
\hline
1）将1000\,ml 10\%DMSO 溶液（TSB）通过配液仪装入一次性灌匀袋（10\,L）中； &
加入体积
% #VALUE_ID: LEX-P0041-V0005
% #FIELD_VALUE: 加入体积
% #HANDWRITTEN: 000.0
\fieldvalue{\handwritten{000.0}}\,ml（g）
\\
\hline
2）细胞分装液（TSB）信息确认

确认组分配液袋包装完整无破损，记录体积/质量： &
细胞分装母液（TSB）：

【1号】体积/质量：
% #VALUE_ID: LEX-P0041-V0006
% #FIELD_VALUE: 1号体积/质量
% #HANDWRITTEN: 479.5
\fieldvalue{\handwritten{479.5}}\,ml（g）；

【2号】体积/质量：
% #VALUE_ID: LEX-P0041-V0007
% #FIELD_VALUE: 2号体积/质量
% #HANDWRITTEN: 300.4
\fieldvalue{\handwritten{300.4}}\,ml（g）；
\\
\hline
3）合并配包

将含有细胞分装母液（TSB）的3\,L 储液袋通过无菌接管机与含有10\%DMSO 溶液（TSB）的一次性配液袋（10\,L）连接，通过自动配液仪将上述一次性配液袋（10\,L）中； &
混合后细胞分装母液（TSB）

体积（V0）
% #VALUE_ID: LEX-P0041-V0009
% #FIELD_VALUE: 体积（V0）
% #HANDWRITTEN: 191.0
\fieldvalue{\handwritten{191.0}}\,ml
\\
\hline
4）在混匀模块中混匀，混匀5\,min，进行模拟取样计数； &
是否符合要求：
% #VALUE_ID: LEX-P0041-V0010
% #FIELD_VALUE: 是
% #HANDWRITTEN: ✓
\fieldvalue{\handwritten{\checkboxfield{checked}}} 是\quad
% #VALUE_ID: LEX-P0041-V0011
% #FIELD_VALUE: 否
\fieldvalue{\checkboxfield{unchecked}} 否
\\
\hline
5）细胞分装液体积应定容至6800\,ml； &
是否符合要求：
% #VALUE_ID: LEX-P0041-V0012
% #FIELD_VALUE: 是
% #HANDWRITTEN: ✓
\fieldvalue{\handwritten{\checkboxfield{checked}}} 是\quad
% #VALUE_ID: LEX-P0041-V0013
% #FIELD_VALUE: 否
\fieldvalue{\checkboxfield{unchecked}} 否
\\
\hline
备注：\quad
% #VALUE_ID: LEX-P0041-V0014
% #TODO #HANDWRITTEN: ／; 手写符号无法确定
\handwritten{／}
&
\\
\hline
\end{tabular}

\vspace{0.5em}

操作人/日期：
% #VALUE_ID: LEX-P0041-V0015
% #FIELD_VALUE: 操作人
% #TODO #HANDWRITTEN: 赵晓玲; 姓名部分笔画较模糊
\fieldvalue{\handwritten{赵晓玲}}
\quad
% #VALUE_ID: LEX-P0041-V0016
% #FIELD_VALUE: 操作日期
% #HANDWRITTEN: 2021.12.13
\fieldvalue{\handwritten{2021.12.13}}
\qquad
复核人/日期：
% #VALUE_ID: LEX-P0041-V0017
% #FIELD_VALUE: 复核人
% #TODO #HANDWRITTEN: 魏晓梅; 姓名部分笔画较模糊
\fieldvalue{\handwritten{魏晓梅}}
\quad
% #VALUE_ID: LEX-P0041-V0018
% #FIELD_VALUE: 复核日期
% #HANDWRITTEN: 2021.12.13
\fieldvalue{\handwritten{2021.12.13}}

\vfill

\begin{center}
第 7 页 共 23 页
\end{center}
% LEXOID_PAGE_COMPLETED: 41/70

\newpage

\begin{center}
记录编号：YZ-VP-DP-APS-012(2502)-07 \quad 版本号：1.0
\end{center}

\begin{tabular}{|p{0.25\linewidth}|p{0.20\linewidth}|p{0.25\linewidth}|p{0.12\linewidth}|}
\hline
品名 & 规格 & 生产批号 & 批量\\
\hline
% #VALUE_ID: LEX-P0042-V0001
% #FIELD_VALUE: 品名
\fieldvalue{艾米克托莫赛注射液(ZZ)无菌工艺模拟试验}
&
% #VALUE_ID: LEX-P0042-V0002
% #FIELD_VALUE: 规格
\fieldvalue{12.5\,$\pm$\,0.5\,ml}
&
% #VALUE_ID: LEX-P0042-V0003
% #FIELD_VALUE: 生产批号
% #TODO #HANDWRITTEN: APS?OS201612?2460; handwriting is partially unclear
\fieldvalue{\handwritten{APS?OS201612?2460}}
&
袋\\
\hline
\end{tabular}

\section{细胞分装}

\subsection{参数设置}

\begin{tabular}{|p{0.64\linewidth}|p{0.28\linewidth}|}
\hline
项目 & 操作记录\\
\hline
1）样品袋（细胞分装液）连接：\newline
① 将一次性袋装液连接管路的管夹关闭；\newline
② 接收细胞分装液，将细胞分装液使用无菌接管机连接至样品管路；\newline
③ 悬挂倒钩：打开配气门和润泡均衡板，将润泡袋与分装袋管路连接，挂到配气门内的挂钩上；将管路左边配气门卡带中，关闭润泡均衡板和配气门；\newline
④ 再次检查确认胶塞、传感器和负压夹均已正确安装；确认无异常后，点击设备上的耗材确认按钮；\newline
⑤ 打开配气袋及对应分装模块管路上的白色荧光管；
&
分装液质量/体积：\newline
% #VALUE_ID: LEX-P0042-V0004
% #FIELD_VALUE: 分装液质量/体积
% #TODO #HANDWRITTEN: 200? g (ml); final digit is unclear
\fieldvalue{\handwritten{200?\,g (ml)}}\newline
是否符合要求：是
% #VALUE_ID: LEX-P0042-V0005
% #FIELD_VALUE: 是
\fieldvalue{\checkboxfield{checked}}
\quad 否
% #VALUE_ID: LEX-P0042-V0006
% #FIELD_VALUE: 否
\fieldvalue{\checkboxfield{unchecked}}\\
\hline
\multicolumn{2}{|l|}{2．参数设置}\\
\hline
\multicolumn{2}{|p{0.92\linewidth}|}{
① 混匀参数设置：\newline
“细胞直径和细胞密度”为“12--20/3E5--1E7”；\newline
“混匀模块类别”选择“3--10L”；\newline
“混匀接触”选择“10L”；\newline
“是否混匀”选择“是”；\newline
设置“混匀时间”；\newline
“样本袋是否放置在混匀仓”选择“是”；\newline
“混匀袋装前停顿”。\newline
② 分装参数设置：\newline
选择分装模块；\newline
“注射器规格”选择“CellBri/50”；\newline
“分装袋规格”；\newline
输入“分装体积”；\newline
输入“总分装袋数”：混匀袋当前体积/12.5（如有小数位则进一位取整）；\newline
“全流程控制”选择“不转气”；\newline
③ 参数设置完成后，进入信息确认及电子签名，签名后点击“确认”；}\\
\hline
& 是否符合要求：是
% #VALUE_ID: LEX-P0042-V0007
% #FIELD_VALUE: 是
\fieldvalue{\checkboxfield{checked}}
\quad 否
% #VALUE_ID: LEX-P0042-V0008
% #FIELD_VALUE: 否
\fieldvalue{\checkboxfield{unchecked}}\\
\hline
备注： & \\
\hline
\end{tabular}

\vspace{0.5em}

操作人/日期：
% #VALUE_ID: LEX-P0042-V0009
% #FIELD_VALUE: 操作人
% #TODO #HANDWRITTEN: 胡丽娟?; handwritten signature is unclear
\fieldvalue{\handwritten{胡丽娟?}}
\quad
% #VALUE_ID: LEX-P0042-V0010
% #FIELD_VALUE: 操作日期
% #HANDWRITTEN: 2025.12.13
\fieldvalue{\handwritten{2025.12.13}}

\hspace{2em}
复核人/日期：
% #VALUE_ID: LEX-P0042-V0011
% #FIELD_VALUE: 复核人
% #TODO #HANDWRITTEN: 张?; handwritten signature is unclear
\fieldvalue{\handwritten{张?}}
\quad
% #VALUE_ID: LEX-P0042-V0012
% #FIELD_VALUE: 复核日期
% #TODO #HANDWRITTEN: 2025.11.12?; date is partially unclear
\fieldvalue{\handwritten{2025.11.12?}}

\begin{center}
第 8 页 共 23 页
\end{center}
% LEXOID_PAGE_COMPLETED: 42/70

\newpage

\noindent 记录编号：YZ-YP-PD-APS-012(2502)-07 \quad 版本号：1.0

\begin{center}
\begin{tabular}{|p{0.20\linewidth}|p{0.20\linewidth}|p{0.27\linewidth}|p{0.18\linewidth}|}
\hline
品名 & 规格 & 生产批号 & 批量\\
\hline
\text{--} & 12.5$\pm$0.5 ml &
% #VALUE_ID: LEX-P0043-V0001
% #FIELD_VALUE: 生产批号
% #TODO #HANDWRITTEN: AS?051?V?（6?2?）>440; 字迹不清
\fieldvalue{\handwritten{AS?051?V?（6?2?）>440}} &
袋\\
\hline
\end{tabular}
\end{center}

\noindent 4.2.\quad 分装模块1\quad
% #VALUE_ID: LEX-P0043-V0002
% #HANDWRITTEN: A1
\handwritten{A1}\quad
% #VALUE_ID: LEX-P0043-V0003
% #TODO #HANDWRITTEN: 分装; 蓝色手写批注字迹不清
\handwritten{分装}\quad
% #VALUE_ID: LEX-P0043-V0004
% #HANDWRITTEN: 2025.12.13
\handwritten{2025.12.13}

\subsubsection*{4.2.1.\quad 分装操作}

\begin{tabular}{|p{0.53\linewidth}|p{0.39\linewidth}|}
\hline
\multicolumn{1}{|c|}{项目} & \multicolumn{1}{c|}{操作记录}\\
\hline
1.分装过程：整个分装过程均由分装模块自动运行，程序自动预推仓、自动判别和掉空管路；\newline
2.分装模块第一轮分装结束后，断开分装模块耗材管路末端的排气阀与设备自带的管路连接，旋转金属卡扣至垂直位置，然后对应管路后取下耗材卡板；\newline
3.点击主机工艺界面中对应模块的“耗材安装1”，分装模块1状态完成后位后，按1.3耗材安装中2种耗材安装步骤中的分装耗材，打开对应管路的白色管夹，点击耗材确认按键后继续分装；\newline
每个模块每次分装进行编号，（如：A号模块 Gentle P-Pac Pro 第1次分装：A-1、A号 Gentle P-Pac Pro 第2次分装：A-2、B号 Gentle P-Pac Pro 第3次分装：B-3等）； &
1.分装开始时间：
% #VALUE_ID: LEX-P0043-V0005
% #FIELD_VALUE: 分装开始时间
% #HANDWRITTEN: 10:24
\fieldvalue{\handwritten{10:24}}\\
&
2.分装模块完成
% #VALUE_ID: LEX-P0043-V0006
% #FIELD_VALUE: 分装模块完成次数
% #HANDWRITTEN: 2
\fieldvalue{\handwritten{2}} 次循环分装；\\
&
3.分装结束时间：
% #VALUE_ID: LEX-P0043-V0007
% #FIELD_VALUE: 分装结束时间
% #HANDWRITTEN: 16:19
\fieldvalue{\handwritten{16:19}}\\
\hline
\multicolumn{2}{|l|}{备注：
% #VALUE_ID: LEX-P0043-V0008
% #FIELD_VALUE: 备注
% #HANDWRITTEN: /
\fieldvalue{\handwritten{/}}}\\
\hline
\end{tabular}

\noindent 操作人/日期：
% #VALUE_ID: LEX-P0043-V0009
% #FIELD_VALUE: 操作人
% #TODO #HANDWRITTEN: 签名; 字迹不清
\fieldvalue{\handwritten{签名}} \quad
% #VALUE_ID: LEX-P0043-V0010
% #FIELD_VALUE: 操作日期
% #HANDWRITTEN: 2025.12.13
\fieldvalue{\handwritten{2025.12.13}}
\hfill
复核人/日期：
% #VALUE_ID: LEX-P0043-V0011
% #FIELD_VALUE: 复核人
% #TODO #HANDWRITTEN: 签名; 字迹不清
\fieldvalue{\handwritten{签名}} \quad
% #VALUE_ID: LEX-P0043-V0012
% #FIELD_VALUE: 复核日期
% #HANDWRITTEN: 2025.12.13
\fieldvalue{\handwritten{2025.12.13}}

\subsubsection*{4.2.2.\quad 冻存袋整合操作}

\begin{tabular}{|p{0.53\linewidth}|p{0.39\linewidth}|}
\hline
1.检查管路无破损、无泄漏，确认标签粘贴良好，冻存袋装置目视无异常；\newline
2.用封管机将冻存袋一侧的 EVA 管路热合三次，用剪刀将第二次热合处位置剪开；\newline
3.检查封口处及细胞冻存袋无泄漏，非标准袋备选原因，并按如下规则确认冷冻标签记录：\newline
\hspace{1em}装置不足的冻存袋标记为B；\newline
\hspace{1em}有破损的细胞冻存袋，标记“×”； &
1.热合开始时间：
% #VALUE_ID: LEX-P0043-V0013
% #FIELD_VALUE: 热合开始时间
% #HANDWRITTEN: 10:38
\fieldvalue{\handwritten{10:38}}\\
&
2.第1轮分装存袋数量：
% #VALUE_ID: LEX-P0043-V0014
% #FIELD_VALUE: 第1轮分装存袋数量
% #HANDWRITTEN: 10
\fieldvalue{\handwritten{10}} 袋；\\
&\quad 分装合格存袋数：
% #VALUE_ID: LEX-P0043-V0015
% #FIELD_VALUE: 第1轮分装合格存袋数
% #HANDWRITTEN: 0
\fieldvalue{\handwritten{0}} 袋；\\
&\quad 标记为B的冻存袋数：
% #VALUE_ID: LEX-P0043-V0016
% #FIELD_VALUE: 第1轮标记为B的冻存袋数
% #HANDWRITTEN: 0
\fieldvalue{\handwritten{0}} 个\\
&\quad 标记为“×”的冻存袋数：
% #VALUE_ID: LEX-P0043-V0017
% #FIELD_VALUE: 第1轮标记为“×”的冻存袋数
% #HANDWRITTEN: 0
\fieldvalue{\handwritten{0}} 个\\
&3.第2轮分装存袋数量：
% #VALUE_ID: LEX-P0043-V0018
% #FIELD_VALUE: 第2轮分装存袋数量
% #HANDWRITTEN: 10
\fieldvalue{\handwritten{10}} 个；\\
&\quad 分装合格存袋数：
% #VALUE_ID: LEX-P0043-V0019
% #FIELD_VALUE: 第2轮分装合格存袋数
% #HANDWRITTEN: 0
\fieldvalue{\handwritten{0}} 个\\
&\quad 标记为B的冻存袋数：
% #VALUE_ID: LEX-P0043-V0020
% #FIELD_VALUE: 第2轮标记为B的冻存袋数
% #HANDWRITTEN: 0
\fieldvalue{\handwritten{0}} 个\\
&\quad 标记为“×”的冻存袋数：
% #VALUE_ID: LEX-P0043-V0021
% #FIELD_VALUE: 第2轮标记为“×”的冻存袋数
% #HANDWRITTEN: 0
\fieldvalue{\handwritten{0}} 个\\
&4.第3轮分装存袋数量：
% #VALUE_ID: LEX-P0043-V0022
% #FIELD_VALUE: 第3轮分装存袋数量
% #HANDWRITTEN: 0
\fieldvalue{\handwritten{0}} 个；\\
&\quad 分装合格存袋数：
% #VALUE_ID: LEX-P0043-V0023
% #FIELD_VALUE: 第3轮分装合格存袋数
% #HANDWRITTEN: 0
\fieldvalue{\handwritten{0}} 个\\
&\quad 标记为B的冻存袋数：
% #VALUE_ID: LEX-P0043-V0024
% #FIELD_VALUE: 第3轮标记为B的冻存袋数
% #HANDWRITTEN: 0
\fieldvalue{\handwritten{0}} 个\\
&\quad 标记为“×”的冻存袋数：
% #VALUE_ID: LEX-P0043-V0025
% #FIELD_VALUE: 第3轮标记为“×”的冻存袋数
% #HANDWRITTEN: 0
\fieldvalue{\handwritten{0}} 个\\
&4.热合结束时间：
% #VALUE_ID: LEX-P0043-V0026
% #FIELD_VALUE: 热合结束时间
% #TODO #HANDWRITTEN: 10:37; 末位字迹不清
\fieldvalue{\handwritten{10:37}}\\
\hline
\multicolumn{2}{|l|}{备注：
% #VALUE_ID: LEX-P0043-V0027
% #FIELD_VALUE: 备注
% #HANDWRITTEN: /
\fieldvalue{\handwritten{/}}}\\
\hline
\end{tabular}

\noindent 操作人/日期：
% #VALUE_ID: LEX-P0043-V0028
% #FIELD_VALUE: 操作人
% #TODO #HANDWRITTEN: 签名; 字迹不清
\fieldvalue{\handwritten{签名}} \quad
% #VALUE_ID: LEX-P0043-V0029
% #FIELD_VALUE: 操作日期
% #HANDWRITTEN: 2025.12.13
\fieldvalue{\handwritten{2025.12.13}}
\hfill
复核人/日期：
% #VALUE_ID: LEX-P0043-V0030
% #FIELD_VALUE: 复核人
% #TODO #HANDWRITTEN: 签名; 字迹不清
\fieldvalue{\handwritten{签名}} \quad
% #VALUE_ID: LEX-P0043-V0031
% #FIELD_VALUE: 复核日期
% #HANDWRITTEN: 2025.12.13
\fieldvalue{\handwritten{2025.12.13}}

\begin{center}
第 9 页 共 23 页
\end{center}
% LEXOID_PAGE_COMPLETED: 43/70

\newpage

\begin{center}
\small
\begin{tabular}{|p{0.23\linewidth}|p{0.20\linewidth}|p{0.25\linewidth}|p{0.13\linewidth}|}
\hline
品名 & 规格 & 生产批号 & 批量\\
\hline
艾米近托赛替康注射液(2)无菌工艺模拟试验
&
12.5$\pm$0.5 ml
&
% #VALUE_ID: LEX-P0044-V0001
% #FIELD_VALUE: 生产批号
% #TODO #HANDWRITTEN: A?24?10(62); handwriting is unclear
\fieldvalue{\handwritten{A?24?10(62)}}
&
% #VALUE_ID: LEX-P0044-V0002
% #FIELD_VALUE: 批量
% #TODO #HANDWRITTEN: 2740; final digit is partially unclear
\fieldvalue{\handwritten{2740}} 瓶
\\
\hline
\end{tabular}
\end{center}

\noindent
记录编号：YZ-VP-PD-APS-012(2502)-07\qquad 版本号：1.0

\medskip

\noindent
4.3.\quad 分装模块2
% #VALUE_ID: LEX-P0044-V0003
% #FIELD_VALUE: 分装模块记录
% #TODO #HANDWRITTEN: A2 [签名] 12.13; middle handwritten portion is unclear
\fieldvalue{\handwritten{A2 [签名] 12.13}}

\medskip

\noindent
4.3.1.\quad 分装操作

\begin{center}
\small
\begin{tabular}{|p{0.55\linewidth}|p{0.37\linewidth}|}
\hline
项目 & 操作记录\\
\hline
1.分装过程：整个分装过程均由分装模块自动运行，程序自动预装带、自动制刷和带套管；
 
2.分装模块第一轮分装结束后，断开分装模块耗材管路末端的排气阀与设备自带的管路连接，旋转金属卡扣至垂直位置，然后对过滤器后段下耗材卡板；

3.点击主机工艺界面中对应模块的“耗材安装1”，分装模块1状态完成复位后，按1.3耗材安装中2)耗材安装步骤安装新的分装耗材，打开对应窗路的白色管夹，点击耗材确认按键继续分装；

每个模块每次分装进行编号。（如：A号模块 Gentle P-Pac Pro 第1次分装：A-1，A号 Gentle P-Pac Pro 第2次分装：A-2，B号 Gentle P-Pac Pro 第3次分装：B-3等）
&
1.分装开始时间：
% #VALUE_ID: LEX-P0044-V0004
% #FIELD_VALUE: 分装开始时间（时）
% #HANDWRITTEN: 10
\fieldvalue{\handwritten{10}}：
% #VALUE_ID: LEX-P0044-V0005
% #FIELD_VALUE: 分装开始时间（分）
% #HANDWRITTEN: 24
\fieldvalue{\handwritten{24}}

2.分装模块完成
% #VALUE_ID: LEX-P0044-V0006
% #FIELD_VALUE: 分装模块完成时间
% #TODO #HANDWRITTEN: 2; handwritten mark is unclear
\fieldvalue{\handwritten{2}}：循环分装；

3.分装结束时间：
% #VALUE_ID: LEX-P0044-V0007
% #FIELD_VALUE: 分装结束时间（时）
% #TODO #HANDWRITTEN: 1; handwriting is faint
\fieldvalue{\handwritten{1}}：
% #VALUE_ID: LEX-P0044-V0008
% #FIELD_VALUE: 分装结束时间（分）
% #HANDWRITTEN: 4
\fieldvalue{\handwritten{4}}
\\
\hline
备注： & \\
\hline
操作人/日期：
% #VALUE_ID: LEX-P0044-V0009
% #FIELD_VALUE: 操作人
% #TODO #HANDWRITTEN: [签名]; signature is illegible
\fieldvalue{\handwritten{[签名]}}
\quad
% #VALUE_ID: LEX-P0044-V0010
% #FIELD_VALUE: 操作日期
% #HANDWRITTEN: 2025.12.13
\fieldvalue{\handwritten{2025.12.13}}
&
复核人/日期：
% #VALUE_ID: LEX-P0044-V0011
% #FIELD_VALUE: 复核人
% #TODO #HANDWRITTEN: [签名]; signature is illegible
\fieldvalue{\handwritten{[签名]}}
\quad
% #VALUE_ID: LEX-P0044-V0012
% #FIELD_VALUE: 复核日期
% #HANDWRITTEN: 2025.12.13
\fieldvalue{\handwritten{2025.12.13}}
\\
\hline
\end{tabular}
\end{center}

\noindent
4.3.2.\quad 冻存袋热合操作

\begin{center}
\small
\begin{tabular}{|p{0.54\linewidth}|p{0.38\linewidth}|}
\hline
\begin{minipage}[t]{0.52\linewidth}
1.检查管路无破损、无泄漏，确认标签贴出完好，冻存袋装量目视无异常；

2.用封管机将冻存袋一侧的EVA管路热合三次，用剪刀将第二次热合处位置剪开；

3.检查封口处及细胞冻存袋无泄漏，非标准装袋原因，并按如下规则确认冻存袋标记记录：

\hspace{1em}装量不足的冻存袋标记为B；

\hspace{1em}有破损的细胞冻存袋，标记“$\times$”；
\end{minipage}
&
1.热合开始时间：
% #VALUE_ID: LEX-P0044-V0013
% #FIELD_VALUE: 热合开始时间（时）
% #HANDWRITTEN: 10
\fieldvalue{\handwritten{10}}：
% #VALUE_ID: LEX-P0044-V0014
% #FIELD_VALUE: 热合开始时间（分）
% #HANDWRITTEN: 28
\fieldvalue{\handwritten{28}}

2.第1轮分装冻存袋数量：
% #VALUE_ID: LEX-P0044-V0015
% #FIELD_VALUE: 第1轮分装冻存袋数量
% #HANDWRITTEN: 10
\fieldvalue{\handwritten{10}}袋；

\hspace{1em}分装合格冻存袋数量：
% #VALUE_ID: LEX-P0044-V0016
% #FIELD_VALUE: 第1轮分装合格冻存袋数量
% #HANDWRITTEN: 0
\fieldvalue{\handwritten{0}}袋；

\hspace{1em}标记为B的冻存袋数：
% #VALUE_ID: LEX-P0044-V0017
% #FIELD_VALUE: 第1轮标记为B的冻存袋数
% #HANDWRITTEN: 0
\fieldvalue{\handwritten{0}}个；

\hspace{1em}标记为“$\times$”的冻存袋数：
% #VALUE_ID: LEX-P0044-V0018
% #FIELD_VALUE: 第1轮标记为“$\times$”的冻存袋数
% #HANDWRITTEN: 0
\fieldvalue{\handwritten{0}}个。

3.第2轮分装冻存袋数量：
% #VALUE_ID: LEX-P0044-V0019
% #FIELD_VALUE: 第2轮分装冻存袋数量
% #HANDWRITTEN: 10
\fieldvalue{\handwritten{10}}个；

\hspace{1em}分装合格冻存袋数量：
% #VALUE_ID: LEX-P0044-V0020
% #FIELD_VALUE: 第2轮分装合格冻存袋数量
% #HANDWRITTEN: 0
\fieldvalue{\handwritten{0}}个；

\hspace{1em}标记为B的冻存袋数：
% #VALUE_ID: LEX-P0044-V0021
% #FIELD_VALUE: 第2轮标记为B的冻存袋数
% #HANDWRITTEN: 0
\fieldvalue{\handwritten{0}}个；

\hspace{1em}标记为“$\times$”的冻存袋数：
% #VALUE_ID: LEX-P0044-V0022
% #FIELD_VALUE: 第2轮标记为“$\times$”的冻存袋数
% #HANDWRITTEN: 0
\fieldvalue{\handwritten{0}}个。
\\
\hline
\begin{minipage}[t]{0.52\linewidth}
\phantom{1.} 
\end{minipage}
&
4.第3轮热合冻存袋数：
% #VALUE_ID: LEX-P0044-V0023
% #FIELD_VALUE: 第3轮热合冻存袋数
% #TODO #HANDWRITTEN: 15; first digit is partly obscured
\fieldvalue{\handwritten{15}}个；

\hspace{1em}分装合格冻存袋数量：
% #VALUE_ID: LEX-P0044-V0024
% #FIELD_VALUE: 第3轮分装合格冻存袋数量
% #HANDWRITTEN: 0
\fieldvalue{\handwritten{0}}个；

\hspace{1em}标记为B的冻存袋数：
% #VALUE_ID: LEX-P0044-V0025
% #FIELD_VALUE: 第3轮标记为B的冻存袋数
% #HANDWRITTEN: 0
\fieldvalue{\handwritten{0}}个；

\hspace{1em}标记为“$\times$”的冻存袋数：
% #VALUE_ID: LEX-P0044-V0026
% #FIELD_VALUE: 第3轮标记为“$\times$”的冻存袋数
% #HANDWRITTEN: 0
\fieldvalue{\handwritten{0}}个。

4.4 热合结束时间：
% #VALUE_ID: LEX-P0044-V0027
% #FIELD_VALUE: 热合结束时间（时）
% #HANDWRITTEN: 16
\fieldvalue{\handwritten{16}}：
% #VALUE_ID: LEX-P0044-V0028
% #FIELD_VALUE: 热合结束时间（分）
% #HANDWRITTEN: 22
\fieldvalue{\handwritten{22}}
\\
\hline
备注： & \\
\hline
操作人/日期：
% #VALUE_ID: LEX-P0044-V0029
% #FIELD_VALUE: 操作人
% #TODO #HANDWRITTEN: [签名]; signature is illegible
\fieldvalue{\handwritten{[签名]}}
\quad
% #VALUE_ID: LEX-P0044-V0030
% #FIELD_VALUE: 操作日期
% #HANDWRITTEN: 2025.12.13
\fieldvalue{\handwritten{2025.12.13}}
&
复核人/日期：
% #VALUE_ID: LEX-P0044-V0031
% #FIELD_VALUE: 复核人
% #TODO #HANDWRITTEN: [签名]; signature is illegible
\fieldvalue{\handwritten{[签名]}}
\quad
% #VALUE_ID: LEX-P0044-V0032
% #FIELD_VALUE: 复核日期
% #HANDWRITTEN: 2025.12.13
\fieldvalue{\handwritten{2025.12.13}}
\\
\hline
\end{tabular}
\end{center}

\begin{center}
第 10 页 共 23 页
\end{center}
% LEXOID_PAGE_COMPLETED: 44/70

\newpage

\begin{center}
\small
\begin{tabular}{|p{0.38\linewidth}|p{0.20\linewidth}|p{0.20\linewidth}|p{0.12\linewidth}|}
\hline
品名 & 规格 & 生产批号 & 批量 \\
\hline
% #VALUE_ID: LEX-P0045-V0001
% #FIELD_VALUE: 品名
\fieldvalue{艾米克坦塞针液(ZZ)无菌工艺模拟试验} &
% #VALUE_ID: LEX-P0045-V0002
% #FIELD_VALUE: 规格
\fieldvalue{12.5$\pm$0.5\,ml} &
% #VALUE_ID: LEX-P0045-V0003
% #FIELD_VALUE: 生产批号
% #HANDWRITTEN: APS2025/12/062
\fieldvalue{\handwritten{APS2025/12/062}} &
% #VALUE_ID: LEX-P0045-V0004
% #FIELD_VALUE: 批量
% #HANDWRITTEN: 2740
\fieldvalue{\handwritten{2740}}袋 \\
\hline
\end{tabular}
\end{center}

\noindent
4.4.\quad 分装模块3\quad
% #VALUE_ID: LEX-P0045-V0005
% #FIELD_VALUE: 分装模块
% #HANDWRITTEN: A2，预调液，2025.12.13
\fieldvalue{\handwritten{A2，预调液，2025.12.13}}

\medskip
\noindent
4.4.1.\quad 分装操作

\begin{center}
\small
\begin{tabular}{|p{0.55\linewidth}|p{0.37\linewidth}|}
\hline
项目 & 操作记录 \\
\hline
1.分装过程：整个分装过程均由分装模块自动运行，程序自动预排、自动刷新和构整合管路。&
1.分装开始时间：
% #VALUE_ID: LEX-P0045-V0006
% #FIELD_VALUE: 分装开始时间
% #HANDWRITTEN: 10：24
\fieldvalue{\handwritten{10：24}}\\
2.分装模块第一轮分装结束后，断开分装模块耗材管路末端的排气阀与设备自带的管路连接，旋转金属卡扣至垂直位置，热合对应该容器后断开卡板；&
2.分装模块完成
% #VALUE_ID: LEX-P0045-V0007
% #FIELD_VALUE: 循环分装次数
% #HANDWRITTEN: 2
\fieldvalue{\handwritten{2}}次循环分装；\\
3.点击主机工艺界面中对应模块的“耗材安装1”，分装模块1状态完成复位后，按1.3耗材安装2；耗材安装步骤安装新的分装耗材，打开对应管路的白色管夹，点击耗材确认按钮继续分装；&
3.分装结束时间：
% #VALUE_ID: LEX-P0045-V0008
% #FIELD_VALUE: 分装结束时间
% #HANDWRITTEN: 12：4
\fieldvalue{\handwritten{12：4}}\\
每个模块每次分装进行编号。（如：A号模块Gentle P-Pac Pro第1次分装：A-1，A号Gentle P-Pac Pro第2次分装：A-2，B号Gentle P-Pac Pro第3次分装：B-3等）&
\\
备注：&\\
\hline
操作人/日期：
% #VALUE_ID: LEX-P0045-V0009
% #FIELD_VALUE: 操作人
% #HANDWRITTEN: 李（字迹不清）
\fieldvalue{\handwritten{李（字迹不清）}}
\quad
% #VALUE_ID: LEX-P0045-V0010
% #FIELD_VALUE: 操作日期
% #HANDWRITTEN: 2025.12.13
\fieldvalue{\handwritten{2025.12.13}}
&
复核人/日期：
% #VALUE_ID: LEX-P0045-V0011
% #FIELD_VALUE: 复核人
% #HANDWRITTEN: 李（字迹不清）
\fieldvalue{\handwritten{李（字迹不清）}}
\quad
% #VALUE_ID: LEX-P0045-V0012
% #FIELD_VALUE: 复核日期
% #HANDWRITTEN: 2025.12.13
\fieldvalue{\handwritten{2025.12.13}}\\
\hline
\end{tabular}
\end{center}

\noindent
4.4.2.\quad 冻存袋热合操作

\begin{center}
\small
\begin{tabular}{|p{0.55\linewidth}|p{0.37\linewidth}|}
\hline
\begin{minipage}[t]{\linewidth}
1.检查管路无破损、无泄漏，确认标签粘贴良好，冻存袋装量目视无异常；\\
2.用封管机将冻存袋一侧的EVA管路热合三次，用剪刀将第二次热合处位置剪开；\\
3.检查封口处及细胞冻存袋无泄漏，非标准装袋原因，并按如下规则确认标签标记记录：\\
\hspace{1em}装量不足的冻存袋标记为B；\\
\hspace{1em}有破损的细胞冻存袋，标记“$\times$”；
\end{minipage}
&
1.热合开始时间：
% #VALUE_ID: LEX-P0045-V0013
% #FIELD_VALUE: 热合开始时间
% #HANDWRITTEN: 10：37
\fieldvalue{\handwritten{10：37}}\newline
2.第1轮分装冻存袋数量：
% #VALUE_ID: LEX-P0045-V0014
% #FIELD_VALUE: 第1轮分装冻存袋数量
% #HANDWRITTEN: 10
\fieldvalue{\handwritten{10}}袋；\newline
\hspace{1em}分装合格冻存袋数量：
% #VALUE_ID: LEX-P0045-V0015
% #FIELD_VALUE: 第1轮分装合格冻存袋数量
% #HANDWRITTEN: 0
\fieldvalue{\handwritten{0}}袋；\newline
\hspace{1em}标记为B的冻存袋数：
% #VALUE_ID: LEX-P0045-V0016
% #FIELD_VALUE: 第1轮标记为B的冻存袋数
% #HANDWRITTEN: 0
\fieldvalue{\handwritten{0}}个；\newline
\hspace{1em}标记为“$\times$”的冻存袋数：
% #VALUE_ID: LEX-P0045-V0017
% #FIELD_VALUE: 第1轮标记为“$\times$”的冻存袋数
% #HANDWRITTEN: 0
\fieldvalue{\handwritten{0}}个。\newline
3.第2轮分装冻存袋数量：
% #VALUE_ID: LEX-P0045-V0018
% #FIELD_VALUE: 第2轮分装冻存袋数量
% #HANDWRITTEN: 10
\fieldvalue{\handwritten{10}}个；\newline
\hspace{1em}分装合格冻存袋数量：
% #VALUE_ID: LEX-P0045-V0019
% #FIELD_VALUE: 第2轮分装合格冻存袋数量
% #HANDWRITTEN: 10
\fieldvalue{\handwritten{10}}个；\newline
\hspace{1em}标记为B的冻存袋数：
% #VALUE_ID: LEX-P0045-V0020
% #FIELD_VALUE: 第2轮标记为B的冻存袋数
% #HANDWRITTEN: 0
\fieldvalue{\handwritten{0}}个；\newline
\hspace{1em}标记为“$\times$”的冻存袋数：
% #VALUE_ID: LEX-P0045-V0021
% #FIELD_VALUE: 第2轮标记为“$\times$”的冻存袋数
% #HANDWRITTEN: 0
\fieldvalue{\handwritten{0}}个。\newline
4.第3轮分装冻存袋数量：
% #VALUE_ID: LEX-P0045-V0022
% #FIELD_VALUE: 第3轮分装冻存袋数量
% #HANDWRITTEN: 10
\fieldvalue{\handwritten{10}}个；\newline
\hspace{1em}分装合格冻存袋数量：
% #VALUE_ID: LEX-P0045-V0023
% #FIELD_VALUE: 第3轮分装合格冻存袋数量
% #HANDWRITTEN: 10
\fieldvalue{\handwritten{10}}个；\newline
\hspace{1em}标记为B的冻存袋数：
% #VALUE_ID: LEX-P0045-V0024
% #FIELD_VALUE: 第3轮标记为B的冻存袋数
% #HANDWRITTEN: 0
\fieldvalue{\handwritten{0}}个；\newline
\hspace{1em}标记为“$\times$”的冻存袋数：
% #VALUE_ID: LEX-P0045-V0025
% #FIELD_VALUE: 第3轮标记为“$\times$”的冻存袋数
% #HANDWRITTEN: 0
\fieldvalue{\handwritten{0}}个。\newline
4.热合结束时间：
% #VALUE_ID: LEX-P0045-V0026
% #FIELD_VALUE: 热合结束时间
% #HANDWRITTEN: 16：31
\fieldvalue{\handwritten{16：31}}
\\
\hline
备注：
% #VALUE_ID: LEX-P0045-V0027
% #FIELD_VALUE: 备注
% #TODO #HANDWRITTEN: /；疑似手写斜线
\fieldvalue{\handwritten{/}}
&
\\
\hline
操作人/日期：
% #VALUE_ID: LEX-P0045-V0028
% #FIELD_VALUE: 操作人
% #HANDWRITTEN: 周（字迹不清）
\fieldvalue{\handwritten{周（字迹不清）}}
\quad
% #VALUE_ID: LEX-P0045-V0029
% #FIELD_VALUE: 操作日期
% #HANDWRITTEN: 2025.12.13
\fieldvalue{\handwritten{2025.12.13}}
&
复核人/日期：
% #VALUE_ID: LEX-P0045-V0030
% #FIELD_VALUE: 复核人
% #HANDWRITTEN: 李（字迹不清）
\fieldvalue{\handwritten{李（字迹不清）}}
\quad
% #VALUE_ID: LEX-P0045-V0031
% #FIELD_VALUE: 复核日期
% #HANDWRITTEN: 2025.12.13
\fieldvalue{\handwritten{2025.12.13}}\\
\hline
\end{tabular}
\end{center}

\begin{center}
第 11 页 共 23 页
\end{center}
% LEXOID_PAGE_COMPLETED: 45/70

\newpage

\noindent
\begin{tabular}{|p{0.22\linewidth}|p{0.23\linewidth}|p{0.25\linewidth}|p{0.16\linewidth}|}
\hline
品名 & 规格 & 生产批号 & 批量 \\
\hline
% #VALUE_ID: LEX-P0046-V0001
% #FIELD_VALUE: 品名
\fieldvalue{艾米近托塞针注射液(22)无菌工艺模拟试验}
&
% #VALUE_ID: LEX-P0046-V0002
% #FIELD_VALUE: 规格
\fieldvalue{12.5$\pm$0.5 ml}
&
% #VALUE_ID: LEX-P0046-V0003
% #FIELD_VALUE: 生产批号
% #TODO #HANDWRITTEN: APS?0?/20?6?24; 字迹不清
\fieldvalue{\handwritten{APS?0?/20?6?24}}
&
袋 \\
\hline
\end{tabular}

\noindent
记录编号：YZ-VP-PD-APS-012(2502)-07 \qquad 版本号：1.0

\section*{4.5\quad 分装模块4}
% #VALUE_ID: LEX-P0046-V0004
% #HANDWRITTEN: 4#隔离器，2025.12.13
% #TODO #HANDWRITTEN: 4#隔离器，2025.12.13; 部分字迹不清
\handwritten{4\#隔离器，2025.12.13}

\subsection*{4.5.1\quad 分装操作}

\noindent
\begin{tabular}{|p{0.56\linewidth}|p{0.35\linewidth}|}
\hline
\textbf{项目} & \textbf{操作记录} \\
\hline
1. 分装过程：整个分装过程均由分装模块自动运行，程序自动预搅拌、自动制剂和排空管路。

2. 分装模块第一轮分装结束后，断开分装模块杜材管路末端的排气阀与设备自带的管路连接，旋转金属卡扣至垂直位置，然后对过管路后取下耗材。

3. 点击主机工艺界面中对应模块的“耗材安装1”，分装模块1状态完成复位后，按1.3耗材安装中2）耗材安装步骤的新分装耗材，打开对应管路的白色管夹，点击耗材确认按键继续分装；

4. 每个模块分装进行编号。（如：A号模块 GentIe P-Pac Pro 第1次分装：A-1、A号模块 GentIe P-Pac Pro 第2次分装：A-2、B号 GentIe P-Pac Pro 第3次分装：B-3等） &
1. 分装开始时间：
% #VALUE_ID: LEX-P0046-V0005
% #FIELD_VALUE: 分装开始时间（小时）
\fieldvalue{\handwritten{10}}
：
% #VALUE_ID: LEX-P0046-V0006
% #FIELD_VALUE: 分装开始时间（分钟）
\fieldvalue{\handwritten{24}}

2. 分装模块完成
% #VALUE_ID: LEX-P0046-V0007
% #FIELD_VALUE: 分装模块完成循环分装次数
\fieldvalue{\handwritten{2}}
次循环分装；

3. 分装结束时间：
% #VALUE_ID: LEX-P0046-V0008
% #FIELD_VALUE: 分装结束时间（小时）
\fieldvalue{\handwritten{10}}
：
% #VALUE_ID: LEX-P0046-V0009
% #FIELD_VALUE: 分装结束时间（分钟）
\fieldvalue{\handwritten{49}} \\
\hline
备注： & \\
\hline
\end{tabular}

\noindent
操作人/日期：
% #VALUE_ID: LEX-P0046-V0010
% #FIELD_VALUE: 操作人
% #HANDWRITTEN: 何...
\fieldvalue{\handwritten{何...}}
\quad
% #VALUE_ID: LEX-P0046-V0011
% #FIELD_VALUE: 操作日期
% #HANDWRITTEN: 2025.12.13
\fieldvalue{\handwritten{2025.12.13}}
\qquad
复核人/日期：
% #VALUE_ID: LEX-P0046-V0012
% #FIELD_VALUE: 复核人
% #HANDWRITTEN: 徐...
\fieldvalue{\handwritten{徐...}}
\quad
% #VALUE_ID: LEX-P0046-V0013
% #FIELD_VALUE: 复核日期
% #HANDWRITTEN: 2025.12.13
\fieldvalue{\handwritten{2025.12.13}}

\subsection*{4.5.2\quad 冻存袋状态}

\noindent
\begin{tabular}{|p{0.56\linewidth}|p{0.35\linewidth}|}
\hline
1. 检查管路无破损、无泄漏，确认标签粘贴良好，冻存袋数量目视无异常；

2. 用封管机将冻存袋一侧的 EVA 管路热合三次，用剪刀将第二次热合处位置剪开；

3. 检查封口处及细胞冻存袋无泄漏，非标准袋备用原因，并按如下规则确认标签记录：

\quad 装置不足的冻存袋标记为 B；

\quad 有破损的细胞冻存袋，标记“$\times$”； &
1. 热合开始时间：
% #VALUE_ID: LEX-P0046-V0014
% #FIELD_VALUE: 热合开始时间（小时）
\fieldvalue{\handwritten{10}}
：
% #VALUE_ID: LEX-P0046-V0015
% #FIELD_VALUE: 热合开始时间（分钟）
\fieldvalue{\handwritten{38}}

2. 第1轮分装存液袋数量：
% #VALUE_ID: LEX-P0046-V0016
% #FIELD_VALUE: 第1轮分装存液袋数量
\fieldvalue{\handwritten{10}}
袋；

\quad 分装合格存液袋数：
% #VALUE_ID: LEX-P0046-V0017
% #FIELD_VALUE: 第1轮分装合格存液袋数
\fieldvalue{\handwritten{0}}
袋；

\quad 标记为B的冻存袋数：
% #VALUE_ID: LEX-P0046-V0018
% #FIELD_VALUE: 第1轮标记为B的冻存袋数
\fieldvalue{\handwritten{0}}
个；

\quad 标记为“$\times$”的冻存袋数：
% #VALUE_ID: LEX-P0046-V0019
% #FIELD_VALUE: 第1轮标记为“$\times$”的冻存袋数
\fieldvalue{\handwritten{0}}
个。

3. 第2轮分装存液袋数量：
% #VALUE_ID: LEX-P0046-V0020
% #FIELD_VALUE: 第2轮分装存液袋数量
\fieldvalue{\handwritten{10}}
个；

\quad 分装合格存液袋数：
% #VALUE_ID: LEX-P0046-V0021
% #FIELD_VALUE: 第2轮分装合格存液袋数
\fieldvalue{\handwritten{0}}
个；

\quad 标记为B的冻存袋数：
% #VALUE_ID: LEX-P0046-V0022
% #FIELD_VALUE: 第2轮标记为B的冻存袋数
\fieldvalue{\handwritten{0}}
个；

\quad 标记为“$\times$”的冻存袋数：
% #VALUE_ID: LEX-P0046-V0023
% #FIELD_VALUE: 第2轮标记为“$\times$”的冻存袋数
\fieldvalue{\handwritten{0}}
个。

4. 第3轮分装存液袋数量：
% #VALUE_ID: LEX-P0046-V0024
% #FIELD_VALUE: 第3轮分装存液袋数量
\fieldvalue{\handwritten{10}}
个；

\quad 分装合格存液袋数：
% #VALUE_ID: LEX-P0046-V0025
% #FIELD_VALUE: 第3轮分装合格存液袋数
\fieldvalue{\handwritten{10}}
个；

\quad 标记为B的冻存袋数：
% #VALUE_ID: LEX-P0046-V0026
% #FIELD_VALUE: 第3轮标记为B的冻存袋数
\fieldvalue{\handwritten{0}}
个；

\quad 标记为“$\times$”的冻存袋数：
% #VALUE_ID: LEX-P0046-V0027
% #FIELD_VALUE: 第3轮标记为“$\times$”的冻存袋数
\fieldvalue{\handwritten{0}}
个。

4. 结合结束时间：
% #VALUE_ID: LEX-P0046-V0028
% #FIELD_VALUE: 结合结束时间（小时）
\fieldvalue{\handwritten{16}}
：
% #VALUE_ID: LEX-P0046-V0029
% #FIELD_VALUE: 结合结束时间（分钟）
\fieldvalue{\handwritten{32}} \\
\hline
备注： & \\
\hline
\end{tabular}

\noindent
操作人/日期：
% #VALUE_ID: LEX-P0046-V0030
% #FIELD_VALUE: 操作人
% #HANDWRITTEN: 何...
\fieldvalue{\handwritten{何...}}
\quad
% #VALUE_ID: LEX-P0046-V0031
% #FIELD_VALUE: 操作日期
% #HANDWRITTEN: 2025.12.13
\fieldvalue{\handwritten{2025.12.13}}
\qquad
复核人/日期：
% #VALUE_ID: LEX-P0046-V0032
% #FIELD_VALUE: 复核人
% #HANDWRITTEN: 徐...
\fieldvalue{\handwritten{徐...}}
\quad
% #VALUE_ID: LEX-P0046-V0033
% #FIELD_VALUE: 复核日期
% #HANDWRITTEN: 2025.12.13
\fieldvalue{\handwritten{2025.12.13}}
% LEXOID_PAGE_COMPLETED: 46/70

\newpage

\begin{center}
\small
\begin{tabular}{|p{0.25\linewidth}|p{0.20\linewidth}|p{0.25\linewidth}|p{0.16\linewidth}|}
\hline
品名 & 规格 & 生产批号 & 批量 \\
\hline
艾克利托莫司注射液(2Z)无菌工艺模拟试验
&
12.5$\pm$0.5\,ml
&
% #VALUE_ID: LEX-P0047-V0001
% #FIELD_VALUE: 生产批号
% #TODO #HANDWRITTEN: APS20251016(?); handwritten batch suffix is unclear
\fieldvalue{\handwritten{APS20251016(?)}} 
&
% #VALUE_ID: LEX-P0047-V0002
% #FIELD_VALUE: 批量
% #TODO #HANDWRITTEN: 720; final digit is partially unclear
\fieldvalue{\handwritten{720}} 袋
\\
\hline
\end{tabular}
\end{center}

记录编号：YZ-VP-PD-APS-012(2502)-07 \quad 版本号：1.0

\section*{4.6\quad 分装模块5}

% #VALUE_ID: LEX-P0047-V0003
% #HANDWRITTEN: B1
\handwritten{B1}

% #VALUE_ID: LEX-P0047-V0004
% #TODO #HANDWRITTEN: 限?; handwritten annotation is unclear
\handwritten{限?}\quad
% #VALUE_ID: LEX-P0047-V0005
% #HANDWRITTEN: 2025.12.13
\handwritten{2025.12.13}

\subsection*{4.6.1\quad 分装操作}

\begin{center}
\small
\begin{tabular}{|p{0.57\linewidth}|p{0.34\linewidth}|}
\hline
项目 & 操作记录 \\
\hline
\begin{minipage}[t]{\linewidth}
1.\ 分装过程：整个分装过程均由分装模块自动运行，程序自动预排气、自动判别和排空管路；

2.\ 分装模块第一轮分装结束后，断开分装模块耗材管路末端的排气阀与设备自带的管路连接，旋转金属卡扣至垂直位置，然后对对应管路段下料并卡板；

3.\ 点击主机工艺界面中对应模块的“耗材安装1”，分装模块1状态完成复位后，按1.3耗材安装2后再对首步耗材安装新的分装耗材，打开对应管路的白色管夹，点击耗材确认按键继续分装；

每个模块每次分装进行编号。（如：A号模块 Gentle P-Pac Pro 第1次分装：A-1，A号模块第2次分装：A-2，B号 Gentle P-Pac Pro 第3次分装：B-3 等）；
\end{minipage}
&
\begin{minipage}[t]{\linewidth}
1.分装开始时间：
% #VALUE_ID: LEX-P0047-V0006
% #FIELD_VALUE: 分装开始时间
% #HANDWRITTEN: 10:24
\fieldvalue{\handwritten{10:24}}

2.分装模块共完成
% #VALUE_ID: LEX-P0047-V0007
% #FIELD_VALUE: 分装模块完成循环次数
% #TODO #HANDWRITTEN: 2; handwriting is faint
\fieldvalue{\handwritten{2}}
次循环分装；

3.装填结束时间：
% #VALUE_ID: LEX-P0047-V0008
% #FIELD_VALUE: 装填结束时间
% #HANDWRITTEN: 12:19
\fieldvalue{\handwritten{12:19}}
\end{minipage}
\\
\hline
备注：
% #VALUE_ID: LEX-P0047-V0009
% #FIELD_VALUE: 备注
% #HANDWRITTEN: /
\fieldvalue{\handwritten{/}}
&
\\
\hline
\end{tabular}
\end{center}

操作人/日期：
% #VALUE_ID: LEX-P0047-V0010
% #TODO #HANDWRITTEN: 张?; handwritten name is unclear
\fieldvalue{\handwritten{张?}}
\quad
% #VALUE_ID: LEX-P0047-V0011
% #HANDWRITTEN: 2025.12.13
\fieldvalue{\handwritten{2025.12.13}}
\qquad
复核人/日期：
% #VALUE_ID: LEX-P0047-V0012
% #TODO #HANDWRITTEN: 陈?; handwritten name is unclear
\fieldvalue{\handwritten{陈?}}
\quad
% #VALUE_ID: LEX-P0047-V0013
% #HANDWRITTEN: 2025.12.13
\fieldvalue{\handwritten{2025.12.13}}

\subsection*{4.6.2\quad 冻存袋热合操作}

\begin{center}
\small
\begin{tabular}{|p{0.55\linewidth}|p{0.37\linewidth}|}
\hline
\begin{minipage}[t]{\linewidth}
1.检查管路无破损、无泄漏，确认标签粘贴完好，冻存袋装量目视异常；

2.用封管机将冻存袋一侧的 EVA 管路热合三次，用剪刀将第二次热合处位置剪开；

3.检查封口处及细胞冻存袋无泄漏，非标准装备经分原因，并按如下规则确认标签标记起止记录：

\begin{itemize}
\item 装量不足的冻存袋标记为B；
\item 有破损的细胞冻存袋，标记“$\times$”；
\end{itemize}
\end{minipage}
&
\begin{minipage}[t]{\linewidth}
1 热合开始时间：
% #VALUE_ID: LEX-P0047-V0014
% #FIELD_VALUE: 热合开始时间
% #HANDWRITTEN: 10:45
\fieldvalue{\handwritten{10:45}}

2.第1轮分装冻存袋数量：
% #VALUE_ID: LEX-P0047-V0015
% #FIELD_VALUE: 第1轮分装冻存袋数量
% #HANDWRITTEN: 10
\fieldvalue{\handwritten{10}} 袋

分装合格冻存袋数量：
% #VALUE_ID: LEX-P0047-V0016
% #FIELD_VALUE: 第1轮分装合格冻存袋数量
% #HANDWRITTEN: 10
\fieldvalue{\handwritten{10}} 袋

标记为B的冻存袋数：
% #VALUE_ID: LEX-P0047-V0017
% #FIELD_VALUE: 第1轮标记为B的冻存袋数
% #HANDWRITTEN: 0
\fieldvalue{\handwritten{0}} 个

标记为“$\times$”的冻存袋数：
% #VALUE_ID: LEX-P0047-V0018
% #FIELD_VALUE: 第1轮标记为“$\times$”的冻存袋数
% #HANDWRITTEN: 0
\fieldvalue{\handwritten{0}} 个

2.第2轮分装冻存袋数量：
% #VALUE_ID: LEX-P0047-V0019
% #FIELD_VALUE: 第2轮分装冻存袋数量
% #HANDWRITTEN: 10
\fieldvalue{\handwritten{10}} 个

分装合格冻存袋数量：
% #VALUE_ID: LEX-P0047-V0020
% #FIELD_VALUE: 第2轮分装合格冻存袋数量
% #HANDWRITTEN: 10
\fieldvalue{\handwritten{10}} 个

标记为B的冻存袋数：
% #VALUE_ID: LEX-P0047-V0021
% #FIELD_VALUE: 第2轮标记为B的冻存袋数
% #HANDWRITTEN: 0
\fieldvalue{\handwritten{0}} 个

标记为“$\times$”的冻存袋数：
% #VALUE_ID: LEX-P0047-V0022
% #FIELD_VALUE: 第2轮标记为“$\times$”的冻存袋数
% #HANDWRITTEN: 0
\fieldvalue{\handwritten{0}} 个

3.第3轮分装冻存袋数量：
% #VALUE_ID: LEX-P0047-V0023
% #FIELD_VALUE: 第3轮分装冻存袋数量
% #HANDWRITTEN: 10
\fieldvalue{\handwritten{10}} 个

分装合格冻存袋数量：
% #VALUE_ID: LEX-P0047-V0024
% #FIELD_VALUE: 第3轮分装合格冻存袋数量
% #HANDWRITTEN: 10
\fieldvalue{\handwritten{10}} 个

标记为B的冻存袋数：
% #VALUE_ID: LEX-P0047-V0025
% #FIELD_VALUE: 第3轮标记为B的冻存袋数
% #HANDWRITTEN: 0
\fieldvalue{\handwritten{0}} 个

标记为“$\times$”的冻存袋数：
% #VALUE_ID: LEX-P0047-V0026
% #FIELD_VALUE: 第3轮标记为“$\times$”的冻存袋数
% #HANDWRITTEN: 0
\fieldvalue{\handwritten{0}} 个

4.热合结束时间：
% #VALUE_ID: LEX-P0047-V0027
% #FIELD_VALUE: 热合结束时间
% #HANDWRITTEN: 16:30
\fieldvalue{\handwritten{16:30}}
\end{minipage}
\\
\hline
\end{tabular}
\end{center}

操作人/日期：
% #VALUE_ID: LEX-P0047-V0028
% #TODO #HANDWRITTEN: 陈?; handwritten name is unclear
\fieldvalue{\handwritten{陈?}}
\quad
% #VALUE_ID: LEX-P0047-V0029
% #HANDWRITTEN: 2025.12.13
\fieldvalue{\handwritten{2025.12.13}}
\qquad
复核人/日期：
% #VALUE_ID: LEX-P0047-V0030
% #TODO #HANDWRITTEN: 潘?; handwritten name is unclear
\fieldvalue{\handwritten{潘?}}
\quad
% #VALUE_ID: LEX-P0047-V0031
% #HANDWRITTEN: 2025.12.13
\fieldvalue{\handwritten{2025.12.13}}

\begin{center}
第 13 页 共 23 页
\end{center}
% LEXOID_PAGE_COMPLETED: 47/70

\newpage

\begin{center}
记录编号：YZ-VP-PD-APS-012(2502)-07 \quad 版本号：1.0
\end{center}

\begin{tabular}{|p{0.25\linewidth}|p{0.20\linewidth}|p{0.25\linewidth}|p{0.18\linewidth}|}
\hline
品名 & 规格 & 生产批号 & 批量 \\
\hline
% #VALUE_ID: LEX-P0048-V0001
% #FIELD_VALUE: 品名
\fieldvalue{艾米沃托萘法液（Z2）无菌工艺模拟试验}
&
% #VALUE_ID: LEX-P0048-V0002
% #FIELD_VALUE: 规格
\fieldvalue{12.5$\pm$0.5 ml}
&
% #VALUE_ID: LEX-P0048-V0003
% #FIELD_VALUE: 生产批号
% #TODO #HANDWRITTEN: APS2025120164?; handwritten batch number is partially unclear
\fieldvalue{\handwritten{APS2025120164?}}
&
% #VALUE_ID: LEX-P0048-V0004
% #FIELD_VALUE: 批量
% #HANDWRITTEN: $\geq$240
\fieldvalue{\handwritten{$\geq$240}}袋 \\
\hline
\end{tabular}

\vspace{0.5em}

4.7．分装模块6
\quad
% #VALUE_ID: LEX-P0048-V0005
% #FIELD_VALUE: 分装模块
% #HANDWRITTEN: B2
\fieldvalue{\handwritten{B2}}
\quad
% #VALUE_ID: LEX-P0048-V0006
% #FIELD_VALUE: 日期
% #HANDWRITTEN: 2025.12.13
\fieldvalue{\handwritten{2025.12.13}}

\subsection*{4.7.1．分装操作}

\begin{tabular}{|p{0.57\linewidth}|p{0.35\linewidth}|}
\hline
项目 & 操作记录 \\
\hline
1．分装过程：整个分装过程均由分装模块自动运行，程序自动预排气、自动判别和排空管路。

2．分装模块第一轮分装结束后，断开分装模块耗材管路末端的排气阀与设备自带的管路连接，旋转金属卡扣至垂直位置，然后对应该容器取样卡板；

3．点击主机工艺界面中对应模块的“耗材安装1”，分装模块1状态完成复位后，按1.3耗材安装中2）耗材安装步骤安装新的分装耗材，打开对应管路的白色管夹，点击耗材确认按钮继续分装；

每个模块每次分装进行编号，（如：A号模块Gentle P-Pac Pro第1次分装：A-1，A号模块第2次分装：A-2，B号Gentle P-Pac Pro第3次分装：B-3等）； &
1．分装开始时间：
% #VALUE_ID: LEX-P0048-V0007
% #FIELD_VALUE: 分装开始时间
% #HANDWRITTEN: 10:24
\fieldvalue{\handwritten{10:24}}

2．分装模块完成
% #VALUE_ID: LEX-P0048-V0008
% #FIELD_VALUE: 分装模块完成时间
% #HANDWRITTEN: 10:19
\fieldvalue{\handwritten{10:19}}

\\
\cline{2-2}
备注：
% #VALUE_ID: LEX-P0048-V0009
% #FIELD_VALUE: 备注
% #HANDWRITTEN: /
\fieldvalue{\handwritten{/}}
&
\\
\hline
\end{tabular}

\vspace{0.3em}

操作人/日期：
% #VALUE_ID: LEX-P0048-V0010
% #FIELD_VALUE: 操作人
% #HANDWRITTEN: 邹凯
\fieldvalue{\handwritten{邹凯}}
\quad
% #VALUE_ID: LEX-P0048-V0011
% #FIELD_VALUE: 操作日期
% #HANDWRITTEN: 2025.12.13
\fieldvalue{\handwritten{2025.12.13}}
\qquad
复核人/日期：
% #VALUE_ID: LEX-P0048-V0012
% #FIELD_VALUE: 复核人
% #HANDWRITTEN: 徐琼
\fieldvalue{\handwritten{徐琼}}
\quad
% #VALUE_ID: LEX-P0048-V0013
% #FIELD_VALUE: 复核日期
% #HANDWRITTEN: 2025.12.13
\fieldvalue{\handwritten{2025.12.13}}

\subsection*{4.7.2．冻存袋热合操作}

\begin{tabular}{|p{0.56\linewidth}|p{0.37\linewidth}|}
\hline
\begin{minipage}[t]{0.54\linewidth}
1．检查管路无破损、无泄漏，确认标签粘贴完好，冻存袋装置目视无异常；

2．用封管机将冻存袋一侧的EVA管路热合三次，用剪刀将第二次热合处位置剪开；

3．检查封口处及细胞冻存袋无泄漏，非标准装备原因，并按如下规则确认冻存袋标记是否正确：

\hspace{1em}装置不足的冻存袋标记为B；

\hspace{1em}有破损的细胞冻存袋，标记“$\times$”；
\end{minipage}
&
1．热合开始时间：
% #VALUE_ID: LEX-P0048-V0014
% #FIELD_VALUE: 热合开始时间
% #HANDWRITTEN: 10:45
\fieldvalue{\handwritten{10:45}}

2．第1轮分装冻存袋数量：
% #VALUE_ID: LEX-P0048-V0015
% #FIELD_VALUE: 第1轮分装冻存袋数量
% #HANDWRITTEN: 10
\fieldvalue{\handwritten{10}}袋；

分装合格冻存袋数量：
% #VALUE_ID: LEX-P0048-V0016
% #FIELD_VALUE: 第1轮分装合格冻存袋数量
% #HANDWRITTEN: 0
\fieldvalue{\handwritten{0}}袋；

标记为B的冻存袋数：
% #VALUE_ID: LEX-P0048-V0017
% #FIELD_VALUE: 第1轮标记为B的冻存袋数
% #HANDWRITTEN: 0
\fieldvalue{\handwritten{0}}个；

标记为“$\times$”的冻存袋数：
% #VALUE_ID: LEX-P0048-V0018
% #FIELD_VALUE: 第1轮标记为“$\times$”的冻存袋数
% #HANDWRITTEN: 0
\fieldvalue{\handwritten{0}}个。

3．第2轮分装冻存袋数量：
% #VALUE_ID: LEX-P0048-V0019
% #FIELD_VALUE: 第2轮分装冻存袋数量
% #HANDWRITTEN: 10
\fieldvalue{\handwritten{10}}个；

分装合格冻存袋数量：
% #VALUE_ID: LEX-P0048-V0020
% #FIELD_VALUE: 第2轮分装合格冻存袋数量
% #HANDWRITTEN: 10
\fieldvalue{\handwritten{10}}个；

标记为B的冻存袋数：
% #VALUE_ID: LEX-P0048-V0021
% #FIELD_VALUE: 第2轮标记为B的冻存袋数
% #HANDWRITTEN: 0
\fieldvalue{\handwritten{0}}个；

标记为“$\times$”的冻存袋数：
% #VALUE_ID: LEX-P0048-V0022
% #FIELD_VALUE: 第2轮标记为“$\times$”的冻存袋数
% #HANDWRITTEN: 0
\fieldvalue{\handwritten{0}}个。

4．第3轮分装冻存袋数量：
% #VALUE_ID: LEX-P0048-V0023
% #FIELD_VALUE: 第3轮分装冻存袋数量
% #HANDWRITTEN: 10
\fieldvalue{\handwritten{10}}个；

分装合格冻存袋数量：
% #VALUE_ID: LEX-P0048-V0024
% #FIELD_VALUE: 第3轮分装合格冻存袋数量
% #HANDWRITTEN: 0
\fieldvalue{\handwritten{0}}个；

标记为B的冻存袋数：
% #VALUE_ID: LEX-P0048-V0025
% #FIELD_VALUE: 第3轮标记为B的冻存袋数
% #HANDWRITTEN: 0
\fieldvalue{\handwritten{0}}个；

标记为“$\times$”的冻存袋数：
% #VALUE_ID: LEX-P0048-V0026
% #FIELD_VALUE: 第3轮标记为“$\times$”的冻存袋数
% #HANDWRITTEN: 0
\fieldvalue{\handwritten{0}}个。

4．热合结束时间：
% #VALUE_ID: LEX-P0048-V0027
% #FIELD_VALUE: 热合结束时间
% #HANDWRITTEN: 16:30
\fieldvalue{\handwritten{16:30}}
\\
\hline
备注：
% #VALUE_ID: LEX-P0048-V0028
% #FIELD_VALUE: 备注
% #HANDWRITTEN: /
\fieldvalue{\handwritten{/}}
&
\\
\hline
\end{tabular}

\vspace{0.3em}

操作人/日期：
% #VALUE_ID: LEX-P0048-V0029
% #FIELD_VALUE: 操作人
% #HANDWRITTEN: 院小娟
\fieldvalue{\handwritten{院小娟}}
\quad
% #VALUE_ID: LEX-P0048-V0030
% #FIELD_VALUE: 操作日期
% #HANDWRITTEN: 2025.12.13
\fieldvalue{\handwritten{2025.12.13}}
\qquad
复核人/日期：
% #VALUE_ID: LEX-P0048-V0031
% #FIELD_VALUE: 复核人
% #HANDWRITTEN: 潘珊珊
\fieldvalue{\handwritten{潘珊珊}}
\quad
% #VALUE_ID: LEX-P0048-V0032
% #FIELD_VALUE: 复核日期
% #HANDWRITTEN: 2025.12.13
\fieldvalue{\handwritten{2025.12.13}}

\begin{center}
第 14 页 共 23 页
\end{center}
% LEXOID_PAGE_COMPLETED: 48/70

\newpage

\noindent
\textbf{记录编号：} YZ-VP-ADPS-012(2502)-07\qquad
\textbf{版本号：}1.0

\begin{center}
\small
\begin{tabular}{|p{0.25\linewidth}|p{0.18\linewidth}|p{0.22\linewidth}|p{0.18\linewidth}|}
\hline
\centering 品名 & \centering 规格 & \centering 生产批号 & \centering 批量\\
\hline
艾米沃托霉素注射液(2.2)无菌工艺模拟试验 &
12.5$\pm$0.5\,ml &
% #VALUE_ID: LEX-P0049-V0001
% #FIELD_VALUE: 生产批号
% #HANDWRITTEN: AP202?12/06?23
\fieldvalue{\handwritten{AP202?12/06?23}} &
% #VALUE_ID: LEX-P0049-V0002
% #FIELD_VALUE: 批量
% #HANDWRITTEN: 7,240
\fieldvalue{\handwritten{7,240}}\ 袋\\
\hline
\end{tabular}
\end{center}

\noindent
4.8.\quad 分装模块7\quad
% #VALUE_ID: LEX-P0049-V0003
% #HANDWRITTEN: B3
\handwritten{B3}\quad
% #VALUE_ID: LEX-P0049-V0004
% #HANDWRITTEN: 溶解定容
\handwritten{溶解定容}\quad
% #VALUE_ID: LEX-P0049-V0005
% #HANDWRITTEN: 2025.12.13
\handwritten{2025.12.13}

\medskip
\noindent
4.8.1.\quad 分装操作

\begin{center}
\small
\begin{tabular}{|p{0.55\linewidth}|p{0.37\linewidth}|}
\hline
\centering 项目 & \centering 操作记录\\
\hline
1.分装过程：整个分装过程均由分装模块自动运行，程序自动模拟、自动判别和抽空管路。& 
1.分装开始时间：
% #VALUE_ID: LEX-P0049-V0006
% #FIELD_VALUE: 分装开始时间
% #HANDWRITTEN: 10:24
\fieldvalue{\handwritten{10:24}}\\[2pt]
2.分装模块第一轮分装结束后，断开分装模块耗材管路末端的排气阀与设备自带的管路连接，旋转金属卡扣至垂直位置，然后对对应容器后取样卡板；&
2.分装模块完成一次循环分装：
% #VALUE_ID: LEX-P0049-V0007
% #FIELD_VALUE: 分装模块完成一次循环分装时间
% #HANDWRITTEN: 10:9?
\fieldvalue{\handwritten{10:9?}}\\[2pt]
3.点击主机工艺界面中对应模块的“耗材安装1”，分装模块1状态完成复位后，按1.3耗材安装（2）耗材安装步骤安装新的分装耗材，打开对应管路的白色管夹，点击耗材确认按键继续分装；&
3.分装结束时间：
% #VALUE_ID: LEX-P0049-V0008
% #FIELD_VALUE: 分装结束时间
% #HANDWRITTEN: 10:19?
\fieldvalue{\handwritten{10:19?}}\\
\hline
每个模块每次分装进行编号，（如：A号模块Gentle P-Pac Pro 第1次分装：A-1，A号Gentle P-Pac Pro第2次分装：A-2，B号Gentle P-Pac Pro第3次分装：B-3等）；&
\\
\hline
备注：&\\
\hline
\end{tabular}
\end{center}

\noindent
操作人/日期：
% #VALUE_ID: LEX-P0049-V0009
% #HANDWRITTEN: 2025.12.13
% #FIELD_VALUE: 操作人/日期
\fieldvalue{\handwritten{2025.12.13}}
\hfill
复核人/日期：
% #VALUE_ID: LEX-P0049-V0010
% #HANDWRITTEN: 李厚强
% #FIELD_VALUE: 复核人
\fieldvalue{\handwritten{李厚强}}
\quad
% #VALUE_ID: LEX-P0049-V0011
% #HANDWRITTEN: 2025.12.13
% #FIELD_VALUE: 复核日期
\fieldvalue{\handwritten{2025.12.13}}

\medskip
\noindent
4.8.2.\quad 冻存袋装操作

\begin{center}
\small
\begin{tabular}{|p{0.56\linewidth}|p{0.36\linewidth}|}
\hline
1.检查管路无破损、无泄漏，确认标签粘贴完好，冻存包装量目视无异常；&
1.装合开始时间：
% #VALUE_ID: LEX-P0049-V0012
% #FIELD_VALUE: 装合开始时间
% #HANDWRITTEN: 10:45
\fieldvalue{\handwritten{10:45}}\\[2pt]
2.用封管机将冻存袋一侧的EVA管路热合三次，用剪刀将第二次热合处位置剪开；&
2.第1轮分装冻存袋数量：
% #VALUE_ID: LEX-P0049-V0013
% #FIELD_VALUE: 第1轮分装冻存袋数量
% #HANDWRITTEN: 10
\fieldvalue{\handwritten{10}}\ 袋\\
& 分装合格冻存袋数量：
% #VALUE_ID: LEX-P0049-V0014
% #FIELD_VALUE: 第1轮分装合格冻存袋数量
% #HANDWRITTEN: 10
\fieldvalue{\handwritten{10}}\ 袋\\
& 标记为B的冻存袋数：
% #VALUE_ID: LEX-P0049-V0015
% #FIELD_VALUE: 第1轮标记为B的冻存袋数
% #HANDWRITTEN: 0
\fieldvalue{\handwritten{0}}\ 个\\
& 标记为“$\times$”的冻存袋数：
% #VALUE_ID: LEX-P0049-V0016
% #FIELD_VALUE: 第1轮标记为“×”的冻存袋数
% #HANDWRITTEN: 0
\fieldvalue{\handwritten{0}}\ 个\\[2pt]
3.检查封口处及细胞冻存袋无泄漏，非标准装备请原因，并按如下规则确认冻存袋标记记录：&
3.第2轮分装冻存袋数量：
% #VALUE_ID: LEX-P0049-V0017
% #FIELD_VALUE: 第2轮分装冻存袋数量
% #HANDWRITTEN: 10
\fieldvalue{\handwritten{10}}\ 个\\
& 分装合格冻存袋数量：
% #VALUE_ID: LEX-P0049-V0018
% #FIELD_VALUE: 第2轮分装合格冻存袋数量
% #HANDWRITTEN: 10
\fieldvalue{\handwritten{10}}\ 个\\
& 标记为B的冻存袋数：
% #VALUE_ID: LEX-P0049-V0019
% #FIELD_VALUE: 第2轮标记为B的冻存袋数
% #HANDWRITTEN: 0
\fieldvalue{\handwritten{0}}\ 个\\
& 标记为“$\times$”的冻存袋数：
% #VALUE_ID: LEX-P0049-V0020
% #FIELD_VALUE: 第2轮标记为“×”的冻存袋数
% #HANDWRITTEN: 0
\fieldvalue{\handwritten{0}}\ 个\\
& \textit{装量不足的冻存袋标记为B；}\\
& \textit{有破损的细胞冻存袋，标记“$\times$”；}\\[2pt]
& 3.第3轮分装冻存袋数量：
% #VALUE_ID: LEX-P0049-V0021
% #FIELD_VALUE: 第3轮分装冻存袋数量
% #HANDWRITTEN: 10
\fieldvalue{\handwritten{10}}\ 个\\
& 分装合格冻存袋数量：
% #VALUE_ID: LEX-P0049-V0022
% #FIELD_VALUE: 第3轮分装合格冻存袋数量
% #HANDWRITTEN: 10
\fieldvalue{\handwritten{10}}\ 个\\
& 标记为B的冻存袋数：
% #VALUE_ID: LEX-P0049-V0023
% #FIELD_VALUE: 第3轮标记为B的冻存袋数
% #HANDWRITTEN: 0
\fieldvalue{\handwritten{0}}\ 个\\
& 标记为“$\times$”的冻存袋数：
% #VALUE_ID: LEX-P0049-V0024
% #FIELD_VALUE: 第3轮标记为“×”的冻存袋数
% #HANDWRITTEN: 0
\fieldvalue{\handwritten{0}}\ 个\\
\hline
\end{tabular}
\end{center}

\noindent
装量不足的冻存袋标记为B；有破损的细胞冻存袋标记“$\times$”。

\noindent
% #VALUE_ID: LEX-P0049-V0025
% #HANDWRITTEN: 冻存袋装合记录2025.12.13
\handwritten{冻存袋装合记录2025.12.13}

\begin{center}
\small
\begin{tabular}{|p{0.56\linewidth}|p{0.36\linewidth}|}
\hline
备注：&
4.共计结束时间：
% #VALUE_ID: LEX-P0049-V0026
% #FIELD_VALUE: 共计结束时间
% #HANDWRITTEN: 16:30
\fieldvalue{\handwritten{16:30}}\\
\hline
\end{tabular}
\end{center}

\noindent
操作人/日期：
% #VALUE_ID: LEX-P0049-V0027
% #HANDWRITTEN: 陈晓霞
% #FIELD_VALUE: 操作人
\fieldvalue{\handwritten{陈晓霞}}
\quad
% #VALUE_ID: LEX-P0049-V0028
% #HANDWRITTEN: 2025.12.13
% #FIELD_VALUE: 操作日期
\fieldvalue{\handwritten{2025.12.13}}
\hfill
复核人/日期：
% #VALUE_ID: LEX-P0049-V0029
% #HANDWRITTEN: 侯向军
% #FIELD_VALUE: 复核人
\fieldvalue{\handwritten{侯向军}}
\quad
% #VALUE_ID: LEX-P0049-V0030
% #HANDWRITTEN: 2025.12.13
% #FIELD_VALUE: 复核日期
\fieldvalue{\handwritten{2025.12.13}}

\begin{center}
第 15 页 共 23 页
\end{center}
% LEXOID_PAGE_COMPLETED: 49/70

\newpage

\begin{center}
\small
\begin{tabular}{|p{0.34\linewidth}|p{0.20\linewidth}|p{0.22\linewidth}|p{0.12\linewidth}|}
\hline
品名 & 规格 & 生产批号 & 批量 \\ \hline
艾米近红霉素注射液(ZZ)无菌工艺模拟试验
&
12.5$\pm$0.5 ml
&
% #VALUE_ID: LEX-P0050-V0001
% #FIELD_VALUE: 生产批号
% #TODO #HANDWRITTEN: AP202?201?6?2>240; handwriting partially illegible
\fieldvalue{\handwritten{AP202?201?6?2>240}}
&
袋
\\ \hline
\end{tabular}
\end{center}

\subsection*{4.9\quad 分装模块8}
% #VALUE_ID: LEX-P0050-V0002
% #FIELD_VALUE: 4.9旁注
% #TODO #HANDWRITTEN: B4 活塞……2025.12.13; handwriting partially illegible
\handwritten{B4\ 活塞……2025.12.13}

\subsubsection*{4.9.1\quad 分装操作}

\begin{center}
\small
\begin{tabular}{|p{0.55\linewidth}|p{0.35\linewidth}|}
\hline
\multicolumn{1}{|c|}{项目} & \multicolumn{1}{c|}{操作记录} \\ \hline
1.\ 分装过程：整个分装过程均由分装模块自动运行，程序自动插袋、自动剔料和料空报警。\newline
2.\ 分装模块第一轮分装结束后，断开分装模块耗材管路末端的排气阀与设备自带的管路连接，旋转金属卡扣至垂直位置，然后对应该管路端下耗材卡板。\newline
3.\ 点击主机工艺界面中对应模块的“耗材安装1”，分装模块1状态完成复位后，按1.3耗材安装+2耗材安装步骤安装新的分装耗材；打开对应管路的白色管夹，点击耗材确认按钮继续分装；\newline
每个模块每次分装进行编号，（如：A号模块 Gentle P-Pac Pro 第1次分装：A-1，A号 Gentle P-Pac Pro 第2次分装：A-2，B号 Gentle P-Pac Pro 第3次分装：B-3等）； & 
1.\ 分装开始时间：
% #VALUE_ID: LEX-P0050-V0003
% #FIELD_VALUE: 分装开始时间
% #HANDWRITTEN: 10:24
\fieldvalue{\handwritten{10:24}}\newline
2.\ 分装模块完成
% #VALUE_ID: LEX-P0050-V0004
% #FIELD_VALUE: 分装模块完成的循环次数
% #HANDWRITTEN: 2
\fieldvalue{\handwritten{2}} 次循环分装：\newline
3.\ 分装结束时间：
% #VALUE_ID: LEX-P0050-V0005
% #FIELD_VALUE: 分装结束时间
% #HANDWRITTEN: 12:19
\fieldvalue{\handwritten{12:19}}
\\ \hline
备注： & \\ \hline
\end{tabular}
\end{center}

\noindent
操作人/日期：
% #VALUE_ID: LEX-P0050-V0006
% #FIELD_VALUE: 操作人
% #TODO #HANDWRITTEN: 赵?; signature is partially illegible
\fieldvalue{\handwritten{赵?}}
\quad
% #VALUE_ID: LEX-P0050-V0007
% #FIELD_VALUE: 操作日期
% #HANDWRITTEN: 2025.12.13
\fieldvalue{\handwritten{2025.12.13}}
\hfill
复核人/日期：
% #VALUE_ID: LEX-P0050-V0008
% #FIELD_VALUE: 复核人
% #TODO #HANDWRITTEN: 张?; signature is partially illegible
\fieldvalue{\handwritten{张?}}
\quad
% #VALUE_ID: LEX-P0050-V0009
% #FIELD_VALUE: 复核日期
% #HANDWRITTEN: 2025.12.13
\fieldvalue{\handwritten{2025.12.13}}

\subsubsection*{4.9.2\quad 冻存袋热合操作}

\begin{center}
\small
\begin{tabular}{|p{0.55\linewidth}|p{0.35\linewidth}|}
\hline
1.\ 检查管路无破损、无泄漏，确认标签粘贴良好，冻存袋装量目视无异常；\newline
2.\ 用封管机将冻存袋袋一侧的 EVA 管路热合三次，用剪刀将第二次热合处位置剪开；\newline
3.\ 检查封口处及细胞冻存袋无泄漏，非标准装备备注原因，并按如下规则确认冻存袋标记记录：\newline
\hspace{1em}装量不足的冻存袋标记为B；\newline
\hspace{1em}有破损的细胞冻存袋，标记“$\times$”； &
1.\ 热合开始时间：
% #VALUE_ID: LEX-P0050-V0010
% #FIELD_VALUE: 热合开始时间
% #HANDWRITTEN: 10:45
\fieldvalue{\handwritten{10:45}}\newline
2.\ 第1轮分装冻存袋数量：
% #VALUE_ID: LEX-P0050-V0011
% #FIELD_VALUE: 第1轮分装冻存袋数量
% #HANDWRITTEN: 10
\fieldvalue{\handwritten{10}} 袋；\newline
\hspace{1em}分装合格冻存袋数：
% #VALUE_ID: LEX-P0050-V0012
% #FIELD_VALUE: 第1轮分装合格冻存袋数
% #HANDWRITTEN: 10
\fieldvalue{\handwritten{10}} 袋；\newline
\hspace{1em}标记为B的冻存袋数：
% #VALUE_ID: LEX-P0050-V0013
% #FIELD_VALUE: 第1轮标记为B的冻存袋数
% #HANDWRITTEN: 0
\fieldvalue{\handwritten{0}} 个；\newline
\hspace{1em}标记为“$\times$”的冻存袋数：
% #VALUE_ID: LEX-P0050-V0014
% #FIELD_VALUE: 第1轮标记为“$\times$”的冻存袋数
% #HANDWRITTEN: 0
\fieldvalue{\handwritten{0}} 个。\\ \hline
\multicolumn{1}{|p{0.55\linewidth}|}{}
&
3.\ 第2轮分装冻存袋数量：
% #VALUE_ID: LEX-P0050-V0015
% #FIELD_VALUE: 第2轮分装冻存袋数量
% #HANDWRITTEN: 10
\fieldvalue{\handwritten{10}} 个；
% #VALUE_ID: LEX-P0050-V0016
% #FIELD_VALUE: 第2轮旁注
% #TODO #HANDWRITTEN: 活塞……2025.12.13; handwriting partially illegible
\handwritten{活塞……2025.12.13}\newline
\hspace{1em}分装合格冻存袋数：
% #VALUE_ID: LEX-P0050-V0017
% #FIELD_VALUE: 第2轮分装合格冻存袋数
% #HANDWRITTEN: 10
\fieldvalue{\handwritten{10}} 个；\newline
\hspace{1em}标记为B的冻存袋数：
% #VALUE_ID: LEX-P0050-V0018
% #FIELD_VALUE: 第2轮标记为B的冻存袋数
% #HANDWRITTEN: 0
\fieldvalue{\handwritten{0}} 个；\newline
\hspace{1em}标记为“$\times$”的冻存袋数：
% #VALUE_ID: LEX-P0050-V0019
% #FIELD_VALUE: 第2轮标记为“$\times$”的冻存袋数
% #HANDWRITTEN: 0
\fieldvalue{\handwritten{0}} 个。\newline
4.\ 第3轮分装冻存袋数量：
% #VALUE_ID: LEX-P0050-V0020
% #FIELD_VALUE: 第3轮分装冻存袋数量
% #HANDWRITTEN: 10
\fieldvalue{\handwritten{10}} 个；
% #VALUE_ID: LEX-P0050-V0021
% #FIELD_VALUE: 第3轮旁注
% #TODO #HANDWRITTEN: 活塞……2025.12.13; handwriting partially illegible
\handwritten{活塞……2025.12.13}\newline
\hspace{1em}分装合格冻存袋数：
% #VALUE_ID: LEX-P0050-V0022
% #FIELD_VALUE: 第3轮分装合格冻存袋数
% #HANDWRITTEN: 10
\fieldvalue{\handwritten{10}} 个；\newline
\hspace{1em}标记为B的冻存袋数：
% #VALUE_ID: LEX-P0050-V0023
% #FIELD_VALUE: 第3轮标记为B的冻存袋数
% #HANDWRITTEN: 0
\fieldvalue{\handwritten{0}} 个；\newline
\hspace{1em}标记为“$\times$”的冻存袋数：
% #VALUE_ID: LEX-P0050-V0024
% #FIELD_VALUE: 第3轮标记为“$\times$”的冻存袋数
% #HANDWRITTEN: 0
\fieldvalue{\handwritten{0}} 个。\newline
4.\ 综合结束时间：
% #VALUE_ID: LEX-P0050-V0025
% #FIELD_VALUE: 综合结束时间
% #HANDWRITTEN: 16:30
\fieldvalue{\handwritten{16:30}}
\\ \hline
\end{tabular}
\end{center}

\noindent
备注：
% #VALUE_ID: LEX-P0050-V0026
% #FIELD_VALUE: 备注
% #HANDWRITTEN: /
\fieldvalue{\handwritten{/}}

\noindent
操作人/日期：
% #VALUE_ID: LEX-P0050-V0027
% #FIELD_VALUE: 操作人
% #TODO #HANDWRITTEN: 1?; signature is partially illegible
\fieldvalue{\handwritten{1?}}
\quad
% #VALUE_ID: LEX-P0050-V0028
% #FIELD_VALUE: 操作日期
% #HANDWRITTEN: 2025.12.13
\fieldvalue{\handwritten{2025.12.13}}
\hfill
复核人/日期：
% #VALUE_ID: LEX-P0050-V0029
% #FIELD_VALUE: 复核人
% #TODO #HANDWRITTEN: 党员?; signature is partially illegible
\fieldvalue{\handwritten{党员?}}
\quad
% #VALUE_ID: LEX-P0050-V0030
% #FIELD_VALUE: 复核日期
% #HANDWRITTEN: 2025.12.13
\fieldvalue{\handwritten{2025.12.13}}
% LEXOID_PAGE_COMPLETED: 50/70

\newpage

\begin{center}
\begin{tabular}{|p{0.36\linewidth}|p{0.18\linewidth}|p{0.18\linewidth}|p{0.16\linewidth}|}
\hline
\multicolumn{4}{r|}{记录编号：YZ-VP-PD-APS-012(2502)-07\quad 版本号：1.0}\\
\hline
品名 & 规格 & 生产批号 & 批量\\
\hline
艾尔迈托寡肽注射液(ZZ)无菌工艺模拟试验
&
12.5$\pm$0.5ml
&
% #VALUE_ID: LEX-P0051-V0001
% #FIELD_VALUE: 生产批号
% #TODO #HANDWRITTEN: Asp.../.../...; handwriting is partially illegible
\fieldvalue{\handwritten{Asp.../.../...}}
&
% #VALUE_ID: LEX-P0051-V0002
% #FIELD_VALUE: 批量
% #HANDWRITTEN: 274
\fieldvalue{\handwritten{274}}\ 袋
\\
\hline
\end{tabular}
\end{center}

\section*{4.10.\quad 分装模块9}

\subsection*{4.10.1.\quad 分装操作}

\begin{center}
\begin{tabular}{|p{0.56\linewidth}|p{0.38\linewidth}|}
\hline
\centering 项目 & \centering 操作记录\\
\hline
1.\ 分装过程：整个分装过程均由分装模块自动运行，程序自动摆料、自动剔料和封签整腔；

2.\ 分装模块第一轮分装结束后，断开分装模块电箱终端的排气阀与设备自带的管路连接，旋转金属卡扣至重置位置，然后对应该管路后取样标板；

3.\ 点击主机工艺界面中对应模块的“耗材安装1”，1状态完成复位后，按1.3耗材安装2和耗材安装3依次进行新的分装耗材，打开对应管路的自启自关，点击耗材确认按键继续分装；

每个模块每次分装进行编号，（如：A号模块Gentle P-Pac Pro第1次分装：A-1，A号Gentle P-Pac Pro第2次分装：A-2，B号Gentle P-Pac Pro第3次分装：B-3等）；
&
1.\ 分装开始时间：\underline{\hspace{1.2cm}}时：\underline{\hspace{1.2cm}}分；

2.\ 分装结束时间：\underline{\hspace{1.2cm}}时：\underline{\hspace{1.2cm}}分；

3.\ 分装结束时间：\underline{\hspace{1.2cm}}时：\underline{\hspace{1.2cm}}分\\
\hline
备注： & \\
\hline
操作人/日期：
% #VALUE_ID: LEX-P0051-V0003
% #FIELD_VALUE: 操作人/日期
% #HANDWRITTEN: /
\fieldvalue{\handwritten{/}}
&
复核人/日期：
% #VALUE_ID: LEX-P0051-V0004
% #FIELD_VALUE: 复核人/日期
% #HANDWRITTEN: /
\fieldvalue{\handwritten{/}}
\\
\hline
\end{tabular}
\end{center}

% #VALUE_ID: LEX-P0051-V0005
% #HANDWRITTEN: [蓝色手写批注，字迹不清]
\handwritten{[蓝色手写批注，字迹不清]}

\subsection*{4.10.2.\quad 冻存袋热合操作}

\begin{center}
\begin{tabular}{|p{0.56\linewidth}|p{0.38\linewidth}|}
\hline
1.\ 检查管路无破损、无泄漏，确认标签粘贴完好，冻存袋装置目视无异常；

2.\ 用封管机将冻存袋一侧的EVA管路热合三次，用剪刀将第二次热合处位置剪开；

3.\ 检查封口处及细胞冻存袋无泄漏，非标准袋备注原因，按如下规则确认冻存袋标记记录：

\hspace{1em}装置不足的冻存袋标记为B；

\hspace{1em}有破损的细胞冻存袋，标记“$\times$”；
&
1.\ 热合开始时间：\underline{\hspace{1.2cm}}时：\underline{\hspace{1.2cm}}分；

2.\ 第1轮分装冻存袋数量：\underline{\hspace{0.8cm}}袋；\newline
\hspace{1em}分装合格冻存袋数：\underline{\hspace{0.8cm}}袋；\newline
\hspace{1em}标记为B的冻存袋数：\underline{\hspace{0.8cm}}个；\newline
\hspace{1em}标记为“$\times$”的冻存袋数：\underline{\hspace{0.8cm}}个；

3.\ 第2轮分装冻存袋数量：\underline{\hspace{0.8cm}}袋；\newline
\hspace{1em}分装合格冻存袋数：\underline{\hspace{0.8cm}}袋；\newline
\hspace{1em}标记为B的冻存袋数：\underline{\hspace{0.8cm}}个；\newline
\hspace{1em}标记为“$\times$”的冻存袋数：\underline{\hspace{0.8cm}}个；

4.\ 第3轮分装冻存袋数量：\underline{\hspace{0.8cm}}袋；\newline
\hspace{1em}分装合格冻存袋数：\underline{\hspace{0.8cm}}袋；\newline
\hspace{1em}标记为B的冻存袋数：\underline{\hspace{0.8cm}}个；\newline
\hspace{1em}标记为“$\times$”的冻存袋数：\underline{\hspace{0.8cm}}个；

4.\ 结合结束时间：\underline{\hspace{1.2cm}}时：\underline{\hspace{1.2cm}}分
\\
\hline
备注： & \\
\hline
操作人/日期：
% #VALUE_ID: LEX-P0051-V0006
% #FIELD_VALUE: 操作人/日期
% #HANDWRITTEN: /
\fieldvalue{\handwritten{/}}
&
复核人/日期：
% #VALUE_ID: LEX-P0051-V0007
% #FIELD_VALUE: 复核人/日期
% #HANDWRITTEN: /
\fieldvalue{\handwritten{/}}
\\
\hline
\end{tabular}
\end{center}

% #VALUE_ID: LEX-P0051-V0008
% #HANDWRITTEN: [蓝色手写批注，字迹不清]
\handwritten{[蓝色手写批注，字迹不清]}

\begin{center}
第 17 页 共 23 页
\end{center}
% LEXOID_PAGE_COMPLETED: 51/70

\newpage

\begin{center}
{\small
\begin{tabular}{|p{0.19\linewidth}|p{0.20\linewidth}|p{0.18\linewidth}|p{0.16\linewidth}|}
\hline
\centering 品名 & \centering 规格 & \centering 生产批号 & \centering 批量 \tabularnewline
\hline
艾杰迈托赛针液(Z2)无菌工艺模拟试验
&
\centering 12.5$\pm$0.5\,ml
&
% #VALUE_ID: LEX-P0052-V0001
% #FIELD_VALUE: 生产批号
% #TODO #HANDWRITTEN: 202?01/06?22; handwriting is partially illegible
\fieldvalue{\handwritten{202?01/06?22}}
&
% #VALUE_ID: LEX-P0052-V0002
% #FIELD_VALUE: 批量
% #TODO #HANDWRITTEN: 7?40; handwriting is partially illegible
\fieldvalue{\handwritten{7?40}}袋
\tabularnewline
\hline
\end{tabular}

\vspace{2pt}
记录编号：YZ-VP-PD-APS-012(2502)-07\qquad 版本号：1.0
}
\end{center}

\section*{4.11\quad 分装模块 10}

\subsection*{4.11.1.\quad 分装操作}

\begin{center}
\small
\begin{tabular}{|p{0.55\linewidth}|p{0.37\linewidth}|}
\hline
\centering 项目 & \centering 操作记录 \tabularnewline
\hline
1.分装过程：整个分装过程均由分装模块自动运行，程序自动预排、自动制剂和样性合管路；\par
2.分装模块第一轮分装结束后，断开分装模块耗材管路末端的排气阀与设备自带的管路连接，旋转金属卡扣至垂直位置，盘合对应部位管路后再卡板卡；
&
1.分装开始时间：\underline{\hspace{0.6cm}}时：\underline{\hspace{0.6cm}}分；\par
2.第1轮分装完成\underline{\hspace{0.7cm}}袋；次循环分装；\par
3.分装结束时间：\underline{\hspace{0.6cm}}时：\underline{\hspace{0.6cm}}分
\tabularnewline
\hline
3.点击主机工艺界面中对应模块的“耗材安装1”，将“模块安装1”状态完成复位后，按1.3耗材安装中2）耗材模块安装新的分装耗材，打开对应管路的自愈合泵头；顺时针闭合该泵头连接线缆分装；
&
\tabularnewline
\hline
每个模块每次分装进行编号。（如：A号模块 Gentle P-Pac Pro 第1次分装：A-1，A号模块 Gentle P-Pac Pro 第2次分装：A-2，B号 Gentle P-Pac Pro 第3次分装：B-3等）
&
\tabularnewline
\hline
备注：
&
\tabularnewline
\hline
\end{tabular}
\end{center}

% #VALUE_ID: LEX-P0052-V0003
% #TODO #HANDWRITTEN: 蓝色斜线/批注，无法辨认；跨越分装操作记录表
% #HANDWRITTEN: 蓝色斜线/批注
\handwritten{\text{蓝色斜线/批注}}

操作人/日期：\underline{\hspace{2.2cm}}
\hfill
% #VALUE_ID: LEX-P0052-V0004
% #HANDWRITTEN: /
\handwritten{/}
复核人/日期：\underline{\hspace{2.2cm}}

\subsection*{4.11.2.\quad 冻存袋热合操作}

\begin{center}
\small
\begin{tabular}{|p{0.55\linewidth}|p{0.37\linewidth}|}
\hline
1.检查管路无破损、无泄漏，确认标签粘贴完成，冻存袋装量目视异常；\par
2.用封管机将冻存袋一侧的EVA管路热合三次，用剪刀将第二次热合处位置剪开；\par
3.检查封口处及细胞冻存袋无泄漏，非标准袋各注原因，并按如下规则对冻存袋标记：\par
\hspace{1em}装量不足的冻存袋标记为B；\par
\hspace{1em}有破损的细胞冻存袋，标记“$\times$”；
&
1.热合开始时间：\underline{\hspace{0.6cm}}时：\underline{\hspace{0.6cm}}分；\par
2.第1轮分装冻存袋数量：\underline{\hspace{0.6cm}}袋；\par
\hspace{1em}分装合格冻存袋数量：\underline{\hspace{0.6cm}}袋；\par
\hspace{1em}标记为B的冻存袋数：\underline{\hspace{0.6cm}}个；\par
\hspace{1em}标记为“$\times$”的冻存袋数：\underline{\hspace{0.6cm}}个；\par
3.第2轮分装冻存袋数量：\underline{\hspace{0.6cm}}袋；\par
\hspace{1em}分装合格冻存袋数量：\underline{\hspace{0.6cm}}袋；\par
\hspace{1em}标记为B的冻存袋数：\underline{\hspace{0.6cm}}个；\par
\hspace{1em}标记为“$\times$”的冻存袋数：\underline{\hspace{0.6cm}}个；\par
4.热合结束时间：\underline{\hspace{0.6cm}}时：\underline{\hspace{0.6cm}}分
\tabularnewline
\hline
备注：
&
\tabularnewline
\hline
\end{tabular}
\end{center}

% #VALUE_ID: LEX-P0052-V0005
% #TODO #HANDWRITTEN: 蓝色手写批注，无法辨认；位于冻存袋热合操作表中部
% #HANDWRITTEN: 蓝色手写批注
\handwritten{\text{蓝色手写批注}}

操作人/日期：\underline{\hspace{2.2cm}}
\hfill
% #VALUE_ID: LEX-P0052-V0006
% #HANDWRITTEN: /
\handwritten{/}
复核人/日期：\underline{\hspace{2.2cm}}

\vfill
\begin{center}
\small 第 18 页 共 23 页
\end{center}
% LEXOID_PAGE_COMPLETED: 52/70

\newpage

\noindent
\begin{tabular}{|p{0.16\linewidth}|p{0.43\linewidth}|p{0.18\linewidth}|p{0.15\linewidth}|}
\hline
\multicolumn{4}{r|}{记录编号：YZ-VP-PD-APS-012(2502)-07 \quad 版本号：1.0}\\[-2pt]
\hline
品名 & 规格 & 生产批号 & 批量\\
\hline
艾尔迈托赛注射液(2Z)无菌工艺模拟试验
&
12.5$\pm$0.5\,ml
&
% #VALUE_ID: LEX-P0053-V0001
% #FIELD_VALUE: 生产批号
% #TODO #HANDWRITTEN: Abs?01/01/0223; 手写批号部分字迹不清
\fieldvalue{\handwritten{Abs?01/01/0223}}
&
% #VALUE_ID: LEX-P0053-V0002
% #FIELD_VALUE: 批量
% #HANDWRITTEN: 2,740
\fieldvalue{\handwritten{2,740}} 袋
\\
\hline
\end{tabular}

\vspace{0.15cm}

4.12.\quad 分装模块 11

\vspace{0.1cm}

4.12.1.\quad 分装操作

\begin{tabular}{|p{0.56\linewidth}|p{0.38\linewidth}|}
\hline
\centering 项目 & \centering 操作记录\\
\hline
1.分装过程：整个分装过程均由分装模块自动运行，程序自动预搅拌、自动刮泡和排管路；\newline
2.分装模块第一轮分装结束后，断开分装模块相背管路末端的排气阀与设备自带的管路连接，旋转金属卡扣至垂直位置，然后对对应管路后下柜卡板；\newline
3.点击主机工艺界面中对应模块的“耗材安装1”，分装模块1状态完成复位后，按1.3耗材安装中2）耗材安装的分装线耗材，打开对应管路的白色管夹，然后顺时针按压连接键继续分装；\newline
每个模块每次分装进行编号：（如，A号模块 Gentle P-Pac Pro 第1次分装：A-1，A号 Gentle P-Pac Pro 第2次分装：A-2，B号 Gentle P-Pac Pro 第3次分装：B-3等）。\newline
备注：\vspace{0.1cm}
&
1.分装开始时间：\underline{\hspace{0.35cm}}时：\underline{\hspace{0.35cm}}分；\newline
2.第1轮分装完成——次循环分装；\newline
分装结束时间：\underline{\hspace{0.35cm}}时：\underline{\hspace{0.35cm}}分；\newline
3.第1轮分装结束后：\vspace{0.05cm}\\
\noalign{\vskip -0.15cm}
\multicolumn{2}{c}{
% #VALUE_ID: LEX-P0053-V0003
% #HANDWRITTEN: /
\handwritten{/}}\\[-0.1cm]
\hline
\end{tabular}

\vspace{0.08cm}

% #VALUE_ID: LEX-P0053-V0004
% #HANDWRITTEN: （手写签名，无法辨认）
\noindent
\textbf{操作人/日期：}
\handwritten{（手写签名，无法辨认）}
\hfill
% #VALUE_ID: LEX-P0053-V0005
% #HANDWRITTEN: /
\textbf{复核人/日期：}\handwritten{/}

\vspace{0.15cm}

4.12.2.\quad 冻存袋热合操作

\begin{tabular}{|p{0.56\linewidth}|p{0.38\linewidth}|}
\hline
1.检查管路无破损、无泄漏，确认标签粘贴良好，冻存袋装量目视异常；\newline
2.用封管机将冻存袋一侧的 EVA 管路热合三次，用剪刀将第二次热合处位置剪开；\newline
3.检查封口处及细胞冻存袋无泄漏，非标准装备注射器，井按如下规则确认标签记录：\newline
装量不足的冻存袋标记为 B；\newline
有破损的细胞冻存袋，标记“$\times$”；
&
1.热合开始时间：\underline{\hspace{0.35cm}}时：\underline{\hspace{0.35cm}}分；\newline
2.第1轮分装冻存袋数量：\underline{\hspace{0.3cm}}袋；\newline
\hspace{0.4cm}分装合格冻存袋数：\underline{\hspace{0.3cm}}袋；\newline
\hspace{0.4cm}标记为B的冻存袋数：\underline{\hspace{0.3cm}}个；\newline
\hspace{0.4cm}标记为“$\times$”的冻存袋数：\underline{\hspace{0.3cm}}个；\newline
3.第2轮分装冻存袋数量：\underline{\hspace{0.3cm}}袋；\newline
\hspace{0.4cm}分装合格冻存袋数：\underline{\hspace{0.3cm}}袋；\newline
\hspace{0.4cm}标记为B的冻存袋数：\underline{\hspace{0.3cm}}个；\newline
\hspace{0.4cm}标记为“$\times$”的冻存袋数：\underline{\hspace{0.3cm}}个；\newline
4.第3轮分装冻存袋数量：\underline{\hspace{0.3cm}}袋；\newline
\hspace{0.4cm}分装合格冻存袋数：\underline{\hspace{0.3cm}}袋；\newline
\hspace{0.4cm}标记为B的冻存袋数：\underline{\hspace{0.3cm}}个；\newline
\hspace{0.4cm}标记为“$\times$”的冻存袋数：\underline{\hspace{0.3cm}}个；\newline
4.结合结束时间：\underline{\hspace{0.35cm}}时：\underline{\hspace{0.35cm}}分；
\\
\hline
\end{tabular}

\vspace{0.08cm}

% #VALUE_ID: LEX-P0053-V0006
% #HANDWRITTEN: /
\noindent
\textbf{备注：}\handwritten{/}

\vspace{0.08cm}

% #VALUE_ID: LEX-P0053-V0007
% #HANDWRITTEN: /
\noindent
\textbf{操作人/日期：}\handwritten{/}
\hfill
% #VALUE_ID: LEX-P0053-V0008
% #HANDWRITTEN: /
\textbf{复核人/日期：}\handwritten{/}

\vfill

\begin{center}
第 19 页 共 23 页
\end{center}
% LEXOID_PAGE_COMPLETED: 53/70

\newpage

% TODO: The small company logo at the upper left is visible, but no image filename is available.
\noindent\textbf{记录编号：}YZ-VP-PD-APS-012(2502)-07\qquad
\textbf{版本号：}1.0

\vspace{2mm}

\begin{center}
\begin{tabular}{|p{0.22\linewidth}|p{0.22\linewidth}|p{0.22\linewidth}|p{0.13\linewidth}|}
\hline
\centering 品名 & \centering 规格 & \centering 生产批号 & \centering 批量 \tabularnewline
\hline
% #VALUE_ID: LEX-P0054-V0001
% #FIELD_VALUE: 品名
\centering \fieldvalue{艾米迈托赛针液(Z2)无菌工艺模拟试验} &
% #VALUE_ID: LEX-P0054-V0002
% #FIELD_VALUE: 规格
\centering \fieldvalue{12.5$\pm$0.5ml} &
% #VALUE_ID: LEX-P0054-V0003
% #FIELD_VALUE: 生产批号
% #TODO #HANDWRITTEN: APS?01?c?b? 2740; handwriting unclear
\centering \fieldvalue{\handwritten{APS?01?c?b? 2740}} &
\centering 袋
\tabularnewline
\hline
\end{tabular}
\end{center}

\noindent 4.13．分装模块 12

\noindent 4.13.1．分装操作

\begin{center}
\begin{tabular}{|p{0.57\linewidth}|p{0.35\linewidth}|}
\hline
\centering 项目 & \centering 操作记录 \tabularnewline
\hline
1.分装过程：整个分装过程均由分装模块自动运行，程序自动预热、自动制剂和排空管路。%
\newline
2.分装模块第一轮分装结束后，断开分装模块耗材管路末端的排气阀与设备自带的管路连接，旋转金属卡扣至垂直位置，然后对贮液器后端取样板卡板；
\newline
3.点击主机工艺界面中对应模块的“耗材安装1”，分装模块1状态完成复位后，按1.3耗材安装中2）耗材安装1新的分装耗材，打开对应管路的自备管夹，点击耗材切换键继续分装；
\newline
每个模块每次分装进行编号。（如：A号模块 Gentle P-Pac Pro 第1次分装：A-1，A号 Gentle P-Pac Pro 第2次分装：A-2，B号 Gentle P-Pac Pro 第3次分装：B-3等） &
1.分装开始时间：\underline{\hspace{1.2cm}}时：\underline{\hspace{1.2cm}}分。
\newline
2.第1分装液完成\underline{\hspace{1.0cm}}次循环分装。
\newline
3.分装结束时间：\underline{\hspace{1.2cm}}时：\underline{\hspace{1.2cm}}分。
\tabularnewline
\hline
备注： & \tabularnewline
\hline
操作人/日期：%
\newline
% #VALUE_ID: LEX-P0054-V0004
% #FIELD_VALUE: 操作人/日期
% TODO #HANDWRITTEN: /; handwritten mark is not legible as a name or date
\fieldvalue{\handwritten{/}} &
复核人/日期：%
\newline
% #VALUE_ID: LEX-P0054-V0005
% #FIELD_VALUE: 复核人/日期
% TODO #HANDWRITTEN: /; handwritten mark is not legible as a name or date
\fieldvalue{\handwritten{/}}
\tabularnewline
\hline
\end{tabular}
\end{center}

\noindent 4.13.2．冻存袋验收操作

\begin{center}
\begin{tabular}{|p{0.57\linewidth}|p{0.35\linewidth}|}
\hline
1.检查管路无破损、无泄漏，确认标签粘贴完整，冻存袋装置无明显异常；
\newline
2.用剪切机将冻存袋一侧的EVA管路热合三次，用剪刀将第二次热合处位置剪开；
\newline
3.检查封口处及细胞冻存袋无泄漏，非标准备各庄原因，并按如下规则确认标签记录：
\newline
装置不是的冻存袋标记为B；
\newline
有破损的细胞冻存袋，标记“×”。 &
1.热合开始时间：\underline{\hspace{1.1cm}}时：\underline{\hspace{1.1cm}}分。
\newline
2.第1轮分装冻存袋数量：\underline{\hspace{0.8cm}}袋；
\newline
分装合格冻存袋数量：\underline{\hspace{0.8cm}}袋；
\newline
标记为B的冻存袋数：\underline{\hspace{0.8cm}}个；
\newline
标记为“×”的冻存袋数：\underline{\hspace{0.8cm}}个。
\newline
3.第2轮分装冻存袋数量：\underline{\hspace{0.8cm}}袋；
\newline
分装合格冻存袋数量：\underline{\hspace{0.8cm}}袋；
\newline
标记为B的冻存袋数：\underline{\hspace{0.8cm}}个；
\newline
标记为“×”的冻存袋数：\underline{\hspace{0.8cm}}个。
\newline
4.结合结束时间：\underline{\hspace{1.1cm}}时：\underline{\hspace{1.1cm}}分。
\tabularnewline
\hline
备注： & \tabularnewline
\hline
操作人/日期：%
\newline
% #VALUE_ID: LEX-P0054-V0006
% #FIELD_VALUE: 操作人/日期
% TODO #HANDWRITTEN: /; handwritten mark is not legible as a name or date
\fieldvalue{\handwritten{/}} &
复核人/日期：%
\newline
% #VALUE_ID: LEX-P0054-V0007
% #FIELD_VALUE: 复核人/日期
% TODO #HANDWRITTEN: /; handwritten mark is not legible as a name or date
\fieldvalue{\handwritten{/}}
\tabularnewline
\hline
\end{tabular}
\end{center}

% #VALUE_ID: LEX-P0054-V0008
% #FIELD_VALUE: 备注/验收记录中的手写批注
% #HANDWRITTEN: 方芳敏
\noindent\fieldvalue{\handwritten{方芳敏}}

\noindent\hfill 第 20 页 共 23 页
% LEXOID_PAGE_COMPLETED: 54/70

\newpage

\begin{center}
\begin{tabular}{|p{0.34\linewidth}|p{0.20\linewidth}|p{0.20\linewidth}|p{0.12\linewidth}|}
\hline
品名 & 规格 & 生产批号 & 批量 \\
\hline
艾咪迈托基注射液(Z2)无菌工艺模拟试验
&
12.5$\pm$0.5ml
&
% #VALUE_ID: LEX-P0055-V0001
% #FIELD_VALUE: 生产批号
% #TODO #HANDWRITTEN: 2024/12/?? 2740; handwriting is partially illegible
\fieldvalue{\handwritten{2024/12/?? 2740}}
&
袋
\\
\hline
\end{tabular}
\end{center}

\section{物料平衡统计}

\subsection{冻存袋平衡统计}

\begin{center}
\small
\begin{tabular}{|p{0.13\linewidth}|p{0.13\linewidth}|p{0.13\linewidth}|p{0.13\linewidth}|p{0.13\linewidth}|p{0.13\linewidth}|p{0.15\linewidth}|}
\hline
支领总量 & 使用总量 & 分装合格袋数 & 分装损耗总数 & 焊接损耗总数 & 剩余总量 & 冻存袋使用平衡 \\
\hline
% #VALUE_ID: LEX-P0055-V0002
% #FIELD_VALUE: 支领总量
% #HANDWRITTEN: 230
\fieldvalue{\handwritten{230}}
&
% #VALUE_ID: LEX-P0055-V0003
% #FIELD_VALUE: 使用总量
% #HANDWRITTEN: 243
\fieldvalue{\handwritten{243}}
&
% #VALUE_ID: LEX-P0055-V0004
% #FIELD_VALUE: 分装合格袋数
% #HANDWRITTEN: 240
\fieldvalue{\handwritten{240}}
&
% #VALUE_ID: LEX-P0055-V0005
% #FIELD_VALUE: 分装损耗总数
% #HANDWRITTEN: 0
\fieldvalue{\handwritten{0}}
&
% #VALUE_ID: LEX-P0055-V0006
% #FIELD_VALUE: 焊接损耗总数
% #HANDWRITTEN: 3
\fieldvalue{\handwritten{3}}
&
% #VALUE_ID: LEX-P0055-V0007
% #FIELD_VALUE: 剩余总量
% #HANDWRITTEN: 7
\fieldvalue{\handwritten{7}}
&
% #VALUE_ID: LEX-P0055-V0008
% #FIELD_VALUE: 冻存袋使用平衡
% #TODO #HANDWRITTEN: 合; handwriting is unclear
\fieldvalue{\handwritten{合}}
\\
\hline
\end{tabular}
\end{center}

公式：冻存袋使用总量=分装合格袋数＋分装损耗总数＋焊接损耗总数

冻存袋使用平衡：（冻存袋使用总量-剩余总量）/冻存袋领总量$\times 100\%$\qquad 限度标准：100\%，。

确认人/日期：
% #VALUE_ID: LEX-P0055-V0009
% #FIELD_VALUE: 确认人
% #TODO #HANDWRITTEN: 陈某某; signature is partially illegible
\fieldvalue{\handwritten{陈某某}}
\qquad
% #VALUE_ID: LEX-P0055-V0010
% #FIELD_VALUE: 确认日期
% #HANDWRITTEN: 2024.12.13
\fieldvalue{\handwritten{2024.12.13}}

\medskip

复核人/日期：
% #VALUE_ID: LEX-P0055-V0011
% #FIELD_VALUE: 复核人
% #TODO #HANDWRITTEN: 张某某; signature is partially illegible
\fieldvalue{\handwritten{张某某}}
\qquad
% #VALUE_ID: LEX-P0055-V0012
% #FIELD_VALUE: 复核日期
% #HANDWRITTEN: 2024.12.13
\fieldvalue{\handwritten{2024.12.13}}

\begin{center}
第 21 页 共 23 页
\end{center}
% LEXOID_PAGE_COMPLETED: 55/70

\newpage

\section*{6.\quad 清场}

\begin{tabular}{|p{0.38\linewidth}|p{0.18\linewidth}|p{0.18\linewidth}|p{0.10\linewidth}|}
\hline
\centering 品名 & \centering 规格 & \centering 生产批号 & \centering 批量\\
\hline
% #VALUE_ID: LEX-P0056-V0001
% #FIELD_VALUE: 品名
\fieldvalue{艾米迈托美芬注射液(ZZ)无菌工艺模拟试验}
&
% #VALUE_ID: LEX-P0056-V0002
% #FIELD_VALUE: 规格
\fieldvalue{12.5$\pm$0.5 ml}
&
% #VALUE_ID: LEX-P0056-V0003
% #FIELD_VALUE: 生产批号
% #TODO #HANDWRITTEN: A?2024/12/13?; handwriting unclear
\fieldvalue{\handwritten{A?2024/12/13?}}
&
袋\\
\hline
\end{tabular}

\vspace{0.5em}

操作依据：

\begin{enumerate}
  \item 《生产区清洁管理规程》\hfill YZ-SMP-PD-03-001
  \item 《生产产、检验状态标识管理规程》\hfill YZ-SMP-QA-04-014
\end{enumerate}

生产工序：
% #VALUE_ID: LEX-P0056-V0004
% #FIELD_VALUE: 生产工序
\fieldvalue{细分装}
\hfill
清场开始时间：
% #VALUE_ID: LEX-P0056-V0005
% #FIELD_VALUE: 清场开始时间
% #HANDWRITTEN: 11时32分
\fieldvalue{\handwritten{11时32分}}

\vspace{0.4em}

清场房间编码：
% #VALUE_ID: LEX-P0056-V0006
% #FIELD_VALUE: 清场房间编码
% #TODO #HANDWRITTEN: 232L（?）?; handwriting unclear
\fieldvalue{\handwritten{232L（?）?}}
\hfill
75\%酒精补号：
% #VALUE_ID: LEX-P0056-V0007
% #FIELD_VALUE: 75\%酒精补号
% #TODO #HANDWRITTEN: 202?12?09; handwriting unclear
\fieldvalue{\handwritten{202?12?09}}

\vspace{0.4em}

\begin{tabular}{|p{0.72\linewidth}|p{0.20\linewidth}|}
\hline
\centering 清场项目及标准 & \centering 执行情况\\
\hline
将本班次使用的各类工具配件恢复原位放置；
&
% #VALUE_ID: LEX-P0056-V0008
% #FIELD_VALUE: 将本班次使用的各类工具配件恢复原位放置；--已执行
\fieldvalue{\checkboxfield{checked}} 已执行
\quad
% #VALUE_ID: LEX-P0056-V0009
% #FIELD_VALUE: 将本班次使用的各类工具配件恢复原位放置；--否
\fieldvalue{\checkboxfield{unchecked}} 否
\\
\hline
将本班次生产有关的其它物品或物料清出生产区域；
&
% #VALUE_ID: LEX-P0056-V0010
% #FIELD_VALUE: 将本班次生产有关的其它物品或物料清出生产区域；--已执行
\fieldvalue{\checkboxfield{checked}} 已执行
\quad
% #VALUE_ID: LEX-P0056-V0011
% #FIELD_VALUE: 将本班次生产有关的其它物品或物料清出生产区域；--否
\fieldvalue{\checkboxfield{unchecked}} 否
\\
\hline
使用75\%酒精消毒设备、操作台面、不锈钢凳子、转移车、传递窗、门把手等；
&
% #VALUE_ID: LEX-P0056-V0012
% #FIELD_VALUE: 使用75\%酒精消毒设备、操作台面、不锈钢凳子、转移车、传递窗、门把手等；--已执行
\fieldvalue{\checkboxfield{checked}} 已执行
\quad
% #VALUE_ID: LEX-P0056-V0013
% #FIELD_VALUE: 使用75\%酒精消毒设备、操作台面、不锈钢凳子、转移车、传递窗、门把手等；--否
\fieldvalue{\checkboxfield{unchecked}} 否
\\
\hline
使用75\%酒精喷洒墙面，然后用抹布擦拭墙面；
&
% #VALUE_ID: LEX-P0056-V0014
% #FIELD_VALUE: 使用75\%酒精喷洒墙面，然后用抹布擦拭墙面；--已执行
\fieldvalue{\checkboxfield{checked}} 已执行
\quad
% #VALUE_ID: LEX-P0056-V0015
% #FIELD_VALUE: 使用75\%酒精喷洒墙面，然后用抹布擦拭墙面；--否
\fieldvalue{\checkboxfield{unchecked}} 否
\\
\hline
按照《生产、检验状态标识管理规程》更换房间、设备状态标识；
&
% #VALUE_ID: LEX-P0056-V0016
% #FIELD_VALUE: 按照《生产、检验状态标识管理规程》更换房间、设备状态标识；--已执行
\fieldvalue{\checkboxfield{checked}} 已执行
\quad
% #VALUE_ID: LEX-P0056-V0017
% #FIELD_VALUE: 按照《生产、检验状态标识管理规程》更换房间、设备状态标识；--否
\fieldvalue{\checkboxfield{unchecked}} 否
\\
\hline
将本班次操作产生的废弃物移出洁净生产区；
&
% #VALUE_ID: LEX-P0056-V0018
% #FIELD_VALUE: 将本班次操作产生的废弃物移出洁净生产区；--已执行
\fieldvalue{\checkboxfield{checked}} 已执行
\quad
% #VALUE_ID: LEX-P0056-V0019
% #FIELD_VALUE: 将本班次操作产生的废弃物移出洁净生产区；--否
\fieldvalue{\checkboxfield{unchecked}} 否
\\
\hline
\end{tabular}

\vspace{0.5em}

清场结束时间：
% #VALUE_ID: LEX-P0056-V0020
% #FIELD_VALUE: 清场结束时间
% #HANDWRITTEN: 10时41分
\fieldvalue{\handwritten{10时41分}}

\vspace{0.4em}

清场人/日期：
% #VALUE_ID: LEX-P0056-V0021
% #FIELD_VALUE: 清场人
% #TODO #HANDWRITTEN: 袁玲玲; handwriting unclear
\fieldvalue{\handwritten{袁玲玲}}
\quad
% #VALUE_ID: LEX-P0056-V0022
% #FIELD_VALUE: 清场日期
% #HANDWRITTEN: 2025.12.13
\fieldvalue{\handwritten{2025.12.13}}

\vspace{0.5em}

\begin{tabular}{|p{0.92\linewidth}|}
\hline
清场检查\\
\hline
% #VALUE_ID: LEX-P0056-V0023
% #FIELD_VALUE: 清场检查--合格
\fieldvalue{\checkboxfield{checked}} 合格\\
\hline
% #VALUE_ID: LEX-P0056-V0024
% #FIELD_VALUE: 清场检查--不合格
\fieldvalue{\checkboxfield{unchecked}} 不合格，具体内容为：\\[1.5em]
\hline
整改情况及整改人：\\[1.5em]
\hline
\end{tabular}

\vspace{0.5em}

确认人/日期：
% #VALUE_ID: LEX-P0056-V0025
% #FIELD_VALUE: 确认人
% #TODO #HANDWRITTEN: 孟梅琴; handwriting unclear
\fieldvalue{\handwritten{孟梅琴}}
\quad
% #VALUE_ID: LEX-P0056-V0026
% #FIELD_VALUE: 确认日期
% #HANDWRITTEN: 2024.12.13
\fieldvalue{\handwritten{2024.12.13}}
% LEXOID_PAGE_COMPLETED: 56/70

\newpage

\begin{center}
\begin{minipage}[t]{0.46\linewidth}
\centering
\textbf{工艺运行报告}

\medskip
\begin{tabular}{|p{0.19\linewidth}|p{0.27\linewidth}|p{0.19\linewidth}|p{0.27\linewidth}|}
\hline
\multicolumn{4}{c}{\textbf{工艺信息}}\\
\hline
产品名称 &
% #VALUE_ID: LEX-P0057-V0001
% #FIELD_VALUE: 产品名称
\fieldvalue{Gentle-PePe Pro} &
产品代码 &
% #VALUE_ID: LEX-P0057-V0002
% #FIELD_VALUE: 产品代码
\fieldvalue{204?8070000}\\
\hline
工艺名称 &
% #VALUE_ID: LEX-P0057-V0003
% #FIELD_VALUE: 工艺名称
\fieldvalue{P-C-P} &
原料名称 &
% #VALUE_ID: LEX-P0057-V0004
% #FIELD_VALUE: 原料名称
\fieldvalue{PG?}\\
\hline
开始时间 &
% #VALUE_ID: LEX-P0057-V0005
% #FIELD_VALUE: 开始时间
\fieldvalue{2025-12-13 10:14} &
用户名称 &
% #VALUE_ID: LEX-P0057-V0006
% #FIELD_VALUE: 用户名称
\fieldvalue{L?}\\
\hline
结束时间 &
% #VALUE_ID: LEX-P0057-V0007
% #FIELD_VALUE: 结束时间
\fieldvalue{6:09?} &
原料代码 &
% #VALUE_ID: LEX-P0057-V0008
% #FIELD_VALUE: 原料代码
\fieldvalue{?}\\
\hline
操作人员 &
% #VALUE_ID: LEX-P0057-V0009
% #FIELD_VALUE: 操作人员
\fieldvalue{?} &
自定义 &
% #VALUE_ID: LEX-P0057-V0010
% #FIELD_VALUE: 自定义
\fieldvalue{?}\\
\hline
质量备注 &
% #VALUE_ID: LEX-P0057-V0011
% #FIELD_VALUE: 质量备注
% #TODO #HANDWRITTEN: 转好; 手写字迹较模糊
\fieldvalue{\handwritten{转好}} &
\multicolumn{2}{l|}{}\\
\hline
\end{tabular}

\medskip
\begin{tabular}{|p{0.24\linewidth}|p{0.18\linewidth}|p{0.18\linewidth}|p{0.18\linewidth}|p{0.18\linewidth}|}
\hline
\multicolumn{5}{c}{\textbf{管路耗材}}\\
\hline
产品名称 & 产品编号 & 产品批号 & 生产日期 & 有效日期\\
\hline
% #VALUE_ID: LEX-P0057-V0012
% #FIELD_VALUE: 管路耗材--产品名称
\fieldvalue{/} &
% #VALUE_ID: LEX-P0057-V0013
% #FIELD_VALUE: 管路耗材--产品编号
\fieldvalue{/} &
% #VALUE_ID: LEX-P0057-V0014
% #FIELD_VALUE: 管路耗材--产品批号
\fieldvalue{/} &
% #VALUE_ID: LEX-P0057-V0015
% #FIELD_VALUE: 管路耗材--生产日期
\fieldvalue{/} &
% #VALUE_ID: LEX-P0057-V0016
% #FIELD_VALUE: 管路耗材--有效日期
\fieldvalue{/}\\
\hline
\end{tabular}

\medskip
\begin{tabular}{|p{0.24\linewidth}|p{0.18\linewidth}|p{0.18\linewidth}|p{0.18\linewidth}|p{0.18\linewidth}|}
\hline
\multicolumn{5}{c}{\textbf{样本信息}}\\
\hline
序号 & 样本名称 & 产品批号 & 样品一编号 & 备注\\
\hline
% #VALUE_ID: LEX-P0057-V0017
% #FIELD_VALUE: 样本信息--序号
\fieldvalue{/} &
% #VALUE_ID: LEX-P0057-V0018
% #FIELD_VALUE: 样本信息--样本名称
\fieldvalue{/} &
% #VALUE_ID: LEX-P0057-V0019
% #FIELD_VALUE: 样本信息--产品批号
\fieldvalue{/} &
% #VALUE_ID: LEX-P0057-V0020
% #FIELD_VALUE: 样本信息--样品一编号
\fieldvalue{/} &
% #VALUE_ID: LEX-P0057-V0021
% #FIELD_VALUE: 样本信息--备注
\fieldvalue{/}\\
\hline
\end{tabular}

\medskip
\begin{tabular}{|p{0.24\linewidth}|p{0.18\linewidth}|p{0.18\linewidth}|p{0.18\linewidth}|p{0.18\linewidth}|}
\hline
\multicolumn{5}{c}{\textbf{制剂信息}}\\
\hline
类别 & 产品名称 & 产品批号 & 样品一编号 & 备注\\
\hline
% #VALUE_ID: LEX-P0057-V0022
% #FIELD_VALUE: 制剂信息--类别
\fieldvalue{/} &
% #VALUE_ID: LEX-P0057-V0023
% #FIELD_VALUE: 制剂信息--产品名称
\fieldvalue{/} &
% #VALUE_ID: LEX-P0057-V0024
% #FIELD_VALUE: 制剂信息--产品批号
\fieldvalue{/} &
% #VALUE_ID: LEX-P0057-V0025
% #FIELD_VALUE: 制剂信息--样品一编号
\fieldvalue{/} &
% #VALUE_ID: LEX-P0057-V0026
% #FIELD_VALUE: 制剂信息--备注
\fieldvalue{/}\\
\hline
\end{tabular}

\medskip
\begin{tabular}{|p{0.24\linewidth}|p{0.70\linewidth}|}
\hline
\multicolumn{2}{c}{\textbf{异常信息记录}}\\
\hline
序号 & 异常信息\\
\hline
% #VALUE_ID: LEX-P0057-V0027
% #FIELD_VALUE: 异常信息记录--序号
\fieldvalue{/} &
% #VALUE_ID: LEX-P0057-V0028
% #FIELD_VALUE: 异常信息记录--异常信息
\fieldvalue{/}\\
\hline
\end{tabular}

\medskip
\begin{tabular}{|p{0.24\linewidth}|p{0.70\linewidth}|}
\hline
\multicolumn{2}{c}{\textbf{工艺详细记录}}\\
\hline
时间 & 详细记录\\
\hline
10:10:14 & 操作开始记录1\\
10:10:14 & 操作开始记录\\
\hline
\end{tabular}

\medskip
% #VALUE_ID: LEX-P0057-V0029
% #FIELD_VALUE: 工艺信息旁手写批次记录
% #TODO #HANDWRITTEN: 生产批次：N? / 数量?; 蓝色手写内容难以辨认
\fieldvalue{\handwritten{生产批次：N? / 数量?}}

% #VALUE_ID: LEX-P0057-V0030
% #FIELD_VALUE: 工艺信息旁手写备注
% #TODO #HANDWRITTEN: 编号?; 蓝色手写内容难以辨认
\fieldvalue{\handwritten{编号?}}

% #VALUE_ID: LEX-P0057-V0031
% #FIELD_VALUE: 工艺信息旁手写备注
% #TODO #HANDWRITTEN: 搅拌?; 蓝色手写内容难以辨认
\fieldvalue{\handwritten{搅拌?}}
\end{minipage}
\hfill
\begin{minipage}[t]{0.46\linewidth}
\centering
\textbf{工艺运行报告}

\scriptsize
\begin{tabular}{|p{0.20\linewidth}|p{0.74\linewidth}|}
\hline
时间 & 操作记录\\
\hline
10:28:31 & 扫描进入原料配料区\\
10:28:31 & 扫描人员、设备及物料信息\\
10:28:31 & 扫描进入原料配料区 1\\
10:28:31 & 扫描人员、设备及物料信息\\
10:28:31 & 扫描进入原料配料区\\
10:28:31 & 扫描人员及设备信息\\
10:28:31 & 扫描进入原料配料区\\
10:28:31 & 扫描人员、设备及物料信息\\
10:28:31 & 生产批次信息校验通过\\
10:28:31 & 原料称量开始\\
10:28:31 & 原料称量完成\\
10:28:31 & 加入原料：500.0 g\\
10:28:31 & 加入原料：200.0 g\\
10:28:31 & 加入原料：50.0 g\\
10:28:31 & 扫描人员及设备信息\\
10:28:31 & 分装设备准备完成\\
10:28:31 & 分装数量：12, 13\\
10:32:32 & 分装设备开始运行\\
10:32:32 & 分装完成：300 min\\
10:32:32 & 分装数量：1, 2, 3, 4, 5, 6, 7, 8, 9, 10\\
10:32:32 & 分装数量：11, 12, 13\\
10:32:32 & 分装完成\\
10:32:32 & 检测设备开始运行\\
10:32:32 & 检测记录生成\\
10:32:32 & 分装数量：12, 13\\
\hline
\end{tabular}

\vfill
\textit{第58页}
\end{minipage}
\end{center}

\medskip

\begin{center}
\begin{minipage}[t]{0.46\linewidth}
\centering
\textbf{工艺运行报告}

\scriptsize
\begin{tabular}{|p{0.20\linewidth}|p{0.74\linewidth}|}
\hline
时间 & 操作记录\\
\hline
10:32:32 & 分装数量：300 min\\
10:32:32 & 分装数量：1, 2, 3, 4, 5, 6, 7, 8, 9, 10\\
10:32:32 & 分装数量：11, 12, 13\\
10:32:32 & 分装完成\\
10:32:33 & 分装设备开始运行\\
10:32:33 & 分装数量：300 min\\
10:32:33 & 分装数量：1, 2, 3, 4, 5, 6, 7, 8, 9, 10\\
10:32:33 & 分装数量：11, 12, 13\\
10:32:33 & 分装完成\\
10:34:37 & 分装数量：300 min\\
10:34:37 & 分装数量：1, 2, 3, 4, 5, 6, 7, 8, 9, 10\\
10:34:37 & 分装数量：11, 12, 13\\
10:34:38 & 分装完成\\
10:34:38 & 分装设备开始运行\\
10:34:38 & 分装数量：300 min\\
10:34:38 & 分装数量：1, 2, 3, 4, 5, 6, 7, 8, 9, 10\\
10:34:38 & 分装数量：11, 12, 13\\
10:34:39 & 分装完成\\
10:34:39 & 分装设备开始运行\\
10:34:39 & 分装数量：300 min\\
10:34:39 & 分装数量：1, 2, 3, 4, 5, 6, 7, 8, 9, 10\\
10:34:39 & 分装数量：11, 12, 13\\
\hline
\end{tabular}

\vfill
\textit{第59页}
\end{minipage}
\hfill
\begin{minipage}[t]{0.46\linewidth}
\centering
\textbf{工艺运行报告}

\scriptsize
\begin{tabular}{|p{0.20\linewidth}|p{0.74\linewidth}|}
\hline
时间 & 操作记录\\
\hline
10:34:39 & 分装数量：300 min\\
10:34:39 & 分装数量：1, 2, 3, 4, 5, 6, 7, 8, 9, 10\\
10:34:39 & 分装数量：11, 12, 13\\
10:34:39 & 分装完成\\
10:34:49 & 分装设备开始运行\\
10:34:49 & 分装数量：200 min\\
10:34:49 & 分装数量：1, 2, 3, 4, 5, 6, 7, 8, 9, 10\\
10:34:49 & 分装数量：11, 12, 13\\
10:34:49 & 分装完成\\
10:34:49 & 分装设备开始运行\\
10:34:49 & 分装数量：200 min\\
10:34:49 & 分装数量：1, 2, 3, 4, 5, 6, 7, 8, 9, 10\\
10:34:49 & 分装数量：11, 12, 13\\
10:35:52 & 分装设备开始运行\\
10:35:52 & 分装数量：200 min\\
10:35:52 & 分装数量：1, 2, 3, 4, 5, 6, 7, 8, 9, 10\\
10:35:52 & 分装数量：11, 12, 13\\
10:35:52 & 分装完成\\
10:35:52 & 分装设备开始运行\\
10:35:52 & 分装数量：200 min\\
10:35:52 & 分装数量：1, 2, 3, 4, 5, 6, 7, 8, 9, 10\\
10:35:52 & 分装数量：11, 12, 13\\
\hline
\end{tabular}

\vfill
% #VALUE_ID: LEX-P0057-V0032
% #HANDWRITTEN: 1/4
\handwritten{1/4}
\end{minipage}
\end{center}
% LEXOID_PAGE_COMPLETED: 57/70

\newpage

\noindent
\begin{minipage}[t]{0.47\linewidth}
\centering
\textbf{翰轩生物}\\[-1mm]
工艺运行报告

\vspace{2mm}
\scriptsize
\begin{tabular}{|p{0.24\linewidth}|p{0.68\linewidth}|}
\hline
时间 & 工艺运行记录 \\ \hline
10:35:55 & 分拣罐3相关运行记录 \\ \hline
10:35:55 & 分拣罐3放料，约2000\,mL/min \\ \hline
10:35:55 & 分拣罐3运行记录 \\ \hline
10:35:57 & 分拣罐3相关运行记录 \\ \hline
10:35:57 & 分拣罐3相关运行记录 \\ \hline
10:35:57 & 分拣罐3运行记录 \\ \hline
10:35:57 & 分拣罐3运行记录 \\ \hline
10:35:57 & 分拣罐3相关运行记录 \\ \hline
10:35:57 & 分拣罐3投料，约200\,mL/min \\ \hline
10:35:57 & 分拣罐3运行记录 \\ \hline
10:35:57 & 分拣罐3运行记录 \\ \hline
10:36:04 & 分拣罐3相关运行记录 \\ \hline
10:36:04 & 分拣罐3相关运行记录 \\ \hline
10:36:07 & 分拣罐3运行记录 \\ \hline
10:36:07 & 分拣罐3运行记录 \\ \hline
10:36:07 & 分拣罐3运行记录 \\ \hline
10:36:07 & 分拣罐3相关运行记录 \\ \hline
10:36:07 & 分拣罐3运行记录 \\ \hline
10:36:07 & 分拣罐3相关运行记录 \\ \hline
10:36:10 & 分拣罐3运行记录 \\ \hline
10:36:10 & 分拣罐3相关运行记录 \\ \hline
10:36:20 & 分拣罐3运行记录 \\ \hline
10:36:20 & 分拣罐3相关运行记录 \\ \hline
10:36:25 & 分拣罐3相关运行记录 \\ \hline
\end{tabular}

\vspace{1mm}

第五页
\end{minipage}
\hfill
\begin{minipage}[t]{0.47\linewidth}
\centering
\textbf{翰轩生物}\\[-1mm]
工艺运行报告

\vspace{2mm}
\scriptsize
\begin{tabular}{|p{0.24\linewidth}|p{0.68\linewidth}|}
\hline
时间 & 工艺运行记录 \\ \hline
10:36:30 & 分拣罐2相关运行记录，约2500\,mL \\ \hline
10:36:30 & 分拣罐2相关运行记录，约2500\,mL \\ \hline
10:36:35 & 分拣罐2相关运行记录，约2500\,mL \\ \hline
10:36:35 & 分拣罐2相关运行记录，约2500\,mL \\ \hline
10:36:41 & 分拣罐2相关运行记录，约2500\,mL \\ \hline
10:36:50 & 分拣罐2相关运行记录，约2500\,mL \\ \hline
10:37:09 & 分拣罐2相关运行记录，约2500\,mL \\ \hline
10:37:17 & 分拣罐2相关运行记录，约2500\,mL \\ \hline
10:37:20 & 分拣罐2相关运行记录，约2500\,mL \\ \hline
10:37:27 & 分拣罐2相关运行记录，约2500\,mL \\ \hline
10:37:27 & 分拣罐2相关运行记录，约2500\,mL \\ \hline
10:37:30 & 分拣罐2相关运行记录，约2500\,mL \\ \hline
10:37:34 & 分拣罐2相关运行记录，约2500\,mL \\ \hline
10:37:40 & 分拣罐2相关运行记录，约2500\,mL \\ \hline
10:37:41 & 分拣罐2相关运行记录，约2500\,mL \\ \hline
10:37:47 & 分拣罐2相关运行记录，约2500\,mL \\ \hline
10:37:48 & 分拣罐2相关运行记录，约2500\,mL \\ \hline
10:38:08 & 分拣罐2相关运行记录，约2500\,mL \\ \hline
10:38:16 & 分拣罐2相关运行记录，约2500\,mL \\ \hline
10:38:19 & 分拣罐2相关运行记录 \\ \hline
\end{tabular}

\vspace{1mm}

第六页
\end{minipage}

\vspace{8mm}

\noindent
\begin{minipage}[t]{0.47\linewidth}
\centering
\textbf{翰轩生物}\\[-1mm]
工艺运行报告

\vspace{2mm}
\scriptsize
\begin{tabular}{|p{0.24\linewidth}|p{0.68\linewidth}|}
\hline
时间 & 工艺运行记录 \\ \hline
10:38:24 & 分拣罐1运行记录 \\ \hline
10:38:24 & 分拣罐1运行记录 \\ \hline
10:38:24 & 分拣罐1运行记录 \\ \hline
10:38:24 & 分拣罐1相关运行记录 \\ \hline
10:38:25 & 分拣罐1运行记录，约500\,mL/min \\ \hline
10:38:25 & 分拣罐1运行记录 \\ \hline
10:38:25 & 分拣罐1相关运行记录 \\ \hline
10:38:25 & 分拣罐1运行记录 \\ \hline
10:38:25 & 分拣罐1运行记录 \\ \hline
10:38:25 & 分拣罐1相关运行记录 \\ \hline
10:38:30 & 分拣罐1运行记录 \\ \hline
10:38:32 & 分拣罐1运行记录 \\ \hline
10:38:32 & 分拣罐1相关运行记录 \\ \hline
10:38:32 & 分拣罐1运行记录 \\ \hline
10:38:32 & 分拣罐1相关运行记录 \\ \hline
10:38:32 & 分拣罐1运行记录 \\ \hline
10:38:32 & 分拣罐1运行记录 \\ \hline
10:38:41 & 分拣罐1运行记录 \\ \hline
10:38:47 & 分拣罐1运行记录 \\ \hline
10:38:47 & 分拣罐1相关运行记录 \\ \hline
\end{tabular}

\vspace{1mm}

第七页
\end{minipage}
\hfill
\begin{minipage}[t]{0.47\linewidth}
\centering
\textbf{翰轩生物}\\[-1mm]
工艺运行报告

\vspace{2mm}
\scriptsize
\begin{tabular}{|p{0.24\linewidth}|p{0.68\linewidth}|}
\hline
时间 & 工艺运行记录 \\ \hline
10:38:47 & 分拣罐2相关运行记录 \\ \hline
10:38:47 & 分拣罐2运行记录，约500\,mL/min \\ \hline
10:38:47 & 分拣罐2运行记录 \\ \hline
10:38:47 & 分拣罐2相关运行记录 \\ \hline
10:38:56 & 分拣罐2运行记录 \\ \hline
10:38:56 & 分拣罐2运行记录 \\ \hline
10:38:56 & 分拣罐2相关运行记录 \\ \hline
10:38:56 & 分拣罐2运行记录 \\ \hline
10:39:40 & 分拣罐2运行记录 \\ \hline
10:39:41 & 分拣罐2相关运行记录 \\ \hline
10:39:41 & 分拣罐2运行记录 \\ \hline
10:39:43 & 分拣罐2运行记录 \\ \hline
10:39:43 & 分拣罐2相关运行记录 \\ \hline
10:39:43 & 分拣罐2运行记录 \\ \hline
10:39:49 & 分拣罐2运行记录 \\ \hline
10:39:51 & 分拣罐2相关运行记录 \\ \hline
10:39:51 & 分拣罐2运行记录 \\ \hline
10:39:55 & 分拣罐2运行记录 \\ \hline
\end{tabular}

\vspace{1mm}

第八页

\vspace{4mm}
% #VALUE_ID: LEX-P0058-V0001
% #HANDWRITTEN: 1/4
\handwritten{1/4}
\end{minipage}

% TODO: The source image contains very small Chinese text in the table descriptions; entries above preserve the visible table structure and readable timestamps, but several detailed descriptions require higher-resolution source text for reliable transcription.
% LEXOID_PAGE_COMPLETED: 58/70

\newpage

\noindent
\begin{minipage}[t]{0.47\linewidth}
\centering
\textbf{工艺运行报告}

\vspace{2mm}
\scriptsize
\begin{tabular}{|p{0.20\linewidth}|p{0.70\linewidth}|}
\hline
时间 & 运行记录 \\ \hline
10:39:58 & 分馏塔4釜再沸器启动 \\ \hline
10:39:58 & 分馏塔4釜温度升至约12℃ \\ \hline
10:39:58 & 分馏塔4釜温度升至约12℃ \\ \hline
10:39:58 & 分馏塔4釜温度约6000mL/min \\ \hline
10:39:58 & 分馏塔4釜温度约6000mL/min \\ \hline
10:39:58 & 分馏塔4釜温度约6000mL/min \\ \hline
10:40:01 & 分馏塔4釜温度约6000mL/min \\ \hline
10:40:11 & 分馏塔4釜温度约6000mL/min \\ \hline
10:40:18 & 分馏塔4釜温度约6000mL/min \\ \hline
10:40:18 & 分馏塔4釜温度约6000mL/min \\ \hline
10:40:22 & 分馏塔4釜温度约6000mL/min \\ \hline
10:40:25 & 工艺参数调整 \\ \hline
10:44:47 & 分馏塔1釜液位调节 \\ \hline
10:44:47 & 分馏塔1釜温度约12℃ \\ \hline
10:44:47 & 分馏塔1釜温度约300mL/min \\ \hline
10:44:47 & 分馏塔1釜温度约300mL/min \\ \hline
10:44:47 & 分馏塔1釜温度约300mL/min \\ \hline
10:44:47 & 分馏塔1釜温度约300mL/min \\ \hline
10:44:47 & 分馏塔1釜温度约300mL/min \\ \hline
10:45:23 & 分馏塔1釜温度约300mL/min \\ \hline
10:45:23 & 分馏塔1釜温度约300mL/min \\ \hline
10:45:23 & 分馏塔1釜温度约300mL/min \\ \hline
10:45:23 & 分馏塔1釜温度约300mL/min \\ \hline
10:45:23 & 分馏塔1釜温度约300mL/min \\ \hline
\end{tabular}

\vspace{1mm}

第9页
\end{minipage}
\hfill
\begin{minipage}[t]{0.47\linewidth}
\centering
\textbf{工艺运行报告}

\vspace{2mm}
\scriptsize
\begin{tabular}{|p{0.20\linewidth}|p{0.70\linewidth}|}
\hline
时间 & 运行记录 \\ \hline
10:46:23 & 分馏塔1釜再沸器约20000Pa \\ \hline
10:46:31 & 分馏塔1釜温度约12℃ \\ \hline
10:46:31 & 分馏塔1釜温度约12℃ \\ \hline
10:46:31 & 分馏塔1釜温度约12℃ \\ \hline
10:46:31 & 分馏塔1釜温度约12℃ \\ \hline
10:46:31 & 分馏塔1釜温度约300mL/min \\ \hline
10:46:31 & 分馏塔1釜温度约300mL/min \\ \hline
10:46:31 & 分馏塔1釜温度约300mL/min \\ \hline
10:46:31 & 分馏塔1釜温度约300mL/min \\ \hline
10:46:31 & 分馏塔1釜温度约300mL/min \\ \hline
10:46:06 & 分馏塔1釜温度约300mL/min \\ \hline
10:46:06 & 分馏塔1釜温度约300mL/min \\ \hline
10:46:06 & 分馏塔1釜温度约300mL/min \\ \hline
10:46:06 & 分馏塔1釜温度约300mL/min \\ \hline
10:46:06 & 分馏塔1釜温度约300mL/min \\ \hline
10:46:36 & 分馏塔1釜温度约20000Pa \\ \hline
10:46:36 & 分馏塔1釜温度约72.5℃ \\ \hline
10:46:36 & 分馏塔1釜温度约300mL/min \\ \hline
10:46:36 & 分馏塔1釜温度约300mL/min \\ \hline
10:46:36 & 分馏塔1釜温度约300mL/min \\ \hline
10:46:41 & 分馏塔1釜温度约300mL/min \\ \hline
10:46:41 & 分馏塔1釜温度约300mL/min \\ \hline
10:46:41 & 分馏塔1釜温度约300mL/min \\ \hline
10:46:41 & 分馏塔1釜温度约300mL/min \\ \hline
\end{tabular}

\vspace{1mm}

第10页
\end{minipage}

\vspace{8mm}

\noindent
\begin{minipage}[t]{0.47\linewidth}
\centering
\textbf{工艺运行报告}

\vspace{2mm}
\scriptsize
\begin{tabular}{|p{0.20\linewidth}|p{0.70\linewidth}|}
\hline
时间 & 运行记录 \\ \hline
10:47:12 & 分馏塔2釜液位调节 \\ \hline
10:47:12 & 分馏塔2釜温度约12℃ \\ \hline
10:47:12 & 分馏塔2釜温度约12℃ \\ \hline
10:47:12 & 分馏塔2釜温度约12℃ \\ \hline
10:47:12 & 分馏塔2釜温度约300mL/min \\ \hline
10:47:12 & 分馏塔2釜温度约300mL/min \\ \hline
10:47:19 & 分馏塔3启动 \\ \hline
10:47:19 & 分馏塔3釜温度约73.5℃ \\ \hline
10:47:19 & 分馏塔3釜温度约300mL/min \\ \hline
10:47:19 & 分馏塔3釜温度约300mL/min \\ \hline
10:47:19 & 分馏塔3釜温度约300mL/min \\ \hline
10:47:19 & 分馏塔3釜温度约300mL/min \\ \hline
10:47:19 & 分馏塔3釜温度约300mL/min \\ \hline
10:47:41 & 分馏塔3釜温度约300mL/min \\ \hline
10:47:41 & 分馏塔3釜温度约300mL/min \\ \hline
10:47:59 & 分馏塔4启动 \\ \hline
10:47:59 & 分馏塔4釜温度约73.5℃ \\ \hline
10:47:59 & 分馏塔4釜温度约300mL/min \\ \hline
10:47:59 & 分馏塔4釜温度约300mL/min \\ \hline
10:47:59 & 分馏塔4釜温度约300mL/min \\ \hline
10:47:59 & 分馏塔4釜温度约300mL/min \\ \hline
10:48:17 & 分馏塔4釜温度约20000Pa \\ \hline
\end{tabular}

\vspace{1mm}

第11页
\end{minipage}
\hfill
\begin{minipage}[t]{0.47\linewidth}
\centering
\textbf{工艺运行报告}

\vspace{2mm}
\scriptsize
\begin{tabular}{|p{0.20\linewidth}|p{0.70\linewidth}|}
\hline
时间 & 运行记录 \\ \hline
10:48:17 & 分馏塔2釜温度约35.3℃ \\ \hline
10:48:17 & 分馏塔2釜温度约300mL/min \\ \hline
10:48:17 & 分馏塔2釜温度约12℃ \\ \hline
10:48:17 & 分馏塔2釜温度约12℃ \\ \hline
10:48:17 & 分馏塔2釜温度约300mL/min \\ \hline
10:48:22 & 分馏塔2釜温度约12℃ \\ \hline
10:48:37 & 分馏塔2釜温度约12℃ \\ \hline
10:48:42 & 分馏塔2釜温度约12℃ \\ \hline
10:48:46 & 分馏塔2釜温度约300mL/min \\ \hline
10:48:56 & 分馏塔2釜温度约300mL/min \\ \hline
10:49:08 & 分馏塔2釜温度约300mL/min \\ \hline
10:49:08 & 分馏塔2釜温度约300mL/min \\ \hline
10:49:33 & 分馏塔2釜温度约300mL/min \\ \hline
10:49:33 & 分馏塔2釜温度约74.5℃ \\ \hline
10:49:33 & 分馏塔2釜温度约300mL/min \\ \hline
10:49:33 & 分馏塔2釜温度约300mL/min \\ \hline
10:49:33 & 分馏塔2釜温度约300mL/min \\ \hline
10:49:33 & 分馏塔2釜温度约300mL/min \\ \hline
10:49:33 & 分馏塔2釜温度约300mL/min \\ \hline
10:49:39 & 分馏塔2釜温度约300mL/min \\ \hline
10:49:39 & 分馏塔2釜温度约300mL/min \\ \hline
10:49:39 & 分馏塔2釜温度约300mL/min \\ \hline
10:49:40 & 分馏塔2釜温度约300mL/min \\ \hline
\end{tabular}

\vspace{1mm}

第12页

% #VALUE_ID: LEX-P0059-V0001
% #HANDWRITTEN: 3/4
\handwritten{3/4}
\end{minipage}

% TODO: The source page contains four densely printed process-log tables. Several individual Chinese log entries are too small to read reliably at the supplied resolution; the visible table structure and legible timestamps have been retained.
% LEXOID_PAGE_COMPLETED: 59/70

\newpage

% TODO: The page contains four reduced-size report pages; some Chinese table text is low-resolution.
\begin{center}
\begin{minipage}[t]{0.47\linewidth}
\centering
\textbf{工艺运行报告}

\scriptsize
\begin{tabular}{|p{0.18\linewidth}|p{0.72\linewidth}|}
\hline
时间 & 运行记录\\ \hline
\multicolumn{2}{|l|}{
% #VALUE_ID: LEX-P0060-V0001
% #FIELD_VALUE: 工艺运行记录
\fieldvalue{10:49:40；分离器1液位调整，11.2}} \\ \hline
\multicolumn{2}{|l|}{
% #VALUE_ID: LEX-P0060-V0002
% #FIELD_VALUE: 工艺运行记录
\fieldvalue{10:49:40；分离器1液位达到200m/min}} \\ \hline
\multicolumn{2}{|l|}{
% #VALUE_ID: LEX-P0060-V0003
% #FIELD_VALUE: 工艺运行记录
\fieldvalue{10:49:40；分离器1液位达到210m/min}} \\ \hline
\multicolumn{2}{|l|}{
% #VALUE_ID: LEX-P0060-V0004
% #FIELD_VALUE: 工艺运行记录
\fieldvalue{10:49:45；分离器1液位达到220m/min}} \\ \hline
\multicolumn{2}{|l|}{
% #VALUE_ID: LEX-P0060-V0005
% #FIELD_VALUE: 工艺运行记录
\fieldvalue{10:49:48；分离器1液位达到230m/min}} \\ \hline
\multicolumn{2}{|l|}{
% #VALUE_ID: LEX-P0060-V0006
% #FIELD_VALUE: 工艺运行记录
\fieldvalue{10:49:52；分离器1液位达到240m/min}} \\ \hline
\multicolumn{2}{|l|}{
% #VALUE_ID: LEX-P0060-V0007
% #FIELD_VALUE: 工艺运行记录
\fieldvalue{10:49:54；分离器1液位达到250m/min}} \\ \hline
\multicolumn{2}{|l|}{
% #VALUE_ID: LEX-P0060-V0008
% #FIELD_VALUE: 工艺运行记录
\fieldvalue{10:49:54；分离器1液位达到260m/min}} \\ \hline
\multicolumn{2}{|l|}{
% #VALUE_ID: LEX-P0060-V0009
% #FIELD_VALUE: 工艺运行记录
\fieldvalue{10:49:54；分离器1液位达到270m/min}} \\ \hline
\multicolumn{2}{|l|}{
% #VALUE_ID: LEX-P0060-V0010
% #FIELD_VALUE: 工艺运行记录
\fieldvalue{10:49:54；分离器1液位达到280m/min}} \\ \hline
\multicolumn{2}{|l|}{
% #VALUE_ID: LEX-P0060-V0011
% #FIELD_VALUE: 工艺运行记录
\fieldvalue{10:49:54；分离器1液位达到290m/min}} \\ \hline
\multicolumn{2}{|l|}{
% #VALUE_ID: LEX-P0060-V0012
% #FIELD_VALUE: 工艺运行记录
\fieldvalue{10:49:54；分离器1液位达到300m/min}} \\ \hline
\multicolumn{2}{|l|}{
% #VALUE_ID: LEX-P0060-V0013
% #FIELD_VALUE: 工艺运行记录
\fieldvalue{10:49:54；分离器1液位达到310m/min}} \\ \hline
\multicolumn{2}{|l|}{
% #VALUE_ID: LEX-P0060-V0014
% #FIELD_VALUE: 工艺运行记录
\fieldvalue{10:49:54；分离器1液位达到320m/min}} \\ \hline
\multicolumn{2}{|l|}{
% #VALUE_ID: LEX-P0060-V0015
% #FIELD_VALUE: 工艺运行记录
\fieldvalue{10:49:54；分离器1液位达到330m/min}} \\ \hline
\multicolumn{2}{|l|}{
% #VALUE_ID: LEX-P0060-V0016
% #FIELD_VALUE: 工艺运行记录
\fieldvalue{10:50:06；分离器1液位达到340m/min}} \\ \hline
\multicolumn{2}{|l|}{
% #VALUE_ID: LEX-P0060-V0017
% #FIELD_VALUE: 工艺运行记录
\fieldvalue{10:50:12；分离器1液位达到350m/min}} \\ \hline
\multicolumn{2}{|l|}{
% #VALUE_ID: LEX-P0060-V0018
% #FIELD_VALUE: 工艺运行记录
\fieldvalue{10:50:15；分离器1液位达到360m/min}} \\ \hline
\multicolumn{2}{|l|}{
% #VALUE_ID: LEX-P0060-V0019
% #FIELD_VALUE: 工艺运行记录
\fieldvalue{10:50:20；分离器1液位达到370m/min}} \\ \hline
\multicolumn{2}{|l|}{
% #VALUE_ID: LEX-P0060-V0020
% #FIELD_VALUE: 工艺运行记录
\fieldvalue{10:50:22；分离器1液位达到380m/min}} \\ \hline
\multicolumn{2}{|l|}{
% #VALUE_ID: LEX-P0060-V0021
% #FIELD_VALUE: 工艺运行记录
\fieldvalue{10:50:31；分离器1液位达到390m/min}} \\ \hline
\multicolumn{2}{|l|}{
% #VALUE_ID: LEX-P0060-V0022
% #FIELD_VALUE: 工艺运行记录
\fieldvalue{10:50:36；分离器1液位达到400m/min}} \\ \hline
\multicolumn{2}{|l|}{
% #VALUE_ID: LEX-P0060-V0023
% #FIELD_VALUE: 工艺运行记录
\fieldvalue{10:50:37；分离器1液位达到410m/min}} \\ \hline
\end{tabular}

\vspace{1mm}
第13页
\end{minipage}
\hfill
\begin{minipage}[t]{0.47\linewidth}
\centering
\textbf{工艺运行报告}

\scriptsize
\begin{tabular}{|p{0.18\linewidth}|p{0.72\linewidth}|}
\hline
时间 & 运行记录\\ \hline
\multicolumn{2}{|l|}{
% #VALUE_ID: LEX-P0060-V0024
% #FIELD_VALUE: 工艺运行记录
\fieldvalue{10:50:37；分离器4达到运行压力}} \\ \hline
\multicolumn{2}{|l|}{
% #VALUE_ID: LEX-P0060-V0025
% #FIELD_VALUE: 工艺运行记录
\fieldvalue{10:50:37；分离器4流量达到25.0 L/min}} \\ \hline
\multicolumn{2}{|l|}{
% #VALUE_ID: LEX-P0060-V0026
% #FIELD_VALUE: 工艺运行记录
\fieldvalue{10:50:52；分离器4液位达到1、2、3、4、5、6、7、8、9、10}} \\ \hline
\multicolumn{2}{|l|}{
% #VALUE_ID: LEX-P0060-V0027
% #FIELD_VALUE: 工艺运行记录
\fieldvalue{10:50:52；分离器4温度达到55.0m/min}} \\ \hline
\multicolumn{2}{|l|}{
% #VALUE_ID: LEX-P0060-V0028
% #FIELD_VALUE: 工艺运行记录
\fieldvalue{10:50:52；分离器4液位达到压力2500.0 mL}} \\ \hline
\multicolumn{2}{|l|}{
% #VALUE_ID: LEX-P0060-V0029
% #FIELD_VALUE: 工艺运行记录
\fieldvalue{10:50:58；分离器4温度达到50.0m/min}} \\ \hline
\multicolumn{2}{|l|}{
% #VALUE_ID: LEX-P0060-V0030
% #FIELD_VALUE: 工艺运行记录
\fieldvalue{10:51:06；分离器4温度达到压力2500.0 mL}} \\ \hline
\multicolumn{2}{|l|}{
% #VALUE_ID: LEX-P0060-V0031
% #FIELD_VALUE: 工艺运行记录
\fieldvalue{10:51:06；分离器4达到运行压力}} \\ \hline
\multicolumn{2}{|l|}{
% #VALUE_ID: LEX-P0060-V0032
% #FIELD_VALUE: 工艺运行记录
\fieldvalue{10:51:06；分离器4流量达到25.0 L/min}} \\ \hline
\multicolumn{2}{|l|}{
% #VALUE_ID: LEX-P0060-V0033
% #FIELD_VALUE: 工艺运行记录
\fieldvalue{10:51:06；分离器4温度达到50.0m/min}} \\ \hline
\multicolumn{2}{|l|}{
% #VALUE_ID: LEX-P0060-V0034
% #FIELD_VALUE: 工艺运行记录
\fieldvalue{10:51:08；分离器4温度达到压力2500.0 mL}} \\ \hline
\multicolumn{2}{|l|}{
% #VALUE_ID: LEX-P0060-V0035
% #FIELD_VALUE: 工艺运行记录
\fieldvalue{10:51:16；分离器4达到运行压力}} \\ \hline
\multicolumn{2}{|l|}{
% #VALUE_ID: LEX-P0060-V0036
% #FIELD_VALUE: 工艺运行记录
\fieldvalue{10:51:16；分离器4流量达到25.0 L/min}} \\ \hline
\multicolumn{2}{|l|}{
% #VALUE_ID: LEX-P0060-V0037
% #FIELD_VALUE: 工艺运行记录
\fieldvalue{10:51:18；分离器4温度达到50.0m/min}} \\ \hline
\multicolumn{2}{|l|}{
% #VALUE_ID: LEX-P0060-V0038
% #FIELD_VALUE: 工艺运行记录
\fieldvalue{10:51:18；分离器4温度达到压力2500.0 mL}} \\ \hline
\multicolumn{2}{|l|}{
% #VALUE_ID: LEX-P0060-V0039
% #FIELD_VALUE: 工艺运行记录
\fieldvalue{10:51:26；分离器4达到运行压力}} \\ \hline
\multicolumn{2}{|l|}{
% #VALUE_ID: LEX-P0060-V0040
% #FIELD_VALUE: 工艺运行记录
\fieldvalue{10:51:26；分离器4流量达到25.0 L/min}} \\ \hline
\multicolumn{2}{|l|}{
% #VALUE_ID: LEX-P0060-V0041
% #FIELD_VALUE: 工艺运行记录
\fieldvalue{10:51:29；分离器4温度达到50.0m/min}} \\ \hline
\multicolumn{2}{|l|}{
% #VALUE_ID: LEX-P0060-V0042
% #FIELD_VALUE: 工艺运行记录
\fieldvalue{10:51:29；分离器4温度达到压力2500.0 mL}} \\ \hline
\multicolumn{2}{|l|}{
% #VALUE_ID: LEX-P0060-V0043
% #FIELD_VALUE: 工艺运行记录
\fieldvalue{10:51:48；分离器4达到运行压力}} \\ \hline
\multicolumn{2}{|l|}{
% #VALUE_ID: LEX-P0060-V0044
% #FIELD_VALUE: 工艺运行记录
\fieldvalue{10:51:48；分离器4流量达到25.0 L/min}} \\ \hline
\multicolumn{2}{|l|}{
% #VALUE_ID: LEX-P0060-V0045
% #FIELD_VALUE: 工艺运行记录
\fieldvalue{10:51:48；分离器4温度达到50.0m/min}} \\ \hline
\multicolumn{2}{|l|}{
% #VALUE_ID: LEX-P0060-V0046
% #FIELD_VALUE: 工艺运行记录
\fieldvalue{10:52:00；分离器4温度达到压力2500.0 mL}} \\ \hline
\multicolumn{2}{|l|}{
% #VALUE_ID: LEX-P0060-V0047
% #FIELD_VALUE: 工艺运行记录
\fieldvalue{10:52:01；分离器4达到运行压力}} \\ \hline
\end{tabular}

\vspace{1mm}
第14页
\end{minipage}

\vspace{5mm}

\begin{minipage}[t]{0.47\linewidth}
\centering
\textbf{工艺运行报告}

\scriptsize
\begin{tabular}{|p{0.18\linewidth}|p{0.72\linewidth}|}
\hline
时间 & 运行记录\\ \hline
\multicolumn{2}{|l|}{
% #VALUE_ID: LEX-P0060-V0048
% #FIELD_VALUE: 工艺运行记录
\fieldvalue{10:52:01；分离器4达到运行压力}} \\ \hline
\multicolumn{2}{|l|}{
% #VALUE_ID: LEX-P0060-V0049
% #FIELD_VALUE: 工艺运行记录
\fieldvalue{10:52:06；分离器4流量达到25.0 L/min}} \\ \hline
\multicolumn{2}{|l|}{
% #VALUE_ID: LEX-P0060-V0050
% #FIELD_VALUE: 工艺运行记录
\fieldvalue{10:52:10；分离器4液位达到1、2、3、4、5、6、7、8、9、10}} \\ \hline
\multicolumn{2}{|l|}{
% #VALUE_ID: LEX-P0060-V0051
% #FIELD_VALUE: 工艺运行记录
\fieldvalue{10:52:10；分离器4温度达到50.0m/min}} \\ \hline
\multicolumn{2}{|l|}{
% #VALUE_ID: LEX-P0060-V0052
% #FIELD_VALUE: 工艺运行记录
\fieldvalue{10:52:10；分离器4温度达到压力2500.0 mL}} \\ \hline
\multicolumn{2}{|l|}{
% #VALUE_ID: LEX-P0060-V0053
% #FIELD_VALUE: 工艺运行记录
\fieldvalue{10:52:10；分离器4达到运行压力}} \\ \hline
\multicolumn{2}{|l|}{
% #VALUE_ID: LEX-P0060-V0054
% #FIELD_VALUE: 工艺运行记录
\fieldvalue{10:52:10；分离器4流量达到25.0 L/min}} \\ \hline
\multicolumn{2}{|l|}{
% #VALUE_ID: LEX-P0060-V0055
% #FIELD_VALUE: 工艺运行记录
\fieldvalue{10:52:10；分离器4温度达到50.0m/min}} \\ \hline
\multicolumn{2}{|l|}{
% #VALUE_ID: LEX-P0060-V0056
% #FIELD_VALUE: 工艺运行记录
\fieldvalue{10:52:10；分离器4温度达到压力2500.0 mL}} \\ \hline
\multicolumn{2}{|l|}{
% #VALUE_ID: LEX-P0060-V0057
% #FIELD_VALUE: 工艺运行记录
\fieldvalue{10:52:16；分离器4达到运行压力}} \\ \hline
\multicolumn{2}{|l|}{
% #VALUE_ID: LEX-P0060-V0058
% #FIELD_VALUE: 工艺运行记录
\fieldvalue{10:52:16；分离器4流量达到25.0 L/min}} \\ \hline
\multicolumn{2}{|l|}{
% #VALUE_ID: LEX-P0060-V0059
% #FIELD_VALUE: 工艺运行记录
\fieldvalue{10:52:16；分离器4温度达到50.0m/min}} \\ \hline
\multicolumn{2}{|l|}{
% #VALUE_ID: LEX-P0060-V0060
% #FIELD_VALUE: 工艺运行记录
\fieldvalue{10:52:16；分离器4温度达到压力2500.0 mL}} \\ \hline
\multicolumn{2}{|l|}{
% #VALUE_ID: LEX-P0060-V0061
% #FIELD_VALUE: 工艺运行记录
\fieldvalue{10:52:16；分离器4达到运行压力}} \\ \hline
\multicolumn{2}{|l|}{
% #VALUE_ID: LEX-P0060-V0062
% #FIELD_VALUE: 工艺运行记录
\fieldvalue{10:52:16；分离器4流量达到25.0 L/min}} \\ \hline
\multicolumn{2}{|l|}{
% #VALUE_ID: LEX-P0060-V0063
% #FIELD_VALUE: 工艺运行记录
\fieldvalue{10:52:16；分离器4液位达到1、2、3、4、5、6、7、8、9、10}} \\ \hline
\multicolumn{2}{|l|}{
% #VALUE_ID: LEX-P0060-V0064
% #FIELD_VALUE: 工艺运行记录
\fieldvalue{10:52:24；分离器4温度达到50.0m/min}} \\ \hline
\multicolumn{2}{|l|}{
% #VALUE_ID: LEX-P0060-V0065
% #FIELD_VALUE: 工艺运行记录
\fieldvalue{10:52:24；分离器4温度达到压力2500.0 mL}} \\ \hline
\multicolumn{2}{|l|}{
% #VALUE_ID: LEX-P0060-V0066
% #FIELD_VALUE: 工艺运行记录
\fieldvalue{10:53:25；分离器4达到运行压力}} \\ \hline
\multicolumn{2}{|l|}{
% #VALUE_ID: LEX-P0060-V0067
% #FIELD_VALUE: 工艺运行记录
\fieldvalue{10:53:25；分离器4流量达到25.0 L/min}} \\ \hline
\multicolumn{2}{|l|}{
% #VALUE_ID: LEX-P0060-V0068
% #FIELD_VALUE: 工艺运行记录
\fieldvalue{10:53:25；分离器4温度达到50.0m/min}} \\ \hline
\multicolumn{2}{|l|}{
% #VALUE_ID: LEX-P0060-V0069
% #FIELD_VALUE: 工艺运行记录
\fieldvalue{10:53:34；分离器4达到运行压力}} \\ \hline
\end{tabular}

\vspace{1mm}
第15页
\end{minipage}
\hfill
\begin{minipage}[t]{0.47\linewidth}
\centering
\textbf{工艺运行报告}

\scriptsize
\begin{tabular}{|p{0.18\linewidth}|p{0.72\linewidth}|}
\hline
时间 & 运行记录\\ \hline
\multicolumn{2}{|l|}{
% #VALUE_ID: LEX-P0060-V0070
% #FIELD_VALUE: 工艺运行记录
\fieldvalue{10:53:34；分离器1液位调整，11.2}} \\ \hline
\multicolumn{2}{|l|}{
% #VALUE_ID: LEX-P0060-V0071
% #FIELD_VALUE: 工艺运行记录
\fieldvalue{10:53:34；分离器1流量达到60.0m/min}} \\ \hline
\multicolumn{2}{|l|}{
% #VALUE_ID: LEX-P0060-V0072
% #FIELD_VALUE: 工艺运行记录
\fieldvalue{10:53:35；分离器1流量达到10.0m/min}} \\ \hline
\multicolumn{2}{|l|}{
% #VALUE_ID: LEX-P0060-V0073
% #FIELD_VALUE: 工艺运行记录
\fieldvalue{10:53:35；分离器1液位达到1、2、3、4、5、6、7、8、9、10}} \\ \hline
\multicolumn{2}{|l|}{
% #VALUE_ID: LEX-P0060-V0074
% #FIELD_VALUE: 工艺运行记录
\fieldvalue{10:53:45；分离器1达到运行压力}} \\ \hline
\multicolumn{2}{|l|}{
% #VALUE_ID: LEX-P0060-V0075
% #FIELD_VALUE: 工艺运行记录
\fieldvalue{10:53:45；分离器1流量达到25.0 L/min}} \\ \hline
\multicolumn{2}{|l|}{
% #VALUE_ID: LEX-P0060-V0076
% #FIELD_VALUE: 工艺运行记录
\fieldvalue{10:53:45；分离器1温度达到50.0m/min}} \\ \hline
\multicolumn{2}{|l|}{
% #VALUE_ID: LEX-P0060-V0077
% #FIELD_VALUE: 工艺运行记录
\fieldvalue{10:53:45；分离器1达到运行压力}} \\ \hline
\multicolumn{2}{|l|}{
% #VALUE_ID: LEX-P0060-V0078
% #FIELD_VALUE: 工艺运行记录
\fieldvalue{16:08:13；分离器1流量达到60.0m/min}} \\ \hline
\multicolumn{2}{|l|}{
% #VALUE_ID: LEX-P0060-V0079
% #FIELD_VALUE: 工艺运行记录
\fieldvalue{16:08:13；分离器1液位调整，11.2}} \\ \hline
\multicolumn{2}{|l|}{
% #VALUE_ID: LEX-P0060-V0080
% #FIELD_VALUE: 工艺运行记录
\fieldvalue{16:08:13；分离器1温度达到300m/min}} \\ \hline
\multicolumn{2}{|l|}{
% #VALUE_ID: LEX-P0060-V0081
% #FIELD_VALUE: 工艺运行记录
\fieldvalue{16:08:13；分离器1温度达到200m/min}} \\ \hline
\multicolumn{2}{|l|}{
% #VALUE_ID: LEX-P0060-V0082
% #FIELD_VALUE: 工艺运行记录
\fieldvalue{16:08:13；分离器1液位达到1、2、3、4、5、6、7、8、9、10}} \\ \hline
\multicolumn{2}{|l|}{
% #VALUE_ID: LEX-P0060-V0083
% #FIELD_VALUE: 工艺运行记录
\fieldvalue{16:08:15；分离器1达到运行压力}} \\ \hline
\multicolumn{2}{|l|}{
% #VALUE_ID: LEX-P0060-V0084
% #FIELD_VALUE: 工艺运行记录
\fieldvalue{16:08:15；分离器1流量达到25.0 L/min}} \\ \hline
\multicolumn{2}{|l|}{
% #VALUE_ID: LEX-P0060-V0085
% #FIELD_VALUE: 工艺运行记录
\fieldvalue{16:08:15；分离器1液位达到1、2、3、4、5、6、7、8、9、10}} \\ \hline
\multicolumn{2}{|l|}{
% #VALUE_ID: LEX-P0060-V0086
% #FIELD_VALUE: 工艺运行记录
\fieldvalue{16:08:16；分离器1流量达到200.0m/min}} \\ \hline
\multicolumn{2}{|l|}{
% #VALUE_ID: LEX-P0060-V0087
% #FIELD_VALUE: 工艺运行记录
\fieldvalue{16:08:16；分离器1液位调整，11.2}} \\ \hline
\multicolumn{2}{|l|}{
% #VALUE_ID: LEX-P0060-V0088
% #FIELD_VALUE: 工艺运行记录
\fieldvalue{16:08:16；分离器1达到运行压力}} \\ \hline
\multicolumn{2}{|l|}{
% #VALUE_ID: LEX-P0060-V0089
% #FIELD_VALUE: 工艺运行记录
\fieldvalue{16:08:16；分离器1液位达到1、2、3、4、5、6、7、8、9、10}} \\ \hline
\multicolumn{2}{|l|}{
% #VALUE_ID: LEX-P0060-V0090
% #FIELD_VALUE: 工艺运行记录
\fieldvalue{16:08:16；分离器1流量达到25.0 L/min}} \\ \hline
\end{tabular}

\vspace{1mm}
第16页
\end{minipage}
\end{center}

% #VALUE_ID: LEX-P0060-V0091
% #HANDWRITTEN: 4/4
\handwritten{4/4}
% LEXOID_PAGE_COMPLETED: 60/70

\newpage

\noindent
\begin{minipage}[t]{0.47\linewidth}
\centering
\textbf{◆ 勤邦生物}\\[2pt]
工艺运行报告\\[-2pt]
\rule{\linewidth}{0.4pt}

\scriptsize
\begin{tabular}{|p{0.19\linewidth}|p{0.73\linewidth}|}
\hline
16:08:16 & 分馏塔1 紧急停机--2000RPM \\ \hline
16:08:17 & 分馏塔1 处置水位--[字迹不清] \\ \hline
16:08:17 & 分馏塔1 紧急停机--[字迹不清] \\ \hline
16:08:17 & 分馏塔1 搅拌器 1、2、15 \\ \hline
16:08:17 & 分馏塔1 负载 30.00r/min \\ \hline
16:08:17 & 分馏塔1 搅拌器 1、2、4、5、6、7、8、9、10 \\ \hline
16:08:17 & 分馏塔1 紧急停机--2000RPM \\ \hline
16:08:17 & 分馏塔1 变频器--[字迹不清] \\ \hline
16:09:21 & 分馏塔1 变频器 06.5 \\ \hline
16:09:21 & 分馏塔1 处置水位--[字迹不清] \\ \hline
16:09:21 & 分馏塔1 搅拌器 1、2 \\ \hline
16:09:21 & 分馏塔1 负载 200mL/min \\ \hline
16:09:21 & 分馏塔1 负载 200mL/min \\ \hline
16:09:24 & 分馏塔1 处置水位--2034Pa \\ \hline
16:09:24 & 分馏塔1 变频器 06.5 \\ \hline
16:09:24 & 分馏塔1 搅拌器 1--[字迹不清] \\ \hline
16:09:24 & 分馏塔1 负载 200mL/min \\ \hline
16:09:24 & 分馏塔1 负载 100mL/min \\ \hline
16:09:38 & 分馏塔1 处置水位--2038Pa \\ \hline
16:09:38 & 分馏塔1 变频器--[字迹不清] \\ \hline
16:09:38 & 分馏塔1 搅拌器 1--[字迹不清] \\ \hline
16:09:38 & 分馏塔1 负载 200mL/min \\ \hline
16:09:38 & 分馏塔1 负载 100mL/min \\ \hline
\end{tabular}

\vspace{2pt}
{\tiny 第197页}
\end{minipage}
\hfill
\begin{minipage}[t]{0.47\linewidth}
\centering
\textbf{◆ 勤邦生物}\\[2pt]
工艺运行报告\\[-2pt]
\rule{\linewidth}{0.4pt}

\scriptsize
\begin{tabular}{|p{0.19\linewidth}|p{0.73\linewidth}|}
\hline
16:09:39 & 分馏塔1 变频器--2037Pa \\ \hline
16:09:39 & 分馏塔1 处置水位 85.6 \\ \hline
16:09:39 & 分馏塔1 搅拌器 1--[字迹不清] \\ \hline
16:09:39 & 分馏塔1 变频器 1--[字迹不清] \\ \hline
16:09:39 & 分馏塔1 负载 200mL/min \\ \hline
16:09:39 & 分馏塔1 搅拌器 1、2、4、5、6、7、8、9、10 \\ \hline
16:10:35 & 分馏塔1 处置水位--[字迹不清] \\ \hline
16:10:35 & 分馏塔1 变频器--[字迹不清] \\ \hline
16:10:35 & 分馏塔1 搅拌器 1--[字迹不清] \\ \hline
16:10:35 & 分馏塔1 负载 200mL/min \\ \hline
16:10:39 & 分馏塔1 处置水位--[字迹不清] \\ \hline
16:10:39 & 分馏塔1 变频器--[字迹不清] \\ \hline
16:10:39 & 分馏塔1 搅拌器 1、2、4、5、6、7、8、9、10 \\ \hline
16:10:39 & 分馏塔1 负载 200mL/min \\ \hline
16:10:39 & 分馏塔1 处置水位 75.0 \\ \hline
16:10:48 & 分馏塔1 变频器 100Hz \\ \hline
16:10:48 & 分馏塔1 负载 200mL/min \\ \hline
16:10:48 & 分馏塔1 负载 200mL/min \\ \hline
16:10:57 & 分馏塔1 处置水位--[字迹不清] \\ \hline
16:10:57 & 分馏塔1 变频器--[字迹不清] \\ \hline
\end{tabular}

\vspace{2pt}
{\tiny 第198页}
\end{minipage}

\vspace{8pt}

\noindent
\begin{minipage}[t]{0.47\linewidth}
\centering
\textbf{◆ 勤邦生物}\\[2pt]
工艺运行报告\\[-2pt]
\rule{\linewidth}{0.4pt}

\scriptsize
\begin{tabular}{|p{0.19\linewidth}|p{0.73\linewidth}|}
\hline
16:10:57 & 分馏塔1 搅拌器 1--[字迹不清] \\ \hline
16:10:58 & 分馏塔1 处置水位--[字迹不清] \\ \hline
16:10:58 & 分馏塔1 变频器--[字迹不清] \\ \hline
16:10:58 & 分馏塔1 负载 200mL/min \\ \hline
16:10:58 & 分馏塔1 负载 200mL/min \\ \hline
16:10:58 & 分馏塔1 搅拌器 1、2、4、5、6、7、8、9、10 \\ \hline
16:11:02 & 分馏塔1 处置水位--[字迹不清] \\ \hline
16:11:02 & 分馏塔1 变频器--[字迹不清] \\ \hline
16:11:02 & 分馏塔1 搅拌器 1--[字迹不清] \\ \hline
16:11:02 & 分馏塔1 负载 200mL/min \\ \hline
16:11:02 & 分馏塔1 负载 200mL/min \\ \hline
16:11:08 & 分馏塔1 处置水位--[字迹不清] \\ \hline
16:11:08 & 分馏塔1 变频器--[字迹不清] \\ \hline
16:11:08 & 分馏塔1 搅拌器 1、2、4、5、6、7、8、9、10 \\ \hline
16:11:15 & 分馏塔1 处置水位--[字迹不清] \\ \hline
16:11:19 & 分馏塔1 变频器--[字迹不清] \\ \hline
16:11:20 & 分馏塔1 负载 200mL/min \\ \hline
16:11:21 & 分馏塔1 处置水位--[字迹不清] \\ \hline
16:11:25 & 分馏塔1 变频器--[字迹不清] \\ \hline
16:11:25 & 分馏塔1 负载 200mL/min \\ \hline
16:11:34 & 分馏塔1 处置水位--[字迹不清] \\ \hline
16:11:40 & 分馏塔1 变频器--[字迹不清] \\ \hline
16:11:45 & 分馏塔1 负载 200mL/min \\ \hline
16:11:54 & 分馏塔1 处置水位--[字迹不清] \\ \hline
\end{tabular}

\vspace{2pt}
{\tiny 第199页}
\end{minipage}
\hfill
\begin{minipage}[t]{0.47\linewidth}
\centering
\textbf{◆ 勤邦生物}\\[2pt]
工艺运行报告\\[-2pt]
\rule{\linewidth}{0.4pt}

\scriptsize
\begin{tabular}{|p{0.19\linewidth}|p{0.73\linewidth}|}
\hline
16:11:55 & 分馏塔1 处置水位--[字迹不清] \\ \hline
16:12:05 & 分馏塔1 变频器--[字迹不清] \\ \hline
16:12:05 & 分馏塔1 负载 200mL/min \\ \hline
16:12:14 & 分馏塔1 处置水位--[字迹不清] \\ \hline
16:12:15 & 分馏塔1 变频器--[字迹不清] \\ \hline
16:12:15 & 分馏塔1 负载 200mL/min \\ \hline
16:12:21 & 分馏塔1 处置水位--[字迹不清] \\ \hline
16:12:25 & 分馏塔1 变频器--[字迹不清] \\ \hline
16:12:25 & 分馏塔1 负载 200mL/min \\ \hline
16:12:31 & 分馏塔1 处置水位--[字迹不清] \\ \hline
16:12:41 & 分馏塔1 变频器--[字迹不清] \\ \hline
16:12:45 & 分馏塔1 负载 200mL/min \\ \hline
16:12:54 & 分馏塔1 处置水位--[字迹不清] \\ \hline
16:12:57 & 分馏塔1 变频器--[字迹不清] \\ \hline
16:13:05 & 分馏塔1 负载 200mL/min \\ \hline
16:13:07 & 分馏塔1 处置水位--[字迹不清] \\ \hline
16:13:10 & 分馏塔1 变频器--[字迹不清] \\ \hline
16:13:10 & 分馏塔1 搅拌器 1--[字迹不清] \\ \hline
16:13:11 & 分馏塔1 负载 200mL/min \\ \hline
16:13:11 & 分馏塔1 搅拌器 1、2、4、5、6、7、8、9、10 \\ \hline
16:13:11 & 分馏塔1 处置水位--[字迹不清] \\ \hline
16:13:15 & 分馏塔1 变频器 1--[字迹不清] \\ \hline
16:13:15 & 分馏塔1 搅拌器 1--[字迹不清] \\ \hline
\end{tabular}

\vspace{2pt}
{\tiny 第200页}
\end{minipage}

\vfill

% #VALUE_ID: LEX-P0061-V0001
% #HANDWRITTEN: 5/14
\handwritten{5/14}
% LEXOID_PAGE_COMPLETED: 61/70

\newpage

\noindent
\begin{minipage}[t]{0.48\linewidth}
\centering
\textbf{康虹生物}

\vspace{0.2em}
工艺运行报告

\vspace{0.4em}
\scriptsize
\begin{tabular}{|p{0.18\linewidth}|p{0.74\linewidth}|}
\hline
16:13:15 & 分装液1（分装液1中间体・针剂）\\ \hline
16:15:30 & 分装液1（分装液1中间体・针剂）\\ \hline
16:18:13 & 分装液1（装液流速50mL/min）\\ \hline
16:18:13 & 分装液1（流速调至50mL/min）\\ \hline
16:18:13 & 分装液1（流速调至50mL/min）\\ \hline
16:19:13 & 分装液1（流速调至2、3、4、5、6、7、8、9、10）\\ \hline
16:19:13 & 分装液1（流速调至50mL/min）\\ \hline
16:25:24 & 分装液2（流速调至50mL/min）\\ \hline
16:26:30 & 分装液2（流速调至50mL/min）\\ \hline
16:30:30 & 分装液2（流速调至1、2）\\ \hline
16:30:30 & 分装液2（流速调至60.0mL/min）\\ \hline
16:30:30 & 分装液2（流速调至50mL/min）\\ \hline
16:30:30 & 分装液2（流速调至2、3、4、5、6、7、8、9、10）\\ \hline
16:30:34 & 分装液2（流速调至50mL/min）\\ \hline
16:40:40 & 分装液4（流速调至50mL/min）\\ \hline
16:40:40 & 分装液4（流速调至50mL/min）\\ \hline
16:40:40 & 分装液4（流速调至50mL/min）\\ \hline
16:40:40 & 分装液4（流速调至50mL/min）\\ \hline
16:40:40 & 分装液4（流速调至2、3、4、5、6、7、8、9、10）\\ \hline
16:40:40 & 分装液4（流速调至50mL/min）\\ \hline
16:41:24 & 分装液3（分装液3中间体・针剂）\\ \hline
16:42:24 & 分装液3（分装液3中间体・针剂）\\ \hline
\end{tabular}

\vspace{0.3em}
第21页
\end{minipage}
\hfill
\begin{minipage}[t]{0.48\linewidth}
\centering
\textbf{康虹生物}

\vspace{0.2em}
工艺运行报告

\vspace{0.4em}
\scriptsize
\begin{tabular}{|p{0.18\linewidth}|p{0.74\linewidth}|}
\hline
16:14:24 & 分装液3（流速调至1、2）\\ \hline
16:14:24 & 分装液3（流速调至60.0mL/min）\\ \hline
16:14:24 & 分装液3（流速调至50mL/min）\\ \hline
16:14:51 & 分装液3（流速调至70mL/min）\\ \hline
16:15:31 & 分装液3（流速调至80mL/min）\\ \hline
16:15:31 & 分装液3（流速调至1、2）\\ \hline
16:15:31 & 分装液3（流速调至1、2）\\ \hline
16:15:31 & 分装液3（流速调至60.0mL/min）\\ \hline
16:14:36 & 分装液3（流速调至50mL/min）\\ \hline
16:14:36 & 分装液3（流速调至50mL/min）\\ \hline
16:14:36 & 分装液3（流速调至50mL/min）\\ \hline
16:14:47 & 分装液3（流速调至1、2）\\ \hline
16:14:47 & 分装液3（流速调至1、2）\\ \hline
16:14:48 & 分装液3（流速调至50mL/min）\\ \hline
16:14:48 & 分装液3（流速调至50mL/min）\\ \hline
16:14:48 & 分装液3（流速调至1、2）\\ \hline
16:14:55 & 分装液4（流速调至1、2）\\ \hline
16:14:55 & 分装液4（流速调至50mL/min）\\ \hline
16:14:55 & 分装液4（流速调至60.0mL/min）\\ \hline
16:14:55 & 分装液4（流速调至50mL/min）\\ \hline
16:14:56 & 分装液4（流速调至1、2）\\ \hline
16:14:58 & 分装液4（流速调至50mL/min）\\ \hline
16:14:58 & 分装液4（流速调至50mL/min）\\ \hline
\end{tabular}

\vspace{0.3em}
第22页
\end{minipage}

\vspace{1.5em}

\noindent
\begin{minipage}[t]{0.48\linewidth}
\centering
\textbf{康虹生物}

\vspace{0.2em}
工艺运行报告

\vspace{0.4em}
\scriptsize
\begin{tabular}{|p{0.18\linewidth}|p{0.74\linewidth}|}
\hline
16:15:08 & 分装液4（流速调至1、2）\\ \hline
16:15:08 & 分装液4（流速调至50mL/min）\\ \hline
16:15:20 & 分装液4（流速调至1、2）\\ \hline
16:15:20 & 分装液4（流速调至50mL/min）\\ \hline
16:19:11 & 工艺空档记录（流速调至1、2）\\ \hline
\end{tabular}

\vspace{1em}
\raggedright
签名：

% #VALUE_ID: LEX-P0062-V0001
% #FIELD_VALUE: 签名
% #TODO #HANDWRITTEN: 张...; signature is partially illegible
\fieldvalue{\handwritten{张...}}

日期：

% #VALUE_ID: LEX-P0062-V0002
% #FIELD_VALUE: 日期
% #HANDWRITTEN: 6月14日
\fieldvalue{\handwritten{6月14日}}

\vspace{1em}
\centering
第23页
\end{minipage}
\hfill
\begin{minipage}[t]{0.48\linewidth}
\centering
\textbf{康虹生物}

\vspace{0.2em}
工艺运行报告

\vspace{0.6em}
\begin{tikzpicture}[x=0.075\linewidth,y=0.045\linewidth]
\draw[->] (0,0) -- (0,10.5) node[above] {流速(mL/min)};
\draw[->] (0,0) -- (10.5,0);
\foreach \y in {1,2,3,4,5,6,7,8,9,10}
  \draw[gray!25] (0,\y) -- (10,\y);
\foreach \x in {1,2,3,4,5,6,7,8,9,10}
  \draw[gray!15] (\x,0) -- (\x,10);
\node[left] at (0,2) {10};
\node[left] at (0,4) {20};
\node[left] at (0,6) {30};
\node[left] at (0,8) {40};
\node[left] at (0,10) {50};
\node[above] at (1,0) {\scriptsize 10:28:31};
\node[above] at (5,0) {\scriptsize 13:28:31};
\node[above] at (9,0) {\scriptsize 16:28:31};
\end{tikzpicture}

\vspace{0.3em}
\scriptsize
\textcolor{red}{\(\bullet\) 流道1}\quad
\textcolor{teal}{\(\bullet\) 流道2}\quad
\textcolor{purple}{\(\bullet\) 流道3}

\vspace{0.5em}
% TODO: The plotted data lines are faint/partially obscured in the source image.
\begin{flushright}
% #VALUE_ID: LEX-P0062-V0003
% #FIELD_VALUE: 日期
% #HANDWRITTEN: 6/14
\fieldvalue{\handwritten{6/14}}
\end{flushright}

\vspace{0.3em}
第24页
\end{minipage}
% LEXOID_PAGE_COMPLETED: 62/70

\newpage

\noindent
\begin{minipage}[t]{0.46\linewidth}
\centering
\textbf{蔚旺生物}\\[0.4em]
工艺运行报告\\[-0.2em]
\rule{\linewidth}{0.4pt}

\vspace{0.8em}
% TODO: Insert the visible temperature chart image; no source filename was provided.
\begin{figure}[h]
\centering
\fbox{\parbox[c][6.5cm][c]{0.92\linewidth}{\centering
温度（℃）\\
纵轴范围：$-10$ 至 $40$\\
图表区域及温度曲线
}}
\caption{图25}
\end{figure}
\end{minipage}
\hfill
\begin{minipage}[t]{0.46\linewidth}
\centering
\textbf{蔚旺生物}\\[0.4em]
工艺运行报告\\[-0.2em]
\rule{\linewidth}{0.4pt}

\vspace{0.8em}
% TODO: Insert the visible pressure chart image; no source filename was provided.
\begin{figure}[h]
\centering
\fbox{\parbox[c][6.5cm][c]{0.92\linewidth}{\centering
压力（Pa）\\
纵轴范围：$-50$ 至 $200$\\
图表区域及压力1、压力2曲线
}}
\caption{图26}
\end{figure}
\end{minipage}

\vfill

\begin{flushright}
% #VALUE_ID: LEX-P0063-V0001
% #HANDWRITTEN: 7/14
\handwritten{7/14}
\end{flushright}
% LEXOID_PAGE_COMPLETED: 63/70

\newpage

\begin{center}
\textbf{工艺运行报告}
\end{center}

\noindent
\begin{minipage}[t]{0.47\linewidth}
\small
\textbf{工艺信息}

\begin{tabular}{|p{0.18\linewidth}|p{0.29\linewidth}|p{0.18\linewidth}|p{0.25\linewidth}|}
\hline
设备名称 &
% #VALUE_ID: LEX-P0064-V0001
% #FIELD_VALUE: 设备名称
\fieldvalue{202401600000} &
设备型号 &
% #VALUE_ID: LEX-P0064-V0002
% #FIELD_VALUE: 设备型号
\fieldvalue{V1.0_Beat}
\\ \hline
工艺名称 &
% #VALUE_ID: LEX-P0064-V0003
% #FIELD_VALUE: 工艺名称
\fieldvalue{P-APC} &
产品名称 &
% #VALUE_ID: LEX-P0064-V0004
% #FIELD_VALUE: 产品名称
\fieldvalue{LW}
\\ \hline
开始时间 &
% #VALUE_ID: LEX-P0064-V0005
% #FIELD_VALUE: 开始时间
\fieldvalue{2025-12-13 09:00:24} &
结束时间 &
% #VALUE_ID: LEX-P0064-V0006
% #FIELD_VALUE: 结束时间
\fieldvalue{2025-12-18 16:18}
\\ \hline
操作员姓名 &
% #VALUE_ID: LEX-P0064-V0007
% #FIELD_VALUE: 操作员姓名
% #HANDWRITTEN: 朴海
\fieldvalue{\handwritten{朴海}} &
审核员 &
\\ \hline
\end{tabular}

\vspace{2mm}
% #TODO #HANDWRITTEN: 蓝色批注内容，疑似“设备编号：12315……”，字迹与分辨率不足以完整辨认。
% #VALUE_ID: LEX-P0064-V0008
% #HANDWRITTEN: 设备编号：12315……
\fieldvalue{\handwritten{设备编号：12315……}}

\vspace{1mm}
% #TODO #HANDWRITTEN: 蓝色批注内容，疑似日期及时间信息，无法完整辨认。
% #VALUE_ID: LEX-P0064-V0009
% #HANDWRITTEN: 日期及时间（无法完整辨认）
\fieldvalue{\handwritten{日期及时间（无法完整辨认）}}

\vspace{1mm}
% #TODO #HANDWRITTEN: 蓝色批注内容无法辨认。
% #VALUE_ID: LEX-P0064-V0010
% #HANDWRITTEN: 无法辨认
\fieldvalue{\handwritten{无法辨认}}

\textbf{管路耗材}

\begin{tabular}{|p{0.18\linewidth}|p{0.20\linewidth}|p{0.20\linewidth}|p{0.20\linewidth}|p{0.12\linewidth}|}
\hline
类别 & 产品名称 & 产品型号 & 规格\slash 货号 & 备注 \\ \hline
 & & & & \\ \hline
\end{tabular}

\vspace{2mm}
\textbf{样本信息}

\begin{tabular}{|p{0.18\linewidth}|p{0.20\linewidth}|p{0.20\linewidth}|p{0.20\linewidth}|p{0.12\linewidth}|}
\hline
类别 & 产品名称 & 产品批号 & 统一标识码 & 备注 \\ \hline
 & & & & \\ \hline
\end{tabular}

\vspace{2mm}
\textbf{制剂信息}

\begin{tabular}{|p{0.18\linewidth}|p{0.20\linewidth}|p{0.20\linewidth}|p{0.20\linewidth}|p{0.12\linewidth}|}
\hline
类别 & 产品名称 & 产品批号 & 统一标识码 & 备注 \\ \hline
 & & & & \\ \hline
\end{tabular}

\vspace{2mm}
\textbf{异常信息记录}

\begin{tabular}{|p{0.25\linewidth}|p{0.65\linewidth}|}
\hline
时间 & 异常信息 \\ \hline
 & \\ \hline
\end{tabular}

\vspace{2mm}
\textbf{工艺详细记录}

\begin{tabular}{|p{0.25\linewidth}|p{0.65\linewidth}|}
\hline
时间 & 详细记录 \\ \hline
% #VALUE_ID: LEX-P0064-V0011
% #FIELD_VALUE: 工艺详细记录时间
\fieldvalue{10:09:22} &
% #VALUE_ID: LEX-P0064-V0012
% #FIELD_VALUE: 工艺详细记录
\fieldvalue{开机待检实验1}
\\ \hline
% #VALUE_ID: LEX-P0064-V0013
% #FIELD_VALUE: 工艺详细记录时间
\fieldvalue{10:09:22} &
% #VALUE_ID: LEX-P0064-V0014
% #FIELD_VALUE: 工艺详细记录
\fieldvalue{样品待检实验1}
\\ \hline
\end{tabular}

\end{minipage}
\hfill
\begin{minipage}[t]{0.47\linewidth}
\scriptsize
\textbf{工艺运行记录}

\begin{tabular}{|p{0.22\linewidth}|p{0.70\linewidth}|}
\hline
时间 & 操作记录 \\ \hline
10:25:51 & [参数输入] 清洗液流量NVR/流程L2-10/36-167 \\ \hline
10:25:51 & [参数输入] 流速设定为5--10 L \\ \hline
10:25:51 & [参数输入] 清洗液流量10 \\ \hline
10:25:51 & [参数输入] 清洗液温度50 \\ \hline
10:25:51 & [参数输入] 清洗液加热功率300 \\ \hline
10:25:51 & [参数输入] 清洗液循环时间200s \\ \hline
10:25:51 & [参数输入] 清洗液排放时间300s \\ \hline
10:25:51 & [参数输入] 清洗液流量8.000 \\ \hline
10:25:51 & [参数输入] 清洗液压力1.2 \\ \hline
10:25:51 & [参数输入] 清洗液循环时间300s \\ \hline
10:25:51 & [参数输入] 清洗液温度50 \\ \hline
10:25:51 & [参数输入] 清洗液压力1.2 \\ \hline
10:25:52 & 分装器2： [0] \\ \hline
10:25:52 & 分装器2： [吸气] - [呼气] \\ \hline
10:25:52 & 分装器2：吸液时间12、13 \\ \hline
10:25:52 & 分装器2：推注速度300mL/min \\ \hline
10:25:52 & 分装器2：进液速度100mL/min \\ \hline
10:25:52 & 分装器2：吸液时间12、3、4、5、6、7、8、9、10 \\ \hline
10:25:52 & 分装器2：吸液时间200s \\ \hline
10:25:52 & 分装器4： [夹持] - [呼气] \\ \hline
10:25:52 & 分装器4：吸液时间12、13 \\ \hline
10:25:52 & 分装器4：推注速度300mL/min \\ \hline
10:25:52 & 分装器4：进液速度100mL/min \\ \hline
10:25:52 & 分装器4：吸液时间200mL/min \\ \hline
\end{tabular}
\end{minipage}

\vspace{8mm}

\begin{center}
\textbf{工艺运行报告}
\end{center}

\noindent
\begin{minipage}[t]{0.47\linewidth}
\scriptsize
\textbf{工艺运行记录（续）}

\begin{tabular}{|p{0.22\linewidth}|p{0.70\linewidth}|}
\hline
时间 & 操作记录 \\ \hline
10:28:52 & 分装器4：吸液时间12、3、4、5、6、7、8、9、10 \\ \hline
10:28:52 & 分装器4：推注速度200mL/min \\ \hline
10:28:52 & 分装器5： [吸液] - [呼气] \\ \hline
10:28:52 & 分装器5：吸液时间12、13 \\ \hline
10:28:52 & 分装器5：推注速度300mL/min \\ \hline
10:28:52 & 分装器5：进液速度100mL/min \\ \hline
10:28:52 & 分装器5：吸液时间12、3、4、5、6、7、8、9、10 \\ \hline
10:28:52 & 分装器5：吸液时间200s \\ \hline
10:28:53 & 分装器6： [吸液] - [呼气] \\ \hline
10:28:53 & 分装器6：吸液时间12、13 \\ \hline
10:28:53 & 分装器6：推注速度300mL/min \\ \hline
10:28:53 & 分装器6：进液速度100mL/min \\ \hline
10:28:53 & 分装器6：吸液时间12、3、4、5、6、7、8、9、10 \\ \hline
10:28:53 & 分装器6：吸液时间200s \\ \hline
10:28:53 & 分装器1： [吸液] - [呼气] \\ \hline
10:28:53 & 分装器1：吸液时间12、13 \\ \hline
10:28:53 & 分装器1：推注速度300mL/min \\ \hline
10:28:53 & 分装器1：进液速度100mL/min \\ \hline
10:28:53 & 分装器1：吸液时间12、3、4、5、6、7、8、9、10 \\ \hline
10:28:53 & 分装器1：吸液时间200s \\ \hline
10:29:50 & 分装器1：推注速度2000mL/min \\ \hline
10:29:50 & 分装器1：吸液速度2500Pa \\ \hline
10:29:50 & 分装器1：吸液时间12、3、4、5、6、7、8、9、10 \\ \hline
10:29:50 & 分装器1：推注速度300mL/min \\ \hline
10:29:50 & 分装器1：进液速度200mL/min \\ \hline
10:29:51 & 分装器1： [吸液] - [呼气] \\ \hline
10:29:51 & 分装器1：吸液时间12、13 \\ \hline
10:29:51 & 分装器1：推注速度300mL/min \\ \hline
10:29:51 & 分装器1：进液速度200mL/min \\ \hline
\end{tabular}
\end{minipage}
\hfill
\begin{minipage}[t]{0.47\linewidth}
\scriptsize
\textbf{工艺运行记录（续）}

\begin{tabular}{|p{0.22\linewidth}|p{0.70\linewidth}|}
\hline
时间 & 操作记录 \\ \hline
10:29:51 & 分装器1：吸液时间200mL/min \\ \hline
10:29:57 & 分装器1：进液速度250Pa \\ \hline
10:29:57 & 分装器1：吸液时间60s \\ \hline
10:29:57 & 分装器1： [吸液] - [呼气] \\ \hline
10:29:58 & 分装器1：吸液时间12、13 \\ \hline
10:29:58 & 分装器1：推注速度300mL/min \\ \hline
10:30:00 & 分装器1：进液速度200mL/min \\ \hline
10:30:00 & 分装器2： [吸液] - [呼气] \\ \hline
10:30:00 & 分装器2：吸液时间12、13 \\ \hline
10:30:00 & 分装器2：推注速度300mL/min \\ \hline
10:30:00 & 分装器2：进液速度100mL/min \\ \hline
10:30:05 & 分装器3： [吸液] - [呼气] \\ \hline
10:30:05 & 分装器3：吸液时间12、13 \\ \hline
10:30:05 & 分装器3：推注速度300mL/min \\ \hline
10:30:05 & 分装器3：进液速度100mL/min \\ \hline
10:30:05 & 分装器3：吸液时间12、3、4、5、6、7、8、9、10 \\ \hline
10:30:08 & 分装器3：吸液时间200s \\ \hline
10:30:08 & 分装器3：推注速度300mL/min \\ \hline
10:30:09 & 分装器3： [吸液] - [呼气] \\ \hline
10:30:09 & 分装器3：吸液时间12、13 \\ \hline
10:30:09 & 分装器3：推注速度300mL/min \\ \hline
\end{tabular}

\vspace{4mm}
% #VALUE_ID: LEX-P0064-V0015
% #HANDWRITTEN: 8/14
\handwritten{8/14}
\end{minipage}

% TODO: The four report panels are continuations of the same visible report; only rows visible on this physical page are included.
% LEXOID_PAGE_COMPLETED: 64/70

\newpage

\begin{center}
\begin{minipage}[t]{0.47\linewidth}
\centering
\textbf{工艺运行报告}

\smallskip
\begin{tabular}{|p{0.18\linewidth}|p{0.72\linewidth}|}
\hline
时间 & 运行记录 \\ \hline
10:30:09 & 分馏塔3 进料流量1000 L/min \\ \hline
10:31:09 & 分馏塔3 分馏段2：分馏器1、2、3、4、5、6、7、8、9、10 \\ \hline
10:31:14 & 分馏段2 启动 \\ \hline
10:31:14 & 分馏段2【启动】1--底1室 \\ \hline
10:31:14 & 分馏塔3 分馏段2温度 \\ \hline
10:31:14 & 分馏塔3 分馏段2压力 \\ \hline
10:31:14 & 分馏段2 压力约12 \\ \hline
10:31:15 & 分馏塔3 进料流量2000 mL/min \\ \hline
10:31:15 & 分馏塔3 进料流量1000 mL/min \\ \hline
10:31:15 & 分馏段2：分馏器1、2、3、4、5、6、7、8、9、10 \\ \hline
10:31:15 & 分馏段2 压力约12 \\ \hline
10:31:15 & 分馏塔3 进料流量约1200 mL/min \\ \hline
10:31:16 & 分馏塔3 进料流量约2000 mL/min \\ \hline
10:31:27 & 分馏塔3 进料流量约1200.500 mL \\ \hline
10:31:26 & 分馏塔3 进料流量约1200.500 mL \\ \hline
10:31:27 & 分馏塔3 进料流量约1200.500 mL \\ \hline
10:31:37 & 分馏塔3 进料流量约1200.500 mL \\ \hline
10:31:37 & 分馏塔3 进料流量约1200.500 mL \\ \hline
10:31:37 & 分馏塔3 进料流量约1200.500 mL \\ \hline
10:31:37 & 分馏塔3 进料流量约1200.500 mL \\ \hline
10:31:46 & 分馏塔3 进料流量约1200.500 mL \\ \hline
\end{tabular}

\smallskip
第65页
\end{minipage}
\hfill
\begin{minipage}[t]{0.47\linewidth}
\centering
\textbf{工艺运行报告}

\smallskip
\begin{tabular}{|p{0.18\linewidth}|p{0.72\linewidth}|}
\hline
时间 & 运行记录 \\ \hline
10:31:47 & 分馏塔3 进料流量约1200.500 mL \\ \hline
10:31:47 & 分馏塔3 进料流量约1200.500 mL \\ \hline
10:31:57 & 分馏塔3 进料流量约1200.500 mL \\ \hline
10:31:57 & 分馏塔3 进料流量约1200.500 mL \\ \hline
10:32:21 & 分馏塔3 进料流量约1200.500 mL \\ \hline
10:32:27 & 分馏塔3 进料流量约1200.500 mL \\ \hline
10:32:34 & 分馏塔3 进料流量约1200.500 mL \\ \hline
10:32:34 & 分馏塔3 进料流量约1200.500 mL \\ \hline
10:32:39 & 分馏塔3 进料流量约1200.500 mL \\ \hline
10:32:42 & 分馏塔3 进料流量约1200.500 mL \\ \hline
10:32:47 & 分馏塔3 进料流量约1200.500 mL \\ \hline
10:32:49 & 分馏塔3 进料流量约1200.500 mL \\ \hline
10:32:55 & 分馏塔3 进料流量约1200.500 mL \\ \hline
10:33:07 & 分馏塔3 进料流量约1200.500 mL \\ \hline
10:33:09 & 分馏塔3 进料流量约1200.500 mL \\ \hline
10:33:09 & 分馏塔3 进料流量约1200.500 mL \\ \hline
10:33:29 & 分馏塔3 进料流量约1200.500 mL \\ \hline
10:33:29 & 分馏塔3 进料流量约1200.500 mL \\ \hline
10:33:33 & 分馏塔3 进料流量约1200.500 mL \\ \hline
10:33:33 & 分馏塔3 进料流量约1200.500 mL \\ \hline
10:33:33 & 分馏段2【停止】1--底室 \\ \hline
\end{tabular}

\smallskip
第66页
\end{minipage}
\end{center}

\vfill

\begin{center}
\begin{minipage}[t]{0.47\linewidth}
\centering
\textbf{工艺运行报告}

\smallskip
\begin{tabular}{|p{0.18\linewidth}|p{0.72\linewidth}|}
\hline
时间 & 运行记录 \\ \hline
10:33:33 & 分馏塔3 进料流量约1200.500 mL \\ \hline
10:33:33 & 分馏塔3 进料流量约1200.500 mL \\ \hline
10:33:33 & 分馏塔3：分馏器1、2、3、4、5、6、7、8、9、10 \\ \hline
10:33:34 & 分馏塔3 进料流量约1200.500 mL \\ \hline
10:33:34 & 分馏塔3 进料流量约1200.500 mL \\ \hline
10:33:37 & 分馏塔3 进料流量约1200.500 mL \\ \hline
10:33:37 & 分馏塔3 进料流量约1200.500 mL \\ \hline
10:33:41 & 分馏塔3【启动】1--顶部 \\ \hline
10:33:41 & 分馏塔3 进料流量约1200.500 mL \\ \hline
10:33:41 & 分馏塔3 进料流量约1200.500 mL \\ \hline
10:33:41 & 分馏塔3 进料流量约1200.500 mL \\ \hline
10:33:44 & 分馏塔3 进料流量约1200.500 mL \\ \hline
10:33:44 & 分馏塔3：分馏器1、2、3、4、5、6、7、8、9、10 \\ \hline
10:33:45 & 分馏塔3 进料流量约1200.500 mL \\ \hline
10:33:46 & 分馏塔3 进料流量约1200.500 mL \\ \hline
10:33:48 & 分馏塔3【停止】1--顶部 \\ \hline
10:33:48 & 分馏塔3 进料流量约1200.500 mL \\ \hline
10:33:48 & 分馏塔3 进料流量约1200.500 mL \\ \hline
10:33:48 & 分馏塔3 进料流量约1200.500 mL \\ \hline
10:33:48 & 分馏塔3：分馏器1、2、3、4、5、6、7、8、9、10 \\ \hline
10:33:49 & 分馏塔3 进料流量约1200.500 mL \\ \hline
\end{tabular}

\smallskip
第67页
\end{minipage}
\hfill
\begin{minipage}[t]{0.47\linewidth}
\centering
\textbf{工艺运行报告}

\smallskip
\begin{tabular}{|p{0.18\linewidth}|p{0.72\linewidth}|}
\hline
时间 & 运行记录 \\ \hline
10:33:49 & 分馏塔3 进料流量约1200.500 mL \\ \hline
10:33:49 & 分馏塔3 进料流量约1200.500 mL \\ \hline
10:33:49 & 分馏塔3 进料流量约1200.500 mL \\ \hline
10:33:49 & 分馏塔3 进料流量约1200.500 mL \\ \hline
10:34:05 & 分馏塔3【启动】1--顶部 \\ \hline
10:34:05 & 分馏塔3【启动】1--顶部 \\ \hline
10:34:05 & 分馏塔3 进料流量约1200.500 mL \\ \hline
10:34:05 & 分馏塔3 进料流量约1200.500 mL \\ \hline
10:34:05 & 分馏塔3 进料流量约1200.500 mL \\ \hline
10:34:53 & 分馏塔3【停止】1--顶部 \\ \hline
10:34:53 & 分馏塔3 进料流量约1200.500 mL \\ \hline
10:34:53 & 分馏塔3 进料流量约1200.500 mL \\ \hline
10:34:53 & 分馏塔3 进料流量约1200.500 mL \\ \hline
10:34:53 & 分馏塔3【停止】1--顶部 \\ \hline
10:35:00 & 分馏塔3 进料流量约1200.500 mL \\ \hline
10:35:00 & 分馏塔3 进料流量约1200.500 mL \\ \hline
10:35:00 & 分馏塔3 进料流量约1200.500 mL \\ \hline
10:35:02 & 分馏塔3 进料流量约1200.500 mL \\ \hline
10:35:02 & 分馏塔3 进料流量约1200.500 mL \\ \hline
10:35:02 & 分馏塔3 进料流量约1200.500 mL \\ \hline
\end{tabular}

\smallskip
第68页
\end{minipage}
\end{center}

\vspace{1em}

% #VALUE_ID: LEX-P0065-V0001
% #TODO #HANDWRITTEN: 9/4; blue handwritten notation, exact characters uncertain
\handwritten{9/4}
% LEXOID_PAGE_COMPLETED: 65/70

\newpage

\begin{center}
\begin{minipage}[t]{0.46\linewidth}
\centering
\textbf{康耐生物}\\[2pt]
工艺运行报告

\vspace{4pt}
\begin{tabular}{|p{0.18\linewidth}|p{0.72\linewidth}|}
\hline
时间 & 运行记录 \\ \hline
10:35:02 & 分装器具及相关运行记录 \\ \hline
10:36:03 & 分装器具及相关运行记录 \\ \hline
10:36:05 & 分装器具及相关运行记录 \\ \hline
10:36:05 & 分装器具及相关运行记录 \\ \hline
10:35:11 & 分装器具及相关运行记录 \\ \hline
10:35:13 & 分装器具及相关运行记录 \\ \hline
10:35:13 & 分装器具及相关运行记录 \\ \hline
10:35:13 & 分装器具及相关运行记录 \\ \hline
10:35:22 & 分装器具及相关运行记录 \\ \hline
10:35:22 & 分装器具及相关运行记录 \\ \hline
10:39:24 & 分装器具及相关运行记录 \\ \hline
10:39:24 & 分装器具及相关运行记录 \\ \hline
10:39:24 & 分装器具及相关运行记录 \\ \hline
10:39:24 & 分装器具及相关运行记录 \\ \hline
10:39:24 & 分装器具及相关运行记录 \\ \hline
10:40:24 & 分装器具及相关运行记录 \\ \hline
10:40:24 & 分装器具及相关运行记录 \\ \hline
10:40:24 & 分装器具及相关运行记录 \\ \hline
10:40:24 & 分装器具及相关运行记录 \\ \hline
10:40:24 & 分装器具及相关运行记录 \\ \hline
10:40:24 & 分装器具及相关运行记录 \\ \hline
10:40:24 & 分装器具及相关运行记录 \\ \hline
10:40:24 & 分装器具及相关运行记录 \\ \hline
10:40:24 & 分装器具及相关运行记录 \\ \hline
10:40:24 & 分装器具及相关运行记录 \\ \hline
10:40:24 & 分装器具及相关运行记录 \\ \hline
10:40:24 & 分装器具及相关运行记录 \\ \hline
10:40:24 & 分装器具及相关运行记录 \\ \hline
10:40:24 & 分装器具及相关运行记录 \\ \hline
10:40:24 & 分装器具及相关运行记录 \\ \hline
10:41:09 & 分装器具及相关运行记录 \\ \hline
10:41:09 & 分装器具及相关运行记录 \\ \hline
\end{tabular}

\vspace{3pt}
第99页
\end{minipage}
\hfill
\begin{minipage}[t]{0.46\linewidth}
\centering
\textbf{康耐生物}\\[2pt]
工艺运行报告

\vspace{4pt}
\begin{tabular}{|p{0.18\linewidth}|p{0.72\linewidth}|}
\hline
时间 & 运行记录 \\ \hline
10:41:09 & 分装器具及相关运行记录 \\ \hline
10:41:09 & 分装器具及相关运行记录 \\ \hline
10:41:09 & 分装器具及相关运行记录 \\ \hline
10:41:09 & 分装器具及相关运行记录 \\ \hline
10:41:38 & 分装器具及相关运行记录 \\ \hline
10:41:38 & 分装器具及相关运行记录 \\ \hline
10:41:39 & 分装器具及相关运行记录 \\ \hline
10:41:39 & 分装器具及相关运行记录 \\ \hline
10:41:39 & 分装器具及相关运行记录 \\ \hline
10:41:50 & 分装器具及相关运行记录 \\ \hline
10:42:00 & 分装器具及相关运行记录 \\ \hline
10:42:09 & 分装器具及相关运行记录 \\ \hline
10:42:21 & 分装器具及相关运行记录 \\ \hline
10:42:27 & 分装器具及相关运行记录 \\ \hline
10:42:27 & 分装器具及相关运行记录 \\ \hline
10:42:27 & 分装器具及相关运行记录 \\ \hline
10:42:27 & 分装器具及相关运行记录 \\ \hline
10:42:27 & 分装器具及相关运行记录 \\ \hline
10:42:27 & 分装器具及相关运行记录 \\ \hline
10:42:27 & 分装器具及相关运行记录 \\ \hline
10:42:27 & 分装器具及相关运行记录 \\ \hline
10:42:27 & 分装器具及相关运行记录 \\ \hline
10:42:27 & 分装器具及相关运行记录 \\ \hline
10:42:27 & 分装器具及相关运行记录 \\ \hline
10:42:27 & 分装器具及相关运行记录 \\ \hline
10:42:27 & 分装器具及相关运行记录 \\ \hline
10:42:27 & 分装器具及相关运行记录 \\ \hline
10:42:27 & 分装器具及相关运行记录 \\ \hline
10:42:27 & 分装器具及相关运行记录 \\ \hline
10:42:27 & 分装器具及相关运行记录 \\ \hline
10:42:27 & 分装器具及相关运行记录 \\ \hline
10:42:27 & 分装器具及相关运行记录 \\ \hline
10:42:27 & 分装器具及相关运行记录 \\ \hline
10:42:27 & 分装器具及相关运行记录 \\ \hline
10:42:27 & 分装器具及相关运行记录 \\ \hline
10:42:27 & 分装器具及相关运行记录 \\ \hline
10:42:27 & 分装器具及相关运行记录 \\ \hline
\end{tabular}

\vspace{3pt}
第100页
\end{minipage}

\vspace{18pt}

\begin{minipage}[t]{0.46\linewidth}
\centering
\textbf{康耐生物}\\[2pt]
工艺运行报告

\vspace{4pt}
\begin{tabular}{|p{0.18\linewidth}|p{0.72\linewidth}|}
\hline
时间 & 运行记录 \\ \hline
10:43:05 & 分装器具及相关运行记录 \\ \hline
10:43:05 & 分装器具及相关运行记录 \\ \hline
10:43:05 & 分装器具及相关运行记录 \\ \hline
10:43:07 & 分装器具及相关运行记录 \\ \hline
10:43:07 & 分装器具及相关运行记录 \\ \hline
10:43:07 & 分装器具及相关运行记录 \\ \hline
10:43:07 & 分装器具及相关运行记录 \\ \hline
10:43:07 & 分装器具及相关运行记录 \\ \hline
10:43:07 & 分装器具及相关运行记录 \\ \hline
10:43:07 & 分装器具及相关运行记录 \\ \hline
10:43:41 & 分装器具及相关运行记录 \\ \hline
10:43:41 & 分装器具及相关运行记录 \\ \hline
10:43:41 & 分装器具及相关运行记录 \\ \hline
10:43:41 & 分装器具及相关运行记录 \\ \hline
10:43:41 & 分装器具及相关运行记录 \\ \hline
10:43:54 & 分装器具及相关运行记录 \\ \hline
10:43:54 & 分装器具及相关运行记录 \\ \hline
10:43:56 & 分装器具及相关运行记录 \\ \hline
10:44:03 & 分装器具及相关运行记录 \\ \hline
10:44:07 & 分装器具及相关运行记录 \\ \hline
10:44:07 & 分装器具及相关运行记录 \\ \hline
10:44:07 & 分装器具及相关运行记录 \\ \hline
10:44:08 & 分装器具及相关运行记录 \\ \hline
\end{tabular}

\vspace{3pt}
第101页
\end{minipage}
\hfill
\begin{minipage}[t]{0.46\linewidth}
\centering
\textbf{康耐生物}\\[2pt]
工艺运行报告

\vspace{4pt}
\begin{tabular}{|p{0.18\linewidth}|p{0.72\linewidth}|}
\hline
时间 & 运行记录 \\ \hline
10:44:08 & 分装器具及相关运行记录 \\ \hline
10:44:08 & 分装器具及相关运行记录 \\ \hline
10:44:12 & 分装器具及相关运行记录 \\ \hline
10:44:16 & 分装器具及相关运行记录 \\ \hline
10:44:21 & 分装器具及相关运行记录 \\ \hline
10:44:21 & 分装器具及相关运行记录 \\ \hline
10:44:21 & 分装器具及相关运行记录 \\ \hline
10:44:21 & 分装器具及相关运行记录 \\ \hline
10:44:21 & 分装器具及相关运行记录 \\ \hline
10:44:21 & 分装器具及相关运行记录 \\ \hline
10:44:21 & 分装器具及相关运行记录 \\ \hline
10:44:24 & 分装器具及相关运行记录 \\ \hline
10:44:24 & 分装器具及相关运行记录 \\ \hline
10:44:24 & 分装器具及相关运行记录 \\ \hline
10:44:24 & 分装器具及相关运行记录 \\ \hline
10:44:24 & 分装器具及相关运行记录 \\ \hline
10:45:08 & 分装器具及相关运行记录 \\ \hline
10:45:17 & 分装器具及相关运行记录 \\ \hline
10:45:22 & 分装器具及相关运行记录 \\ \hline
10:45:22 & 分装器具及相关运行记录 \\ \hline
10:45:22 & 分装器具及相关运行记录 \\ \hline
10:45:22 & 分装器具及相关运行记录 \\ \hline
\end{tabular}

\vspace{3pt}
第102页
\end{minipage}
\end{center}

% TODO: The detailed Chinese event descriptions in the four densely printed tables are too small to be transcribed reliably from the provided page image.

% #VALUE_ID: LEX-P0066-V0001
% #HANDWRITTEN: 10/14
\handwritten{10/14}
% LEXOID_PAGE_COMPLETED: 66/70

\newpage

\noindent
\begin{minipage}[t]{0.48\linewidth}
\centering
\textbf{工艺运行报告}

\small
\begin{tabular}{|p{0.20\linewidth}|p{0.72\linewidth}|}
\hline
时间 & 运行记录\\ \hline
\multicolumn{2}{|l|}{
% #VALUE_ID: LEX-P0067-V0001
% #FIELD_VALUE: 工艺运行记录
\fieldvalue{\texttt{10:45:22}\quad 记录内容（原图文字过小）}} \\ \hline
\multicolumn{2}{|l|}{
% #VALUE_ID: LEX-P0067-V0002
% #FIELD_VALUE: 工艺运行记录
\fieldvalue{\texttt{10:45:22}\quad 记录内容（原图文字过小）}} \\ \hline
\multicolumn{2}{|l|}{
% #VALUE_ID: LEX-P0067-V0003
% #FIELD_VALUE: 工艺运行记录
\fieldvalue{\texttt{10:45:27}\quad 记录内容（原图文字过小）}} \\ \hline
\multicolumn{2}{|l|}{
% #VALUE_ID: LEX-P0067-V0004
% #FIELD_VALUE: 工艺运行记录
\fieldvalue{\texttt{10:45:28}\quad 记录内容（原图文字过小）}} \\ \hline
\multicolumn{2}{|l|}{
% #VALUE_ID: LEX-P0067-V0005
% #FIELD_VALUE: 工艺运行记录
\fieldvalue{\texttt{10:45:33}\quad 记录内容（原图文字过小）}} \\ \hline
\multicolumn{2}{|l|}{
% #VALUE_ID: LEX-P0067-V0006
% #FIELD_VALUE: 工艺运行记录
\fieldvalue{\texttt{10:45:34}\quad 记录内容（原图文字过小）}} \\ \hline
\multicolumn{2}{|l|}{
% #VALUE_ID: LEX-P0067-V0007
% #FIELD_VALUE: 工艺运行记录
\fieldvalue{\texttt{10:45:34}\quad 记录内容（原图文字过小）}} \\ \hline
\multicolumn{2}{|l|}{
% #VALUE_ID: LEX-P0067-V0008
% #FIELD_VALUE: 工艺运行记录
\fieldvalue{\texttt{10:45:34}\quad 记录内容（原图文字过小）}} \\ \hline
\multicolumn{2}{|l|}{
% #VALUE_ID: LEX-P0067-V0009
% #FIELD_VALUE: 工艺运行记录
\fieldvalue{\texttt{10:45:39}\quad 记录内容（原图文字过小）}} \\ \hline
\multicolumn{2}{|l|}{
% #VALUE_ID: LEX-P0067-V0010
% #FIELD_VALUE: 工艺运行记录
\fieldvalue{\texttt{10:45:40}\quad 记录内容（原图文字过小）}} \\ \hline
\multicolumn{2}{|l|}{
% #VALUE_ID: LEX-P0067-V0011
% #FIELD_VALUE: 工艺运行记录
\fieldvalue{\texttt{10:45:40}\quad 记录内容（原图文字过小）}} \\ \hline
\multicolumn{2}{|l|}{
% #VALUE_ID: LEX-P0067-V0012
% #FIELD_VALUE: 工艺运行记录
\fieldvalue{\texttt{10:45:40}\quad 记录内容（原图文字过小）}} \\ \hline
\multicolumn{2}{|l|}{
% #VALUE_ID: LEX-P0067-V0013
% #FIELD_VALUE: 工艺运行记录
\fieldvalue{\texttt{10:45:40}\quad 记录内容（原图文字过小）}} \\ \hline
\multicolumn{2}{|l|}{
% #VALUE_ID: LEX-P0067-V0014
% #FIELD_VALUE: 工艺运行记录
\fieldvalue{\texttt{10:45:44}\quad 记录内容（原图文字过小）}} \\ \hline
\multicolumn{2}{|l|}{
% #VALUE_ID: LEX-P0067-V0015
% #FIELD_VALUE: 工艺运行记录
\fieldvalue{\texttt{10:45:44}\quad 记录内容（原图文字过小）}} \\ \hline
\multicolumn{2}{|l|}{
% #VALUE_ID: LEX-P0067-V0016
% #FIELD_VALUE: 工艺运行记录
\fieldvalue{\texttt{10:45:51}\quad 记录内容（原图文字过小）}} \\ \hline
\multicolumn{2}{|l|}{
% #VALUE_ID: LEX-P0067-V0017
% #FIELD_VALUE: 工艺运行记录
\fieldvalue{\texttt{10:45:52}\quad 记录内容（原图文字过小）}} \\ \hline
\multicolumn{2}{|l|}{
% #VALUE_ID: LEX-P0067-V0018
% #FIELD_VALUE: 工艺运行记录
\fieldvalue{\texttt{10:45:52}\quad 记录内容（原图文字过小）}} \\ \hline
\multicolumn{2}{|l|}{
% #VALUE_ID: LEX-P0067-V0019
% #FIELD_VALUE: 工艺运行记录
\fieldvalue{\texttt{10:45:55}\quad 记录内容（原图文字过小）}} \\ \hline
\multicolumn{2}{|l|}{
% #VALUE_ID: LEX-P0067-V0020
% #FIELD_VALUE: 工艺运行记录
\fieldvalue{\texttt{10:45:58}\quad 记录内容（原图文字过小）}} \\ \hline
\multicolumn{2}{|l|}{
% #VALUE_ID: LEX-P0067-V0021
% #FIELD_VALUE: 工艺运行记录
\fieldvalue{\texttt{10:46:02}\quad 记录内容（原图文字过小）}} \\ \hline
\multicolumn{2}{|l|}{
% #VALUE_ID: LEX-P0067-V0022
% #FIELD_VALUE: 工艺运行记录
\fieldvalue{\texttt{10:46:06}\quad 记录内容（原图文字过小）}} \\ \hline
\multicolumn{2}{|l|}{
% #VALUE_ID: LEX-P0067-V0023
% #FIELD_VALUE: 工艺运行记录
\fieldvalue{\texttt{10:46:08}\quad 记录内容（原图文字过小）}} \\ \hline
\end{tabular}

\vspace{2pt}
第13页
\end{minipage}
\hfill
\begin{minipage}[t]{0.48\linewidth}
\centering
\textbf{工艺运行报告}

\small
\begin{tabular}{|p{0.20\linewidth}|p{0.72\linewidth}|}
\hline
时间 & 运行记录\\ \hline
\multicolumn{2}{|l|}{
% #VALUE_ID: LEX-P0067-V0024
% #FIELD_VALUE: 工艺运行记录
\fieldvalue{\texttt{10:46:11}\quad 记录内容（原图文字过小）}} \\ \hline
\multicolumn{2}{|l|}{
% #VALUE_ID: LEX-P0067-V0025
% #FIELD_VALUE: 工艺运行记录
\fieldvalue{\texttt{10:46:12}\quad 记录内容（原图文字过小）}} \\ \hline
\multicolumn{2}{|l|}{
% #VALUE_ID: LEX-P0067-V0026
% #FIELD_VALUE: 工艺运行记录
\fieldvalue{\texttt{10:46:13}\quad 记录内容（原图文字过小）}} \\ \hline
\multicolumn{2}{|l|}{
% #VALUE_ID: LEX-P0067-V0027
% #FIELD_VALUE: 工艺运行记录
\fieldvalue{\texttt{10:46:13}\quad 记录内容（原图文字过小）}} \\ \hline
\multicolumn{2}{|l|}{
% #VALUE_ID: LEX-P0067-V0028
% #FIELD_VALUE: 工艺运行记录
\fieldvalue{\texttt{10:46:13}\quad 记录内容（原图文字过小）}} \\ \hline
\multicolumn{2}{|l|}{
% #VALUE_ID: LEX-P0067-V0029
% #FIELD_VALUE: 工艺运行记录
\fieldvalue{\texttt{10:46:13}\quad 记录内容（原图文字过小）}} \\ \hline
\multicolumn{2}{|l|}{
% #VALUE_ID: LEX-P0067-V0030
% #FIELD_VALUE: 工艺运行记录
\fieldvalue{\texttt{10:46:14}\quad 记录内容（原图文字过小）}} \\ \hline
\multicolumn{2}{|l|}{
% #VALUE_ID: LEX-P0067-V0031
% #FIELD_VALUE: 工艺运行记录
\fieldvalue{\texttt{10:46:20}\quad 记录内容（原图文字过小）}} \\ \hline
\multicolumn{2}{|l|}{
% #VALUE_ID: LEX-P0067-V0032
% #FIELD_VALUE: 工艺运行记录
\fieldvalue{\texttt{10:46:40}\quad 记录内容（原图文字过小）}} \\ \hline
\multicolumn{2}{|l|}{
% #VALUE_ID: LEX-P0067-V0033
% #FIELD_VALUE: 工艺运行记录
\fieldvalue{\texttt{10:46:51}\quad 记录内容（原图文字过小）}} \\ \hline
\multicolumn{2}{|l|}{
% #VALUE_ID: LEX-P0067-V0034
% #FIELD_VALUE: 工艺运行记录
\fieldvalue{\texttt{10:46:58}\quad 记录内容（原图文字过小）}} \\ \hline
\multicolumn{2}{|l|}{
% #VALUE_ID: LEX-P0067-V0035
% #FIELD_VALUE: 工艺运行记录
\fieldvalue{\texttt{10:47:01}\quad 记录内容（原图文字过小）}} \\ \hline
\multicolumn{2}{|l|}{
% #VALUE_ID: LEX-P0067-V0036
% #FIELD_VALUE: 工艺运行记录
\fieldvalue{\texttt{10:47:01}\quad 记录内容（原图文字过小）}} \\ \hline
\multicolumn{2}{|l|}{
% #VALUE_ID: LEX-P0067-V0037
% #FIELD_VALUE: 工艺运行记录
\fieldvalue{\texttt{10:47:20}\quad 记录内容（原图文字过小）}} \\ \hline
\multicolumn{2}{|l|}{
% #VALUE_ID: LEX-P0067-V0038
% #FIELD_VALUE: 工艺运行记录
\fieldvalue{\texttt{10:47:25}\quad 记录内容（原图文字过小）}} \\ \hline
\multicolumn{2}{|l|}{
% #VALUE_ID: LEX-P0067-V0039
% #FIELD_VALUE: 工艺运行记录
\fieldvalue{\texttt{10:47:25}\quad 记录内容（原图文字过小）}} \\ \hline
\multicolumn{2}{|l|}{
% #VALUE_ID: LEX-P0067-V0040
% #FIELD_VALUE: 工艺运行记录
\fieldvalue{\texttt{10:47:25}\quad 记录内容（原图文字过小）}} \\ \hline
\multicolumn{2}{|l|}{
% #VALUE_ID: LEX-P0067-V0041
% #FIELD_VALUE: 工艺运行记录
\fieldvalue{\texttt{10:47:25}\quad 记录内容（原图文字过小）}} \\ \hline
\multicolumn{2}{|l|}{
% #VALUE_ID: LEX-P0067-V0042
% #FIELD_VALUE: 工艺运行记录
\fieldvalue{\texttt{10:47:25}\quad 记录内容（原图文字过小）}} \\ \hline
\multicolumn{2}{|l|}{
% #VALUE_ID: LEX-P0067-V0043
% #FIELD_VALUE: 工艺运行记录
\fieldvalue{\texttt{10:47:35}\quad 记录内容（原图文字过小）}} \\ \hline
\multicolumn{2}{|l|}{
% #VALUE_ID: LEX-P0067-V0044
% #FIELD_VALUE: 工艺运行记录
\fieldvalue{\texttt{10:47:35}\quad 记录内容（原图文字过小）}} \\ \hline
\multicolumn{2}{|l|}{
% #VALUE_ID: LEX-P0067-V0045
% #FIELD_VALUE: 工艺运行记录
\fieldvalue{\texttt{10:47:42}\quad 记录内容（原图文字过小）}} \\ \hline
\multicolumn{2}{|l|}{
% #VALUE_ID: LEX-P0067-V0046
% #FIELD_VALUE: 工艺运行记录
\fieldvalue{\texttt{10:47:44}\quad 记录内容（原图文字过小）}} \\ \hline
\multicolumn{2}{|l|}{
% #VALUE_ID: LEX-P0067-V0047
% #FIELD_VALUE: 工艺运行记录
\fieldvalue{\texttt{10:47:52}\quad 记录内容（原图文字过小）}} \\ \hline
\end{tabular}

\vspace{2pt}
第14页
\end{minipage}

\vspace{8pt}

\noindent
\begin{minipage}[t]{0.48\linewidth}
\centering
\textbf{工艺运行报告}

\small
\begin{tabular}{|p{0.20\linewidth}|p{0.72\linewidth}|}
\hline
时间 & 运行记录\\ \hline
\multicolumn{2}{|l|}{
% #VALUE_ID: LEX-P0067-V0048
% #FIELD_VALUE: 工艺运行记录
\fieldvalue{\texttt{10:47:56}\quad 记录内容（原图文字过小）}} \\ \hline
\multicolumn{2}{|l|}{
% #VALUE_ID: LEX-P0067-V0049
% #FIELD_VALUE: 工艺运行记录
\fieldvalue{\texttt{10:47:56}\quad 记录内容（原图文字过小）}} \\ \hline
\multicolumn{2}{|l|}{
% #VALUE_ID: LEX-P0067-V0050
% #FIELD_VALUE: 工艺运行记录
\fieldvalue{\texttt{10:47:56}\quad 记录内容（原图文字过小）}} \\ \hline
\multicolumn{2}{|l|}{
% #VALUE_ID: LEX-P0067-V0051
% #FIELD_VALUE: 工艺运行记录
\fieldvalue{\texttt{10:47:56}\quad 记录内容（原图文字过小）}} \\ \hline
\multicolumn{2}{|l|}{
% #VALUE_ID: LEX-P0067-V0052
% #FIELD_VALUE: 工艺运行记录
\fieldvalue{\texttt{10:47:56}\quad 记录内容（原图文字过小）}} \\ \hline
\multicolumn{2}{|l|}{
% #VALUE_ID: LEX-P0067-V0053
% #FIELD_VALUE: 工艺运行记录
\fieldvalue{\texttt{10:47:56}\quad 记录内容（原图文字过小）}} \\ \hline
\multicolumn{2}{|l|}{
% #VALUE_ID: LEX-P0067-V0054
% #FIELD_VALUE: 工艺运行记录
\fieldvalue{\texttt{10:47:56}\quad 记录内容（原图文字过小）}} \\ \hline
\multicolumn{2}{|l|}{
% #VALUE_ID: LEX-P0067-V0055
% #FIELD_VALUE: 工艺运行记录
\fieldvalue{\texttt{10:47:56}\quad 记录内容（原图文字过小）}} \\ \hline
\multicolumn{2}{|l|}{
% #VALUE_ID: LEX-P0067-V0056
% #FIELD_VALUE: 工艺运行记录
\fieldvalue{\texttt{10:48:05}\quad 记录内容（原图文字过小）}} \\ \hline
\multicolumn{2}{|l|}{
% #VALUE_ID: LEX-P0067-V0057
% #FIELD_VALUE: 工艺运行记录
\fieldvalue{\texttt{10:48:05}\quad 记录内容（原图文字过小）}} \\ \hline
\multicolumn{2}{|l|}{
% #VALUE_ID: LEX-P0067-V0058
% #FIELD_VALUE: 工艺运行记录
\fieldvalue{\texttt{10:48:09}\quad 记录内容（原图文字过小）}} \\ \hline
\multicolumn{2}{|l|}{
% #VALUE_ID: LEX-P0067-V0059
% #FIELD_VALUE: 工艺运行记录
\fieldvalue{\texttt{10:48:09}\quad 记录内容（原图文字过小）}} \\ \hline
\multicolumn{2}{|l|}{
% #VALUE_ID: LEX-P0067-V0060
% #FIELD_VALUE: 工艺运行记录
\fieldvalue{\texttt{10:48:10}\quad 记录内容（原图文字过小）}} \\ \hline
\multicolumn{2}{|l|}{
% #VALUE_ID: LEX-P0067-V0061
% #FIELD_VALUE: 工艺运行记录
\fieldvalue{\texttt{10:48:10}\quad 记录内容（原图文字过小）}} \\ \hline
\multicolumn{2}{|l|}{
% #VALUE_ID: LEX-P0067-V0062
% #FIELD_VALUE: 工艺运行记录
\fieldvalue{\texttt{10:48:10}\quad 记录内容（原图文字过小）}} \\ \hline
\multicolumn{2}{|l|}{
% #VALUE_ID: LEX-P0067-V0063
% #FIELD_VALUE: 工艺运行记录
\fieldvalue{\texttt{10:48:10}\quad 记录内容（原图文字过小）}} \\ \hline
\multicolumn{2}{|l|}{
% #VALUE_ID: LEX-P0067-V0064
% #FIELD_VALUE: 工艺运行记录
\fieldvalue{\texttt{10:48:10}\quad 记录内容（原图文字过小）}} \\ \hline
\multicolumn{2}{|l|}{
% #VALUE_ID: LEX-P0067-V0065
% #FIELD_VALUE: 工艺运行记录
\fieldvalue{\texttt{10:49:07}\quad 记录内容（原图文字过小）}} \\ \hline
\multicolumn{2}{|l|}{
% #VALUE_ID: LEX-P0067-V0066
% #FIELD_VALUE: 工艺运行记录
\fieldvalue{\texttt{10:49:07}\quad 记录内容（原图文字过小）}} \\ \hline
\multicolumn{2}{|l|}{
% #VALUE_ID: LEX-P0067-V0067
% #FIELD_VALUE: 工艺运行记录
\fieldvalue{\texttt{10:49:07}\quad 记录内容（原图文字过小）}} \\ \hline
\multicolumn{2}{|l|}{
% #VALUE_ID: LEX-P0067-V0068
% #FIELD_VALUE: 工艺运行记录
\fieldvalue{\texttt{10:49:07}\quad 记录内容（原图文字过小）}} \\ \hline
\multicolumn{2}{|l|}{
% #VALUE_ID: LEX-P0067-V0069
% #FIELD_VALUE: 工艺运行记录
\fieldvalue{\texttt{10:49:18}\quad 记录内容（原图文字过小）}} \\ \hline
\multicolumn{2}{|l|}{
% #VALUE_ID: LEX-P0067-V0070
% #FIELD_VALUE: 工艺运行记录
\fieldvalue{\texttt{10:49:20}\quad 记录内容（原图文字过小）}} \\ \hline
\end{tabular}

\vspace{2pt}
第15页
\end{minipage}
\hfill
\begin{minipage}[t]{0.48\linewidth}
\centering
\textbf{工艺运行报告}

\small
\begin{tabular}{|p{0.20\linewidth}|p{0.72\linewidth}|}
\hline
时间 & 运行记录\\ \hline
\multicolumn{2}{|l|}{
% #VALUE_ID: LEX-P0067-V0071
% #FIELD_VALUE: 工艺运行记录
\fieldvalue{\texttt{10:49:20}\quad 记录内容（原图文字过小）}} \\ \hline
\multicolumn{2}{|l|}{
% #VALUE_ID: LEX-P0067-V0072
% #FIELD_VALUE: 工艺运行记录
\fieldvalue{\texttt{10:49:21}\quad 记录内容（原图文字过小）}} \\ \hline
\multicolumn{2}{|l|}{
% #VALUE_ID: LEX-P0067-V0073
% #FIELD_VALUE: 工艺运行记录
\fieldvalue{\texttt{10:49:21}\quad 记录内容（原图文字过小）}} \\ \hline
\multicolumn{2}{|l|}{
% #VALUE_ID: LEX-P0067-V0074
% #FIELD_VALUE: 工艺运行记录
\fieldvalue{\texttt{10:49:21}\quad 记录内容（原图文字过小）}} \\ \hline
\multicolumn{2}{|l|}{
% #VALUE_ID: LEX-P0067-V0075
% #FIELD_VALUE: 工艺运行记录
\fieldvalue{\texttt{10:49:27}\quad 记录内容（原图文字过小）}} \\ \hline
\multicolumn{2}{|l|}{
% #VALUE_ID: LEX-P0067-V0076
% #FIELD_VALUE: 工艺运行记录
\fieldvalue{\texttt{10:49:27}\quad 记录内容（原图文字过小）}} \\ \hline
\multicolumn{2}{|l|}{
% #VALUE_ID: LEX-P0067-V0077
% #FIELD_VALUE: 工艺运行记录
\fieldvalue{\texttt{10:49:27}\quad 记录内容（原图文字过小）}} \\ \hline
\multicolumn{2}{|l|}{
% #VALUE_ID: LEX-P0067-V0078
% #FIELD_VALUE: 工艺运行记录
\fieldvalue{\texttt{10:49:27}\quad 记录内容（原图文字过小）}} \\ \hline
\multicolumn{2}{|l|}{
% #VALUE_ID: LEX-P0067-V0079
% #FIELD_VALUE: 工艺运行记录
\fieldvalue{\texttt{10:49:32}\quad 记录内容（原图文字过小）}} \\ \hline
\multicolumn{2}{|l|}{
% #VALUE_ID: LEX-P0067-V0080
% #FIELD_VALUE: 工艺运行记录
\fieldvalue{\texttt{10:49:32}\quad 记录内容（原图文字过小）}} \\ \hline
\multicolumn{2}{|l|}{
% #VALUE_ID: LEX-P0067-V0081
% #FIELD_VALUE: 工艺运行记录
\fieldvalue{\texttt{10:49:40}\quad 记录内容（原图文字过小）}} \\ \hline
\multicolumn{2}{|l|}{
% #VALUE_ID: LEX-P0067-V0082
% #FIELD_VALUE: 工艺运行记录
\fieldvalue{\texttt{16:07:52}\quad 记录内容（原图文字过小）}} \\ \hline
\multicolumn{2}{|l|}{
% #VALUE_ID: LEX-P0067-V0083
% #FIELD_VALUE: 工艺运行记录
\fieldvalue{\texttt{16:07:52}\quad 记录内容（原图文字过小）}} \\ \hline
\multicolumn{2}{|l|}{
% #VALUE_ID: LEX-P0067-V0084
% #FIELD_VALUE: 工艺运行记录
\fieldvalue{\texttt{16:07:52}\quad 记录内容（原图文字过小）}} \\ \hline
\multicolumn{2}{|l|}{
% #VALUE_ID: LEX-P0067-V0085
% #FIELD_VALUE: 工艺运行记录
\fieldvalue{\texttt{16:07:52}\quad 记录内容（原图文字过小）}} \\ \hline
\multicolumn{2}{|l|}{
% #VALUE_ID: LEX-P0067-V0086
% #FIELD_VALUE: 工艺运行记录
\fieldvalue{\texttt{16:07:52}\quad 记录内容（原图文字过小）}} \\ \hline
\multicolumn{2}{|l|}{
% #VALUE_ID: LEX-P0067-V0087
% #FIELD_VALUE: 工艺运行记录
\fieldvalue{\texttt{16:07:52}\quad 记录内容（原图文字过小）}} \\ \hline
\multicolumn{2}{|l|}{
% #VALUE_ID: LEX-P0067-V0088
% #FIELD_VALUE: 工艺运行记录
\fieldvalue{\texttt{16:07:52}\quad 记录内容（原图文字过小）}} \\ \hline
\multicolumn{2}{|l|}{
% #VALUE_ID: LEX-P0067-V0089
% #FIELD_VALUE: 工艺运行记录
\fieldvalue{\texttt{16:07:52}\quad 记录内容（原图文字过小）}} \\ \hline
\multicolumn{2}{|l|}{
% #VALUE_ID: LEX-P0067-V0090
% #FIELD_VALUE: 工艺运行记录
\fieldvalue{\texttt{16:07:52}\quad 记录内容（原图文字过小）}} \\ \hline
\multicolumn{2}{|l|}{
% #VALUE_ID: LEX-P0067-V0091
% #FIELD_VALUE: 工艺运行记录
\fieldvalue{\texttt{16:07:53}\quad 记录内容（原图文字过小）}} \\ \hline
\multicolumn{2}{|l|}{
% #VALUE_ID: LEX-P0067-V0092
% #FIELD_VALUE: 工艺运行记录
\fieldvalue{\texttt{16:07:53}\quad 记录内容（原图文字过小）}} \\ \hline
\multicolumn{2}{|l|}{
% #VALUE_ID: LEX-P0067-V0093
% #FIELD_VALUE: 工艺运行记录
\fieldvalue{\texttt{16:07:53}\quad 记录内容（原图文字过小）}} \\ \hline
\multicolumn{2}{|l|}{
% #VALUE_ID: LEX-P0067-V0094
% #FIELD_VALUE: 工艺运行记录
\fieldvalue{\texttt{16:07:33}\quad 记录内容（原图文字过小）}} \\ \hline
\end{tabular}

\vspace{2pt}
第16页
\end{minipage}

\vfill

% #VALUE_ID: LEX-P0067-V0095
% #HANDWRITTEN: 1/14
\handwritten{1/14}
% LEXOID_PAGE_COMPLETED: 67/70

\newpage

\begin{center}
\scriptsize
\begin{tabular}{|p{0.12\linewidth}|p{0.32\linewidth}|}
\hline
\multicolumn{2}{c}{\textbf{工艺运行报告}}\\
\hline
16:07:33 & 分馏塔1 进料流量~30000Pa\\ \hline
16:07:34 & 分馏塔1 进料压力~500.0kPa\\ \hline
16:07:34 & 分馏塔1 进料温度~97.2℃\\ \hline
16:07:34 & 分馏塔1 回流量~12.1\\ \hline
16:07:34 & 分馏塔1 塔顶压力~101.2kPa\\ \hline
16:07:34 & 分馏塔1 塔顶温度~300.0℃\\ \hline
16:07:34 & 分馏塔1 塔底温度~300.0℃\\ \hline
16:07:34 & 分馏塔1 塔顶回流比~1、2、4、5、6、7、8、9、10\\ \hline
16:07:34 & 分馏塔1 塔顶冷凝器~1000mL\\ \hline
16:07:35 & 分馏塔1 塔底温度~300.0℃\\ \hline
16:08:33 & 分馏塔1 进料流量~50.2\\ \hline
16:08:33 & 分馏塔1 回流比~3000mL\\ \hline
16:08:33 & 分馏塔1 塔底压力~1\\ \hline
16:08:33 & 分馏塔1 塔顶温度~200.0mL\\ \hline
16:08:33 & 分馏塔1 塔顶温度~200.0mL\\ \hline
16:08:40 & 分馏塔1 进料流量~200.0\\ \hline
16:08:40 & 分馏塔1 塔顶压力~3000.0Pa\\ \hline
16:08:40 & 分馏塔1 回流比~6.5\\ \hline
16:08:40 & 分馏塔1 塔顶温度~300.0℃\\ \hline
16:08:40 & 分馏塔1 塔顶回流比~1、2、4、5、6、7、8、9、10\\ \hline
16:08:40 & 分馏塔1 塔顶冷凝器~1000mL\\ \hline
16:08:47 & 分馏塔1 塔底温度~200.0mL\\ \hline
16:08:47 & 分馏塔1 塔顶压力~300.0Pa\\ \hline
16:08:47 & 分馏塔1 回流比~1\\ \hline
16:08:47 & 分馏塔1 塔顶温度~100.0mL\\ \hline
16:08:47 & 分馏塔1 塔底温度~200.0mL\\ \hline
\end{tabular}
\hfill
\begin{tabular}{|p{0.12\linewidth}|p{0.32\linewidth}|}
\hline
\multicolumn{2}{c}{\textbf{工艺运行报告}}\\
\hline
16:08:53 & 分馏塔2 塔底压力~200150Pa\\ \hline
16:08:53 & 分馏塔2 塔顶压力~500.0kPa\\ \hline
16:08:53 & 分馏塔2 塔顶温度~97.2℃\\ \hline
16:08:53 & 分馏塔2 回流比~1\\ \hline
16:08:53 & 分馏塔2 塔顶温度~300.0℃\\ \hline
16:08:53 & 分馏塔2 塔底温度~100.0℃\\ \hline
16:08:53 & 分馏塔2 塔顶回流比~1、2、4、5、6、7、8、9、10\\ \hline
16:08:53 & 分馏塔2 塔顶冷凝器~1000mL\\ \hline
16:09:50 & 分馏塔2 进料流量~200.0\\ \hline
16:09:50 & 分馏塔2 回流比~1\\ \hline
16:09:51 & 分馏塔2 塔底温度~200.0mL\\ \hline
16:09:51 & 分馏塔2 塔顶温度~100.0mL\\ \hline
16:09:51 & 分馏塔2 塔顶回流比~1、2、4、5、6、7、8、9、10\\ \hline
16:09:53 & 分馏塔2 塔顶冷凝器~1000mL\\ \hline
16:09:53 & 分馏塔2 回流比~72.59\\ \hline
16:09:53 & 分馏塔2 塔底压力~1\\ \hline
16:09:53 & 分馏塔2 塔顶温度~300.0℃\\ \hline
16:09:53 & 分馏塔2 塔顶温度~200.0mL\\ \hline
16:09:53 & 分馏塔2 塔顶回流比~1、2、4、5、6、7、8、9、10\\ \hline
16:10:02 & 分馏塔2 塔底温度~75.23\\ \hline
16:10:02 & 分馏塔2 回流比~1\\ \hline
16:10:03 & 分馏塔2 塔顶温度~500.0mL\\ \hline
16:10:03 & 分馏塔2 塔底温度~1\\ \hline
16:10:03 & 分馏塔2 回流比~1\\ \hline
\end{tabular}

\vspace{0.8cm}

\begin{tabular}{|p{0.12\linewidth}|p{0.32\linewidth}|}
\hline
\multicolumn{2}{c}{\textbf{工艺运行报告}}\\
\hline
16:10:03 & 分馏塔3 塔底温度~200.0mL\\ \hline
16:10:03 & 分馏塔3 塔顶温度~100.0mL\\ \hline
16:10:05 & 分馏塔3 塔顶温度~500.0mL\\ \hline
16:10:10 & 分馏塔3 塔底压力~300.0Pa\\ \hline
16:10:10 & 分馏塔3 塔底温度~1\\ \hline
16:10:10 & 分馏塔3 塔顶温度~1\\ \hline
16:10:10 & 分馏塔3 塔顶温度~100.0mL\\ \hline
16:10:10 & 分馏塔3 塔顶回流比~1、2、4、5、6、7、8、9、10\\ \hline
16:10:14 & 分馏塔3 塔顶温度~500.0mL\\ \hline
16:10:14 & 分馏塔3 塔顶冷凝器~500.0mL\\ \hline
16:10:22 & 分馏塔3 塔顶冷凝器~500.0mL\\ \hline
16:10:23 & 分馏塔3 塔顶温度~500.0mL\\ \hline
16:10:25 & 分馏塔3 塔顶温度~500.0mL\\ \hline
16:10:32 & 分馏塔3 塔顶温度~500.0mL\\ \hline
16:10:33 & 分馏塔3 塔顶冷凝器~500.0mL\\ \hline
16:10:34 & 分馏塔3 塔顶冷凝器~500.0mL\\ \hline
16:10:43 & 分馏塔3 塔顶温度~500.0mL\\ \hline
16:10:44 & 分馏塔3 塔顶温度~500.0mL\\ \hline
16:10:55 & 分馏塔3 塔顶温度~500.0mL\\ \hline
16:11:09 & 分馏塔3 塔顶温度~500.0mL\\ \hline
16:11:19 & 分馏塔3 塔顶冷凝器~500.0mL\\ \hline
\end{tabular}
\hfill
\begin{tabular}{|p{0.12\linewidth}|p{0.32\linewidth}|}
\hline
\multicolumn{2}{c}{\textbf{工艺运行报告}}\\
\hline
16:11:20 & 分馏塔4 塔顶冷凝器~500.0mL\\ \hline
16:11:21 & 分馏塔4 塔顶冷凝器~500.0mL\\ \hline
16:11:28 & 分馏塔4 塔顶冷凝器~500.0mL\\ \hline
16:11:29 & 分馏塔4 塔顶冷凝器~500.0mL\\ \hline
16:11:30 & 分馏塔4 塔顶冷凝器~500.0mL\\ \hline
16:11:31 & 分馏塔4 塔顶冷凝器~500.0mL\\ \hline
16:11:39 & 分馏塔4 塔顶冷凝器~500.0mL\\ \hline
16:11:39 & 分馏塔4 塔顶冷凝器~500.0mL\\ \hline
16:11:40 & 分馏塔4 塔顶冷凝器~499.9mL\\ \hline
16:11:42 & 分馏塔4 塔顶冷凝器~500.0mL\\ \hline
16:11:52 & 分馏塔4 塔顶冷凝器~500.0mL\\ \hline
16:11:59 & 分馏塔4 塔顶冷凝器~500.0mL\\ \hline
16:12:08 & 分馏塔4 塔顶冷凝器~500.0mL\\ \hline
16:12:09 & 分馏塔4 塔顶冷凝器~500.0mL\\ \hline
16:12:10 & 分馏塔4 塔顶冷凝器~500.0mL\\ \hline
16:12:10 & 分馏塔4 塔顶回流比~1、2、4、5、6、7、8、9、10\\ \hline
16:12:23 & 分馏塔4 塔顶温度~500.0mL\\ \hline
16:12:23 & 分馏塔4 塔底压力~1\\ \hline
16:12:23 & 分馏塔4 回流比~1、2、4、5、6、7、8、9、10\\ \hline
16:12:25 & 分馏塔4 塔顶冷凝器~500.0mL\\ \hline
16:12:33 & 分馏塔4 塔顶温度~500.0mL\\ \hline
16:12:35 & 分馏塔4 塔顶温度~500.0mL\\ \hline
16:12:35 & 分馏塔4 塔底压力~1\\ \hline
16:12:38 & 分馏塔4 塔顶温度~500.0mL\\ \hline
16:12:39 & 分馏塔4 塔底温度~1\\ \hline
\end{tabular}
\end{center}

% #VALUE_ID: LEX-P0068-V0001
% #HANDWRITTEN: 1/14
\handwritten{1/14}
% LEXOID_PAGE_COMPLETED: 68/70

\newpage

\begin{center}
\begin{minipage}[t]{0.47\linewidth}
\centering
\textbf{蔚蓝生物}

\vspace{0.2em}
工艺运行报告

\vspace{0.3em}
\begin{tabular}{|p{0.18\linewidth}|p{0.75\linewidth}|}
\hline
16:12:23 & 分装器1 流量约为11.12 \\ \hline
16:12:23 & 分装器1 流量为400.00mL/min \\ \hline
16:12:23 & 分装器1 流量约为500.00mL/min \\ \hline
16:12:23 & 分装器1 流量约为1、2、3、4、5、6、7、8、9、10 \\ \hline
16:12:23 & 分装器1 流量约为500.00mL/min \\ \hline
16:12:33 & 分装器1 流量约为500.00mL/min \\ \hline
16:12:33 & 分装器1 流量约为500.00mL/min \\ \hline
16:12:39 & 分装器4 流量约为500.00mL/min \\ \hline
16:12:39 & 分装器4 流量约为11.12 \\ \hline
16:12:39 & 分装器4 流量为400.00mL/min \\ \hline
16:12:39 & 分装器4 流量约为500.00mL/min \\ \hline
16:12:39 & 分装器4 流量约为1、2、3、4、5、6、7、8、9、10 \\ \hline
16:12:39 & 分装器4 流量约为500.00mL/min \\ \hline
16:12:43 & 分装器2 流量约为500.00mL/min \\ \hline
16:12:43 & 分装器2 流量约为11.12 \\ \hline
16:12:43 & 分装器2 流量为400.00mL/min \\ \hline
16:12:43 & 分装器2 流量约为500.00mL/min \\ \hline
16:12:43 & 分装器2 流量约为1、2、3、4、5、6、7、8、9、10 \\ \hline
16:12:43 & 分装器2 流量约为500.00mL/min \\ \hline
16:12:36 & 分装器1 流量约为500.00mL/min \\ \hline
16:12:36 & 分装器1 流量约为11.12 \\ \hline
16:12:36 & 分装器1 流量为400.00mL/min \\ \hline
16:12:36 & 分装器1 流量约为500.00mL/min \\ \hline
16:12:36 & 分装器1 流量约为1、2、3、4、5、6、7、8、9、10 \\ \hline
16:12:36 & 分装器1 流量约为500.00mL/min \\ \hline
\end{tabular}

\vspace{0.4em}
第21页
\end{minipage}
\hfill
\begin{minipage}[t]{0.47\linewidth}
\centering
\textbf{蔚蓝生物}

\vspace{0.2em}
工艺运行报告

\vspace{0.3em}
\begin{tabular}{|p{0.18\linewidth}|p{0.75\linewidth}|}
\hline
16:13:36 & 分装器1 流量约为500.00mL/min \\ \hline
16:13:37 & 分装器1 流量约为11.12 \\ \hline
16:13:37 & 分装器1 流量为400.00mL/min \\ \hline
16:13:37 & 分装器3 流量约为500.00mL/min \\ \hline
16:13:37 & 分装器3 流量约为1、2、3、4、5、6、7、8、9、10 \\ \hline
16:13:37 & 分装器3 流量约为500.00mL/min \\ \hline
16:13:48 & 分装器3 流量约为500.00mL/min \\ \hline
16:13:48 & 分装器3 流量约为11.12 \\ \hline
16:13:48 & 分装器3 流量为400.00mL/min \\ \hline
16:13:50 & 分装器4 流量约为500.00mL/min \\ \hline
16:13:50 & 分装器4 流量约为1、2、3、4、5、6、7、8、9、10 \\ \hline
16:13:50 & 分装器4 流量约为500.00mL/min \\ \hline
16:13:56 & 分装器2 流量约为500.00mL/min \\ \hline
16:13:56 & 分装器2 流量约为11.12 \\ \hline
16:13:56 & 分装器2 流量为400.00mL/min \\ \hline
16:13:56 & 分装器2 流量约为500.00mL/min \\ \hline
16:13:56 & 分装器2 流量约为1、2、3、4、5、6、7、8、9、10 \\ \hline
16:13:56 & 分装器2 流量约为500.00mL/min \\ \hline
16:14:02 & 分装器4 流量约为500.00mL/min \\ \hline
16:14:02 & 分装器4 流量约为11.12 \\ \hline
16:14:06 & 分装器4 流量为400.00mL/min \\ \hline
16:14:06 & 分装器4 流量约为500.00mL/min \\ \hline
16:14:06 & 分装器4 流量约为1、2、3、4、5、6、7、8、9、10 \\ \hline
\end{tabular}

\vspace{0.4em}
第22页
\end{minipage}
\end{center}

\vspace{1em}

\begin{center}
\begin{minipage}[t]{0.47\linewidth}
\centering
\textbf{蔚蓝生物}

\vspace{0.2em}
工艺运行报告

\vspace{0.3em}
\begin{tabular}{|p{0.18\linewidth}|p{0.75\linewidth}|}
\hline
16:14:10 & 分装器1 流量约为500.00mL/min \\ \hline
16:14:15 & 分装器2 流量约为500.00mL/min \\ \hline
16:14:16 & 工艺结束 \\ \hline
16:15:23 & 停止运行，设备温度为37℃ \\ \hline
\end{tabular}

\vspace{1em}
审核：

% #VALUE_ID: LEX-P0069-V0001
% #FIELD_VALUE: 审核
% #TODO #HANDWRITTEN: 张俊杰?; handwriting is partly illegible
\fieldvalue{\handwritten{张俊杰?}}

日期：

% #VALUE_ID: LEX-P0069-V0002
% #FIELD_VALUE: 日期
% #TODO #HANDWRITTEN: 202?年12月3日; first digit of the year is unclear
\fieldvalue{\handwritten{202?年12月3日}}

\vspace{0.4em}
第23页
\end{minipage}
\hfill
\begin{minipage}[t]{0.47\linewidth}
\centering
\textbf{蔚蓝生物}

\vspace{0.2em}
工艺运行报告

\vspace{0.5em}
\begin{center}
\begin{tikzpicture}[x=0.075\linewidth,y=0.055\linewidth]
\draw[gray] (0,0) rectangle (9,9);
\draw[gray] (0,1.8) -- (9,1.8);
\draw[gray] (0,3.6) -- (9,3.6);
\draw[gray] (0,5.4) -- (9,5.4);
\draw[gray] (0,7.2) -- (9,7.2);
\draw[gray] (1.2,0) -- (1.2,9);
\draw[gray] (4.5,0) -- (4.5,9);
\draw[gray] (7.8,0) -- (7.8,9);
\draw[->] (0,0) -- (9.3,0);
\draw[->] (0,0) -- (0,9.3);
\node[anchor=south west] at (0,9) {流速(mL/min)};
\node[anchor=east] at (0,0) {0};
\node[anchor=east] at (0,1.8) {10};
\node[anchor=east] at (0,3.6) {20};
\node[anchor=east] at (0,5.4) {30};
\node[anchor=east] at (0,7.2) {40};
\node[anchor=east] at (0,9) {50};
\node[anchor=north] at (1.2,0) {10:25:31};
\node[anchor=north] at (4.5,0) {13:25:31};
\node[anchor=north] at (7.8,0) {16:25:31};
\end{tikzpicture}
\end{center}

\begin{itemize}
\item[] \textcolor{red}{\(\bullet\)} 流速1
\item[] \textcolor{teal}{\(\bullet\)} 流速2
\item[] \textcolor{purple}{\(\bullet\)} 流速3
\end{itemize}

\vspace{0.3em}
第24页

% #VALUE_ID: LEX-P0069-V0003
% #HANDWRITTEN: 13/14
\fieldvalue{\handwritten{13/14}}
\end{minipage}
\end{center}

% TODO: The page is a four-up scan. Several small table entries are blurred and may require verification against the source PDF.
% LEXOID_PAGE_COMPLETED: 69/70

\newpage

\begin{center}
\begin{minipage}[t]{0.46\linewidth}
\begin{center}
工艺运行报告
\end{center}

\begin{center}
\fbox{\parbox[c][7cm][c]{0.92\linewidth}{%
\centering
% TODO: 插入本页左侧温度趋势图。\\[0.5em]
纵轴：温度（$^\circ$C）\\
图例：温度\\
纵轴范围：$-10$ 至 $40$\\
横轴刻度：10:23:51，13:23:51，16:23:51
}}
\end{center}

\begin{center}
第25页
\end{center}
\end{minipage}
\hfill
\begin{minipage}[t]{0.46\linewidth}
\begin{center}
工艺运行报告
\end{center}

\begin{center}
\fbox{\parbox[c][7cm][c]{0.92\linewidth}{%
\centering
% TODO: 插入本页右侧压力趋势图。\\[0.5em]
纵轴：压力（Pa）\\
图例：压力1，压力2，压力3，压力4\\
横轴刻度：10:23:51，13:23:51，16:23:51
}}
\end{center}

\begin{center}
第26页
\end{center}
\end{minipage}
\end{center}

\vfill

\begin{flushright}
% #VALUE_ID: LEX-P0070-V0001
% #TODO #HANDWRITTEN: illegible; blue handwritten name or signature
\handwritten{\text{（姓名，难辨）}}\\[0.4em]
% #VALUE_ID: LEX-P0070-V0002
% #HANDWRITTEN: 2024.11.13
\handwritten{2024.11.13}\\[0.4em]
% #VALUE_ID: LEX-P0070-V0003
% #TODO #HANDWRITTEN: 14/14; first character and date format are uncertain
\handwritten{14/14}
\end{flushright}

\end{document}
% LEXOID_PAGE_COMPLETED: 70/70
